% 日本語 LuaLaTeX。2026-11-22 / v1.62
% エンジン、出力先、SyncTeX は repo ルートの .latexmkrc で指定する。
\RequirePackage{fix-cm}
\documentclass[a4paper,11pt]{ltjsarticle}
\usepackage{amsmath,amssymb,amsthm,mathtools,mathrsfs}
% rsfs の小さな添字も、代替サイズではなく設計フォントの拡縮で組む。
\DeclareFontFamily{U}{rsfs}{\skewchar\font127}
\DeclareFontShape{U}{rsfs}{m}{n}{<-6> rsfs5 <6-8> rsfs7 <8-> rsfs10}{}
\usepackage[margin=25mm]{geometry}
\usepackage{booktabs,enumitem}
\usepackage[unicode,pdfusetitle,colorlinks=true,
 linkcolor=blue,citecolor=blue,urlcolor=blue]{hyperref}
\numberwithin{equation}{section}
\theoremstyle{plain}
\newtheorem{theorem}{定理}[section]
\newtheorem{proposition}[theorem]{命題}
\newtheorem{lemma}[theorem]{補題}
\newtheorem{corollary}[theorem]{系}
\theoremstyle{definition}
\newtheorem{definition}[theorem]{定義}
\newtheorem{example}[theorem]{例}
\newtheorem{problem}[theorem]{問題}
\theoremstyle{remark}
\newtheorem{remark}[theorem]{注意}
\renewcommand{\proofname}{証明}
\newcommand{\kk}{k}
\newcommand{\OO}{\mathcal O}
\newcommand{\cC}{\mathcal C}
\newcommand{\PP}{\mathbb P}
\newcommand{\HH}{\mathcal H_2}
\newcommand{\HHstack}{\mathfrak H_2}
\DeclareMathOperator{\Hom}{Hom}
\DeclareMathOperator{\End}{End}
\DeclareMathOperator{\Ext}{Ext}
\DeclareMathOperator{\coh}{coh}
\DeclareMathOperator{\qgr}{qgr}
\DeclareMathOperator{\gldim}{gldim}
\DeclareMathOperator{\Aut}{Aut}
\DeclareMathOperator{\Pic}{Pic}
\DeclareMathOperator{\Sym}{Sym}
\DeclareMathOperator{\Gr}{Gr}
\DeclareMathOperator{\GL}{GL}
\DeclareMathOperator{\PGL}{PGL}
\DeclareMathOperator{\rank}{rank}
\DeclareMathOperator{\ch}{ch}
\DeclareMathOperator{\td}{td}
\DeclareMathOperator{\length}{length}
\title{構造対象を指定した非可換二次曲面：\\平方根と二次 AS 正則表示の忘却}
\author{}
\date{2026年11月22日\quad v1.62（完全分類を中心とする原稿の再編）}

\begin{document}
\maketitle
\begin{abstract}
構造対象 $\OO$ のみを指定した強い非可換二次曲面
$(\cC,\OO)$ を pointed exact 同値で分類する。
$\OO$ は Gabriel--Krull 商で rank-one torsion-free exceptional とし、
反 Serre 関手の exact ample な平方根と、その軌道環の
型 $(1,4,6,4,1)$ の二次 AS 正則代数による正則正規二次商表示を要求する。
平方根・その同型・ambient は標識として固定しない。

体上の全対象は、linear データと、被約・非被約を含む
Van den Bergh 許容四つ組から構成される cubic 領域、
および中央二重平方表示を持つ noncubic 領域に分かれる。
前者では四生成・七二次関係の有限代数的判定によって、
選んだ平方根を忘れた pointed exact 同値を区別する。
後者は $\PP^1$ 上の二つの被約二点集合の非順序対の
同時射影同値で分類できる。全退化型を含み、
未知の AS 条件や圏同値の解決を分類条件には残さない。

さらに楕円四つ組の指定範囲では、構造対象の
$\mathbb Z^2$ 格子と $\OO$-標識忘却の有限／可算無限ファイバーを記述する。
相対モジュライの有限型Artin性・noncubic 境界の幾何は
別稿 \cite{QuadricModuli2026} の対象とする。
既知の cubic 分類と本稿の pointed 同値分類を区別し、
包括的な新規性・外部査読の確定は主張しない。
\end{abstract}
\tableofcontents

\section{Introduction}
\label{sec:introduction}

滑らかな二次曲面 $Q=\PP^1\times\PP^1$ の偏極
$H=\OO_Q(1,1)$ は $H^{\otimes2}=\omega_Q^{-1}$ を満たす。
従って $\coh Q$ の Serre 関手 $S$ に対し
$s=-\otimes H$ は $s^2\simeq S^{-1}[2]$ という平方根である。
本稿ではこの性質と、非可換射影三次元空間内の二次超曲面としての表示を用い、
構造対象付き圏 $(\cC,\OO)$ の分類を目標とする。
以下の $\HH$ の添字 $2$ は超曲面の次数を表す。
可換例の反標準次数は $(-K_Q)^2=8$ であり、次数二の del Pezzo 曲面とは異なる。

三次超曲面の切断環では、二次関係だけから ambient を回復できる場合がある。
しかし二次超曲面では曲面を切り出す式自体が二次関係に含まれる。
切断環 $B$ の二次関係は七次元、ambient の二次関係は六次元となるため、
通常の quadratic envelope をそのまま用いることはできない。
適切な強い条件は、正則正規元 $g\in A_2$ による表示 $B\simeq A/(g)$ の存在である。
この修正は ambient の選択を新たな marking として加えるものではない。

\subsection{分類の中心結果}
\label{sec:intro-classification}

本稿の主定理は定理~\ref{thm:complete-h2-classification} である。
許容データから対象を構成し、任意の強い構造対象付き二次曲面から
cubic/noncubic のいずれかのデータを回復する。
その際、選択した Serre 平方根、helix foundation や ambient を
最終的な同値関係に含めず、四生成環の二次関係と根対配置によって
pointed exact 同値を判定する。
幾何的四つ組の同型や通常の Zhang twist と同一視しない。
既知の cubic $\mathbb Z$-algebra の分類は
Van den Bergh \cite{VdB2011} を入力とするが、
構造対象と平方根の忘却の扱いは別に証明する。
Normal helix、適合二周期の Serre 平方関係、AS4 被覆、
最初の軌道対象の内在的認識は、この分類を支える主要な道具である。

\subsection{相対モジュライ別稿との分担}

局所有限型基底の内在的相対 $\mathfrak H_2$ の有限型smooth Artin性、
cubic/noncubic 族の相対効果性、framed 実現、
被約noncubic境界の群商・慣性・good/coarse moduli と形式局所商は
別稿 \cite{QuadricModuli2026} で扱う。
体上の分類だけから非被約基底の相対結論は従わない。
現在は原稿再編の途中であり、本稿の後半には参照のため
相対定理の主張を残している。分類側の依存を検査した後、
不要な相対節の \verb|\input| と重複した説明を除く予定である。
本稿の分類主定理の証明と、別稿の代数性の公表前監査は区別する。

\subsection{先行研究と本稿の範囲}

相対 tails の導来認識では、Bondal--Van den Bergh
\cite[\S4.1, Lemma 4.2.2]{BondalVanDenBergh2003} の torsion 局所化の方法を用いる。
その体上の定理を一般基底へそのまま適用せず、
有限 projective 係数の cubic 解消から acyclicity と有限 cohomological dimension を
証明し直す（第\ref{sec:relative-cubic-tails}節）。
noncubic 境界では極小 Betti 数が跳ぶため、五頂点の非極小 bar 解消へ置き換える
（第\ref{sec:noncubic-local-tails}節）。
有限 perfect 複体の特殊 fiber の集中から同じ局所環上の全前向き Hom を制御する。
この段階には perfect 複体の標準的な fiber 判定\cite[Tag 0BDK]{StacksPerfectFibers} を用いる。
これにより任意基底での canonical な非有界導来比較を得る。
Noetherian 性は別に証明する必要があり、第\ref{sec:relative-cubic-noetherian}節で
相対 point curve と正規族を scheme-theoretic に構成して解決する。
曲線商の Noetherian 性には Keeler \cite[Theorem 1.5, 付属 Erratum \S2]{Keeler2003}
の一般 Noetherian 基底上の消滅を用いる。
linear 部分の非被約な Segre 表示と、
第\ref{sec:formal-noetherian-patching}節の次数別形式的貼合せによって、
単一の基底全体へ結論を移す。
Noetherian 性の下では、有限解消と有限 presentation から
一様 Serre 消滅と全相対条件を導ける。
一方、非可換射影直線束の候補対角列を使う比較には具体的な反例がある
（相対モジュライ別稿~\cite{QuadricModuli2026}）。
既知の束の Noetherian 性 \cite[Theorem 1.2]{VanDenBergh2012P1} は保たれるが、
その候補列は非零 coherent 対象を生成しない。


Artin--Zhang の再構成理論 \cite{AZ1994} は ample pair から切断環と圏を対応させる。
本稿では ampleness と再構成の正確な形として
Mori--Ueyama \cite[Definition 2.5, Theorem 2.6]{MU2021} を用いる。
Van den Bergh \cite[Theorem 5.5.10]{VdB2011} は cubic 三次元
AS 正則 $\mathbb Z$-algebra を geometric quintuple 等で分類し、
同論文の Theorem 6.1、Proposition 6.6、Corollary 6.7 は
二次の正則正規元を用いた四次元二次 AS 正則環による表示を与える。
ここで番号は arXiv v3 の番号である。
その quadratic hull の一意性には次数 $-4$ の自己同型の固定が含まれる。
単に切断環だけを固定した ambient 一意性とは異なる。

Smith--Van den Bergh \cite{SVdB2013} は quantum $\PP^3$ 内の
中心二次切断に対し、Clifford 型代数、smoothness、ruling の関係を研究した。
その standard 性に関する記述には
Mori--Ueyama \cite[Section 5, Appendix]{MU2021} による訂正がある。
同じ tails 圏でも Zhang twist により standard でない表示が生じるため、
表示の性質を無条件に $(\cC,\OO)$ の不変量とはしない。
Mori--Ueyama--Wu \cite[Theorems 4.10, 4.12, 5.5]{MUW2026} は、
二つの ruling の作用による Zhang twist の standard 性判定と、
滑らかな中心 Sklyanin quadric の standard 性を証明した。
同論文の一般的な定義には ambient の Auslander 正則性と
Cohen--Macaulay 性という追加仮定があるため、型 $(1,4,6,4,1)$ だけから
その全結果を本稿の全対象へ適用しない。
また同論文の Introduction は、正規切断を graded Morita 同値で
常に中心切断へ置き換えられるかを未解決としている。
Mori--Nyman \cite[Theorem 6.13, arXiv v1]{MN2025} は、中心切断、domain、
graded isolated singularity、standard という仮定の下で
Gorenstein parameter 四の三次元 AS 正則 $\mathbb Z$-algebra への比較を与える。
本稿の命題~\ref{prop:standard-pointed} はこの条件下で構造対象と平方根の像を記述する。
定理~\ref{thm:normal-helix} は、二つの MCM 型を使い分けることで、中心性・standard 性・ambient の条件 $(*)$ を要しない比較を与える。

Okawa--Ueda \cite[Theorem 1.1]{OU2014} の geometric quintuple の GIT 商
$\PP(2,4,4,6)$ と conifold potential の GIT 商 $\PP(1,2,3,4)$、
および後者から前者への次数四の射は、既存の重要なモジュライ記述である。
しかしその tensor や helix の標識を忘れた同値関係まで自動的に与えるわけではない。
本稿の可換点における四つの平方根と、その次数四の射を同一視する主張も行わない。
Abdelgadir--Okawa--Ueda \cite[Definition 3.5, Proposition 3.6]{AOU2014} は、
有限非巡回箙の関係式を固定次元の Grassmannian と道代数の自己同型群からなる
有限型スタックとして構成する。Mori--Okawa--Ueda \cite[\S6]{MOU2019} は、
sheaf bimodule、関係式、曲線と束のスタックを比較する。
本稿の有限 foundation の構成はこの有限代数による方法に対応するが、
標識の忘却と全 $\HH$ の相対再構成を別の問題として扱う。
固定した dg 圏の対象についての To\"en--Vaqui\'e の代数性定理
\cite[Theorem 3.6, Corollaries 3.15, 3.24]{TV2007} も、
全ての変動する $\HH$ 圏のモジュライの代数性を直ちには与えない。


変形論の入力は Lowen--Van den Bergh
\cite[Theorem 3.1; §§7.7--7.9, author preprint]{LowenVanDenBergh2005} と
\cite{LowenVanDenBergh2006}、および De Deken--Lowen
\cite[Theorems 5.13, 5.15]{DeDekenLowen2011} の既知定理である。
本稿では normal helix の compact 生成性と geometric Hom の条件を確認した上で適用する。
これにより全 Artin Abel 変形の canonical 反 Serre 偏極を再構成し、
関係式変形から Abel 変形への忘却の接核を
$\operatorname{HH}^1$ と同定する。
また、可逆双加群の曲率障害と近恒等群の平方根補正を使って Serre 根を持ち上げ、
可換点では全 Artin 変形を相対 $\mathfrak H_2$ に実現する。
被覆の構成には Bauwens--Van den Bergh \cite{BauwensVanDenBergh1998} と
Van den Bergh \cite[\S6]{VdB2011} の canonical hull を参照するが、
二周期の相対 tensor 恒等式、正則性、Koszul 双対の平坦性は別途証明する。
定理~\ref{thm:universal-cubic-hull} はこれを任意の geometric 二周期族へ拡張し、
逆 coevaluation の核から全次数 projective な AS4 被覆を構成する。
全 row・column の普遍 Noetherian 性には二周期自体も不要である
（定理~\ref{thm:relative-geometric-cubic-noetherian}）。
被覆の存在と、同じ二周期が Serre 平方根になることは区別する。
一次 Jacobian 被覆は全九次元の具体式を与える。
Toda \cite[Theorem 4.7, arXiv v3]{Toda2009} の一次 Fourier--Mukai 持上げ定理は
選択済みの四根の一次条件に用いるが、全 Artin 段階の結論は別途証明する。
To\"en--Vaqui\'e \cite[Corollaries 3.22, 3.24]{TV2007} の一般基礎環上の
自己同値モジュライに、標準根の全 Artin 存在・一意性を適用すると、
可換点の foundation 表示の \'etale 近傍に実際の Morita 根を構成できる。
さらに標準 foundation の二項移動と四項窓を代数化し、三次関係と
四次重なりを有限局所自由な tensor として認識する
（命題~\ref{prop:quadric-etale-finite-window}）。
可換点の局所環では全次数の三次提示・標準 cubic 解消と
偶数環の相対 Koszul 性まで得る
（定理~\ref{thm:quadric-local-all-degree-cubic}）。
幾何点の cubic AS 正則性は開近傍上で確認できる
（命題~\ref{prop:quadric-geometric-quintuple-open}）。
全次数の一様な平坦性と内在 Hom の同定は
命題~\ref{prop:quadric-uniform-cubic-hom-open} で得られる。
さらに一般 canonical hull により相対 AS4 被覆を同じ開上に構成できる
（命題~\ref{prop:quadric-zariski-homological-as-cover}）。
この構成に canonical tails 比較を合わせ、
定理~\ref{thm:quadric-finite-type-effectivity} で
有限型開上の内在的相対実現を得る。
この族の全閉点で De Deken--Lowen の helix 持上げを適用し、
指定された pointed 同型まで延長することで、
系~\ref{cor:quadric-effective-smooth-neighborhood} の滑らかな Artin 開近傍を構成する。
有限表示な各基底変更への Artin 判定と smooth groupoid の商には
\cite[Tags 02HY, 04X0, 04TK]{StacksSmoothAtlas} を用いる。
これらの適用を含む包括的な先行研究上の新規性は、なお確認を要する。

Presotto \cite[Theorem 2.8; Appendix A]{Presotto2018} は、一般の cubic quadric が
$\mathbb Z$-domain であることを証明した。この定理に整性・被約性の仮定はなく、
本稿では退化曲線から作る even part の domain 性に適用する。
一方、二周期の存在だけからその平方が反 Serre 関手であるとは結論できない。
本稿は四次テンソルの巡回写像を双対化双加群の作用と直接比較し、
全ての被約・非被約 $(2,2)$ 曲線について、束の貼合せと切断上の二周期を記述する。

主結果は定理~\ref{thm:complete-h2-classification} である。
既知の cubic データの分類を入力としつつ、全 $\HH$ のうち cubic 表示を
持たない領域も回復し、その領域を二点集合の対で分類する。
さらに、最初の軌道対象の内在的認識により、選んだ根を保持しない pointed exact 同値を
有限の代数的条件へ帰着する。この判定は未知の AS 正則性や
未知の自己同値群の計算を入力とせず、特殊な退化点も含む。
既存の cubic 分類、$\mathbb Z$-domain 定理、embedding 定理を本稿の新しい成果とはしない。
本稿の直接証明に関する包括的な先行研究調査と新規性の確定は、なお公表前の課題である。

可換対象の認識、四つの適合根、具体的な Sklyanin・Ore・多パラメータ族は、
定義と忘却を確認する独立した基準例として保持する。
可換認識には一次生成性と $\dim B_1=4,\dim B_2=9$ が主要な入力であり、
楕円型の具体例には \cite[arXiv v3, Proposition 7.1]{BSW2010}、
\cite[arXiv v2, Lemma 1.5]{RRZ2014}、\cite[\S 10.5]{SVdB2013} を用いる。
これらの例に複素数体や無限位数を仮定する箇所があっても、主定理は
標数零の任意の代数閉体上で成り立つ。

八本・十本・十二本矢の表示モジュライと反標準 avatar の関係は、別稿
\cite{QuiverAvatarModuli2026} で扱う。同稿では二つの隣接乗法を保持する
有限型 foundation スタックを構成し、二種類の十本矢開部分の交わりを
八本矢部分として記述する。また、直積型 collection の mutation、
$\Sigma_2$ を含む avatar の族、内在的忘却に必要な相対条件を比較する。
この有限表示の開性は、本稿の $\mathfrak H_2$ 実現部分の開性とは別である。

\section{構造対象と存在条件}
\label{sec:definition}

以下、$k$ は標数零の代数閉体とする。
圏は本質的に小さい $k$-線形圏とし、Abel 圏が noetherian であるとは
その全対象が noetherian であることをいう。
次数付き環については右加群を用いる。

\begin{definition}[Gabriel--Krull rank one torsion-free]
\label{def:gk-rank}
noetherian Abel 圏 $\cC$ の有限長対象からなる Serre 部分圏を $\cC_0$、
$\cC/\cC_0$ で有限長となる対象の逆像を $\cC_1$ とし、
$q:\cC\to\cC/\cC_1$ を商関手とする。
$\OO\in\cC$ が Gabriel--Krull rank one torsion-free であるとは、
\[
 q\OO\text{ が非零単純であり、}\qquad
 \Hom(T,\OO)=0\quad(T\in\cC_1)
\]
であることをいう。これは ``GK dimension が $1$'' という条件ではない。
さらに $\OO$ が exceptional であるとは
$\End(\OO)=k$ および $\Ext^i(\OO,\OO)=0$（$i>0$）をいう。
\end{definition}

exact 自己同値は有限長性と Serre 商の有限長性を保つので、
$\cC_0,\cC_1$ を保つ。従ってこの rank と torsion-free 条件は exact 同値で不変である。

\begin{definition}[Artin--Zhang ample pair]
\label{def:ample}
$s$ を $\cC$ の exact 自己同値とする。$(\OO,s)$ が ample であるとは、
次の二条件を満たすことをいう。
\begin{enumerate}[label=(A\arabic*)]
\item 任意の $M\in\cC$ に対し、ある $n_1,\ldots,n_r\ge0$ と全射
      $\bigoplus_{j=1}^r s^{-n_j}\OO\twoheadrightarrow M$ が存在する。
\item 任意の全射 $M\twoheadrightarrow N$ に対し、十分大きい全ての $n$ で
      $\Hom(s^{-n}\OO,M)\to\Hom(s^{-n}\OO,N)$ が全射である。
\end{enumerate}
切断環を
\[
 B(\cC,\OO,s)=\bigoplus_{n\ge0}\Hom(\OO,s^n\OO),
 \qquad a*b=s^n(a)\circ b\quad(a\in B_m,\ b\in B_n)
\]
と定める。この積は右加群の再構成と整合する規約である。
\end{definition}

\begin{definition}[型 $(1,4,6,4,1)$]
\label{def:as14641}
$A$ は局所有限、連結、次数一生成、関係イデアルが二次で生成され、左右 noetherian とする。
$A$ がここでいう型 $(1,4,6,4,1)$ とは、Koszul であり、
global dimension と Gorenstein parameter がともに $4$ の AS 正則環で、
自明加群の極小分解が
\[
 0\to A(-4)\to A(-3)^4\to A(-2)^6\to A(-1)^4\to A\to k\to0
\]
となることをいう。AS 条件は左右で要求する。
特に $\Ext_A^i(k,A)$ は $i=4$ 以外で零、$i=4$ で $k(4)$ であり、
$H_A(t)=(1-t)^{-4}$ である。
元 $g$ が正則正規であるとは、左右の零因子でなく、$gA=Ag$ となることをいう。
\end{definition}

\begin{definition}[$\OO$-marked 非可換二次曲面の強いクラス]
\label{def:h2}
$\HH$ の対象は対 $(\cC,\OO)$ であり、次を要求する。
\begin{enumerate}[label=(Q\arabic*)]
\item $\cC$ は $k$-線形 noetherian Abel 圏で、$D^b(\cC)$ は Hom 有限かつ
      Serre 関手 $S$ を持つ。
\item $\OO$ は定義~\ref{def:gk-rank} の rank one torsion-free exceptional 対象である。
\item exact 自己同値 $s$ が存在し、その導来自己同値について
      $s^2\simeq S^{-1}[2]$ かつ $(\OO,s)$ が Artin--Zhang ample である。
\item \emph{同じ $s$} に対し、型 $(1,4,6,4,1)$ の環 $A$ と
      正則正規元 $g\in A_2$ が存在し、全次数で
      $B(\cC,\OO,s)\simeq A/(g)$ という次数付き環同型がある。
\end{enumerate}
$s$、根の同型 $s^2\simeq S^{-1}[2]$、$A$、$g$ および商表示は指定構造ではない。
対象間の同値は $k$-線形 exact 同値 $F:\cC\xrightarrow{\sim}\cC'$ と
$F(\OO)\simeq\OO'$ によるものとする。
\end{definition}

\begin{remark}
条件 (Q4) は一次生成性と Hilbert 関数を含意する。
これらを平方根と豊富性のみの帰結とはしない。
また $g$ の中心性は定義に含めない。中心性を課した表示だけを先に調べる場合も、
それが全対象を覆うかは別途証明する必要がある。
単なる $\qgr A/(g)\simeq\cC$ だけでなく、指定 $\OO$ と選んだ $s$ に対する
\emph{全切断環} の表示を要求していることにも注意する。
\end{remark}

平方根を保持した補助対象 $(\cC,\OO,s)$、さらに ambient 表示を保持した補助対象から、
最終対象 $(\cC,\OO)$ への忘却を考える。
平方根付きの同値では $Fs\simeq s'F$ も要求する。
根同型を固定した場合の整合性や慣性は、本稿では分類しない。
この区別により、表示の重複を異なる二次曲面として数えることを避ける。

% !TEX root = ../nc-quadric-surface.tex
\section{全対象の分類定理}
\label{sec:classification}

この節では分類のデータ、許容条件、同値関係を先に述べる。
構成と被覆の証明は後続の各節に置く。
平方根と ambient は構成に使うが、最終的な同値関係では保持しない。

\subsection{cubic 側の幾何データと四生成環}

\begin{definition}[許容四つ組]
\label{def:classification-admissible}
$P_{\mathrm{cub}}$ は一つの記号 $\mathrm{lin}$ と、次の四つ組
$U=(C,L_0,L_1,L_2)$ 全体との直和とする。
$C$ は射影曲線、$L_i$ は直線束であり、次を満たす。
\begin{enumerate}[label=(D\arabic*)]
\item $h^0(C,L_i)=2$ で、$L_i$ は大域切断で生成される。
      隣接する組 $(L_0,L_1)$ と $(L_1,L_2)$ による射はそれぞれ
      $C$ を $\PP^1\times\PP^1$ の二次数 $(2,2)$ の Cartier 因子として閉埋め込む。
\item 各既約成分 $C_j\subset C_{\mathrm{red}}$ 上で
      $\deg(L_0|_{C_j})=\deg(L_2|_{C_j})$ である。
\item $L_0\not\simeq L_2$ である。
\end{enumerate}
これは \cite[arXiv v3, Proposition 5.1.2]{VdB2011} の許容四つ組の条件である。
$C$ の整性、被約性、滑らかさは要求しない。
記号 $\mathrm{lin}$ は $(\coh(\PP^1\times\PP^1),\OO)$ を表す。
\end{definition}

四つ組から $L_{i+3}=L_i^{-1}\otimes L_{i+1}\otimes L_{i+2}$ により
両側の helix を作り、
\[
 V_i=H^0(C,L_i),\qquad
 R_i=\ker\bigl(V_i\otimes V_{i+1}\otimes V_{i+2}
       \longrightarrow H^0(C,L_i\otimes L_{i+1}\otimes L_{i+2})\bigr)
\]
と置く。$\dim V_i=\dim R_i=2$ である。
これらを隣接空間と三次関係とする $\mathbb Z$-algebra を $D(U)$ とする。
定理~\ref{thm:all-admissible-inverse} の構成に従い、
反 Serre の四周期 $\theta$ と $\beta^2=\theta$ を対角型まで満たす
二周期 $\beta:D(U)\to D(U)(2)$ を一つ選ぶ。
この選択の存在は全ての許容データで証明されている。
具体的には曲線の自己同型 $\sigma$ と束の同型
$\rho_i:\sigma^*L_i\xrightarrow\sim L_{i+2}$ を取り、
$\beta_i=H^0(\rho_i)\circ\sigma^*$ とする。

$V=V_0\otimes V_1$ とし、$\beta_{01}=\beta_0\otimes\beta_1:
V\to V_2\otimes V_3$ と書く。有限次元の空間
\begin{equation}
 K_U=(1_V\otimes\beta_{01}^{-1})
       (R_0\otimes V_3+V_0\otimes R_1)\subset V\otimes V,
 \qquad B_U=T(V)/(K_U)
 \label{eq:classification-cubic-kernel}
\end{equation}
を作る。$\dim V=4$、$\dim K_U=7$ である。
ここで $R_0\otimes V_3$ と $V_0\otimes R_1$ の交わりは
四次テンソルの張る一次元空間 $W_0$ である。
三次関係を偶数頂点間でまとめれば、二次環 $B_U$ は
\[
 (B_U)_n=D(U)_{0,2n},\qquad a*b=a\beta^m(b)
 \quad(a\in(B_U)_m)
\]
と一致する。これは曲線 $C$ の切断環ではない。
全次数での二次表示は、偶数長の道に現れる $R_j$ の開始添字が偶数なら
右に一文字、奇数なら左に一文字を加え、上の二種類の四文字関係にまとめることで従う。
特に次数四の全切断乗法の核を課すと曲線への追加商となるので、
式~\eqref{eq:classification-cubic-kernel} の七次元空間を用いる。
linear データでは、$B_{\mathrm{lin}}$ を通常の Segre 環
$k[x_{00},x_{01},x_{10},x_{11}]/(x_{00}x_{11}-x_{01}x_{10})$ とする。

\subsection{cubic データ上の有限同値判定}

各 $V$ の基底を取り、上で得た二次関係を
$K,K'\subset k^4\otimes k^4$ とする。
$G=\GL_4(k)$ と、方程式で与えられる閉部分集合
\[
 E(K,K')=\{(g,h)\in G\times G\mid(g\otimes h)K=K'\}
\]
を使い、次の反復を定める。
\begin{equation}
 \begin{split}
 C_0&=G,\\
 C_{n+1}&=\operatorname{pr}_1(E\cap(G\times C_n))
       \cap\operatorname{pr}_2(E\cap(C_n\times G)).
 \end{split}
 \label{eq:classification-finite-relation}
\end{equation}
命題~\ref{prop:first-orbit-finite-equivalence} により、この下降列は有限段で安定する。
安定集合が非空のとき $U\sim_{\mathrm{cub}}U'$ と定める。
行列式を可逆とした座標環で、射影の消去と根基イデアルの比較だけを用いる判定である。
標数零の抽象的な代数閉体上では有限の代数的判定を意味し、
機械上で実行するには係数体が有効に与えられていることを前提とする。
基底、束の表示、二周期 $\beta$ の選び方によらないことは、次の定理で示す。

\subsection{noncubic データと主定理}

$P_{\mathrm{nc}}$ を、$\PP^1$ 上の被約な二点集合 $Z_1,Z_2$ の非順序対の全体とする。
同値関係 $\sim_{\mathrm{nc}}$ は $\PGL_2$ の同時作用である。
それぞれの零点が $Z_i$ となる非退化二次形式 $f_i(u,v)$ を選び、
\[
 C_{Z_1,Z_2}=k[u,v]\langle c,d\rangle/
       (c^2-f_1(u,v),\ d^2-f_2(u,v))
\]
と置く。全ての生成元の次数は一で、$u,v$ は中央とする。
この表示自体の次数シフトが適合根であるとは要求しない。
適合根と AS 正則被覆は定理~\ref{thm:ore-anticommuting-central} で構成する。

\begin{theorem}[全 $\HH$ の pointed exact 分類]
\label{thm:complete-h2-classification}
標数零の代数閉体 $k$ 上で、次の対応は全単射である。
\[
 (P_{\mathrm{cub}}/{\sim_{\mathrm{cub}}})
 \ \sqcup\ (P_{\mathrm{nc}}/{\sim_{\mathrm{nc}}})
 \ \xrightarrow{\ \sim\ }\ 
 \HH/{\text{pointed exact equivalence}}.
\]
cubic 側では $U$ を $(\qgr B_U,\pi B_U)$ へ、
noncubic 側では $\{Z_1,Z_2\}$ を
$(\qgr C_{Z_1,Z_2},\pi C_{Z_1,Z_2})$ へ送る。
両領域の像は交わらず、前者は pointed cubic 表示を持つ対象の全体、
後者はそのような表示を持たない対象の全体である。

noncubic 側の商は、二対が一致する一類、一点を共有する一類と、
四点が異なる場合の $\lambda\in k\setminus\{0,1\}$ を
$\lambda\sim\lambda^{-1}$ で割った一パラメータ族からなる。
この領域で non-standard な適合表示を持つのは、二対が一致する一類だけである。
\end{theorem}

\begin{proof}
まず全許容四つ組の $D(U)$ は cubic AS 正則 $\mathbb Z$-algebra である
\cite[arXiv v3, Theorems 5.5.7--5.5.10]{VdB2011}。
定理~\ref{thm:all-admissible-inverse} は、全曲線型で適合二周期を構成し、
even part の全切断環としての表示、型 $(1,4,6,4,1)$ の被覆、
正則正規元、Gabriel--Krull 条件、exceptionality、AZ ampleness を確認する。
従って cubic 側の写像は定義される。
linear データは通常の二次曲面であり、$-\otimes\OO(1,1)$ を使えばよい。
noncubic 側の構成と同値判定は定理~\ref{thm:noncubic-classification} による。

任意の $(\cC,\OO)\in\HH$ を取る。
定理~\ref{thm:cubic-coverage} の階数判定により、pointed cubic 表示を持つか、
二つの MCM 型の階数がともに三であるかに分かれる。
前者では \cite[arXiv v3, Theorem 5.5.10]{VdB2011} の分類により、
linear データまたは定義~\ref{def:classification-admissible} の四つ組から得られる。
後者では定理~\ref{thm:noncubic-classification} により中央二重平方環へ帰着する。
これで全被覆と二領域の非交差を得た。

cubic 側の各 $B_U$ は、構成した同じ pointed 圏の適合切断環である。
命題~\ref{prop:first-orbit-finite-equivalence} は二つの適合切断環について、
反復~\eqref{eq:classification-finite-relation} の非空性が
pointed exact 同値と必要十分であることを示す。
従って $\sim_{\mathrm{cub}}$ は同値関係であり、写像は単射である。
一つの四つ組から異なる基底・束同型・適合二周期を選んでも、
いずれも $(\qgr D(U),\pi e_0D(U))$ と pointed 同値なので、同じ判定で同一視される。
具体的には $\check B_{ij}=B_{j-i}$ から $D(U)_{2i,2j}$ への写像
$a\mapsto\beta^i(a)$ が積と整合するので、偶数頂点の圏と添字零の対象を保つ。
未知の自己同値群や全ての平方根の列挙を、判定の前提として残してはいない。
noncubic 側の最後の具体的記述も定理~\ref{thm:noncubic-classification} による。
\end{proof}

\begin{remark}[分類の範囲]
主定理は全対象の同値類の分類である。
cubic 側の同値を幾何的な少数の正規形へさらに簡約すること、
全自己同値、平方根と ambient の各忘却ファイバー、
族の代数性と相対モジュライは別の問題である。
Sklyanin pencil の局所的な候補数の評価に現れる追加の単射性条件は、
式~\eqref{eq:classification-finite-relation} では使わず、
主定理の許容領域を狭めない。
既知入力と本稿の直接証明は区別しており、後者の包括的な新規性は未確認である。
\end{remark}


\section{内在的条件からの基本的帰結}
\label{sec:categorical-consequences}

\begin{proposition}[大域次元]
\label{prop:gldim-two}
条件 (Q1), (Q3) の平方根が存在し、$\cC\ne0$ ならば $\gldim\cC=2$ である。
\end{proposition}
\begin{proof}
$S\simeq s^{-2}[2]$ なので、Serre 双対性から
\[
 \Ext^i(M,N)\simeq\Hom_{D^b(\cC)}(s^2N,M[2-i])^* .
\]
$i>2$ では heart の対象間の負次数 Hom であるため零である。
非零 $N$ に対し $\Ext^2(s^2N,N)\simeq\End(s^2N)^*\ne0$ なので
大域次元は正確に $2$ である。
\end{proof}

\begin{lemma}[非零切断は単射]
\label{lem:sections-monic}
$\OO$ が rank one torsion-free で $s$ が exact 自己同値ならば、
任意の整数 $n$ と非零射 $f:\OO\to s^n\OO$ に対し $f$ は単射である。
従って非負次数の切断環 $B(\cC,\OO,s)$ は domain である。
\end{lemma}
\begin{proof}
$s$ は $\cC_1$ を保つため $q\OO$ と $q(s^n\OO)$ は非零単純であり、
$s^n\OO$ も torsion-free である。
$qf=0$ ならば $\operatorname{im}f\in\cC_1$ となり、
$\operatorname{im}f\hookrightarrow s^n\OO$ は零となるので $f=0$ に矛盾する。
よって $qf$ は単純対象間の非零射、従って同型であり、$\ker f\in\cC_1$ となる。
$\ker f\hookrightarrow\OO$ と torsion-free 性から $\ker f=0$ を得る。
切断環の非零斉次元の積は単射の合成であり零にならない。
一般の非零元については最高次数の斉次成分を取れば domain 性が従う。
\end{proof}

\begin{proposition}[切断環と再構成]
\label{prop:reconstruction}
$\HH$ の対象と条件を満たす平方根 $s$ に対して、切断環 $B$ は二次 domain であり、
\[
 H_B(t)=\frac{1-t^2}{(1-t)^4},\qquad \dim B_n=(n+1)^2\quad(n\ge0).
\]
負次数では $\Hom(\OO,s^{-m}\OO)=0$（$m>0$）であり、
\[
 (\cC,\OO,s)\simeq(\qgr B,\pi B,(1))
\]
という algebraic triple の同値がある。
\end{proposition}
\begin{proof}
Hilbert 級数と二次性は二次の正則元による商の完全列から従う
（命題~\ref{prop:quadric-hilbert}）。domain 性は補題~\ref{lem:sections-monic} による。
もし非零射 $\OO\to s^{-m}\OO$ があればそれは単射であり、
$s^{n+m}$ を作用させてから $\Hom(\OO,-)$ を取ると
$B_{n+m}\hookrightarrow B_n$ を得る。$n\ge0$ では
$(n+m+1)^2>(n+1)^2$ となり矛盾する。
最後は Hom 有限性と Artin--Zhang ampleness に対する再構成
\cite[Theorem 2.6(2)]{MU2021} の適用である。
ここでは $B$ が noetherian なので coherent module と finitely generated module の
区別は生じず、tails は通常の $\qgr$ となる。
\end{proof}

\begin{lemma}[一次生成と偶数部分列の ampleness]
\label{lem:ample-square-from-generation}
$\mathcal C$ を Hom 有限な $k$ 線型 Abel 圏，$E\in\mathcal C$，
$s$ を完全自己同値とする．$(E,s)$ が Artin--Zhang ample であり，
右規約の orbit algebra
$B=\bigoplus_{n\ge0}\operatorname{Hom}(E,s^nE)$ が一次生成であるとする．
このとき評価射
\[
 \operatorname{Hom}(s^{-1}E,E)\otimes_k s^{-1}E\longrightarrow E
 \tag{E}\label{eq:degree-one-evaluation}
\]
は全射であり，$(E,s^2)$ も Artin--Zhang ample である．
\end{lemma}
\begin{proof}
AZ の生成条件を $sE$ に適用してから $s^{-1}$ を作用させると，
$E$ は $s^{-m_i}E$ の有限直和の商になり，すべての $m_i$ を正に取れる．
これは生成条件の記述で指数 $0$ を許していても成り立つ．
一次生成より，$m>0$ に対する任意の射 $s^{-m}E\to E$ は
\eqref{eq:degree-one-evaluation}を経由する．
実際，対応する $B_m$ の元を次数一の元の積の線型結合として書き，
積 $a*b=s^{\deg b}(a)\circ b$ を $s^{-m}$ で移すと，
最後の射は常に $s^{-1}E\to E$ となる．
従って上記の全射は評価射を経由し，評価射自身も全射である．

任意の対象を生成する負の $s$-shift の直和のうち，
奇数指数の項 $s^{-m}E$ は，評価射を $s^{-m}$ で移すことにより
$s^{-(m+1)}E$ の有限直和で置き換えられる．
このため偶数指数だけで生成条件を満たす．
必要なら生成条件を先に $s^N F$ に適用してから $s^{-N}$ を作用させ，
すべての指数を任意の下限より大きくできる．
全射に対する Hom の最終的全射性は，元の AZ 条件を偶数部分列へ制限すればよい．
\end{proof}

\begin{corollary}
$\HH$ の全対象に対し、$(\OO,s^2)=(\OO,S^{-1}[2])$ も
Artin--Zhang ample である。
\end{corollary}
\begin{proof}
条件 (Q4) により切断環が一次生成なので、直前の補題を適用できる。
\end{proof}

\begin{remark}[Gabriel--Krull 次元]
条件 (Q2) と (A1) により、$\cC/\cC_1$ の全対象は有限長である。
実際、任意の対象は有限個の $s^{-n}\OO$ の直和の商であり、その像は単純対象の
有限直和の商となる。$q\OO\ne0$ なので Gabriel--Krull のこの規約で
$\cC$ は次元二である。さらに $B_1\ne0$ から $q\OO\simeq q(s\OO)$ となり、
商圏の単純対象は $q\OO$ の一種類に限られる。
従って rank-one 条件を次数付き環の GK dimension と取り違える必要はない。
\end{remark}

\begin{remark}[Nakayama twist の点検]
\label{rem:nakayama}
例えば \cite[Section 5]{MU2021} の中心二次切断の仮定の下では、
右加群規約で canonical bimodule は $M\mapsto M_\nu(-2)$ と書ける。
従って $S^{-1}[2]$ は $M\mapsto M_{\nu^{-1}}(2)$ である。
通常の次数シフト $s=(1)$ が平方根になるには、誘導される
$M\mapsto M_{\nu^{-1}}$ が tails 上で自然に恒等であることが必要である。
この性質を環自己同型 $\nu=\mathrm{id}$ と同一視しない。
また単に一部の対象が $\nu$ で不変であるだけでは、関手の自然同型は得られない。
一般の正規切断へ公式を使う場合は双対化加群の仮定と規約を個別に照合する。
\end{remark}

% !TeX root = ../nc-quadric-classification.tex
% 2026-11-21 / v0.1
\section{二次包絡の障害と二次超曲面表示}
\label{sec:quadric-envelope-obstruction}

以下、$k$ は代数閉体とし、代数は $k$ 上の局所有限な連結次数付き代数とする。
次数一で生成される代数 $B$ に対し、$V=B_1$ と置き、通常の二次包絡を
\[
 \operatorname{QEnv}(B)
 :=T(V)/\bigl(\ker(V\otimes V\longrightarrow B_2)\bigr)
\]
と定める。括弧は両側イデアルの生成を表す。
これは $B$ の二次関係をすべて課し、三次以上の独立な関係を忘れる構成である。
したがって $B$ 自身が二次代数ならば
$\operatorname{QEnv}(B)=B$ である。

型 $(1,4,6,4,1)$ は定義~\ref{def:as14641}の意味で用いる。

\begin{proposition}[次数二の正則元による商の数値]
\label{prop:quadric-hilbert}
$A$ を上記の型の代数、$g\in A_2$ を非零の両側正則正規元とし、
$B=A/(g)$ と置く。このとき $B$ は二次代数であり、
\[
 H_B(t)=\frac{1-t^2}{(1-t)^4}
       =\frac{1+t}{(1-t)^3},\qquad
 \dim_k B_n=(n+1)^2\quad(n\geq0).
\]
特に $\dim B_1=4$、$\dim B_2=9$ であり、
$\dim\ker(B_1\otimes B_1\to B_2)=7$ である。
\end{proposition}

\begin{proof}
正規性から $(g)=gA=Ag$ であり、左から $g$ を掛ける写像によって
右加群の完全列
\[
 0\longrightarrow A(-2)\xrightarrow{\,g\cdot\,}A
 \longrightarrow B\longrightarrow0
\]
を得る。したがって Hilbert 級数は主張の通りである。
係数は $\binom{n+3}{3}-\binom{n+1}{3}=(n+1)^2$ である
（第二項は $n<2$ なら零とする）。
$A$ が二次代数で $g$ が次数二なので、$B$ も二次代数である。
最後の次元は $4^2-9=7$ から従う。
\end{proof}

\begin{corollary}[通常の二次包絡を用いる定義の障害]
\label{cor:quadric-envelope-fails}
命題~\ref{prop:quadric-hilbert} の商 $B$ について
$\operatorname{QEnv}(B)=B$ であり、これは型 $(1,4,6,4,1)$ の
AS 正則代数ではない。
\end{corollary}

\begin{proof}
$B$ は二次代数なので二次包絡と一致する。
その次数二成分の次元は $9$ であるが、指定された型の AS 正則代数では
$10$ である。
\end{proof}

とくに $Q=\mathbb P^1\times\mathbb P^1$、
$\mathcal O=\mathcal O_Q$、$s=-\otimes\mathcal O_Q(1,1)$ に対しては、
$s^2\simeq S^{-1}[2]$ であり、$(\mathcal O,s)$ は Artin--Zhang ample である。
しかしその切断代数は
\[
 B=\bigoplus_{n\geq0}H^0(Q,\mathcal O_Q(n,n))
   \simeq k[x,y,u,v]/(xy-uv)
\]
なので、通常の二次包絡の AS 正則性を要求するとこの基本例を除外してしまう。
従って quadric における ambient 条件は、二次包絡の正則性から、
\emph{次数二の正則正規元を核に持つ AS 正則な二次被覆の存在}へ置き換える。
$s$ も被覆も、$\mathcal O$-marked object の指定データには含めず、
存在条件の証人として扱う。

\section{固定された切断代数の ambient 被覆}
\label{sec:ambient-incidence}

ここでは $B$ を二次代数とし、$V=B_1$、
$K=\ker(V\otimes V\to B_2)$ と置く。
さらに $\dim V=4$ および
$H_B(t)=(1-t^2)/(1-t)^4$ を仮定するので、$\dim K=7$ である。

\begin{lemma}[正規性と Hilbert 級数による正則性の判定]
\label{lem:normal-hilbert-regular}
$A$ は局所有限な連結次数付き代数、$0\ne g\in A_d$ は正規元とし、
$B=A/(g)$ とする。もし
$H_B(t)=(1-t^d)H_A(t)$ ならば、$g$ は左右の零因子でない。
\end{lemma}

\begin{proof}
右加群の完全列
\[
 0\longrightarrow\ker(g\cdot:A(-d)\to A)
 \longrightarrow A(-d)\xrightarrow{\,g\cdot\,}A
 \longrightarrow B\longrightarrow0
\]
を考える。正規性により像は商の核 $(g)$ に等しい。
Hilbert 級数を取ると、仮定は左端の核の Hilbert 級数が零であることを意味する。
局所有限性により核は零である。右から掛ける写像についても、
$Ag=(g)$ を用いて同じ議論を適用できる。
\end{proof}

\begin{proposition}[被覆の超平面表示]
\label{prop:ambient-hyperplane}
固定された $B$ に対する、次数一の同型 $A_1\to B_1=V$ を持つ
型 $(1,4,6,4,1)$ の AS 正則二次被覆
$A\twoheadrightarrow B$ で、その核が次数二の正則正規元で生成されるものは、
\[
 R\in\operatorname{Gr}(6,K)\simeq\mathbb P(K^*)
\]
のうち、次の二条件を満たすものによって分類される。
\begin{enumerate}
 \item $A_R:=T(V)/(R)$ は型 $(1,4,6,4,1)$ の AS 正則代数である。
 \item $(K/R)\subset(A_R)_2$ を張る任意の非零元 $g_R$ は正規元である。
\end{enumerate}
ここで被覆の同型は $B$ への写像と可換な次数付き代数同型とし、
核の生成元自体のスカラー倍は区別しない。
第二条件の下では $g_R$ の両側正則性は自動的に従う。
第一条件の下で、第二条件は二つの線形写像
\[
 V\longrightarrow(A_R)_3,\quad v\longmapsto g_Rv,
 \qquad
 V\longrightarrow(A_R)_3,\quad v\longmapsto vg_R
\]
がともに階数 $4$ を持つことと同値である。
\end{proposition}

\begin{proof}
与えられた被覆から
$R=\ker(V\otimes V\to A_2)$ を取ると、$\dim A_2=10$ より
$\dim R=6$ であり、合成写像から $R\subset K$ である。
二次性により $A=T(V)/(R)$ である。
核の次数二成分は $K/R$ であり、仮定された正則正規元の張る直線に等しい。

逆に $R$ が二条件を満たすとする。
$B=T(V)/(K)$ なので、$A_R/(g_R)=B$ である。
Hilbert 級数の仮定と第一条件から
$H_B(t)=(1-t^2)H_{A_R}(t)$ であり、
補題~\ref{lem:normal-hilbert-regular} によって $g_R$ は両側正則である。
この二構成は互いに逆である。実際、$B$ への写像と可換な被覆の同型は
次数一の同一視 $A_1=V$ を保存しなければならず、$A$ は次数一で生成されるので
同じ $R$ を与える。

階数による判定について、$J=\ker(A_R\to B)$ と置くと
$\dim J_3=\dim(A_R)_3-\dim B_3=20-16=4$ である。
二つの写像の像はともに $J_3$ に含まれるので、ともに階数 $4$ ならば
$g_RV=J_3=Vg_R$ である。次数一生成性を用いて生成元を一つずつ動かすと
$g_RA_R=A_Rg_R$ が従う。
逆に正規性が成り立てば、既に示した両側正則性から両写像は単射であり、
その階数は $4$ である。
\end{proof}

\begin{remark}[有限次数の判定と未解決部分]
$A_R$ は次数一で生成されるため、
$g_RV=Vg_R$ が $(A_R)_3$ の部分空間として成立することと
$g_R$ の正規性は同値である。
一方の包含から生成元を一つずつ動かして $g_RA_R\subset A_Rg_R$ を得て、
逆の包含も同様に示せるからである。
これは正規性の有限次元線形代数による判定であるが、
有限次数での Hilbert 関数の一致だけで AS 正則性まで証明できるとは主張しない。
また、この許容集合が開集合である、滑らかである、あるいは一点であるという
主張はここには含まれない。
\end{remark}

\section{同じ可換 quadric を持つ非同一の ambient 被覆}
\label{sec:ore-ambient-family}

\begin{proposition}[正規被覆を忘れる写像の非自明なファイバー]
\label{prop:ore-ambient-family}
$a\in k^{\times}$ に対し
\[
 A_a=k\langle x,y,u,v\rangle/
 \bigl([u,v],[u,x],[u,y],[v,x],[v,y],
       yx-a\,xy-(1-a)uv\bigr),\qquad g=xy-uv
\]
と定める。このとき次が成り立つ。
\begin{enumerate}
 \item $A_a$ は noetherian な二次 AS 正則 domain であり、型は
       $(1,4,6,4,1)$ である。
 \item $g$ は次数二の正則正規元であり、
       $gx=a xg$、$gy=a^{-1}yg$、$gu=ug$、$gv=vg$ を満たす。
 \item 生成元を同一視する標準同型
       \[
        A_a/(g)\simeq k[x,y,u,v]/(xy-uv)=:B
       \]
       がある。
 \item $B_1$ の同一視を固定した被覆としては、$a\ne b$ ならば
       $A_a\twoheadrightarrow B$ と $A_b\twoheadrightarrow B$ は同型でない。
       $A_1$ は可換であり、$a\ne1$ ならば $A_a$ は非可換で $g$ は非中心である。
\end{enumerate}
\end{proposition}

\begin{proof}
$R_0=k[u,v,x]$ とし、
\[
 \sigma(u)=u,\quad\sigma(v)=v,\quad\sigma(x)=ax,
 \qquad
 \delta(u)=\delta(v)=0,\quad\delta(x)=(1-a)uv
\]
と置く。$a\ne0$ なので $\sigma$ は次数を保つ自己同型である。
$\delta$ は次数一の $\sigma$-derivation に延びる。
たとえば可換関係 $xu-ux$ に対して
\[
 \delta(xu-ux)=(1-a)uvu-u(1-a)uv=0
\]
であり、$xv-vx$、$uv-vu$ についても同様である。
従って表示から
$A_a=R_0[y;\sigma,\delta]$ である。

この Ore 拡大は、多項式環の自明な Zhang twist を基礎環とする
次数一の Ore 拡大であるため、
\cite[Proposition~2.2]{LV2026}（$n=3$、twist は恒等写像）の仮定を満たす。
同命題により $A_a$ は noetherian AS 正則 domain、global dimension $4$、
Koszul である。この適用に必要な条件は直前で確認した。
Ore 拡大の基底
\[
 \{u^i v^j x^r y^s\mid i,j,r,s\geq0\}
\]
から $H_{A_a}(t)=(1-t)^{-4}$ を得る。
Koszul 性により極小自由分解の第 $i$ 項は次数 $i$ で生成されるので、
Hilbert 級数の逆数
\[
 (1-t)^4=1-4t+6t^2-4t^3+t^4
\]
から分解の型が $(1,4,6,4,1)$ であると分かる。

基底より $g=xy-uv\ne0$ である。$A_a$ は domain なので $g$ は両側正則である。
$u,v$ の中心性と表示関係から
\[
 \begin{split}
 gx&=xyx-uvx
     =x\bigl(a\,xy+(1-a)uv\bigr)-xuv
     =a x(xy-uv)=a xg,\\
 yg&=yxy-yuv
     =\bigl(a\,xy+(1-a)uv\bigr)y-uvy
     =a(xy-uv)y=a gy
 \end{split}
\]
を得る。これは主張された正規性の関係を与える。

商で $g=0$ を課すと $xy=uv$ であり、表示関係は
$yx=a uv+(1-a)uv=uv=xy$ となる。
残る関係はすべて可換関係なので、商はちょうど上記の $B$ である。

最後に、$W\subset V\otimes V$ を表示中の最初の五つの可換関係の張る部分空間、
$r_a=y\otimes x-a x\otimes y-(1-a)u\otimes v$ と置く。
$R_a=W+kr_a$ が被覆の二次関係空間である。
もし $R_a=R_b$ なら、$y\otimes x$ の係数を比較して
$r_a-r_b=(b-a)(x\otimes y-u\otimes v)\in W$ となる。
右辺の括弧内の元は $W$ に含まれないので $a=b$ である。
命題~\ref{prop:ambient-hyperplane} の同型に関する議論により、
これは固定された $B$ 上の被覆の非同型性を意味する。
$a\ne1$ なら $[y,x]=(a-1)g\ne0$ であり、
$gx-xg=(a-1)xg\ne0$ なので、非可換性と非中心性も従う。
\end{proof}

\begin{remark}[この命題が分類について述べる範囲]
上の族は、固定された切断代数 $B$ から ambient 正規被覆が一意に復元される
という主張への反例である。同じ
$(\operatorname{coh}(\mathbb P^1\times\mathbb P^1),\mathcal O,s)$ を与えるので、
異なる $a$ を異なる $\mathcal O$-marked quadric と数えてはならない。
ただし、基底の代数 $B$ の自己同型まで許す場合の $a$ の同値関係を
この命題は分類していない。次数二の中心元だけを許す被覆という別の問題も、
上記の正規被覆の問題と区別する必要がある。
この具体例の文献上の新規性は未確認であり、新分類定理として扱わない。
\end{remark}


% !TEX root = ../nc-quadric-surface.tex
\subsection{二次被覆の有限判定}
\label{subsec:finite-ambient-test}

\begin{theorem}
\label{thm:finite-ambient-test}
$B=T(V)/(K)$ を左右 noetherian な連結二次 Koszul 環とし、
$H_B(z)=(1+z)/(1-z)^3$ とする。
超平面 $J\subset K$ に対して $A=T(V)/(J)$ と置き、
$K/J$ の非零元を $g$ とする。次は同値である。
\begin{enumerate}
\item $A$ は型 $(1,4,6,4,1)$ の環であり、$g$ は左右正則正規である。
\item 次の有限次元の条件が成り立つ。
\begin{enumerate}
\item $A_3$ で $gV=Vg$ である。
\item $A_i\to A_{i+2}$ の $g$ の左右乗法は $i=1,2$ で単射である。
\item 二次双対 $A^!=T(V^*)/(J^\perp)$ の次数零から五までの次元は
$1,4,6,4,1,0$ である。
\item $A^!_i\otimes A^!_{4-i}\to A^!_4$ の積の組合せが $i=1,2,3$ で非退化である。
\end{enumerate}
\end{enumerate}
従って $B$ の AS Gorenstein 性を先に仮定せず、適合被覆の存在は、
$\operatorname{Gr}(6,K)$ 上で有限個の階数条件を満たす点の有無で判定できる。
\end{theorem}

\begin{proof}
必要性は正規性と Koszul 双対の Hilbert 級数および Frobenius 性から従う。
逆向きに、一次生成性から $gA=Ag$ である。
$\operatorname{Tor}$ の環変換スペクトル系列を用いる。
右 $A$-加群 $B$ に対し $A(-2)\xrightarrow{g\cdot}A\to B\to0$ は極小表示である。
左乗法に核があると仮定し、その最小次数を $d$（$A$ の元の次数）とする。
$g\ne0$ と単射の条件から $d\ge3$ である。
極小分解の第二項から、$\operatorname{Tor}^A_2(B,k)$ の
最小非零内部次数は $d+2$ となる。
また $\operatorname{Tor}^A_1(B,k)\simeq k(-2)$ は
内部次数二だけに集中する左 $B$-加群である。

\[
 E^2_{p,q}=\operatorname{Tor}^B_p(k,\operatorname{Tor}^A_q(B,k))
 \ \Longrightarrow\ \operatorname{Tor}^A_{p+q}(k,k)
\]
において、$E^2_{0,2}$ の内部次数 $d+2$ は非零である。
実際、この最小次数には $B_{\ge1}$ による積が入らない。
この部分に届く可能性のある微分は $(2,1)$ および $(3,0)$ からであるが、
$B$ が Koszul なのでそれぞれ内部次数四および三に集中する。
$(0,2)$ から出る微分はなく、内部次数 $d+2\ge5$ の元は残る。
これは二次環 $A$ の $\operatorname{Tor}^A_2(k,k)$ が
内部次数二に集中することに反する。
反対環にも同じ議論を適用して、$g$ の左右正則性を得る。

従って $B$ の右 $A$-自由分解は長さ一であり、環変換スペクトル系列には
$q=0,1$ の二行のみが残る。$b_p=\dim(B^!)_p$ とすると、
$E^2_{p,0}$ は内部次数 $p$ に $b_p$ 次元、
$E^2_{p,1}$ は内部次数 $p+2$ に $b_p$ 次元ある。
全次数 $n$ の対角内部次数 $n$ に寄与するのは
$E^\infty_{n,0}=\ker(d_2:E^2_{n,0}\to E^2_{n-2,1})$ のみである。
二次環の対角 Tor の次元は二次双対の次元と一致する。
$A^!$ は一次生成なので、次数五の消滅からそれより上も消滅する。
よって仮定により
\[
 \dim E^\infty_{n,0}
 = [z^n](1+z)^4
 = b_n-b_{n-2},
\]
ただし負添字の $b_p$ は零とする。
最後の等号は $H_{B^!}(z)=(1+z)^3/(1-z)$ から従う。
従って全ての $d_2:E^2_{n,0}\to E^2_{n-2,1}$ は全射であり、
非対角の Tor は残らない。これで $A$ は Koszul となる。
同じ対角次元を持つ反対環にも適用できる。

正則正規元で割った環 $B$ が左右 noetherian なので、$A$ も左右 noetherian である。
実際、$(g)$ 進次数環は $B$ 上の skew polynomial 環であり、
斉次イデアルでは固定内部次数でフィルトレーションが有限なので
初期イデアルから有限個の生成元を持ち上げられる。
一般のイデアルには内部次数による最高次成分を用いる。
二次双対は最高次数四で積の組合せが完全なので graded Frobenius である。
Koszul 複体を双対化すると、この組合せによって左右の複体が
内部次数四と複体次数四の移動を除いて互いに同定される。
従って $\Ext_A^j(k,A)$ および反対環の Ext は、$j=4$ の $k(4)$ を除き零となる。
Koszul 性と二次双対の次元により自明加群の極小分解は
長さ四で所要の型となる。任意の有限生成次数付き加群の極小自由分解の長さも、
反対側の $k$ の分解を用いて Tor を計算すれば四以下である。
非次数付き有限生成加群には生成元からの次数フィルトレーションを入れ、
付随次数加群の自由分解を持ち上げると同じ上界を得る。
従って巡回加群に対する射影次元の上界から大域次元も四となる。
正則二次切断から $H_A(z)=(1-z)^{-4}$ が従う。
\end{proof}

% !TeX root = ../nc-quadric-classification.tex
% 2026-11-21 / v0.1
\section{可換二次曲面と平方根の非一意性}
\label{sec:commutative-roots}

この節では $k$ を標数 $0$ の代数閉体，
$Q=\mathbb P^1_k\times\mathbb P^1_k$ とする．
$\mathcal O_Q(a,b)$ を通常の線束とし，
$S_Q=(-)\otimes\mathcal O_Q(-2,-2)[2]$ と書く．
ここで分類する平方根は，$k$ 線型完全自己同値 $s$ であって
$s^2\cong S_Q^{-1}[2]$ を満たすものの同型類である．
この同型そのものは指定せず，Serre 双対性の trace と両立する
追加の巡回的 coherence 条件も課さない．
根の同値は，$F(\mathcal O_Q)\cong\mathcal O_Q$ を満たす
完全自己同値 $F$ と $Fs\cong s'F$ の存在を意味する．
この節は族や慣性の分類を主張しない．

\begin{lemma}[可換標識対象]
\label{lem:commutative-marking-line}
$Q$ 上の rank-one torsion-free exceptional coherent sheaf $E$ は線束である．
したがって $(\operatorname{coh}Q,E)$ は
$(\operatorname{coh}Q,\mathcal O_Q)$ と pointed exact 同値である．
\end{lemma}
\begin{proof}
$Q$ は滑らかな曲面なので $L=E^{\vee\vee}$ は線束であり，
$T=L/E$ は有限長である．$n=\operatorname{length}T$ と置く．
$E$ は完全複体として扱え，Riemann--Roch により
\[
 \chi(E,E)
 =\int_Q\operatorname{ch}(E^{\vee}_{\mathrm{der}})
             \operatorname{ch}(E)\operatorname{td}(Q)
 =\chi(\mathcal O_Q)-2n=1-2n.
\]
実際，$D=c_1(L)$ とすれば
$\operatorname{ch}(E)=e^D(1-n[\mathrm{pt}])$ であり，
導来双対の Chern character は $e^{-D}(1-n[\mathrm{pt}])$ である．
exceptionality から $\chi(E,E)=1$ なので $n=0$，従って $E=L$ である．
$(-)\otimes L^{-1}$ が求める pointed exact 同値を与える．
\end{proof}

\begin{theorem}[四つの平方根]
\label{thm:commutative-four-roots}
$\operatorname{coh}Q$ の $k$ 線型完全自己同値 $s$ が
\[
 s^2\cong(-)\otimes\mathcal O_Q(2,2),
 \qquad \dim_k\operatorname{Hom}(\mathcal O_Q,s\mathcal O_Q)=4
 \tag{$*$}\label{eq:commutative-root-conditions}
\]
を満たすとする．このような $s$ は pointed exact 同値を除いてちょうど四種類であり，
\[
 s_\varphi(F)=\varphi^*F\otimes\mathcal O_Q(1,1)
\]
において $\varphi$ を次のいずれかに取ったものである．
\begin{enumerate}
 \item 恒等写像．
 \item 一方の因子の非自明な対合 $\iota\times\mathrm{id}$．
 \item 両因子の非自明な対合 $\iota\times\iota$．
 \item 二因子の交換 $\epsilon$．
\end{enumerate}
ここで $\iota([u_0:u_1])=[u_0:-u_1]$ と取れる．
四者はいずれも $(\mathcal O_Q,s_\varphi)$ が Artin--Zhang ample であり，
対応する右規約の orbit algebra は $(14641)$ 型 quadratic AS-regular domain の
正則正規な二次斉次元による商である．
\end{theorem}

\begin{proof}
まず完全自己同値の形を決める．
Bondal--Orlov の定理 \cite[Theorem~3.1]{BondalOrlov2001} は，
反標準束が ample な滑らかな射影多様体の導来自己同値を
自己同型，線束による tensor，shift で記述する．
Abel 圏の完全自己同値から得られる導来自己同値は標準 heart を保つため shift は零である．
従って $s=\varphi^*(-)\otimes L$ と書ける．
平方根の条件から
\[
 \varphi^2=\mathrm{id},\qquad
 L\otimes\varphi^*L\cong\mathcal O_Q(2,2).
 \label{eq:commutative-root-norm}
\]
$\varphi$ が二つの ruling をそれぞれ保つ場合，Picard 群への作用は恒等なので
$L=\mathcal O_Q(1,1)$ である．
$\varphi$ が ruling を交換する場合は
$L=\mathcal O_Q(a,2-a)$ である．
このとき $h^0(L)=4$ はまず $0\le a\le2$ を強制し，さらに
\[
 (a+1)(3-a)=4
\]
から $a=1$ を強制する．従って，いずれの場合も $L=\mathcal O_Q(1,1)$ である．

$F(\mathcal O_Q)\cong\mathcal O_Q$ を保つ完全自己同値 $F$ は，
同じ定理により $\psi^*$ の形に限られる．
すべての $\psi\in\operatorname{Aut}(Q)$ は
$\mathcal O_Q(1,1)$ の同型類を保つので，求める同値類は
$\operatorname{Aut}(Q)$ における対合（恒等元を含む）の共役類である．
\[
 \operatorname{Aut}(Q)
 =\bigl(\operatorname{PGL}_2(k)\times\operatorname{PGL}_2(k)\bigr)
       \rtimes\langle\epsilon\rangle.
\]
$\operatorname{PGL}_2(k)$ の非自明な対合は，標数が $2$ でなく $k$ が代数閉体なので
すべて $\iota$ に共役である．
恒等成分に属する対合は二因子の対合の組であり，
因子交換も許すと，非自明な因子の数 $0,1,2$ の三種類になる．
他方，恒等成分の外にある対合は $(g,g^{-1})\epsilon$ と書け，
$(g,\mathrm{id})\epsilon(g,\mathrm{id})^{-1}=(g,g^{-1})\epsilon$
よりすべて $\epsilon$ に共役である．
前の三種類は ruling の作用と各因子の作用で区別でき，
最後の種類は ruling を交換するのでこれらと共役ではない．

残りの条件を確認する．各 $\varphi$ は $\mathcal O_Q(1,1)$ を保ち，
\[
 s_\varphi^{-n}\mathcal O_Q\cong\mathcal O_Q(-n,-n)
 \qquad(n\in\mathbb Z)
\]
である．$\mathcal O_Q(1,1)$ の ampleness による大域生成と Serre 消滅から，
この列は Artin--Zhang の生成条件と全射に対する Hom の最終的全射性を満たす．
以下の命題が orbit algebra の AS-regular cover を与える．
\end{proof}

\begin{proposition}[四平方根の明示的な cover]
\label{prop:commutative-twist-covers}
$V=H^0(Q,\mathcal O_Q(1,1))$，$P=\operatorname{Sym}(V)$ とし，
Segre 方程式を
$q=x_{00}x_{11}-x_{01}x_{10}\in P_2$ と書く．
定理\ref{thm:commutative-four-roots}の各 $\varphi$ について，
その $V$ 上の持上げ $\sigma$ を $\sigma^2=1$ となるように選ぶ．
$P$ と同じ次数付きベクトル空間に積
\[
 f\star g=\sigma^{\deg g}(f)g
 \qquad(f,g\text{ homogeneous})
\]
を入れた代数を $A_\sigma$ とする．
すると $A_\sigma$ は左右 noetherian な quadratic Koszul AS-regular domain であり，
解消型は $(14641)$ である．
$P_2=(A_\sigma)_2$ というベクトル空間の同一視で得られる元 $q$ は正則正規で，
右規約
$a*b=s_\varphi^{\deg b}(a)\circ b$
による orbit algebra は $A_\sigma/(q)$ に同型である．
\end{proposition}

\begin{proof}
上記の積は，$\sigma$ が次数付き自己同型であることから結合的である．
$\sigma$ の固有基底 $z_1,\ldots,z_4$ を取り，固有値を $e_i\in\{1,-1\}$ とすると，
\[
 z_i\star z_j=(e_i/e_j)z_j\star z_i.
\]
従って $A_\sigma$ は四変数の skew polynomial ring である．
順序付き単項式は PBW 基底をなし，自己同型による反復 Ore 拡大として
左右 noetherian かつ domain である．
通常の Koszul 複体の交換係数を上記の符号へ変更すると，
正規正則列 $z_1,\ldots,z_4$ に沿う標準的な skew Koszul 複体が得られ，
\[
 0\longrightarrow A_\sigma(-4)
 \longrightarrow A_\sigma(-3)^4
 \longrightarrow A_\sigma(-2)^6
 \longrightarrow A_\sigma(-1)^4
 \longrightarrow A_\sigma\longrightarrow k\longrightarrow0
\]
は完全である．双対複体のコホモロジーは最高次数の $k(4)$ だけであり，
反復 Ore 拡大の大域次元は $4$，Hilbert 級数は $(1-t)^{-4}$ である．
これが主張した AS-regularity と解消型を与える．

$\varphi$ は $Q$ を保つので $\sigma(q)=c q$ であり，$c\in\{1,-1\}$ である．
$x\in V$ に対して
\[
 q\star x=\sigma(q)x=c qx,
 \qquad x\star q=\sigma^2(x)q=xq
\]
だから $q\star x=c(x\star q)$ となる．
従って $q$ は正規であり，domain 性から正則でもある．
二側イデアル $(q)$ は通常の $qP$ と同じ次数付きベクトル空間なので，
$A_\sigma/(q)$ は Segre 切断環 $P/(q)$ の同じ積による twist である．
$\varphi^*\mathcal O_Q(1,1)\cong\mathcal O_Q(1,1)$ の同型を
$\sigma$ と整合的に選ぶと，orbit algebra の積はまさに
$a*b=\sigma^{\deg b}(a)b$ になる．
\end{proof}

\begin{remark}
上の四代表で，$\sigma$ はそれぞれ恒等作用，
$x_{10},x_{11}$ の符号変更，$x_{01},x_{10}$ の符号変更，
$x_{01}$ と $x_{10}$ の交換として取れる．
$q$ の固有値 $c$ は順に $1,-1,1,1$ である．
従って一因子対合のこの cover では $q$ は各一次元と反可換になり，中心元ではない．
これは正則正規な二次斉次元を許す定義を点検する具体例になる．
ここで $q$ は $P_2=(A_\sigma)_2$ の同一視で指定した元であり，
$\star$ 積で表した式の係数は通常の可換積の表示と異なり得る．
\end{remark}

\begin{example}[平方根と ampleness だけでは一次元数は決まらない]
\label{ex:commutative-unbalanced-root}
二因子の交換 $\epsilon$ を用いて
\[
 t(F)=\epsilon^*F\otimes\mathcal O_Q(2,0)
\]
と置くと $t^2\cong(-)\otimes\mathcal O_Q(2,2)$ であり，
$(\mathcal O_Q,t)$ は Artin--Zhang ample である．
実際，$t^{-2m}\mathcal O_Q\cong\mathcal O_Q(-2m,-2m)$，
$t^{-(2m+1)}\mathcal O_Q\cong\mathcal O_Q(-2m,-2m-2)$ である．
偶数部分列による生成と，両部分列に対する Serre 消滅から ampleness が従う．
しかし
\[
 \dim_k\operatorname{Hom}(\mathcal O_Q,t\mathcal O_Q)
 =h^0(\mathcal O_Q(2,0))=3.
\]
従って，その orbit algebra は四つの一次生成元を持つ AS-regular 環の
二次斉次元による商ではない．
平方根の存在と ampleness だけから $(14641)$ 型 cover の数値条件を結論することはできない．
\end{example}

\begin{remark}[対象と根の区別]
補題\ref{lem:commutative-marking-line}により，heart を
$\operatorname{coh}(\mathbb P^1\times\mathbb P^1)$ に固定した可換対象には，
pointed exact 同値を除いて一つの rank-one torsion-free exceptional 標識がある．
一方，定理\ref{thm:commutative-four-roots}はその同じ対象上で許容平方根が四つの同値類を持つことを示す．
この四種類を四つの $\mathcal O$-marked 曲面として数えてはならない．
この節の結果は既知の自己同値記述と明示的計算による基礎例であり，
独立した新規性は主張しない．
\end{remark}

% !TeX root = ../nc-quadric-classification.tex
% 2026-11-21 / v0.2。引用番号の版は本文で指定する。
\section{楕円型 Sklyanin 二次曲面の例}
\label{sec:sklyanin-examples}

この節では $k=\mathbb C$ とする。
$T$ を \cite[\S 7, (7.1), (7.2)]{BSW2010} の通常の四次元 Sklyanin 環とし、
\cite[\S 10.1]{SVdB2013} の楕円型データ $(E,L,\tau)$ に対応すると仮定する。
ここで $E$ は滑らかな楕円曲線、$\deg L=4$、$\tau$ は無限位数の平行移動である。
原点は同文献の規約に合わせ、$L$ による埋込みで四点が同一平面上にあることと
その和が零であることが同値になるように選ぶ。
したがって平行移動 $\tau$ は $E$ の点ともみなす。
この仮定は modified Sklyanin 環や退化した点スキームを含めない。

次数二の中央元からなる二次元空間 $\mathbb C\Omega_1+\mathbb C\Omega_2\subset T_2$ の
非零元による商は、$z\in E$ を用いて
\[
 R_z=T/(g_z),\qquad Q_z=\operatorname{Proj}R_z,
 \qquad Q_z=Q_{-z-2\tau}
\]
と表示される \cite[\S\S 10.2--10.4]{SVdB2013}。
同文献の \S 2 は右加群で定義される一方、\S 10 の幾何的記述には
$T/TW$ という左加群の表示が使われる。
この幾何的表示を本稿の右加群規約に合わせるため、
\[
 A=T^{\mathrm{op}},\qquad B_z=R_z^{\mathrm{op}}
       =A/(g_z),\qquad
 \cC_z=\qgr B_z,\qquad \OO_z=\pi B_z
\]
と置く。以下の反自己同型により、左右いずれの規約による結果も
構造対象と次数シフトを保って移送できる。
反対環への移行は以下で用いる左右 AS 正則性、Koszul 性、domain 性、
中央正則元、Calabi--Yau 性を保つ。

\begin{lemma}[左右規約の比較]
\label{lem:sklyanin-opposite}
通常の Sklyanin 座標において、
$\theta(x_0)=-x_0$、$\theta(x_i)=x_i$（$1\le i\le3$）は
$T$ の次数付き反自己同型を定め、中央二次 pencil を各点で固定する。
従って $R_z\simeq R_z^{\mathrm{op}}$ であり、左右の coherent module 圏の間には
構造対象と次数シフトを保つ $\mathbb C$-線形 exact 同値がある。
\end{lemma}
\begin{proof}
\cite[\S 7]{BSW2010} の関係式を $r_i,s_i$ と書くと、
積の順序を逆転して上記の符号を代入する写像は $r_i$ を $r_i$、$s_i$ を $-s_i$ に送る。
その二乗は恒等なので反自己同型である。
\cite[\S 1.1, pp.~38--39]{LS1993} の中央二次元空間は
$x_0^2,x_1^2,x_2^2,x_3^2$ の線型結合で張られ、各元は $\theta$ で固定される。
同文献の仮定 $\tau\notin E[4]$ は無限位数の仮定から従う。
よって $\theta$ は $g_z$ を固定して商へ降りる。
左加群 $M$ に右作用 $m\cdot r=\theta(r)m$ を入れることが所要の同値を与える。
\end{proof}

\begin{proposition}[滑らかな楕円型の対象]
\label{prop:sklyanin-h2}
上の仮定の下で、$z+\tau\notin E[2]$ ならば
$(\cC_z,\OO_z)\in\HH$ である。
存在条件の証人として $s=(1)$ と次数二の中央正則元による表示 $B_z=A/(g_z)$ を取れる。
特に切断環は高次数で一致するだけでなく、全非負次数で
\[
 B(\cC_z,\OO_z,(1))\cong B_z
\]
である。ここで $E[2]$ は $E$ の二捩れ部分群を表す。
\end{proposition}

\begin{proof}
以下 $B=B_z$、$\cC=\cC_z$、$\OO=\OO_z$ と略記する。
\cite[\S\S 10.1--10.2]{SVdB2013} により、$A$ は左右 noetherian な
四次元二次 AS 正則 Koszul domain で、$H_A(t)=(1-t)^{-4}$ である。
従って極小分解は型 $(1,4,6,4,1)$ となる。
$g_z$ は中央正則であり、$B$ も noetherian domain である。
商は AS Gorenstein 次元 $3$、Gorenstein parameter $2$ を持ち、
\[
 H_B(t)=\frac{1-t^2}{(1-t)^4}=\frac{1+t}{(1-t)^3}.
\]
この商の AS Gorenstein 性と次数については
\cite[\S 5, Definition 5.1 直後]{MU2021} も参照できる。
\cite[Theorem 10.2]{SVdB2013} により、仮定 $z+\tau\notin E[2]$ の下で
$\cC$ は滑らかであり、$\gldim\cC=2$ である。
また $B$ は深さ $3$ の Cohen--Macaulay 環であり、両側の条件 $\chi$ を満たす
\cite[\S\S 2.5--2.8]{SVdB2013}。
$\chi$ による coherent cohomology の有限性、構造対象の shift による生成、
および $\gldim\cC=2$ から $D^b(\cC)$ は Hom 有限である。

まず Serre 関手の自己同型による捩れを確認する。
\cite[arXiv v3, Proposition 7.1]{BSW2010} は $T$ が非捩れの
Calabi--Yau 環であることを示す。従って $A$ の Nakayama 自己同型は恒等である。
\cite[arXiv v2, Lemma 1.5]{RRZ2014} を中央正則元 $g_z$ に適用できる。
同補題の $\mu_A$-normal 条件は $\mu_A(g_z)=g_z$ により満たされ、
正規化自己同型も恒等なので、$\mu_B=\mathrm{id}$ となる。
従って $B$ の canonical bimodule は両側加群として $B(-2)$ である。
\cite[Example 3.6(2) および \S 5 の Proposition 5.2 直後]{MU2021} により
\[
 \mathbb S_{\cC}\simeq(-2)[2],\qquad
 s^2=(2)\simeq\mathbb S_{\cC}^{-1}[2].
\]
ここでは hypersurface 環 $B$ 自体が有限 global dimension を持つとは主張していない。

$\mathfrak m=B_{\ge1}$ とする。
\cite[(2.2)]{SVdB2013} と $H^0_{\mathfrak m}(B)=H^1_{\mathfrak m}(B)=0$ により
\[
 B\xrightarrow{\sim}
 \bigoplus_{n\in\mathbb Z}\Hom_{\cC}(\OO,\OO(n))
\]
である。この写像は本稿の切断環の積と整合する。
従って全次数での飽和性が成り立ち、$\End(\OO)=\mathbb C$ を得る。
同文献 (2.3) と深さの条件から
\[
 \Ext^1_{\cC}(\OO,\OO)=H^2_{\mathfrak m}(B)_0=0.
\]
さらに Serre 双対性と飽和性により
\[
 \Ext^2_{\cC}(\OO,\OO)
 \cong\Hom_{\cC}(\OO,\OO(-2))^*
 \cong B_{-2}^*=0.
\]
次数 $i>2$ の Ext も消えるので $\OO$ は exceptional である。
有限生成加群の十分高い切断に生成元を取ると、任意の対象は有限個の
$\OO(-n)$、$n\ge0$、の直和の商になる。
また $\chi$ と (2.3) は任意の $K\in\cC$ に対し
$\Ext^1(\OO(-n),K)=0$ を $n\gg0$ で与える。
全射の核にこれを適用すれば Artin--Zhang ample の第二条件も従う。

最後に定義~\ref{def:gk-rank} の Gabriel--Krull rank 条件を確認する。
これは $B$ の Gelfand--Kirillov dimension や商体上の rank の計算だけでは済まない。
有限生成次数付き $B$-加群 $M$ は有限生成 $A$-加群でもあり、
$A$ 上の有限次数付き自由分解を持つので、Hilbert 関数は最終的に多項式 $P_M(n)$ となる。
$M$ は有限個の $B$ の shift の商なので $\deg P_M\le2$ である。
$P_M$ は $\pi M$ のみで決まり、短完全列に対して加法的である。
零多項式となることと $\pi M=0$ は同値である。

$\mathcal D_0\subset\cC$ を Hilbert 多項式が定数である対象、
$\mathcal D_1\subset\cC$ を次数が高々 $1$ である対象の部分圏とする
（いずれも零対象を含める）。両者は Serre 部分圏である。
非零な定数 Hilbert 多項式の値は正整数なので、
$\mathcal D_0$ の対象の任意の真の部分商列の長さはその定数で抑えられる。
従って $\mathcal D_0\subset\cC_0$ である。
同様に、次数 $1$ の Hilbert 多項式の一次係数は正整数であり、
$\cC/\mathcal D_0$ において $\mathcal D_1$ の対象の長さを抑える。
exact な商関手 $\cC/\mathcal D_0\to\cC/\cC_0$ は有限長性を保つから
\[
 \mathcal D_1\subset\cC_1.
 \tag{\dagger}\label{eq:sklyanin-polynomial-filtration}
\]
ここではこれらの包含を等号とする必要はない。

$g_z$ と線形独立な次数二の中央元 $h$ を同じ pencil から取る。
$B$ におけるその像も $h$ と書く。
$h$ は非零で正則であり、$B/(h)$ は楕円曲線 $E$ の twisted homogeneous
coordinate ring の反対環である \cite[\S 10.2]{SVdB2013}。
従って $\qgr B/(h)\simeq\coh E$ は $\cC$ に exact かつ fully faithful に埋め込まれ、
$\pi(B/(h))$ とその次数 shift は $E$ 上の可逆層に対応する。
可逆層は、点 $p\in E$ における零点の位数を増やす部分可逆層の下降列を持つので、
この埋込み後も有限長ではない。
そこで $\OO$ の下降列
\[
 \OO\supset\pi(hB)\supset\pi(h^2B)\supset\cdots
\]
を考える。隣接する商は $\pi(B/(h))(-2j)$ なので $\cC_0$ に属さない。
この列は $\cC/\cC_0$ でも真の下降列であり、その中の $\OO$ の像は有限長でない。
従って $\OO\notin\cC_1$ である。

$0\ne N\subset\OO$ を任意の部分対象とする。
Serre 商における部分対象の持ち上げにより $N=\pi I$ と書ける。
ここで $I\subset B$ は非零な次数付き右イデアルである。
次数 $d$ の非零元 $f\in I$ を取ると、domain 性により
$B(-d)\xrightarrow{f\cdot}I$ は単射なので、$\OO(-d)\subset N$ である。
shift は $\cC_1$ を保つから $N\notin\cC_1$ となる。
一方、$B/I$ は $B/fB$ の商であり、
\[
 H_{B/fB}(t)=(1-t^d)H_B(t)
\]
の Hilbert 多項式は次数高々 $1$ である。
従って \eqref{eq:sklyanin-polynomial-filtration} により
$\OO/N\in\cC_1$ となる。
$q:\cC\to\cC/\cC_1$ とすると、$q\OO$ は非零であり、
その部分対象は $\OO$ の部分対象から持ち上がるので、以上から $q\OO$ は単純である。
また $\cC_1$ に属する非零部分対象を $\OO$ は持たないので
$\Hom(T,\OO)=0$（$T\in\cC_1$）である。
これで定義~\ref{def:gk-rank} と $\HH$ の全条件が従う。
\end{proof}

\begin{remark}[standard 性と既存分類への入口]
\label{rem:sklyanin-standard}
上の滑らかな Sklyanin 商は \cite[Theorem 5.5]{MUW2026} により standard である。
反対環の規約は補題~\ref{lem:sklyanin-opposite} で処理できる。
従って命題~\ref{prop:standard-pointed} により、
三次元 AS 正則 $\mathbb Z$-algebra への比較で標識と平方根を追跡できる。
\end{remark}

\begin{proposition}[可換二次曲面と同値でない対象]
\label{prop:sklyanin-genuinely-noncommutative}
上と同じ $T$ に対して $z=\tau$ と取ると、
$(\cC_\tau,\OO_\tau)\in\HH$ であり、
$\cC_\tau$ は $\coh(\PP^1\times\PP^1)$ と $\mathbb C$-線形 exact 同値でない。
\end{proposition}

\begin{proof}
$\tau$ が無限位数なので $2\tau\notin E[2]$ であり、
命題~\ref{prop:sklyanin-h2} が適用できる。
注意~\ref{rem:sklyanin-standard} の standard 性により、以下では
一般の中心二次商についての誤った standard 性の主張を用いない。
\cite[\S 10.5, p.~854]{SVdB2013} では、
$\omega\in E[2]$、$i\in\mathbb N$ に対して、$Q_{\omega+i\tau}$ 上の
単純対象 $F_{\omega+i\tau}$ を構成し、
同文献 Proposition 9.3 の Euler 形式を用いて
$\chi(F_{\omega+i\tau},F_{\omega+i\tau})=2$ を計算している。
同文献の (10.1) でこの単純対象は coherent な line object の商として与えられるので、
$\coh Q_{\omega+i\tau}$ に属する。
ここで $\omega=0$、$i=1$ とすれば $\cC_\tau$ に単純対象 $F$ が存在して
\[
 \chi(F,F):=\sum_{j=0}^2(-1)^j\dim_{\mathbb C}\Ext^j(F,F)=2
\]
となる。
一方、$\coh(\PP^1\times\PP^1)$ の単純対象は閉点の skyscraper 層 $k(p)$ であり、
滑らかな二次元局所環の Koszul 分解からその自己 Ext の次元は $(1,2,1)$ である。
従って $\chi(k(p),k(p))=0$ となる。
exact 同値は単純性と全 Ext を保つので、両圏は同値でない。
\end{proof}

この例は可換二次曲面の同じ Abel 圏に異なる平方根や twisted coordinate ring を
選ぶことでは説明されない。
一方、この命題は全ての滑らかな Sklyanin pencil の成員について非同値性を主張するものでも、
楕円型の全対象を分類するものでもない。
引用した Sklyanin 環・pencil・単純対象の構成は既知の結果であり、
本節はそれらが本稿の標識付き定義、とくに全次数の切断環と
Gabriel--Krull rank 条件を満たすことを確認したものである。

% !TEX root = ../nc-quadric-surface.tex
\section{反 Serre 環による同値判定}
\label{sec:antiserre-ring}

\begin{theorem}[反 Serre 環と pointed exact 同値]
\label{thm:antiserre-invariant}
$(\cC,\OO)\in\HH$ に対し、$t=S^{-1}[2]$ の heart 上の代表を用い、
\[
 R(\cC,\OO)=\bigoplus_{n\ge0}\Hom(\OO,t^n\OO)
\]
に定義~\ref{def:ample} の右規約で積を入れる。
その次数付き同型類は代表、Serre 関手および平方根の選択に依存しない。
任意の適合根の切断環 $B$ に対し $R\simeq B^{(2)}$ であり、
$(\cC,\OO,t)\simeq(\qgr R,\pi R,(1))$ となる。
二つの $\HH$ 対象が pointed exact 同値であるための必要十分条件は
反 Serre 環が次数付き $k$-代数として同型となることである。
\end{theorem}

\begin{proof}
$t$ は導来圏上で $s^2$ と自然同型なので heart を保ち、制限は exact 自己同値である。
heart から導来圏への埋込みの充満忠実性により自然同型を制限できる。
一般に $F:\cC\to\cC'$ が $k$-線形同値で、
$u:F\OO\xrightarrow{\sim}\OO'$ および $\eta:Ft\xrightarrow{\sim}t'F$
があるとする。反復同型 $\eta_n:Ft^n\to t'^nF$ を
$\eta_0=1$、$\eta_{n+1,X}=t'(\eta_{n,X})\circ\eta_{t^nX}$ で定める。
この定義から
$\eta_{m+n,X}=t'^n(\eta_{m,X})\circ\eta_{n,t^mX}$ となる。
次数 $m$ で
\[
 a\longmapsto t'^m(u)\circ\eta_{m,\OO}\circ F(a)\circ u^{-1}
\]
と置くと、$\eta_n$ の自然性を $a:\OO\to t^m\OO$ に適用して
積 $t^n(a)\circ b$ と整合することが分かる。
$F$ の充満忠実性により各次数で線形同型である。

この構成を恒等関手と自然同型 $s^2\simeq t$ に適用すると
$R\simeq B^{(2)}$ を得る。Serre 関手の一意性から代表の変更にも同じ構成を使える。
pointed exact 同値 $F$ の導来同値は Serre 関手を自然同型まで移すので、
heart 上で $Ft\simeq t'F$ となり、必要性が従う。

$B^{(2)}$ は左右 noetherian である。実際、右イデアル $I$ を $B$ へ拡張すると
$IB\cap B^{(2)}=I$ である。積の偶数成分への射影は $I$ の元との左乗法と可換であり、
$IB$ の任意の元の偶数成分は $I$ に入る。左イデアルにも同様に適用できる。
補題~\ref{lem:ample-square-from-generation} により $(\OO,t)$ は ample なので、
\cite[Theorem 2.6(2)]{MU2021} が主張した triple の再構成を与える。
逆に次数付き環同型は右加群の移送を通じて $\pi R$ を保つ $\qgr$ の exact 同値を与える。
\end{proof}

\begin{proposition}[反 Serre 環の二次関係]
\label{prop:antiserre-quadratic}
$B=T(V)/(K)$ を適合根に対応する切断環とし、
$U=B_2=(V\otimes V)/K$ と書く。このとき
\[
 R\simeq T(U)/(J),\qquad
 J=\ker(U\otimes U\to B_4)
   =\operatorname{im}(V\otimes K\otimes V\to U\otimes U).
\]
$\dim U=9$、$\dim J=56$ であり、
$\dim R_n=(2n+1)^2$、$H_R(z)=(1+6z+z^2)/(1-z)^3$ となる。
\end{proposition}
\begin{proof}
$T(V)^{(2)}=T(V\otimes V)$ と同一視する。
次数 $2n$ での関係イデアルは $2n$ 文字の中に $K$ を一箇所挿入した部分空間の和である。
$K$ の前の文字数が偶数なら $K$ は二文字一組の内部にあり、一次関係 $K$ が生成する。
奇数なら挿入箇所の前後に各一文字を加えて二組を取り出せるため、
$V\otimes K\otimes V$ が生成する。最初の関係で割った後は $T(U)$ となり、
残る生成関係は次数二の上記の像である。
命題~\ref{prop:reconstruction} より $\dim B_2=9$、$\dim B_4=25$ だから
$\dim J=81-25=56$。同じ Hilbert 関数の偶数項を足せば級数を得る。
\end{proof}

従って二つの適合表示から計算した関係空間 $J,J'$ について、
$(h\otimes h)J=J'$ を満たす線形同型 $h:U\to U'$ の存在が
pointed exact 同値の必要十分条件である。
この判定では $B$ の基底や被覆を保持する必要はない。
体上で実現する軌道は定理~\ref{thm:complete-h2-classification} で全て記述する。
これを族のモジュライへ拡張するには、関係空間の実現部分を底変更と整合する
代数的 locus として構成し、相対再構成を別途証明する必要がある。

\begin{corollary}[導来圏と構造対象からの heart の認識]
\label{cor:pointed-derived-exact}
二つの $\HH$ 対象の有界導来圏の $k$-線形三角同値が
$\OO$ を $\OO'$ に送るなら、その同値は標準の heart に関して $t$-exact である。
\end{corollary}
\begin{proof}
$t=S^{-1}[2]$ とする。任意の $M\in\cC$ に対し、
$t^{-n}\OO$ から $M$ への正次数 Ext は $n\gg0$ で消える。
実際、豊富性により $M$ への全射の始域を有限個の $t^j\OO$ の直和に取れる。
従って $\Hom(M,t^{-n-1}\OO)=0$ は負切断の消滅から従い、
Serre 双対性で $\Ext^2(t^{-n}\OO,M)=0$ を得る。
$\Ext^1$ の消滅は二次の正則正規商 $B$ の AS Gorenstein 性を用いる
\cite[Lemma 2.4]{MU2022}。
深さ三により $\Ext^1_{\qgr B}(\pi B(i),\pi B(j))=0$ である
（\cite[Lemma 5.9(2)]{MU2021} と同じ打切りの議論）。
上の全射の核を $K$ とすると、$\Ext^1(t^{-n}\OO,M)$ は
$\Ext^2(t^{-n}\OO,K)$ に埋まり、後者は既に十分大きい $n$ で消える。

従って $L\in D^b(\cC)$ に対し、十分大きい $n$ では
$\Hom(t^{-n}\OO,L[i])\simeq\Hom(t^{-n}\OO,H^i(L))$ である。
これは標準打切りに伴う Ext のスペクトル系列から従う。
豊富性の生成条件により、非零な $H^i(L)$ への非零射は
任意に大きい $n$ のある $t^{-n}\OO$ から存在する。
従って $L$ が heart に属する条件は、全ての $i\ne0$ について
$\Hom(t^{-n}\OO,L[i])$ が $n\gg0$ で消えることである。
三角同値は Serre 関手とシフトを保つので、$\OO$ を保てば
この条件を双方向に保つ。
\end{proof}

\subsection{可換 quadric と楕円型 Sklyanin pencil の判別}

可換 quadric の四つの根が与える反 Serre 環はいずれも
$\bigoplus_{n\ge0} H^0(Q,\OO_Q(2n,2n))$ である。
これは $t=-\otimes\OO_Q(2,2)$ から直接従い、
命題~\ref{prop:ore-ambient-family} の被覆の選択にも依存しない。

\begin{proposition}[滑らかな Sklyanin 例の内在的な楕円商]
\label{prop:sklyanin-antiserre-center}
命題~\ref{prop:sklyanin-h2} の範囲で $R=B_z^{(2)}$ とする。
pencil の残りの中心二次元 $h$ の像を $R_1$ で考えると
$Z(R)=k[h]$ であり、$R/(h)=(B_z/(h))^{(2)}$ となる。
特にこの範囲の全対象は可換 quadric と pointed exact 同値でなく、
商に現れる楕円曲線は pointed exact 同値類から回復できる。
\end{proposition}
\begin{proof}
$h$ は中心正則である。偶数成分を取ると主張した商が得られる。
既知の楕円商表示（命題~\ref{prop:sklyanin-h2} の証明）により
この商は楕円曲線の twisted homogeneous coordinate ring の反対環である。
偶数次だけを取ると、平行移動 $\sigma$ は $\sigma^2$ に、
次数一の線束 $L$ は $L\otimes\sigma^*L$ に替わる。
$\sigma^2$ は無限位数で、この線束の次数は八である。

このような楕円環の中心は $k$ である。実際、線束を関数体 $k(E)$ 上で
自明化すると、環は skew polynomial ring $k(E)[u;\sigma^2]$ に入る
（反対環を取っても中心は同じである）。
中心に非零斉次元 $f u^m$（$m>0$）があれば、
次数一の任意の二つの非零切断 $a u,b u$ との可換性を割り算して
$(\sigma^2)^m(a/b)=a/b$ を得る。
次数八の線束は very ample なので、その比で関数体が生成される。
これは $\sigma^2$ の無限位数性に反する。
従って $R/(h)$ の中心は $k$ であり、$R$ の中心にある任意の正次数斉次元は
$h$ の倍数となる。$h$ が中心正則なので割った元も中心にあり、次数による帰納法で
$Z(R)=k[h]$ が従う。

従って $R$ は非可換であり、定理~\ref{thm:antiserre-invariant} と
可換 quadric の $R$ の表示から、pointed exact 非同値である。
さらに $Z(R)_1$ の生成するイデアルで割る操作は内在的であり、
商の非零斉次元を逆にした次数付き商環の零次体は $k(E)$ となる。
実際、商は skew Laurent 環（またはその反対環）に埋め込まれ、
次数一切断の比が $k(E)$ を生成する。
斉次 Ore 局所化は noetherian domain 性から存在する。
斉次元同士の共通倍数は非零な斉次成分を取って選べる。
従って $k(E)$ は同型不変であり、その滑らかな射影模型 $E$ も同型まで定まる。
さらに商環の次数一の可逆元による零次体上の共役作用は、
零次体が可換なので元の選択に依存しない。この作用は上の表示で
$\sigma^2$（反対環の規約ではその逆）に対応する。
\end{proof}

% !TEX root = ../nc-quadric-surface.tex
\section{中心 standard 表示との比較}
\label{sec:standard-comparison}

\begin{proposition}[標識とシフトの移送]
\label{prop:standard-pointed}
$B=A/(g)$ を $\HH$ 対象の適合表示とし、$g$ は中心元、
$B$ は \cite{MN2025} の意味で standard とする。
$X$ を二つの ruling に対応する極大 Cohen--Macaulay 加群の一つとする。
すると \cite[Theorem 6.13]{MN2025} の三次元 AS 正則
$\mathbb Z$-algebra $D$（Gorenstein parameter 四）への同値 $F$ は
\[
 F(\pi B(i))\simeq\pi e_{-2i}D,\qquad
 F(\pi X(i))\simeq\pi e_{-2i-1}D
\]
を満たす。特に構造対象は $\pi e_0D$ に移る。
$D$ の二添字の周期移動から得る自己同値 $T$ に対して
$F(1)\simeq TF$ であり、$T^2\simeq S_D^{-1}[2]$ となる。
\end{proposition}

\begin{proof}
命題~\ref{prop:reconstruction} により $B$ は domain であり、
命題~\ref{prop:gldim-two} により $\qgr B$ の大域次元は有限である。
よって中心性と standard 性を加えれば
\cite[arXiv v1, Theorem 6.13]{MN2025} の仮定を満たす。
その構成は \cite[Theorem 5.16]{MU2021} の列
$E_{2i}=\pi B(i)$、$E_{2i+1}=\pi X(i)$ を用いる。
$D_{ij}=\Hom(E_{-j},E_{-i})$ の積は合成であり、
$j$ 成分が $\Hom(E_{-j},M)$ の右 $D$-加群を tails に送ると $F$ になる
\cite[arXiv v1, Definition 5.1, Theorems 5.8 and 6.4]{MN2025}。
従って $F(E_{-i})\simeq\pi e_iD$ である。

シフト $s=(1)$ は $s E_i\simeq E_{i+2}$ を与える。
これらの同型を使い、射に $s$ を作用させることで
$\alpha:D_{ij}\xrightarrow{\sim}D_{i-2,j-2}$ を得る。
右加群上で $(TN)_j=N_{j+2}$ とし、作用を $\alpha^{-1}$ で引き戻す。
写像
\[
 \Hom(E_{-j-2},M)\longrightarrow\Hom(E_{-j},sM),
 \qquad f\longmapsto s(f)
\]
（始域には上記の同一視を用いる）は合成と可換であり、
$TF\simeq Fs$ を与える。導来同値による Serre 関手の移送から
$T^2\simeq S_D^{-1}[2]$ が従う。

具体的に $p\le q$ に対して $B_{q-p}\to D_{2p,2q}$ は
$a\mapsto s^{-q}(a)$ であり、$p\le q\le r$ での合成は
$s^{-q}(a)\circ s^{-r}(b)=s^{-r}(s^{r-q}(a)\circ b)$ となる。
従って偶数添字部分は右規約の $\check B$ であり、
$R_n$ は $D_{0,4n}$ に対応する。積は周期移動と合成で与えられる。
\end{proof}

命題~\ref{prop:sklyanin-h2} の滑らかな Sklyanin 例では
\cite[Theorem 5.5]{MUW2026} が standard 性を保証するので、
命題~\ref{prop:standard-pointed} を適用できる。
両方の ruling を別々に用いて得た比較は、どちらも同じ $(\qgr B,\pi B)$ を回復する。
逆向きに任意の cubic $\mathbb Z$-algebra から出発した場合、
二添字の周期移動だけから $T^2\simeq S_D^{-1}[2]$ を結論することはできない。

% !TeX root = ../nc-quadric-surface-classification.tex

\subsection{三次正則表示を与える ruling}
\label{subsec:cubic-ruling}

\begin{proposition}[三次型の判定]
\label{prop:cubic-ruling}
命題~\ref{prop:standard-pointed} の中央標準表示 $B$ に対し、
二つの次数零生成 MCM 加群の一方 $X$ を選び、他方を $Y$ とする。
$E_{2n}=\pi B(n)$、$E_{2n+1}=\pi X(n)$ から作る
$D_{ij}=\Hom(E_{-j},E_{-i})$ は、次の積が同型である場合にちょうど
Van den Bergh の三次三次元 AS 正則 $\mathbb Z$-algebra となる。
\[
 \mu_X:\Hom(\pi X,\pi B(1))\otimes_k\Hom(\pi B,\pi X)
       \longrightarrow \Hom(\pi B,\pi B(1)),\qquad h\otimes x\longmapsto h\circ x .
\]
\end{proposition}

\begin{proof}
必要性は \cite[arXiv v3, \S 5.1, (Q2)--(Q3)]{VdB2011} の
距離二の積に関する記述から従う。三つの Hom はそれぞれ次元 $2,2,4$ である
\cite[Lemmas 5.6, 5.9]{MU2021}。

十分性を示す。$X(n)$ を標的とする射について、
\[
 0\longrightarrow Y(n-1)\longrightarrow B(n)^2\longrightarrow X(n)
   \longrightarrow0
\]
を使う \cite[Proposition 5.5(1)]{MU2021}。
tails で $E_j$（$j<2n+1$）から $\Hom$ を取ると、
$\Ext^1(E_j,\pi Y(n-1))=0$ なので、各射は
直前の $\pi B(n)$ を経由する射の和となる。消滅は
\cite[Lemmas 5.6(8), 5.9(2), 5.11]{MU2021} による。
$B$ の一次生成性と $\mu_X$ の全射性から、$\pi B(a)$ から $\pi B(b)$
（$b>a$）への射は、隣接する helix 対象を介した合成で張られる。
残る $\pi X(a)$ から $\pi B(b)$（$b>a$）への射に対しては、もう一方の列
$0\to X(-1)\to B^2\to Y\to0$ を $\Hom_B(-,B)$ で双対化する。
$Y$ は MCM なので $B^2\to\Hom_B(X,B)(1)$ は全射となる。
ここで次数付き内部 Hom を用いる。従って $X$ から $B$ への内部 Hom は、
次数一の二つの射で左 $B$-加群として生成され、その左作用は
$B$ のシフト間の射との合成である。tails との Hom の比較
\cite[Lemma 5.9(1)]{MU2021} を使えば、全ての射が隣接する helix 対象を
介した合成で張られる。

よって $D$ は隣接成分で生成され、各隣接成分は二次元である。
右単純加群の極小分解の初めは $P_{i+1}^2\to P_i$ であり、
左単純加群でも同じ隣接生成性を使える。
Mori--Nyman の AS 条件で大域次元 $3$、Gorenstein parameter $4$ なので、
$P_i=e_iD$ とすると右単純加群の最後の項は $P_{i+4}$ である。
双対複体が添字 $i+4$ の左単純加群の極小分解になることから、
残る項は $P_{i+3}^2$ である
\cite[arXiv v1, Theorem 4.5 and Remark 4.6]{MN2025}。
$D$ の成分次元は \cite[Lemmas 5.6, 5.9]{MU2021} の
Hilbert 級数から添字の差に関して二次で抑えられる。
これは \cite[arXiv v3, Definitions 4.1.1--4.1.2]{VdB2011} の三次正則型を与える。
\end{proof}

\begin{proposition}[Sklyanin pencil での判定]
\label{prop:sklyanin-cubic-ruling}
命題~\ref{prop:sklyanin-h2} の滑らかな対象に対し、
ruling のラベルを $w\in\{z,-z-2\tau\}$ と書く。
この ruling に属する線加群を $\operatorname{coker}(X(-1)\to B_z)$
と表示して得られる $X$ に対し、
\[
 \rank\mu_X=
 \begin{cases}
  4 & w\notin E[2],\\
  3 & w\in E[2].
 \end{cases}
\]
従って少なくとも一方の選択で、標識と平方根を保持した三次正則表示が得られる。
\end{proposition}

\begin{proof}
命題~\ref{prop:standard-pointed} の仮定は
\cite[Lemma 5.1 and Theorem 5.5]{MUW2026} と命題~\ref{prop:sklyanin-h2} により満たされる。
右加群での ruling と secant の対応には
\cite[Lemma 3.20, Proposition 3.21 and Lemmas 5.3--5.4]{MUW2026} を用いる。
ここでは補題~\ref{lem:sklyanin-opposite} の次数付き同型で
$B_z$ を通常の右加群表示 $T/(g_z)$ に移してから $w$ を読む。
固定した ruling は、$p+q=w$ を満たす secant $pq$ の線加群の族である。
実際、同じラベルの二つの線加群は Lemma 5.4 の列で共通の補線を取り、
異なるラベルの二つには同じ補題の列を直接適用できる。
標準性と Proposition 3.21 により、同じラベルと同じ ruling が対応する。

異なる一般の二つの secant $pq,p'q'$ を取り、四点は相異なるとする。
対応する線加群を定める $X(-1)\to B_z$ の二射は
二次元の Hom 空間の異なる一次元部分空間を与えるので、基底をなす。
それぞれの次数一像 $W,W'\subset (B_z)_1$ は secant 上で消える
線型形式の二次元空間であり、$\operatorname{im}\mu_X=W+W'$ である。
$W\cap W'\ne0$ は四点が同一平面上にあることと同値である。
楕円曲線の次数四の埋込みでは、四点が同一平面上にあることと
$p+q+p'+q'=0$ が同値である \cite[\S 10.1]{SVdB2013}。
従って共通部分が非零となるのは $2w=0$ の場合であり、
二つの secant は相異なるため、その場合の共通部分は一次元となる。
両ラベルが $E[2]$ に属すると $4\tau=0$ となり、
$\tau$ の無限位数に反する。
命題~\ref{prop:cubic-ruling} と命題~\ref{prop:standard-pointed} から結論を得る。
\end{proof}

$z=0$ は $\tau$ が無限位数のとき滑らかな例であり、
$w=0$ の選択では距離二の積が階数三、$w=-2\tau$ の選択では階数四である。
同じ構造対象と平方根から出発しても、どちらの ruling を取るかで
三次正則型の判定結果が変わる。


\subsection{制約を保つ Zhang twist}

\begin{lemma}[次数付き自己同型が tails 上で自明となる条件]
\label{lem:scalar-tails}
適合切断環 $B$ の次数付き自己同型 $\sigma$ に対し、
右作用の引き戻し $(-)_\sigma$ が $\qgr B$ 上で恒等と自然同型となるための
必要十分条件は、ある $\lambda\in k^\times$ に対して
$\sigma(b)=\lambda^d b$ ($b\in B_d$) となることである。
この次数スカラーを $\xi_\lambda$ と書く。
\end{lemma}
\begin{proof}
$\pi B(i)_\sigma\to\pi B(i)$ の基準同型は $b\mapsto\sigma^{-1}(b)$ である。
$\End(\pi B(i))=k$ より、恒等への任意の自然同型の $\pi B(i)$ 成分は
この基準同型の非零定数 $c_i$ 倍である。
$a\in B_1$ の左乗法 $\pi B(i)\to\pi B(i+1)$ に対する自然性から
$c_{i+1}\sigma^{-1}(a)=c_i a$ を得る。ここで左乗法は
$B_1\to\Hom(\pi B(i),\pi B(i+1))$ で単射である。
実際、非零 $a$ の左乗法は domain 性により非零の十分高い次数においても単射であり、
tails で零射にならない。一次生成により $\sigma$ は次数スカラーとなる。
逆に $\sigma=\xi_\lambda$ なら次数 $j$ の成分に $\lambda^{-j}$ を掛けることで
$M_\sigma\to M$ の自然同型が得られる。
\end{proof}

\begin{proposition}[固定表示での中心化と平方根性]
\label{prop:central-twist-obstruction}
適合表示 $B=A/(g)$ に対し、$ag=g\nu(a)$ で正規化自己同型 $\nu$ を定める。
$\sigma\in\Aut_{\mathrm{gr}}(A)$ が $\sigma(g)=c g$ を満たすとし、
積 $a\star b=a\sigma^{\deg a}(b)$ の Zhang twist を取る。
twist 後に $g$ が中心元となる条件は $\sigma^2=\xi_c\nu$ である。
一方 twist 後の商の次数シフトが反 Serre の平方根となる条件は
$\bar\sigma^2$ が $B$ の次数スカラーとなることである。
従って両条件が成り立つなら $\nu=\xi_1$ または $\nu=\xi_{-1}$ である。
\end{proposition}
\begin{proof}
$g$ の正則正規性が $\nu$ の存在と一意性を与える。
$a\in A_d$ に対し $a\star g=c^d g\nu(a)$、$g\star a=g\sigma^2(a)$
なので、中心性の判定は左から $g$ を消去すればよい。
これは $g$ の直線を保つ場合の
\cite[Remark 4.4]{MUW2026} の正規化公式である。

右加群の twist 同値 $F$ は作用を、次数 $j$ の $m$ に対して
$m\star b=m\bar\sigma^j(b)$ と変える。これは regular module を保つ。
$F(M)(1)$ の作用では $M(1)$ における次数 $j$ の元に元の次数 $j+1$ を使うので、
$F(M)(1)\simeq F((M(1))_{\bar\sigma})$ となる。
従って新しいシフトを $\qgr B$ に戻すと $s(-)_{\bar\sigma}$ であり、
その二乗は $s^2(-)_{\bar\sigma^2}$ となる。
$s^2\simeq S^{-1}[2]$ と Serre 関手の移送、および
補題~\ref{lem:scalar-tails} により平方根性の条件を得る。
両条件下で $\nu$ は $B$ 上で次数スカラー、従って
$A_1=B_1$ と一次生成性から $A$ 上でも次数スカラーである。
正規化の定義から $\nu(g)=g$ なので、その一次のスカラーは $\pm1$ である。
\end{proof}

\begin{example}
\label{ex:ore-central-twist}
命題~\ref{prop:ore-ambient-family} の $A_a$ では
$\nu(x)=a^{-1}x$、$\nu(y)=ay$、$\nu(u)=u$、$\nu(v)=v$ である。
$b^2=a$ を満たす $b$ を選べば
$(x,y,u,v)\mapsto (b^{-1}x,by,u,v)$ は関係と $g$ を保ち、
二乗が $\nu$ なので中心化できる。しかし $a\ne1$ なら $\nu$ は次数スカラーでなく、
この表示からはどの次数付き自己同型 twist でも
中心性とシフトの平方根性を同時には得られない。
同じ標識対象には $a=1$ の中心表示が存在するため、
中心表示の存在一般に対する反例は与えない。
\end{example}

\subsection{二添字周期が平方根となる条件}

\begin{proposition}
\label{prop:periodicity-gauge}
$D$ を命題~\ref{prop:standard-pointed} で構成した $\mathbb Z$-algebra とする。
より一般に \cite[Theorem 6.8]{MN2025} の coherent $\mathrm{ASF}^{++}$ 正則性の仮定を満たす
次元三・Gorenstein parameter 四の連結な $D$ で、
$D_{ij}\to\Hom(\pi e_jD,\pi e_iD)$ が同型である場合にも以下は成立する。
Nakayama 同型を $\nu:D_{ij}\to D_{i-4,j-4}$ とする。
二添字移動 $\alpha:D_{ij}\to D_{i-2,j-2}$ から得る自己同値 $T_\alpha$ が
$T_\alpha^2\simeq S_D^{-1}[2]$ を満たす必要十分条件は、
ある $c_i\in k^\times$ に対し
\[
 (\nu^{-1}\alpha^2)(a)=c_i c_j^{-1}a\qquad (a\in D_{ij})
\]
となることである。
また $\pi e_0D$ の反 Serre 環は次数 $n$ で $D_{-4n,0}$、
積は $a*b=\nu^n(a)b$ ($a\in D_{-4m,0}$、$b\in D_{-4n,0}$) となる。
\end{proposition}
\begin{proof}
命題~\ref{prop:standard-pointed} の $D$ に対する仮定は
\cite[Theorems 6.4, 6.13 and Proposition 4.11]{MN2025} による。
\cite[Theorem 6.8]{MN2025} の右規約では
$S_D^{-1}[2](N)=N(4)_{\nu^{-1}}$ であり、
$\pi e_iD$ を $\pi e_{i-4}D$ に移し、左乗法の射には $\nu$ を作用させる。
$T_\alpha^2$ は同じ添字移動を行い、射には $\alpha^2$ を作用させる。
二つの関手の比は添字を固定する自己同型 $\delta=\nu^{-1}\alpha^2$ による
作用の引き戻しである。
$\End(\pi e_iD)=k$ と標準 Hom 同型の下で
各左乗法に自然性を適用すると、$\delta$ が成分ごとに $c_i c_j^{-1}$ 倍となる。
逆にこの形なら、作用を $\delta^{-1}$ で引き戻した右加群から元の加群への同型は
次数 $j$ の成分で $c_j^{-1}$ を掛けて与えられる。従って tails 上でも自然同型を得る。

積は $S_D^{-1}[2]$ を $n$ 回左乗法に作用させてから合成するので、
表示された公式となる。次数 $n$ ごとに $\nu^{-n}$ で移すと
$D_{0,4n}$ 上で積は $a*b=a\nu^{-m}(b)$ となる。
特に適合表示が存在する場合、二次関係空間は
$D_{0,4}\otimes D_{0,4}\to D_{0,8}$、$a\otimes b\mapsto a\nu^{-1}(b)$
の核である。
\end{proof}

% !TEX root = ../nc-quadric-surface.tex
\section{正規切断の標識付き比較}
\label{sec:normal-comparison}

\begin{theorem}[正規切断からの四周期 helix]
\label{thm:normal-helix}
$\HH$ の任意の適合表示を $B=A/(g)$ とする。
二つの次数零生成非自由極大 Cohen--Macaulay 型を $X,Y$ とし、
ラベル集合 $\{X,Y\}$ の交換を $c$、$Z\mapsto\Omega Z(1)$ の置換を $p$ と書く。
任意の $X_0\in\{X,Y\}$ から $X_{i-1}=pc(X_i)$ ($i\in\mathbb Z$) と定めると、
\begin{equation}
 E_{2i}=\pi B(i),\qquad E_{2i+1}=\pi X_i(i)
\end{equation}
は ample な四周期の full geometric helix をなす。
従って $D_{ij}=\Hom(E_{-j},E_{-i})$ ($i\le j$) から
右 noetherian な次元三・Gorenstein parameter 四の AS 正則
$\mathbb Z$-algebra を得る。標識は $\pi e_0D$ に移り、
反 Serre 関手は $s^2$ から誘導される四添字の移動に移る。
standard の場合は $X_i$ が一定であり、non-standard の場合は二型が交互に現れる。
\end{theorem}

\begin{proof}
正則正規切断に対する Rees の補題により $B$ は次元三・parameter 二の
AS Gorenstein 環である
\cite[Lemma 2.4]{MU2022}。parameter の変化は切断次数二による。
$B$ は domain なので $g$ は二つの次数一元の積に分解できず、
\cite[Definitions 4.10 and 5.9]{MU2022} の $n=4$ の high rank 条件を満たす。
$A$ は同論文の quantum polynomial algebra であり、$\qgr B$ の大域次元は二である。
従って \cite[Theorem 5.5, Lemma 5.13(4)]{MU2022} により
$C(B)\simeq M_2(k)\times M_2(k)$ となる。
ここで用いる定義は AS 正則性と Hilbert 級数を要求し、
Auslander 正則性と Cohen--Macaulay 性の追加条件を課さない。
同論文の Lemma 4.13(5) により、
安定圏の次数零生成部分の二つの非零単純型に対応する非自由加群を $X,Y$ と取れる。
Lemma 5.11 の証明で、対応する単純加群の次元二は
線形行列分解の階数二に等しい。この分解から極小自由分解の
各項が $B(-j)^2$ となり、$\Omega Z(1)$ も次数零生成の非自由 MCM 型である。
$\Omega(1)$ は安定圏の自己同値なので二型を置換する。
自由直和因子を除いた極小代表を取っているため、これは加群の同型型の置換 $p$ を与える。
$p^2=1$ から $X_{i+2}=X_i$ である。
$Z\in\{X,Y\}$ の極小自由被覆は
\begin{equation}
 0\longrightarrow p(Z)(-1)\longrightarrow B^2\longrightarrow Z
 \longrightarrow0
 \label{eq:normal-ruling-evaluation}
\end{equation}
を与える。
$Z$ と $p(Z)$ を交換した同様の列を双対化すると、
内部 $\Hom_B(Z,B)$ は次数一に二元で生成され、特に次数零以下は零である。
ここでは $\Ext_B^1(p(Z),B)=0$ を使った。

\cite[Lemma 4.13(3)--(5)]{MU2022} の安定圏の導来同値は、
二型を半単純環の二つの単純加群に送る。
安定圏で $\Omega$ は懸垂の逆であり、内部シフト $(1)$ は
$p$ と懸垂 $[1]$ の合成となる。従って二型の安定 Hom は
内部シフト零でのみ非零になり、同じ型の場合に $k$ である。
次数零以下での $Z$ から自由加群を経由する射は、
$\Hom_B(Z,B(j))=0$ ($j\le0$) と $Z'$ の負次数の消滅から零となる。
従って未シフトの二型の次数零 Hom も $\delta_{ZZ'}k$ である。
また $\Ext_B^1(Z,Z'(d))$ は安定 Hom
$\underline{\Hom}_B(p(Z),Z'(d+1))$ であり、
$d=-1$ かつ $Z'=p(Z)$ の場合にのみ非零となる。

自由加群と二型はいずれも深さ三である。実際、
\eqref{eq:normal-ruling-evaluation} を $p(Z)$ にも適用して
$Z$ を自由加群へ埋め込む列を得る。その余核も二型のシフトである。
この操作を繰り返すと、自由加群の深さ三と長完全列から
$\Ext^j_B(k,Z)=0$ ($0\le j<3$) が従う。
従って $\Hom$ と $\Ext^1$ は、\cite[Lemma 5.9(1)--(2)]{MU2021}
の打切りの議論により graded 加群から tails に移しても変わらない。
$\Ext^1$ はどちらかが自由加群なら零である。
奇数項同士では前向きの射の $\Ext^1$ は零であり、
四連続項の中の逆向きでは、$X_{i-1}=pc(X_i)\ne p(X_i)$
により零である。$\Hom$ の逆向き消滅は次数と
$\Hom_B(Z,B)$ の次数一以上への集中から従う。
$\Ext^2(E_i,E_j)$ は Serre 双対性により $\Hom(E_j,E_{i-4})$ の双対であり、
$i\le j$ または $0<i-j<4$ のとき逆向きの Hom なので零となる。
高次は大域次元から零である。
この結果と $E_{i+4}\simeq s^2E_i$ が geometric helix の条件を与える。
ここで $s^{-2}$ は導来圏上で $S[-2]$ となる exact 自己同値なので、
この圏の canonical bimodule を与える。

完全性を示す。Orlov の分解と MCM 安定圏の記述により、
$D^b(\qgr B)$ は $\pi B(-1),\pi B,\pi X,\pi Y$ で生成される
\cite[Theorem 2.5(i)]{Orlov2009}。
実際、$\Hom_B(Z,B(j))=0$ ($j\le0$) なので、
Orlov の埋込みは二型の特異圏での像を $\pi X,\pi Y$ に送る
\cite[Lemma 4.1]{MU2016Stable}。二型は半単純な安定圏の導来表示を生成する。
これを任意に次数シフトする。$Y_i=c(X_i)$ と書くと、
\eqref{eq:normal-ruling-evaluation} の $Y_i(i)$ の核が $X_{i-1}(i-1)$ なので、
生成列 $(\pi B(i-1),\pi B(i),\pi Y_i(i),\pi X_i(i))$
の中の隣接二項を対応する三角で入れ替えて、偶数添字で始まる四項の完全性を得る。
奇数添字で始まる場合は、補型 $c(X_i)$ を同じ方法で選んだ
偶数添字から始まる四項
$(\pi B(i+1),\pi Y_{i+1}(i+1),\pi B(i+2),\pi Y_{i+2}(i+2))$
で、二つの隣接対をそれぞれ同じ三角で置き換える。

豊富性の生成条件は偶数項から従う。
残る条件に必要な $\Ext^1(\pi Z(-n),K)=0$ ($n\gg0$) を示す。
任意の $K\in\qgr B$ について、
$0\to L\to P\to K\to0$ を短完全列とし $P$ を有限個の $\pi B$ のシフトの直和とする。
$\Ext^1(\pi Z(-n),P)=0$ なので $\Ext^1(\pi Z(-n),K)$ は
$\Ext^2(\pi Z(-n),L)$ に埋まる。
同様の自由被覆 $P'\twoheadrightarrow L$ と大域次元二を用いると、
$\Ext^2(\pi Z(-n),L)$ は $\Ext^2(\pi Z(-n),P')$ の商となる。
Serre 双対性と負次数の消滅から
$\Ext^2(\pi Z(-n),P')=\Hom(P',\pi Z(-n-2))^*=0$ ($n\gg0$) となる。
よって十分負の奇数項からの $\Hom$ は全射を保つ。
以上で \cite[Definition 5.6]{MN2025} の豊富性が従う。

\cite[Theorems 5.8, 6.4 and 6.5]{MN2025} を適用する。
構造対象と反 Serre 関手の移送は命題~\ref{prop:standard-pointed}
と同じ Hom の構成で得られる。non-standard の場合、
ここで使う表示に二添字の移動があることは要求していない。
\end{proof}

\begin{corollary}[三次正則型からのずれ]
\label{cor:normal-cubic-defect}
定理~\ref{thm:normal-helix} の $D$ は隣接成分と、距離二の偶数添字間の成分で生成される。
選んだ列について
\[
 r_i=\rank\bigl(
 \Hom(\pi X_i,\pi B(1))\otimes\Hom(\pi B,\pi X_i)
 \longrightarrow\Hom(\pi B,\pi B(1))
 \bigr)
\]
と置くと、$r_i\in\{3,4\}$ であり、距離二の偶数添字間で
追加の生成元を $4-r_i$ 個要する。$D$ が Van den Bergh の
三次正則型となる条件は全ての $r_i$ が四となることである。
また $r_i$ は $i$ に依存しない。
\end{corollary}
\begin{proof}
$\pi B(i)$ を始域かつ $\pi B(j)$ を終域とする射は、$B$ の一次生成性から
距離二の積で生成される。
奇数項を終域とする射で始域がそれより前にある場合は、
\eqref{eq:normal-ruling-evaluation} のシフトに $\Hom$ を適用して
直前の偶数項を経由する射で張られる。
実際、核のシフトは始域のシフトより小さくない
（始域が同じシフトなら自由加群である）ので、核への $\Ext^1$ は零である。
奇数項を始域とし偶数項を終域とする射については、内部 $\Hom_B(X_i,B)$ の
次数一生成性を使う。よって主張した距離で生成される。
各隣接空間は二次元である。

下限 $r_i\ge3$ を示す。$Z=X_i$ とする。$Z$ と $p(Z)$ は
\eqref{eq:normal-ruling-evaluation} からともに自由加群へ埋め込める。
noetherian domain $B$ の商斜体上で正の rank を持ち、
同じ完全列から両者の rank の和は二なので、各 rank は一である。
従って任意の非零斉次射 $Z(-1)\to B$ は単射で、
その像は次数一の二次元部分空間で生成される。
一次独立な二射の像の次数一成分が一致したとすると、
次数一生成性から像全体が一致する。このとき二射は
同じ像への同型であり、$\End_{\mathrm{gr}}(Z)=k$ に反する。
従って二つの二次元部分空間の和の次元は三以上となる。
この和が上の積の像である。
残る距離二の偶数添字間の生成元の個数は、その積の余次元 $4-r_i$ となる。

距離二の奇数項間の Hom は四次元である。
実際、$X_{i+1}(i+1)$ の被覆列の核 $p(X_{i+1})(i)$ に対し、
$\Hom(\pi X_i(i),-)$ を適用する。$X_i\ne p(X_{i+1})$ による
次数零 Hom の消滅と $\Ext^1$ の消滅から、
直前の偶数項を経由する積が同型であることが分かる。

$D$ の添字 $j$ から始まる距離二の生成元の個数を $d_j$ とする
（奇数 $j$ では零）。右単純加群の極小分解は
\[
 0\longrightarrow P_{j+4}\longrightarrow
 P_{j+3}^{2}\oplus P_{j+2}^{d_{j+2}}\longrightarrow
 P_{j+1}^{2}\oplus P_{j+2}^{d_j}\longrightarrow P_j\longrightarrow S_j
 \longrightarrow0
\]
となる。右側の第一項は生成元から分かり、双対極小分解では同じ議論が
左単純加群に使える。左側から見た距離二の生成元の個数も、
同じ隣接積の余次元なので $d_{j+2}$ である。
分解の次数 $j+2$ 成分の Euler 標数を取ると、
$4-2\cdot2-d_j+d_{j+2}=0$、すなわち $d_j=d_{j+2}$ を得る。
従って $r_i$ は $i$ に依存しない。特に non-standard では
二つの MCM 型の階数が等しい。

全ての $r_i$ が四なら分解は三次型の形になる。
成分次元の多項式上界は、
$B(i)^2\twoheadrightarrow X_i(i)$ と各 $X_j(j)$ の線形自由分解による
Hilbert 級数 $2(1-t)^{-3}$ から従う。
逆に三次正則型なら距離二の積は全射である。
\end{proof}

\begin{remark}[階数三の行列表示]
\label{rem:rank-three-matrix}
階数三の場合、$\Hom_B(X_i,B(1))$ の基底と $(X_i)_0$ の基底を選ぶと、
合成の四成分は $A_1=B_1$ の元を成分とする行列 $\Phi$ になる。
双対評価列と線形行列分解において、これは残りの行列 $\Psi$ と
$\Phi\Psi=gI$ を満たす側の行列として取れる。
四成分の間の唯一の線形関係を表す $2\times2$ 係数行列は階数二である。
実際、もし階数一なら基底変換で $\Phi$ の一成分を零にできる。
この場合、零成分を含む行と積 $\Phi\Psi$ の対角成分を見ると
$g$ が二つの次数一元の積となり、$B$ の domain 性に反する。
従って行と列の基底変換でこの関係を trace が零という形にできる。
そのように取ると
\[
 \Phi=\begin{pmatrix} a&b\\ c&-a\end{pmatrix},
 \qquad a,b,c\in A_1\ \text{は線形独立}
\]
となる。階数三の領域では、この形と対応する $\Psi$ の積の式が
残りの階数および別表示との比較の出発点となる。
\end{remark}

\begin{corollary}[non-standard 型の正規化条件]
\label{cor:nonstandard-square-matrix}
$p(Z)=Z$ のとき、対応する評価行列を $M\in M_2(A_1)$ として
\[
 M^2=gC,\qquad \nu(M)=CMC^{-1},\qquad C\in\GL_2(k)
\]
を満たす形に取れる。ここで $\nu$ は $ag=g\nu(a)$ の規約の正規化自己同型である。
特に $g$ が中心なら $C$ はスカラーであり、$M$ を定数倍して $M^2=gI$ とできる。
\end{corollary}
\begin{proof}
\eqref{eq:normal-ruling-evaluation} の射を
$\iota:Z(-1)\to B^2$ と $q:B^2\to Z$ とする。
積 $\iota q(-1):B(-1)^2\to B^2$ を $\bar M$ と書くと、
$q\iota=0$ により $\bar M^2=0$ となる。
これは $Z$ の極小線形表示であり、
それを与える行列分解の第一行列は、自由加群の基底変換の後、
$\bar M$ の $A_1$ への一意な持上げ $M$ と一致する。
従って $M$ は $A$ 上で単射である。
$\bar M^2=0$ と $\ker(A_2\to B_2)=kg$ から $M^2=gC$ となる。
$C$ が非可逆なら非零の定数列 $v$ で $Cv=0$ があり、
$M$ の単射性に反して $M^2v=0$ となる。
最後に $M(gC)=(gC)M$ に $Mg=g\nu(M)$ を代入し、
正則元 $g$ を左から消去して共役の式を得る。
中心切断では $C$ は $M$ の各係数行列と可換である。
系~\ref{cor:normal-cubic-defect} により係数行列の張る空間は三次元以上だが、
非スカラーな二次正方行列の中心化環は二次元なので、$C$ はスカラーである。
基礎体は代数閉なので必要な平方根を取り、$M$ を定数倍できる。
\end{proof}

% !TEX root = ../nc-quadric-surface.tex
\subsection{二型がともに階数三となる非標準表示}
\label{sec:rank-three-example}

\begin{proposition}
\label{prop:nonstandard-rank-three-example}
標数零の代数閉体上に、二つの MCM 型の行列成分がともに三次元を張る
non-standard な $\HH$ の適合表示が存在する。
この表示と定理~\ref{thm:normal-helix} のいずれの初期型を用いても、
得られる四周期表示は Van den Bergh の cubic 型にならない。
\end{proposition}

\begin{proof}
$k$ の元 $\mathrm i$ を $\mathrm i^2=-1$ となるように選び、
\[
 A=k[u,v]\otimes_k k\langle x,y\rangle/(yx+xy),\qquad
 g=\tfrac14(x^2+y^2)-uv,\qquad B=A/(g)
\]
とする。$A$ は四変数の skew polynomial ring であり、
命題~\ref{prop:commutative-twist-covers} で用いた skew Koszul 複体と同様に
型 $(1,4,6,4,1)$ である。$g$ は中央正則である。
$B$ は $k[u,v,x]$ 上で $y$ に関する Ore 拡大を取り、
モニックな関係 $y^2=4uv-x^2$ を課した環なので、
左 $k[u,v,x]$ 加群として基底 $1,y$ を持つ。
特に $u$ は正則である。$u$ を可逆にすると $v$ を消去でき、
\[
 B[u^{-1}]\simeq k[u,u^{-1}]\otimes_k k\langle x,y\rangle/(yx+xy)
\]
となるため、$B$ は domain である。

滑らかさは $B$ の Koszul 双対から判定できる。
$U,V,X,Y$ を $u,v,x,y$ の双対基底とすると、二次関係の直交補空間は
\[
\begin{gathered}
 U^2=V^2=0,\quad X^2=Y^2=w,\quad UV+VU=-4w,\quad XY=YX,\\
 UX=-XU,\quad UY=-YU,\quad VX=-XV,\quad VY=-YV
\end{gathered}
\]
を与える。ここで $w=X^2$ は中央であり、$B^!/(w)=A^!$ となる
正則元の代表である \cite[Lemma 4.11(2)]{MU2022}。$P=(U+V)/(2\mathrm i)$、$Q=(U-V)/2$ と置けば
$P^2=Q^2=w$、$PQ=-QP$ である。
$C(B)=B^![w^{-1}]_0$ は $P X^{-1},Q X^{-1},Y X^{-1}$ によって生成される。
最初の二つはそれぞれ平方が $-1$ で互いに反交換し、最後は中央で平方が $1$ である。
従って $M_2(k)\times M_2(k)$ から $C(B)$ への全射を得る。
\cite[Lemma 4.13(1)]{MU2022} により $C(B)$ は八次元なのでこの全射は同型であり、
同論文 Theorem 5.5 によって $\qgr B$ は滑らかである。

$B$ は正則二次切断なので次元三・parameter 二の AS Gorenstein 環である
\cite[Lemma 2.4]{MU2022}。canonical bimodule を与える右作用の twist を計算する。
\cite[arXiv v2, Example 5.5]{RRZ2014} により、$A$ の Nakayama 自己同型は
$u,v$ を固定し $x,y$ に符号 $-1$ を掛ける。
これは $g$ を固定するため同論文 Lemma 1.5 を適用でき、
$B$ の Nakayama 自己同型もその誘導 $\mu$ となる。
$\mu^{-1}=\mu$ なので、\cite[Example 3.6(2)]{MU2021} の
右 canonical bimodule の規約では
$S_{\qgr B}^{-1}[2]\simeq(-)_\mu(2)$ となる。

ここで graded 自己同型
\[
 \sigma(u)=-u,\qquad \sigma(v)=v,\qquad
 \sigma(x)=\mathrm i x,\qquad \sigma(y)=\mathrm i y
\]
を取る。$\sigma(g)=-g$ なので $B$ に誘導され、
その二乗は $\mu$ である。命題~\ref{prop:central-twist-obstruction} の証明中の
作用比較を用いると、Zhang 同値 $F:\qgr B\to\qgr B^\sigma$ は $\pi B$ を
$\pi B^\sigma$ に送り、新しい次数シフトを旧圏では $(-)_\sigma(1)$ として実現する。
その二乗は $(-)_\mu(2)$ なので、新しい次数シフトは反 Serre の平方根である。
$A^\sigma$ は型 $(1,4,6,4,1)$ を保ち \cite[Lemma 4.1]{MUW2026}、
$g$ はその中で左右正則正規であり、$B^\sigma=A^\sigma/(g)$ である。

$\pi B^\sigma$ に対する切断と豊富性を確認する。
正規正則切断の AS Gorenstein 性と深さ三から、
全次数での飽和性、$\operatorname{Ext}^1(\pi B^\sigma,\pi B^\sigma)=0$ が従う。
すなわち次数付き有限長加群 $L$ について $\operatorname{Ext}^j(L,B^\sigma)=0$
($j<3$) を用い、自由加群の打切りと Hom の極限を取る
\cite[Lemma 5.9(1)--(2)]{MU2021}。
$\operatorname{Ext}^2$ は Serre 双対性により負次数の切断の双対なので零である。
高次も大域次元二により零であり、$\End(\pi B^\sigma)=k$ と合わせて exceptionality を得る。
任意の対象は有限個の $\pi B^\sigma(-n)$ の直和の商であり、
その指数は打切りにより任意に大きく取れる。
全射の核 $K$ について $\operatorname{Ext}^1(\pi B^\sigma(-n),K)=0$
($n\gg0$) を得るには、$K$ の自由被覆を取って、
その核への $\operatorname{Ext}^2$ に埋め込む。
その核の自由被覆と大域次元二から、最後の群は
自由加群の像への $\operatorname{Ext}^2$ の商となり、Serre 双対性と
負次数の切断の消滅により零である。
従って $(\pi B^\sigma,(1))$ は ample であり、その切断環は $B^\sigma$ である。

Gabriel--Krull 条件は同値で不変なので $\qgr B$ で確認する。
命題~\ref{prop:sklyanin-h2} と同じく $A$ 上の有限自由分解により
Hilbert 関数は最終的に多項式となる。定数なら有限長、次数が高々一なら
$\cC_1$ に入ることは \eqref{eq:sklyanin-polynomial-filtration} と同じ議論で従う。
一方 $B/uB$ 上で $v$ は正則であり、$B/(u,v)$ は無限次元である。
従って $v$ の冪による列は $\pi(B/uB)$ の真の下降列となり、
この対象は有限長でない。さらに $u$ の冪による列の各相続商は
$\pi(B/uB)$ のシフトであるから、$\pi B\notin\cC_1$ である。
$\pi B$ の任意の非零部分対象は、非零斉次元 $f\in B$ に対する
$\pi(fB)\simeq\pi B(-\deg f)$ を含む。
$B/fB$ の Hilbert 多項式は次数が高々一なので、非零部分対象による商は
$\cC_1$ に入る。よって $\pi B$ は rank one torsion-free である。

最後に $c=(x+y)/2$、$d=(x-y)/2$ とし、$r=c,d$ に対し
\[
 M_r=\begin{pmatrix}r&u\\ -v&-r\end{pmatrix}
\]
と置く。$c^2=d^2=(x^2+y^2)/4$ により $M_r^2=gI$ であり、
これらは reduced な線形行列分解である。
成分の張る空間はそれぞれ $\langle u,v,c\rangle$ と
$\langle u,v,d\rangle$ である。これらは異なり、
行と列の定数基底変換で保たれるので、二つの余核は同型でない。
各分解では二つの連続行列が同じなので、$\Omega(1)$ は対応する余核を固定する。
この二つが全 MCM 型であることは $C(B)\simeq M_2(k)\times M_2(k)$ と
\cite[Lemma 4.13(4)--(5)]{MU2022} による。
実際、余核は二元生成で、表示行列は非零の最小関係を与えるため自由ではない。
非自由直和因子は少なくとも二元で生成されるので、各余核は直既約である。

$D=\operatorname{diag}(\mathrm i,1)$ とすると
$\sigma(M_r)=\mathrm i D M_r D^{-1}$ なので、右作用の $\sigma$ による引戻しは
両方の型を固定する。
Zhang 同値は極小自由被覆とその核を移送し、
$F(M)(1)\simeq F(M(1)_\sigma)$ であるから、
$B^\sigma$ でも $\Omega(1)$ は両方の型を固定する。
最小表示行列の成分は twist 後も同じ次元の空間を張る。
実際、$F(B(-1))\simeq B^\sigma(-1)$ を $m\mapsto\sigma^{-1}(m)$ で同定すると、
次数一成分が $a$ の左乗法は $a\star(-)$ となる。
従って $B^\sigma$ は non-standard であり、双方の成分空間が三次元である。
系~\ref{cor:normal-cubic-defect} により、どちらの初期型を選んでも
cubic 型は得られない。
\end{proof}

任意の別表示を含めた cubic 型の排除は、
\ref{thm:cubic-coverage} で証明する。

\begin{corollary}[この skew polynomial 環の中央切断の許容域]
\label{cor:skew-rank-three-central-locus}
上の $A$ の非零中央二次元 $g$ を
\[
 g=a x^2+b y^2+q(u,v)
\]
と書く。$A/(g)$ が domain かつ滑らかとなるための必要十分条件は
$ab\ne0$ かつ二次形式 $q$ が非退化となることである。
この領域では全ての中央切断およびその Zhang twist は
同じ pointed exact 同値類を与える。この類は $\HH$ に属し、
定理~\ref{thm:cubic-coverage} の cubic 表示を持たない類である。
\end{corollary}

\begin{proof}
次数二の中心は表示された五つの二次元で張られる。
これは skew polynomial の標準単項式を用い、
$x,y$ と交換する元の $x,y$ の次数がともに偶数となることから従う。
$ab\ne0$ かつ $q$ 非退化なら、$x,y$ のスカラー変更と
$u,v$ の線形基底変更で命題~\ref{prop:nonstandard-rank-three-example} の
$g$ に移る。この変更は $A$ の自己同型である。
Zhang 同値は構造対象を保存するので、この条件が十分である。

必要性を示す。$a=b=0$ なら非零の $q$ は一次元の積に分解し、
二因子は商で非零なので商は domain でない。
残りでは $x,y$ を交換して $a\ne0$ としてよい。
$B=A/(g)$ とし、$X,Y$ を $x,y$ の双対基底とする。
$B^!$ において $w=X^2$ は中央正則元の代表であり、
$XY=YX$、$Y^2=(b/a)w$ となる。
$u,v$ の双対に由来する二つの生成元は $X,Y$ と反交換し、
その平方と相互の反交換子は $w$ に $q$ の係数を掛けたものである。
これらはいずれも $B$ の二次関係空間との直交性を評価して得られる。
従って $C(B)=B^![w^{-1}]_0$ は
$Y X^{-1}$ と、$u,v$ の双対元に $X^{-1}$ を掛けた二元で生成される。

$b=0$ なら $Y X^{-1}$ は非零の中央平方零元なので $C(B)$ は
半単純でない。非零性は $w$ の正則性による。
また $q$ が退化しているなら、$u,v$ の双対側でその二次形式の根基から
非零元 $L$ を取れる。$L$ は平方が零であり、他の双対二次形式の
生成元との反交換子も零である。
よって $L X^{-1}$ は非零の平方零元で、
他の生成元と交換または反交換する。従ってそれが生成する
両側イデアルは非零で平方零となる。
いずれの場合も \cite[Theorem 5.5]{MU2022} により $\qgr B$ は滑らかでない。
\end{proof}

\begin{theorem}[四変数符号型の対角切断]
\label{thm:sign-quadric-diagonal}
$\epsilon_{ij}\in\{1,-1\}$ とし、
$A_\epsilon=k\langle x_1,\dots,x_4\rangle/
(x_jx_i-\epsilon_{ij}x_ix_j)_{i<j}$ とする。
非零中央元 $g=\sum_i a_i x_i^2$ による切断を考える。
$A_\epsilon/(g)$ が滑らかな domain となるのは、全ての $a_i$ が非零であり、
\[
 \epsilon_{ij}\epsilon_{ik}\epsilon_{jk}=-1
 \qquad(i<j<k)
\]
となる三つ組の個数が零または二の場合にちょうど限る。
前者は可換二次曲面の標識類、後者は
命題~\ref{prop:nonstandard-rank-three-example} の標識類を与える。
この二類はいずれも $\HH$ に属し、pointed exact 同値でない。
\end{theorem}

\begin{proof}
$a_j\ne0$ を一つ選ぶ。
二次双対において $w=X_j^2$ は中央正則で、
$X_i^2=(a_i/a_j)w$、$X_iX_l=-\epsilon_{il}X_lX_i$ ($i<l$) である。
$C(B)$ は $X_iX_j^{-1}$ ($i\ne j$) で生成され、
これらは相互に符号付きで交換する。
$a_i=0$ なら $X_iX_j^{-1}$ は零でない平方零元であり、
これが生成する両側イデアルも平方零である。
従って滑らかさに必要な半単純性が失われる
\cite[Theorem 5.5]{MU2022}。

以後 $a_i\ne0$ とする。対角符号自己同型 $x_i\mapsto s_i x_i$ による
Zhang twist は $\epsilon_{ij}$ を $s_i s_j\epsilon_{ij}$ に替える。
$g$ の新しい各平方に対する係数は非零のままなので、
変数のスカラー変更で正規化できる。
$g$ は自己同型で固定され、商の twist は構造対象を保つ。
また domain 性は斉次元の積に関してこの twist で不変である。

この符号変更で特定の頂点と残りの全ての間の符号を正にできる。
その後に残る三辺の符号は、その頂点を含む三つ組の積で決まる。
従って三つ組の四つの符号が各ラベルで一致するなら
二つの符号行列は twist で移り合う。
四符号の積は一であり、頂点の置換は三つ組の置換を全て実現する。
従って負の三つ組の個数は零、二、四のいずれかであり、
それぞれ全辺正、一辺だけ負、全辺負を代表として取れる。

零の場合の中央切断は滑らかな可換二次曲面である。
二の場合は可換変数二つと反交換する変数二つを持つ
前の系の非退化切断である。
四の場合には $\sum_i a_i x_i^2=(\sum_i \sqrt{a_i}x_i)^2$ となる。
一次の元は商で非零であり、平方が零なので domain ではない。
二つの残る標識類の非同値は、
可換類が cubic 表示を持ち、後者が持たないこと
（定理~\ref{thm:cubic-coverage}）から従う。
\end{proof}

% !TEX root = ../nc-quadric-surface-classification.tex
\subsection{符号型 skew polynomial 環の全正規切断}
\label{subsec:sign-normal-quadrics}

\begin{theorem}
\label{thm:sign-quadric-normal}
四変数符号型 skew polynomial 環 $A_\epsilon$ の非零正規二次元 $g$ を考える。
$B=A_\epsilon/(g)$ と置く。$(\qgr B,\pi B)$ が $\HH$ に属する場合の
pointed exact 同値類は、次の三つである。
\begin{enumerate}
\item 滑らかな可換二次曲面。
\item 命題~\ref{prop:nonstandard-rank-three-example} の非 cubic 型。
\item $g=x_1x_2+x_3x_4$、
$(\epsilon_{12},\epsilon_{34})=(-1,1)$、
$(\epsilon_{13},\epsilon_{14},\epsilon_{23},\epsilon_{24})=(1,-1,1,-1)$ の切断。
この類は cubic 型であり、反 Serre 環は非可換である。
\end{enumerate}
\end{theorem}

\begin{proof}
まず中央切断は定理~\ref{thm:sign-quadric-diagonal} に帰着する。
実際、$\epsilon_{ii}=1$ として符号行列を対称に書くと、
二次単項式 $x_i x_j$ ($i\ne j$) は第 $i,j$ 行が一致する場合に限り中央である。
各変数との交換は標準単項式ごとに検査できる。
同じ行を持つ変数は相互に可換であり、他の変数との交換符号も一様である。
従って同じ行の変数内での任意の線形基底変更は自己同型を与える。
中央二次形式はこれらの可換ブロック上の二次形式の和なので、
ブロックごとに対角化すればよい。

次に非中央の正規 $g$ を扱う。$x_i g=g\nu(x_i)$ と書く。
符号型の標準単項式の任意の辞書式順序で最高次項を比較すると、
積の最高指数は二因子の最高指数の和である。
従って $\nu(x_i)$ の最高単項式はどの順序でも $x_i$ であり、
$\nu(x_i)=\chi_i x_i$ となる。ここで $\chi_i=\pm1$ である。
$g$ の各単項式は同じ交換指標 $\chi=(\chi_i)_i$ を持つ。
$\chi\ne1$ なので平方は現れず、$x_i x_j$ が現れる条件は
第 $i,j$ 行の成分ごとの積が $\chi$ となることである。
同じ行のブロックとそれに $\chi$ を掛けたブロックは相異なる。
従って $g$ は互いに素なブロック対に属する双線形形式の和である。
それぞれを左右の基底変更で階数標準形にすると、四変数では
一つの積または二つの互いに素な積の和となる。
一つの積の商は domain でないので、
$g=x_1x_2+x_3x_4$ と取る。

$p=\epsilon_{12}$、$q=\epsilon_{34}$、$r=\epsilon_{13}$ とすると
両単項式の指標が一致するため
\[
 (\epsilon_{13},\epsilon_{14},\epsilon_{23},\epsilon_{24})
       =(r,pr,qr,pqr),\qquad
 \chi=(p,p,q,q).
\]
$p=q$ なら全ての三つ組の符号積は正である。
全ての交換符号を正にする対角符号 twist は、
両方の積を同じ符号 $p$ 倍するので $g$ の直線を保つ。
その可換商は滑らかな Segre 二次曲面を与える。
$p\ne q$ なら二組の交換と、必要なら第二組内の交換で
$(p,q,r)=(-1,1,1)$ を選べる。項の係数は変数のスカラー変更で戻せる。
これが (3) に掲げた切断である。

以下 (3) の $B=A_\epsilon/(g)$ を確認する。
$x_1$ は正規であり、$A_\epsilon/(x_1)$ は domain で $g$ の像も非零である。
$x_1 b=ga$ とすれば $a\in (x_1)$ となり、
$A_\epsilon$ 上で $x_1$ を消去して $b\in(g)$ を得る。
従って $x_1$ は $B$ 上で正則である。
$x_1$ を可逆にすると $x_2=-x_1^{-1}x_3x_4$ を代入でき、
残る三変数の skew Laurent 拡大に埋まるので $B$ は domain である。

二つの線形行列
\[
 \Phi=\begin{pmatrix}x_1&x_3\\x_4&x_2\end{pmatrix},\qquad
 \Psi=\begin{pmatrix}x_2&x_3\\x_4&-x_1\end{pmatrix}
\]
は $\Phi\Psi=gI$、$\Psi\nu(\Phi)=gI$ を満たす。
従ってどちらから始めても reduced な正規行列分解を得る。
可換多項式と見なした二つの行列式は
$x_1x_2-x_3x_4$ と $-x_1x_2-x_3x_4$ であり、比例しない。
従って定数の行・列変換で互いに移らず、
極小表示が異なる二つの二元生成 MCM 型を得る。
$g$ は一次元の積に分解しないので
\cite[Lemma 5.11]{MU2022} により、
対応する二次元の $C(B)$-加群は単純である。
ここで使う階数と次元の対応は同補題の証明および
同論文 Lemma 4.13 による。
$\dim C(B)=8$ なので半単純商の二つの四次元成分で全次元を使い、
$C(B)=M_2(k)\times M_2(k)$ となる。
従って $\qgr B$ は滑らかである。
両行列の成分空間は四次元である。

この符号行列では各行の符号積が正なので
$A_\epsilon$ の Nakayama 自己同型は恒等である
\cite[Example 5.5]{RRZ2014}。
$B$ の canonical twist は $\nu=\nu^{-1}$ である
\cite[Lemma 1.5]{RRZ2014}。
canonical bimodule と Serre 関手との対応には \cite[Example 3.6(2)]{MU2021} を用いる。
$\sigma$ を各変数に順に $\mathrm i,-\mathrm i,1,1$ を掛ける対角自己同型とすると、
$\sigma(g)=g$、$\sigma^2=\nu$ である。
従って $B^\sigma$ のシフトは旧圏では $(1)(-)_{\sigma}$ であり、
その平方が反 Serre となる。$A_\epsilon^\sigma$ は型 $(1,4,6,4,1)$ を保ち、
$B^\sigma$ はその正則正規二次商である。
飽和性、例外性、豊富性は命題~\ref{prop:nonstandard-rank-three-example} と同じ
正則切断および Serre 双対性の議論で従う。

Gabriel--Krull 条件については $B/(x_1)$ が
$B/(x_1,x_3)\simeq k_{-1}[x_2,x_4]$ に全射を持つ。
後者の正則元 $x_2$ の冪による下降列は、
隣接商が $k[x_4]$ のシフトなので tails で真に下降する。
従って $\pi(B/x_1B)$ は有限長でなく、
$x_1$ の冪による列により $\pi B\notin\cC_1$ となる。
domain 性と Hilbert 多項式から、命題~\ref{prop:nonstandard-rank-three-example}
と同様に rank one torsion-free 性を得る。
従って (3) は $\HH$ に属する。
twist 後も極小表示の成分空間の次元は保たれるので、
定理~\ref{thm:cubic-coverage} により (3) は cubic 型である。

最後に $B^\sigma$ の偶数 Veronese において、
$x_1x_3$ と $x_2x_3$ は反交換する。
旧積で既に反交換し、偶数の新しい積では
両方に同じ $\sigma^2$ の符号 $-1$ が働くからである。
domain 性から積は非零なので反 Serre 環は非可換である。
従って (3) は可換類とも非 cubic 類とも異なる。

逆に、どの切断の標識対も、$\HH$ に入るなら滑らかである。
また正則正規二次切断の AS Gorenstein 性により
元の切断環は飽和している。
元の次数シフトの $\pi B$ の軌道も rank one torsion-free なので、
切断の非零積は単射の合成となり、元の $B$ は domain である。
このため上で排除した非 domain または非滑らかな切断は、
別の平方根を選んでも新たな $\HH$ 類を与えない。
\end{proof}

% !TEX root = ../nc-quadric-classification.tex
\section{多パラメータ skew polynomial 環上の正規二次切断}
\label{sec:multiparameter-quadrics}

ここでは $q_{ij}\in k^\times$、$q_{ji}=q_{ij}^{-1}$、$q_{ii}=1$ とし、
$A=k\langle x_1,x_2,x_3,x_4\rangle/(x_jx_i-q_{ij}x_ix_j)_{i<j}$ を考える。
$Q_r$ ($r\in k^\times$) を、次のパラメータの skew polynomial 環の
$g=x_1x_2+x_3x_4$ による商とする。
\[
 q_{12}=q_{34}=1,\qquad (q_{13},q_{14},q_{23},q_{24})=(r,r^{-1},r^{-1},r).
\]

\begin{theorem}[多パラメータ切断の分類]
\label{thm:multiparameter-quadric-classification}
上の $A$ の正規二次切断から得られる $\HH$ の pointed exact 同値類は、
命題~\ref{prop:nonstandard-rank-three-example} の非 cubic 型一つと、
$Q_r$ の族である。$Q_r$ はすべて cubic 型であり、
\[
 (\qgr Q_r,\pi Q_r)\simeq(\qgr Q_{r'},\pi Q_{r'})
 \quad\Longleftrightarrow\quad
 r'^2=r^2\ \text{または}\ r'^2=r^{-2}
\]
が成り立つ。可換類は $r^2=1$ に対応する。
\end{theorem}

\begin{proof}
正規切断に対してその元の切断環が domain かつ滑らかであることは、
定理~\ref{thm:sign-quadric-normal} の末尾と同じ理由で必要である。
正規化自己同型が対角となる最高単項式の議論も、交換係数が非零ならそのまま使える。
各変数の交換指標を $\rho_i=(q_{ij})_j$ と書く。
$g$ の各単項式は同じ指標 $\chi$ を持ち、
$x_i x_j$ の指標は $\rho_i\rho_j$ である。
等しい指標を持つ変数のブロック内では変数は可換であり、
ブロックごとの基底変更は $A$ の自己同型である。
$\rho\mapsto\chi/\rho$ が固定するブロックには二次形式が現れ、
別のブロックへ移す場合にはブロック対に双線形形式が現れる。
それぞれ対角化または階数標準形にすることで、$g$ は
重複のない変数の平方と、互いに素な変数対の積の和になる。

この形で $g$ に現れない変数があれば商は滑らかでない。
実際、その変数の双対 $X$ は $(A/(g))^!$ で平方零であり、
他の双対生成元と非零スカラーを介して交換する。
このため $(X)$ は非零平方零イデアルである。
$w$ の代表は、$g$ に現れる一つの平方の双対平方、
または一つの積に対応する双対の二項和として取れる。
いずれも同じ変数の積からなるため、$Xw$ と $wX$ はスカラー倍で一致する。
\cite[Lemma 4.11]{MU2022} の正則正規二次元 $w$ を逆にしても消えず、
その局所化の次数零部分にも非零の平方零イデアルを与える。
実際 $Xw\ne0$ なので、$w$ を二文字の和に展開すると
ある双対生成元 $Y$ に対して $XY\ne0$ となり、$XYw^{-1}$ がその部分に入る。
従って $C(B)$ は半単純でなく、\cite[Theorem 5.5]{MU2022} により
商は滑らかでない。

全変数が現れる場合、可能な形は四平方、二平方と一対の積、二対の積である。
四平方なら各成分で $\chi_j=q_{jj}^2=1$ なので全交換係数が符号となり、
定理~\ref{thm:sign-quadric-normal} が適用できる。
二平方の場合に $g=x_1^2+x_2^2+x_3x_4$ とする。
指標の等式から $a=q_{13}$ とすれば
$q_{14}=a^{-1}$、$q_{23}=\pm a$、$q_{24}=q_{23}^{-1}$、
$q_{34}=a^{-2}$、$q_{12}=\pm1$ である。
対角自己同型 $(1,1,a,a^{-1})$ による Zhang twist は
$g$ の直線を保存し、交換係数をすべて符号にする。
従ってこの場合も符号型の分類に帰着する。

二対の積の場合は $g=x_1x_2+x_3x_4$ と取る。
$p=q_{12}$、$q=q_{34}$、$a=q_{13}$ とすると
\[
 q_{14}=p/a,\qquad q_{23}=1/(qa),\qquad q_{24}=qa/p
\]
である。$b^2p=c^2q=1$ を満たす $b,c$ を選び、
対角自己同型 $(b,b^{-1},c,c^{-1})$ で twist すれば、
新しい $p,q$ はともに一になり、$a$ は $r=ba/c$ となる。
二つの積の係数は非零なので変数のスカラー変更で $g$ を正規化できる。
これにより $Q_r$ が得られる。
符号型の混合符号の第三類もこの操作により $r^2=-1$ の $Q_r$ へ移る。

次に全ての $Q_r$ が $\HH$ に属することを示す。
$x_1$ は正則正規であり、それを逆にすると
$x_2=-x_1^{-1}x_3x_4$ を消去できるので、商は domain である。
正則性と消去は定理~\ref{thm:sign-quadric-normal} の議論に従い、
交換係数に上の $r$ を入れて確かめられる。
$r^2=1$ の場合は符号 twist で可換例になる。
$r^2\ne1$ の場合、
\[
 \Phi=\begin{pmatrix}x_1&x_3\\x_4&-r x_2\end{pmatrix},\qquad
 \Psi=\begin{pmatrix}x_2&x_3\\x_4&-r^{-1}x_1\end{pmatrix}
\]
は積が両方向とも $gI$ となる線形行列分解を与える。
可換多項式の行列式は比例しないため、二つの MCM 型は相異なる。
$C(Q_r)$ が $M_2(k)\times M_2(k)$ となることと滑らかさは、
定理~\ref{thm:sign-quadric-normal} と同じ次元の議論で従う。
両方の行列の成分は四次元を張る。
各変数の交換係数の積が一で $g$ は中央なので、
Nakayama 自己同型は $Q_r$ でも恒等である
\cite[Example 5.5 and Lemma 1.5]{RRZ2014}。
従って通常シフトが反 Serre の平方根となる。
飽和性、例外性および豊富性は正則切断の議論で確認できる。
$x_1$ で割った後さらに $x_3$ で割ると二変数の skew polynomial 環が得られ、
その tails の構造対象に無限下降列がある。
従って同じ $x_1$ の冪の議論で Gabriel--Krull 条件も従う。
以上で $\HH$ に入り、定理~\ref{thm:cubic-coverage} から cubic 型となる。

最後に同値を判定する。$r\mapsto-r$ は $g$ の直線を保つ
対角符号 twist で実現でき、$r\mapsto r^{-1}$ は $x_1,x_2$ の交換で実現できる。
必要性には命題~\ref{prop:first-orbit-finite-equivalence} の
添字付き Hom 環の同型を使う。
関係の双線形化が定める点対 $(p,p')\in\PP^3\times\PP^3$ では、
交換関係から $p,p'$ の非零座標の集合が一致する。
さらに三座標が非零なら対応する三角形の交換比の積が一でなければならない。
この積は $r^2$ または $r^{-2}$ なので、$r^2\ne1$ なら非零座標は高々二つである。
$g$ の双線形化は座標対 $12,34$ を排除するため、
各射影の像は座標対 $13,14,23,24$ の四直線からなる正方形である。
$r^2=1$ の可換類では像は滑らかな二次曲面なので、
二つの領域は添字付き線形同型で交わらない。

両方が $r^2\ne1$ の領域にあれば、その同型の各添字の射影線形変換は
四直線の交点、すなわち四座標点を置換する。
点対が各座標点を同じ座標点へ結ぶので、この置換 $\pi$ は添字に依存しない。
従って第 $i$ 添字で $x_a$ を $c_{i,a} x_{\pi(a)}$ へ送ると書ける。
隣接二添字で交換関係を送ったとき、任意の異なる $a,b$ について
$q'_{\pi(a)\pi(b)}/q_{ab}$ は
$c_{i+1,b}/c_{i,b}$ を $c_{i+1,a}/c_{i,a}$ で割った値である。
これは反対辺の対 $12,34$ についても成り立つ。
その二つの単項式は相互に比例するが零ではなく、
各交換関係の像は一方の対の単項式だけを含むからである。
よって三角形を回ると定数倍の比は相殺し、
交換比の三角積が保存される。三角積は $r^2$ と $r^{-2}$ なので
$r'^2=r^2$ または $r'^2=r^{-2}$ が必要となる。
\end{proof}

% !TEX root = ../nc-quadric-surface-classification.tex
\section{交差項を持つ Ore 環の正規二次切断}
\label{sec:ore-normal-quadrics}

$u,v$ を中央の次数一元、$q\in k^\times$、$0\ne h\in k[u,v]_2$ とし、
\[
 A=k[u,v]\langle x,y\rangle/(yx-qxy-h)
\]
を考える。これは $k[u,v,x]$ 上の Ore 拡大であり、型 $(1,4,6,4,1)$ を持つ
（命題~\ref{prop:ore-ambient-family} と同じ Ore 拡大定理を用いる）。

\begin{proposition}
\label{prop:ore-normal-quadrics}
$A$ の非零正規二次元は、スカラー倍を除いて次の形に限る。
\begin{enumerate}
\item $q=1$ の場合は $k[u,v]_2$ の元。
\item $q\ne 1,-1$ の場合は $k[u,v]_2$ の元または $xy+h/(q-1)$。
\item $q=-1$ の場合は $a x^2+b y^2+H(u,v)$ または $xy-h/2$。
\end{enumerate}
従って $q\ne-1$ のとき、これらの切断が与える $\HH$ の標識類は、
$q\ne1$ かつ $h$ が非退化である場合の可換二次曲面一つに限る。
$q=-1$ の非中央切断も同じ条件で同じ標識類を与える。
\end{proposition}

\begin{proof}
$g$ を正規二次元とし $ag=g\nu(a)$ と書く。中央元 $u,v$ は $\nu$ で固定される。
$x,y$ の全次数を先に比較する次数フィルトレーションを取ると、
その付随次数環は $yx=qxy$ の skew polynomial 環である。
その中で $x,y$ の辞書順序をどちらから始めてもよく、
積の最高単項式では指数が加算される。
$xg=g\nu(x)$ と $yg=g\nu(y)$ の最高指数を比較すると、
$\nu(x)=\lambda x+L(u,v)$、$\nu(y)=\mu y+M(u,v)$
となり、$\lambda,\mu\ne0$ である。

$q\ne1$ なら $\nu$ を Ore 関係に適用する。
$x,y$ の一次の係数から $(1-q)\lambda M=(1-q)\mu L=0$ となり、
$L=M=0$ である。また $h\ne0$ なので $\lambda\mu=1$ である。
$g$ の $x,y$ 次数が最大の層にある単項式 $x^m y^n$ を取ると、
$x,y$ との交換から $\lambda=q^{-n}$、$\mu=q^m$ となる。
従って $q^{m-n}=1$ である。
$q\ne1,-1$ なら $m+n\le2$ と合わせて $m=n$ となる。
最大層が $xy$ の場合、$x,y$ 次数二の比較で
$u,v$ を係数とする一次の項が消える。
実際、$x$ との交換で $x$ 項は消え、$y$ との交換で $y$ 項は消える。
$g=a xy+H(u,v)$ と書くと
$xg=gq^{-1}x$ は $(q-1)H=ah$ と同値である。
最大層が次数零なら $g\in k[u,v]_2$ である。

$q=-1$ の場合、最大層が $xy$ なら同じ計算を適用する。
他の最大層では $m-n$ が偶数で、平方の層または次数零しかなく、
$\lambda=\mu=1$ となる。$x^2,y^2$ は中央元であり、
残る $x,y$ の一次の項は中心に入らないので、
中央二次式は表示した形になる。

$q=1$ なら交換子は $x,y$ 次数を下げる。
最高層の比較で $\lambda=\mu=1$ となり、
$[x,g]=gL$、$[y,g]=gM$ の次数を比較すれば $L=M=0$ である。
従って $g$ は中央である。
PBW 表示で $[y,-]$ は $h\,\partial/\partial x$、
$[x,-]$ は $-h\,\partial/\partial y$ として作用する。
標数零かつ $h\ne0$ なので中心は $k[u,v]$ である。

$k[u,v]_2$ の非零元は一次式の積に分かれ、商は domain でない。
一方 $g=xy+h/(q-1)$ で割ると $xy=yx=h/(1-q)$ となり、
商は四変数可換多項式環の二次切断である。
これは $h$ が非退化である場合にちょうど滑らかで、
その場合は可換二次曲面を与える。
元の切断環の domain 性と tails の滑らかさは、
定理~\ref{thm:sign-quadric-normal} の末尾と同じく $\HH$ への所属に必要である。
従って上の標識類以外は、この範囲では現れない。
\end{proof}

\subsection{反交換 Ore 環の中央切断}
\label{subsec:ore-central-anticommuting}

\begin{theorem}
\label{thm:ore-anticommuting-central}
$q=-1$ の中央切断 $g=a x^2+b y^2+H(u,v)$ が
$\HH$ の標識類を与えるための必要十分条件は、
$a,b$ がともに非零であり、以下の二つの二次形式 $f_1,f_2$ が非退化となることである。
$x,y$ をスカラー変更して $a=b=1$ とし、その規約で $f_1=-H+h$、$f_2=-H-h$ とする。
許容される標識類は、$\PP^1$ 上の相異なる二点の非順序対を二つ取り、
対の交換と $\PGL_2$ で割ったデータで判定される
（二つの対は互いに一致してもよい）。どの類も標識付き cubic 表示を持たない。
\end{theorem}

\begin{proof}
$a=b=0$ なら商は domain でない。
一方だけ零、例えば $b=0$ なら $B^!$ で $w=X^2$ は中央正則で、
$Y^2=0$、$XY=YX$ である。$X,Y$ は $u,v$ の双対生成元と反交換する。
従って $C(B)$ 中で $YX^{-1}$ は非零中央平方零元であり、
この場合は滑らかでない。

次に $a=b=1$ とする。$c=x+y,d=x-y$ と置けば
\[
 B=k[u,v]\langle c,d\rangle/(c^2-f_1,d^2-f_2).
\]
交互に並ぶ $c,d$ の語が $k[u,v]$ 上の基底である。
書換え $c^2\mapsto f_1$、$d^2\mapsto f_2$ は長さを下げ、
三文字の同じ文字の重なりは係数の中心性で解消するからである。
元の Ore 環の関係は $c^2-d^2=f_1-f_2$ である。
双対で $CD=DC=0$、$w=C^2+D^2$ は正則正規二次元の代表である。
$C^2,D^2$ は中央なので、$C(B)$ は直交冪等元 $C^2/w,D^2/w$ により分解する。
第一成分では $C$ が対応する局所化の角で可逆となり、
$UC^{-1},VC^{-1}$ が二次形式が $f_1$ に比例する二元 Clifford 環を生成する。
第二成分では $D$ と $f_2$ を用いた同じ計算になる。
実際、$C,D$ は $U,V$ と反交換し、
$U,V$ の平方と反交換子は $f_1,f_2$ の係数を掛けた $C^2,D^2$ の和である。
従って各成分は高々四次元で、$\dim C(B)=8$ により両方が四次元である。
標数零なので二元 Clifford 環は非退化の場合にちょうど半単純であり、
退化の場合には根基のベクトルが非零平方零イデアルを生成する。
従って滑らかさは $f_1,f_2$ がともに非退化であることと同値である
\cite[Lemma 4.13 and Theorem 5.5]{MU2022}。

以後 $f_1,f_2$ を非退化とする。$B$ は domain である。
$F=k(u,v)$ に移し、中央元 $z=cd+dc$ を取ると、
交互語の基底から $F[z]$ 上で $1,c,d,cd$ が自由基底になる
（$z$ の冪の最高長さの交互語を比較する）。
さらに $F(z)$ に移し $e=d-zc/(2f_1)$ とすれば、
$c^2=f_1$、$ec=-ce$、$e^2=-(z^2-4f_1 f_2)/(4f_1)$ の四元数環となる。
$f_1,f_2$ はともに $F$ の平方でない。
$z^2-4f_1 f_2$ の任意の既約因子の素点では、$f_1$ は剰余体でも平方ではない。
$f_1 f_2$ が平方なら剰余体は $F$ であり、そうでなければ $F(\sqrt{f_1 f_2})$ である。
後者で $f_1$ が平方なら $F$ で $f_1$ または $f_2$ が平方になってしまう。
従ってこの素点における $\sqrt{f_1}$ の拡大は非分岐で剰余次数二となる。
ノルムの付値は偶数だが $e^2$ の付値は一である。
よって $e^2$ はノルムでなく、この四元数環は斜体である。
自由性により元の $B$ も domain となる。

適合根と被覆を構成する。
$f_1,f_2$ の二つの根の対を互いに交換する $\PP^1$ の対合を取れる。
根の対が一致するときは恒等を使う。
一根を共有するとき、座標で二対を $\{0,\infty\},\{0,1\}$ とすれば
$t\mapsto t/(t-1)$ が働く。四根が相異なるときは
$\{0,\infty\},\{1,\lambda\}$ とし、$t\mapsto\lambda(t-1)/(t-\lambda)$ を使う。
二乗が恒等の線形持上げ $L$ を選び、
$Lf_1=\eta f_2,Lf_2=\eta^{-1}f_1$ とする。
$\alpha^2=\eta$ を選び、$\sigma$ を $u,v$ 上で $L$、
$c\mapsto\alpha d$、$d\mapsto-\alpha^{-1}c$ とする。
$r_c=c^2-f_1,r_d=d^2-f_2$ と書けば $\sigma$ は
$r_c-\eta r_d$ の直線と $r_c+\eta r_d$ の直線を保つ。
$A'$ を中央 $u,v$ と第一の関係 $r_c-\eta r_d$ だけで定め、
第二の元 $g'=r_c+\eta r_d$ で切断すれば $B$ が得られる。
$c\pm\sqrt\eta\,d$ を新しい変数に取ると、
$A'$ は反交換 Ore 環である。型 $(1,4,6,4,1)$ を持ち、
$c^2,d^2$ が中央であるため $g'$ は中央正則である。

$A'$ の Nakayama 自己同型は $u,v$ を固定し $c,d$ に負号を掛ける。
Ore 表示の基礎環での作用は \cite[Theorem 3.3]{LWW2014} により、
残る変数の作用は Ore 関係との整合から従う。
Ore 関係の $k[u,v]$ 項が非零なら、基礎環の非中央変数に関する次数を比較して
Ore 関係を保つために平行移動項は零となり、係数は $-1$ である。
Ore 関係の $k[u,v]$ 項が零なら符号型の Nakayama 公式を使える。
従って中央切断の公式 \cite[Lemma 1.5]{RRZ2014} により、$B$ の Nakayama 自己同型も
$u,v$ を固定し $c,d$ を反転する。これは $\sigma^2$ と一致する。
よって $B^\sigma=A'^\sigma/(g')$ の通常シフトが適合根となる。
飽和性、例外性、豊富性は前の符号型の場合と同じ正則切断の議論による。

GK 条件には、$B$ が $k[u,v]$ 上自由で、
$B/(u,v)=k\langle c,d\rangle/(c^2,d^2)$ が無限次元であることを使う。
$B/uB$ 上の $v$ の冪の列は tails で真に下降するので、
$B/uB$ の像は有限長でない。
$u$ の冪の列と domain 性を使うと、前と同様に rank one torsion-free 性が従う。

$f_1=-\ell_1\ell_2$ と一次式の積に書けば、
行列 $\left(\begin{smallmatrix}c&\ell_1\\ \ell_2&-c\end{smallmatrix}\right)$ の二乗は
$r_c I$ であり、$A'$ 上で $g'$ の非零スカラー倍になる。
$f_2$ と $d$ に対しても同様である。
二型の成分空間はそれぞれ $\langle u,v,c\rangle$ と
$\langle u,v,d\rangle$ で、両方が階数三である。
twist は階数を保つので標識付き cubic 表示は存在しない。

最後に二つの許容切断を比較する。
適合環の添字付き代数は、それぞれ twist 前の $B$ の添字付き代数と
同型である。従って命題~\ref{prop:first-orbit-finite-equivalence} を使える。
上の $B$ の双線形化関係は、$(u,v)\ne(0,0)$ では
\[
 E=\{c^2=f_1(u,v),\ d^2=f_2(u,v)\}\subset\PP^3
\]
の上の対角対応を与える。$u,v$ のいずれかを使う交換関係から
次の点が比例するためである。$(u,v)=(0,0)$ では、
$c,d$ の二つの座標点が互いに移る。これらは $E$ には含まれない。
二つの多項式の根の対が異なれば $E$ は整で $\PP^3$ を張る。
実際、$\PP^1$ の関数体上で独立な二平方根を付加する被覆であり、
式は基底 $1,c,d,cd$ の有限自由被覆を与える。
根の対が一致すれば $E$ は異なる二平面にある二つの滑らかな円錐曲線であり、
二点で交わる。いずれの場合も $E$ を点ごとに固定する射影線形変換は恒等である。
整の場合は固有空間のいずれか一つに全体が入ることから、
二成分の場合は二平面上のスカラーが交点で一致することから従う。

従って添字付き同型の連続する射影線形変換は $E$ 上で一致するので、
スカラー倍まで同じである。よって同値は $B$ の次数付き同型に帰着する。
次数一の中心は $\langle u,v\rangle$ である。
$c$ の像を $a c'+b d'+\ell(u,v)$ と書くと、平方が中央変数の多項式であるため
交互語の基底から $ab=0,\ell=0$ を得る。$d$ に対しても同様である。
従って次数付き同型は $u,v$ の基底変更と、$c,d$ の交換およびスカラー変更
に尽きる。これはちょうど二つの根の対に対する主張した同値関係である。
逆に任意の二対を定める非退化な二次形式を取り、
必要なら一方をスカラー倍して $f_1-f_2\ne0$ とすれば、この Ore 族から実現できる。
\end{proof}

\begin{corollary}
多パラメータ skew polynomial の切断と、この節の Ore 環の切断を合わせると、
$\HH$ に属する標識類は、$r^2$ を逆元で割った cubic 族と、
定理~\ref{thm:ore-anticommuting-central} の非 cubic 族である。
後者は根の対が一致する一点、根を一つ共有する一点、
および四根が異なる場合の $\lambda\in k\setminus\{0,1\}$ を
$\lambda\leftrightarrow\lambda^{-1}$ で割った族からなる。
\end{corollary}
\begin{proof}
前者は定理~\ref{thm:multiparameter-quadric-classification} による。
同定の際、多パラメータ切断の非 cubic の一点は
命題~\ref{prop:nonstandard-rank-three-example} を $c,d$ の表示で書けば
両平方が同じ二次形式となる場合である。
後者の四根が異なる場合は二対を $\{0,\infty\},\{1,\lambda\}$ へ移せる。
各対内の交換は交比を逆にし、二対そのものの交換は交比を保存する。
一致と一根共有の場合は、それぞれ二点または三点に対する
$\PGL_2$ の推移性から一点となる。
\end{proof}

% !TEX root = ../nc-quadric-surface.tex
\subsection{反 Serre Hilbert 多項式による最初の軌道対象の認識}
\label{sec:first-orbit-object}

\begin{proposition}
\label{prop:first-orbit-object}
$(\cC,\OO)\in\HH$ とし、$t=S^{-1}[2]$、適合根を $s$ とする。
Gabriel--Krull rank one torsion-free な exceptional 対象 $M$ が
\[
 \chi(\OO,t^n M)=(2n+2)^2\qquad(n\gg0)
\]
を満たせば、$M\simeq s\OO$ である。
従って、どの適合根 $s'$ を用いても $s'\OO\simeq s\OO$ である。
\end{proposition}

\begin{proof}
定理~\ref{thm:normal-helix} の四連続項
$(E_0,E_1,E_2,E_3)=(\OO,\pi X_0,s\OO,\pi X_1(1))$
の類は $K_0(\cC)$ の基底である。Euler 形式の行列は
\[
 G=\begin{pmatrix}
 1&2&4&6\\
 0&1&2&4\\
 0&0&1&2\\
 0&0&0&1
 \end{pmatrix}.
\]
実際、$\Hom(\OO,\pi X_i)$ の次元は二元生成から $2$ であり、
$\Hom(\pi X_0,s\OO)$ の次元も双対評価列から $2$ である。
同じ評価列を次数一で取ると $\dim (X_1)_1=2\dim B_1-\dim(p(X_1))_0=6$ であり、
$\dim B_1=4$ である。
残る $\Hom(\pi X_0,\pi X_1(1))$ は
系~\ref{cor:normal-cubic-defect} の奇数項間の積の同型から四次元である。
他の成分は full geometric helix の Ext 消滅と例外性から従う。

$G$ は $Q=\PP^1\times\PP^1$ 上の
$\OO_Q,\OO_Q(1,0),\OO_Q(1,1),\OO_Q(2,1)$ の Euler 行列と一致する。
これらは $K_0(Q)$ の整数基底である。例えば積の標準基底
$\OO_Q,\OO_Q(1,0),\OO_Q(0,1),\OO_Q(1,1)$ から、
$[\OO_Q(2,1)]=2[\OO_Q(1,1)]-[\OO_Q(0,1)]$ によって移る。
従って上記の順序で $K_0(\cC)$ を $K_0(Q)$ と Euler 形式を保って同定できる。
Serre 関手の $K_0$ 上の作用は
$\chi(F,G)=\chi(G,SF)$ と Euler 形式の非退化性で一意に定まる。
従ってこの同定で $t$ の作用は $\OO_Q(2,2)$ をテンソルする作用である。

$H_1,H_2$ を $Q$ の二つの因子から来る超平面類とし、
$[M]$ に対応する類の Chern 指標を
$r+aH_1+bH_2+cH_1H_2$ と書く。
$Q$ 上の Riemann--Roch により
\[
 \chi(\OO,t^n M)=r(2n+1)^2+(a+b)(2n+1)+c,\qquad
 \chi(M,M)=r^2+2rc-2ab.
\]
第一式の係数比較から $r=1$、$a+b=2$、$c=1$ を得る。
例外性から第二式は $1$ なので $ab=1$、従って $a=b=1$ である。
Chern 指標の有理同型と $K_0(Q)$ の無捩れ性から $[M]=[s\OO]$ が従う。

$L=s\OO$ と置く。$\chi(L,M)=1$ より $\Hom(L,M)$ と
$\Ext^2(L,M)$ の少なくとも一方は非零である。
後者が非零なら Serre 双対性により $M\to t^{-1}L$ に非零射がある。
両対象は rank one torsion-free なので、この射は単射である。
実際、\ref{lem:sections-monic} の証明と同じく、商圏で非零となるため
商圏で同型となり、核は torsion-free 性によって零である。
従って $\Hom(\OO,t^n M)\hookrightarrow\Hom(\OO,t^{n-1}L)$ を得る。
ここで $\dim\Hom(\OO,t^n M)=\chi(\OO,t^n M)$ は $n\gg0$ で成り立つ。
定理~\ref{thm:normal-helix} の証明末尾の自由被覆の議論を偶数項に適用すれば
$\Ext^1(t^{-n}\OO,M)=0$ を得る。
また $\Ext^2(t^{-n}\OO,M)=\Hom(M,t^{-n-1}\OO)^*=0$ は、
$M$ の自由被覆と負次数切断の消滅から従う。
よって上の単射は $(2n+2)^2\le (2n)^2$ を与え、矛盾する。
従って $\Hom(L,M)$ に非零射があり、同じ理由でこれは単射である。
余核の $K_0$ 類は零である。
$B=A/(g)$ 上の有限生成加群の Hilbert 多項式は、
$A$ 上の有限自由分解により存在し、tails 上で加法的である。
非零の tails の対象は零 Hilbert 多項式を持たないので、余核は零となる。

適合根 $s'$ に対しては
$\dim\Hom(\OO,t^n s'\OO)=(2n+2)^2$ が切断環の Hilbert 関数から従い、
前段の Ext 消滅により漸近的に Euler 特性と一致する。
$s'\OO$ は $\OO$ と同じ rank one torsion-free 性と例外性を持つため、
第一の主張を適用できる。
\end{proof}

\begin{corollary}[異なる適合根における階数の保存]
\label{cor:first-orbit-rank-invariance}
同じ標識対象の二つの適合切断環について、二つの次数零生成 MCM 型の
最小線形表示行列の成分空間の次元は、二つの値の多重集合として一致する。
特に命題~\ref{prop:nonstandard-rank-three-example} の例では、
どの適合根から定理~\ref{thm:normal-helix} の方法を用いても
二つの階数はともに三である。
\end{corollary}
\begin{proof}
$s^2\simeq (s')^2\simeq t$ と命題~\ref{prop:first-orbit-object} から
$s^j\OO\simeq(s')^j\OO$ が全整数 $j$ について従う。
これらの対象の添字を保つ同型を選ぶと、合成を積とする
$\Hom(s^{-j}\OO,s^{-i}\OO)$ の $\mathbb Z$-algebra が同型になる。
各切断環 $B$ とその $\mathbb Z$-algebra の間では、
添字 $(i,j)$ 成分を $B_{j-i}$ と同定する
$a\mapsto s^{-j}(a)$ が合成を保つ。
従って得られる次数付き右加群圏の同値は
全添字の自由加群 $B(i)$ を対応する自由加群に移し、
最小自由被覆、Ext、次数零生成性を保存する。
特に自由加群の全シフトを係数に取った Ext 消滅で定まる MCM 型を移送する。
最小線形表示に現れる $B(-1)\to B$ の Hom 空間が線形同型で移るので、
その中で行列成分が張る空間の次元は保たれる。
\end{proof}

\begin{proposition}[四生成元による有限同値判定]
\label{prop:first-orbit-finite-equivalence}
二つの $\HH$ 対象の適合切断環を、
次数一の基底を選んでそれぞれ $B=T(k^4)/(K)$、$B'=T(k^4)/(K')$
とする。$G=\GL_4(k)$ とし、
\[
 E=\{(g,h)\in G\times G\mid (g\otimes h)K=K'\}
\]
を作る。$C_0=G$ から始め、射影 $\operatorname{pr}_i$ を用いて
\[
 C_{n+1}=\operatorname{pr}_1\bigl(E\cap(G\times C_n)\bigr)
         \cap\operatorname{pr}_2\bigl(E\cap(C_n\times G)\bigr)
\]
と定める。この下降列は有限段で安定し、安定集合が非空であることと
二つの対象の pointed exact 同値性は同値である。
各段は $\GL_4$ の座標環での有限の代数的消去により定まる。
機械的な計算可能性を述べる場合には、係数体が有効に与えられていることを前提とする。
\end{proposition}

\begin{proof}
命題~\ref{prop:first-orbit-object} から、pointed exact 同値は
各添字の軌道対象を保ち、系~\ref{cor:first-orbit-rank-invariance} の証明と同じ
indexed Hom の構成により $\check B\simeq\check B'$ を与える。
逆にこのような添字を保つ $\mathbb Z$-algebra 同型は、
各次数成分を各添字の加群成分と同定することにより、
次数付き右加群の圏の間の同値を誘導する。全ての自由シフトと
次数ごとのベクトル空間を保つため、有限生成性と有限次元性を保つ。
$B,B'$ は右 noetherian なので有限生成次数付き加群の圏は Abel 圏であり、
その torsion 対象は有限次元である。この部分圏を割ると exact 同値を得て、
自由加群の添字零は $\OO$ に対応する。

$B,B'$ は二次環なので、$\check B$ と $\check B'$ はそれぞれ
各隣接成分を $k^4$ とするテンソル経路環から、各二段の関係 $K$、$K'$ を課して得られる。
従って添字を保つ同型は、各隣接成分上の $g_i\in G$ により、
$(g_i,g_{i+1})\in E$ ($i\in\mathbb Z$) として必要十分に記述される。

$G\times G$ は $V^{\otimes2}$ の部分空間の Grassmannian に作用する。
$E$ は $K$ の $K'$ への transporter なので、空集合か、
閉安定化部分群 $H_E\subset G\times G$ の剰余類である。
ここで有限停止には、一般の構成可能集合の下降列ではなく
\emph{剰余類の構造が再帰で保たれる}ことが重要である。
実際、$C_n$ が非空な閉部分群 $H_n\subset G$ の剰余類なら、
$E\cap(G\times C_n)$ と $E\cap(C_n\times G)$ は
それぞれ非空な場合には
$H_E\cap(G\times H_n)$ と $H_E\cap(H_n\times G)$ の剰余類である。
それらの $\operatorname{pr}_1,\operatorname{pr}_2$ による像は、
代数群準同型による閉部分群の像の剰余類だから閉集合である。
二つの剰余類の共通部分も、非空なら閉部分群の剰余類である。
$C_0=G$ からの帰納法により、全ての $C_n$ は空集合または
閉代数部分群の剰余類となる。
従って $C_n$ は閉集合の下降列であり、$G$ の座標環の noetherian 性から
有限段で安定する。ここでは任意の閉集合の射影が閉とは主張しない。

また $g\in C_n$ となることは、$g$ を中央とする長さ $2n$ の有限列を
作れることと同値である。両側への延長可能性が再帰式を与える。
一度 $C_{n+1}=C_n$ となれば再帰式により以後は同じである。
安定集合が非空なら、その中の任意の元から左右に安定集合内の隣を選び続けられるので、
両側無限列を構成できる。逆向きには両側無限列の中央項がすべての $C_n$ に入る。

上記の剰余類に限れば射影像は閉である。
従って行列式の逆元も変数として含めた座標環で消去を行い、
各像と共通部分の根基イデアルを求めればよい。
根基の一致および全環になるかどうかを比較すると停止と非空性を判定できる。
\end{proof}

\begin{corollary}[二つの行列による pointed 同値判定]
\label{cor:two-matrix-pointed-equivalence}
前命題の仮定の下で、pointed exact 同値が存在することは
\[
 \exists\,g,h\in\GL_4(k),\qquad
 (g\otimes h)K=K',\qquad(h\otimes g)K=K'
\]
と同値である。
\end{corollary}
\begin{proof}
pointed exact 同値 $F$ を取り、両側の適合根を $s,s'$、Serre 関手を $S,S'$ とする。
$D^b(F)$ は Serre 双対を移すので、表現対象の一意性から自然同型
$D^b(F)S\simeq S'D^b(F)$ を得る。$s^{-2}\simeq S[-2]$ と
$(s')^{-2}\simeq S'[-2]$ を使い、heart に制限して
$\alpha:F s^{-2}\xrightarrow{\sim}(s')^{-2}F$ を一つ固定する。
$P_i=s^{-i}\OO,\ P'_i=(s')^{-i}\OO'$ と置く。pointing と
命題~\ref{prop:first-orbit-object} から $u_0:F P_0\simeq P'_0$ と
$u_{-1}:F P_{-1}\simeq P'_{-1}$ を選べる。
各偶奇の列において
\[
 u_{i+2}=(s')^{-2}(u_i)\circ\alpha_{P_i}
\]
で正負両方向へ延長する。逆向きにも $(s')^{-2}$ は全忠実なので一意に延長できる。

$a\in B_1$ に対し $a_i=s^{-(i+1)}(a):P_{i+1}\to P_i$ と書き、
$a'_i$ も $B'_1$ に対して同様に定める。一次成分の同型 $g_i$ は
$(g_i a)'_i=u_iF(a_i)u_{i+1}^{-1}$ で与えられる。$\alpha$ の自然性から
\[
 u_{i+2}F(a_{i+2})u_{i+3}^{-1}
   =(s')^{-2}\bigl(u_iF(a_i)u_{i+1}^{-1}\bigr)
   =(g_i a)'_{i+2}.
\]
従って $g_{i+2}=g_i$ である。右加群規約では
$a_i b_{i+1}=s^{-(i+2)}(s(a)\circ b)$ なので、
二段の関係を保つ条件は順に $(g_0\otimes g_1)K=K'$ と
$(g_1\otimes g_0)K=K'$ となる。
逆に解 $(g,h)$ から偶数辺に $g$、奇数辺に $h$ を置けば全ての二次関係を保つ。
逆行列も逆向きの関係を保つので、添字付き環の同型となり、
前命題の tails 構成で pointed exact 同値を得る。
$\alpha$ は存在証明のための選択であり、根を保つ同値を仮定する必要はない。
\end{proof}

\begin{example}[二行列を一行列に制限できない例]\label{ex:two-matrix-segre}
Segre 環 $B=k[x_{00},x_{01},x_{10},x_{11}]/(\det x)$ で、第一行の符号を変える
次数付き対合 $\gamma$ を取る。右 twist $B'=B^\gamma$ の積を
$a\star b=a\gamma^{\deg a}(b)$ とすると $K'=(1\otimes\gamma)K$ である。
$\gamma\otimes\gamma$ は $K$ を保つので $(1,\gamma)$ は上の二条件を満たす。
両辺とも適合表示である。実際、$\gamma$ を四変数多項式環へ同じ対合として
持ち上げると、twist 後も型 $(1,4,6,4,1)$ の AS 正則被覆を持ち、
$\det x$ の直線は安定なので正則正規二次切断を与える。
右加群の twist 同値は自由加群を自由加群に送り、$\gamma^2=1$ より
二重シフトと自然に可換する（右作用 $m\star b=m\gamma^{\deg m}(b)$ を比較すればよい）。
ゆえに通常シフトの平方は両側で反 Serre 移動となる。
一方、$x_{00}\star x_{10}=x_{00}x_{10}=-x_{10}\star x_{00}\ne0$ である。
従って可換な $B$ から $B'$ への次数付き環同型はなく、$g=h$ の解は存在しない。
\end{example}

% !TEX root = ../nc-quadric-classification.tex
\section{cubic 型による被覆の限界}
\label{sec:cubic-coverage}

\begin{theorem}[cubic 表示の判定]
\label{thm:cubic-coverage}
$\HH$ 対象 $(\cC,\OO)$ の任意の適合表示から
定理~\ref{thm:normal-helix} の二つの初期 MCM 型を取る。
次の条件は同値である。
\begin{enumerate}
\item $(\cC,\OO)$ はある Van den Bergh の三次三次元 AS 正則
$\mathbb Z$-algebra $D$ に対する $(\qgr D,\pi e_0D)$ と pointed exact 同値である。
\item 二つの初期型の少なくとも一方で、偶数添字間の隣接合成が階数四になる。
\end{enumerate}
特に命題~\ref{prop:nonstandard-rank-three-example} の対象は
いかなる cubic 表示でも添字零の構造対象を保って実現できない。
\end{theorem}

\begin{proof}
十分性は定理~\ref{thm:normal-helix} と系~\ref{cor:normal-cubic-defect}
から従う。必要性のため、cubic 表示があると仮定し、
$\pi e_iD$ の像を $Q_i\in\cC$ と書く。
cubic 正則性と \cite[Proposition 4.11, Theorem 6.8]{MN2025} により、
$tQ_i\simeq Q_{i-4}$、および
$\Hom(Q_j,Q_i)=D_{ij}$ である。
この適用に必要な正則性と周期性は \cite[arXiv v3, Definition 4.1.1,
Proposition 5.6.1]{VdB2011} と
\cite[Theorem 4.19]{MN2025} から従う。

$M=Q_{-2}$ と置く。二添字移動は $D$ の自己同型として実現されるので
$Q_{-2}$ は $\OO=Q_0$ の exact 自己同値による像である。
従って $M$ は rank one torsion-free exceptional である。
さらに cubic 型の成分次元から
$\Hom(\OO,t^nM)=D_{-4n-2,0}$ は $(2n+2)^2$ 次元である。命題~\ref{prop:first-orbit-object} の証明の Ext 消滅により、
$n\gg0$ では Euler 特性も同じ値を取る。
命題~\ref{prop:first-orbit-object} により $M\simeq s\OO$ となる。
従ってすべての整数 $j$ で $Q_{-2j}\simeq s^j\OO$ であり、
これらの同型は Hom と合成で定めた二つの添字付き環の同型
$D_{\mathrm{even}}\simeq\check B$ を与える（偶数添字は二で割る）。
ここでは cubic 側の二添字移動そのものが反 Serre の平方根であることは使わない。

$P_i=e_iD$ の偶数成分への制限を $P_i^{\mathrm{ev}}$ と書く。
上の同型により次数付き右 $B$-加群とみなせる。
奇数 $i$ に対する右単純加群の極小分解を偶数成分に制限すると
\[
 0\longrightarrow P_{i+4}^{\mathrm{ev}}
 \longrightarrow (P_{i+3}^{\mathrm{ev}})^2
 \longrightarrow (P_{i+1}^{\mathrm{ev}})^2
 \longrightarrow P_i^{\mathrm{ev}}\longrightarrow0
\]
となる。偶数の添字の制限は $B$ の自由シフトである。
$i$ を前後に四ずつ変えると、$P_i^{\mathrm{ev}}$ は自由加群への埋込みを
繰り返し持ち、その余核は次の自由加群へ埋め込める。
従ってこの列の二回の次元移動を反復し、
$B$ の有限単射次元を使えば $\Ext_B^j(P_i^{\mathrm{ev}},B)=0$ ($j>0$) となる。

$X=P_{-1}^{\mathrm{ev}}$ は次数零に二元生成され、
$H_X(z)=2/(1-z)^3$ であり、自由加群ではない。
定理~\ref{thm:normal-helix} の二型の記述により、非自由 MCM の
直和因子は少なくとも二元を必要とする。従って $X$ は
定理~\ref{thm:normal-helix} の二型の一方である。
$D_{-2,-1}$ は左乗法と制限により $X\to B(1)$ の射を与え、
$D_{-1,0}$ は $B\to X$ の射を与える。
それらの積は $D_{-2,0}\simeq B_1$ を張る。
実際、cubic 型では距離二に関係がなく、この隣接積は同型である。
これらは tails への移行後も同じ Hom と合成になるので、
$X$ の偶数項間の合成は階数四である。
\end{proof}

% !TEX root = ../nc-quadric-surface-classification.tex
\subsection{階数三の四周期表示の生成元と関係}
\label{subsec:rank-three-relations}

\begin{proposition}
\label{prop:rank-three-relations}
定理~\ref{thm:normal-helix} の階数三の場合に、
矢を $E_i\to E_j$ の向きで書き、道の積を合成で定める。
隣接する二元の矢に加えて、$E_{2i}\to E_{2i+2}$ ごとに一つの追加矢を取る。
最小関係は以下の三種類である。
\begin{enumerate}
\item $E_{2i}\to E_{2i+2}$ に一つ。二つの隣接矢の空間を通る合成の核で、
その二因子テンソルの階数は二である。
\item $E_{2i}\to E_{2i+3}$ に二つ。
$X_{i+1}$ の極小評価列の核の次数一部分を、
$\Hom(\pi B(i),\pi B(i+1))$ と $\Hom(\pi B(i+1),\pi X_{i+1}(i+1))$
を経由する道として表す。
\item $E_{2i+1}\to E_{2i+4}$ に二つ。
$X_i$ の双対評価列の核の最初の非零次数部分を、
$\Hom(\pi X_i(i),\pi B(i+1))$ と $\Hom(\pi B(i+1),\pi B(i+2))$
を経由する道として表す。
\end{enumerate}
各距離三の関係は、距離二の関係を先に課した後の道の空間から
一意な二次元部分空間として得られる。
\end{proposition}

\begin{proof}
四周期表示の単純加群の極小分解
（系~\ref{cor:normal-cubic-defect} の証明）により、最小関係の個数は
偶数項間の距離二で一、距離三で各二であり、他にはない。
距離二の偶数項間では階数三の合成に第四の基底を加えてあるので、
唯一の関係の追加矢における係数は零である。
この関係が階数二になることは
注意~\ref{rem:rank-three-matrix} の行列分解の議論による。

次に $E_{2i}\to E_{2i+3}$ の道を考える。
距離二の関係を先に課すと、道の空間は
$\Hom(\pi B(i),\pi B(i+1))\otimes\Hom(\pi B(i+1),\pi X_{i+1}(i+1))$
と同定できる。実際、距離三には三つの隣接矢からなる道と
追加矢を一つ含む道があり、前者に距離二の関係を適用すれば
最初の二矢の積と追加矢を合わせて $B_1$ の基底が得られる。
従って合成の核を求めるには、$X_{i+1}$ の評価列
\eqref{eq:normal-ruling-evaluation} を次数一で取ればよい。
その核は $(p(X_{i+1}))_0$ であり、二次元である。

$E_{2i+1}\to E_{2i+4}$ についても、
後二つの隣接矢の積と追加矢を合わせると $B_1$ の基底が得られる。
ここで $B$ への内部 Hom を取った評価列
\[
 0\to\Hom_B(p(X_i),B)\to B^2
       \to \Hom_B(X_i,B)(1)\to0
\]
を使うと、合成の核は左端の次数一成分である。
これも二次元であり、主張した全ての距離三の関係を与える。
\end{proof}

比較先の四周期環を逆に構成するとき、
この表示で距離二の関係を階数二に正規化してから
距離三の各二次元部分空間を変化させられる。
ここから得る候補の全次数の許容条件は、
極小分解が実際に完全であることと、周期偶数側が
命題~\ref{prop:first-orbit-object} の軌道を与えることを含む。
これらは有限次元の関係の個数からは従わない。

% !TeX root = ../nc-quadric-classification.tex
\section{階数三支持超平面の共通二変数環}
\label{sec:common-plane}

基礎体 $k$ は標数零の代数閉体とする。
次数一で生成された連結二次次数付き domain $B=T(V)/(K)$ について
\[
 \dim V=4,\qquad K\subset V\otimes V,\qquad \dim K=7
\]
を仮定する。従って $\dim B_2=9$ である。
線形行列 $M\in\operatorname{Mat}_2(V)$ の成分が張る部分空間を $U(M)$ と書く。

\begin{lemma}[行列積の四成分の独立性]
\label{lem:w-matrix-products}
$L_1,L_2\subset\operatorname{Mat}_2(k)$ を三次元部分空間とする。
trace pairing に関する $L_2^\perp$ が可逆行列で生成されるなら、
行列の乗法
\[
 L_1\otimes L_2\longrightarrow\operatorname{Mat}_2(k)
\]
は全射である。
従って $P\in\operatorname{Mat}_2(U_1)$ と $Q\in\operatorname{Mat}_2(U_2)$ の係数行列が
それぞれ $L_1,L_2$ の基底をなすとき、$PQ$ の四成分は
形式的テンソル空間 $U_1\otimes U_2$ で一次独立である。
\end{lemma}
\begin{proof}
全射でないなら、非零行列 $D$ があって
$\operatorname{tr}(DXY)=0$ が全ての $X\in L_1,Y\in L_2$ に対して成立する。
従って $DL_1\subset L_2^\perp=kC$ であり、$C$ は可逆である。
$D$ が可逆なら $\dim DL_1=3$ なので矛盾する。
$\operatorname{rank} D=1$ なら $DX$ は全て特異行列であるため $DL_1=0$ となるが、
左乗法 $X\mapsto DX$ の核は二次元であり、三次元の $L_1$ を含めない。
後半はこの全射の双対を取ったものである。
\end{proof}

\begin{lemma}[平方零行列の成分の関係]
\label{lem:w-invertible-annihilator}
$M\in\operatorname{Mat}_2(V)$ が $B$ 上で $M^2=0$ を満たし、$\dim U(M)=3$ とする。
四成分の唯一の一次関係を表す二次正方係数行列は可逆である。
従って $M^2$ の四成分を $V\otimes V$ の元として解釈すると、
それらは四次元部分空間 $R(M)\subset K\cap U(M)^{\otimes2}$ を張る。
\end{lemma}
\begin{proof}
関係の係数行列の階数が一なら、非零の定数行ベクトル $a$ と
列ベクトル $b$ によって $aMb=0$ と書ける。
$ab\ne0$ なら同時基底変換、すなわち $M$ の共役変換によって
この式を $M_{11}=0$ にできる。
このとき $M=\left(\begin{smallmatrix}0&x\\y&z\end{smallmatrix}\right)$
と書け、$M^2=0$ の第一対角成分から $xy=0$ となる。
domain 性により $x=0$ または $y=0$ で、成分空間の次元三に反する。
$ab=0$ なら同時基底変換で $M_{12}=0$ にできる。
このとき $M=\left(\begin{smallmatrix}x&0\\y&z\end{smallmatrix}\right)$
であり、二つの対角成分から $x^2=z^2=0$ となって同様に矛盾する。
関係の係数行列の転置は係数行列空間の trace-annihilator を生成するので、
補題~\ref{lem:w-matrix-products} を適用できる。
\end{proof}

\begin{lemma}[二次元部分空間にある関係]
\label{lem:w-binary-relations}
二次元部分空間 $W\subset V$ に対し、$\dim(K\cap W^{\otimes2})\le1$ である。
等号が成り立つ場合、非零の関係は二因子テンソルとして階数二である。
\end{lemma}
\begin{proof}
階数一の非零テンソル $w\otimes w'$ が $K$ にあれば、
$B$ で $ww'=0$ となって domain 性に反する。
もし $K\cap W^{\otimes2}$ に独立な二つの元 $r_0,r_1$ があれば、
$\det(sr_0+tr_1)$ は $(s,t)$ の斉次二次式である。
代数閉性により、この式が恒等的に零の場合も含めて、
非零の $(s,t)$ で行列式が零になる。
$sr_0+tr_1$ は非零なので階数一となり、矛盾する。
\end{proof}

\begin{theorem}[non-standard 二型の共通二変数環]
\label{thm:w-common-plane}
$M_1,M_2\in\operatorname{Mat}_2(V)$ が $B$ 上で $M_i^2=0$ を満たすとする。
$U_i=U(M_i)$ は三次元で、$U_1\ne U_2$ と仮定する。
$W=U_1\cap U_2$ および $R_i=R(M_i)$ と置けば、
\[
 \dim W=2,\qquad
 R_1\cap R_2=K\cap W^{\otimes2}=kr,\qquad
 \operatorname{rank} r=2,\qquad K=R_1+R_2.
\]
特に二つの平方零行列は $B$ の全ての二次関係を与える。
\end{theorem}
\begin{proof}
補題~\ref{lem:w-invertible-annihilator} により $\dim R_i=4$ である。
$R_i\subset K$ と $\dim K=7$ から $\dim(R_1\cap R_2)\ge1$ となる。
$U_1,U_2$ に適合する $V$ の基底を取れば
\[
 (U_1\otimes U_1)\cap(U_2\otimes U_2)
   =(U_1\cap U_2)\otimes(U_1\cap U_2)=W^{\otimes2}
\]
である。従って $R_1\cap R_2\subset K\cap W^{\otimes2}$ であり、
補題~\ref{lem:w-binary-relations} によって両辺は一次元に等しい。
さらに $\dim(R_1+R_2)=4+4-1=7=\dim K$ となる。
\end{proof}

\paragraph{正規切断への適用.}
系~\ref{cor:nonstandard-square-matrix} により、
正規切断 $B=A/(g)$ で $p(X)=X,p(Y)=Y$ のとき、評価行列は
$M_X^2=gC_X$, $M_Y^2=gC_Y$ ($C_X,C_Y\in\operatorname{GL}_2(k)$) と取れる。
$B$ に移せば平方零なので、双方が階数三かつ支持が異なる場合には
定理~\ref{thm:w-common-plane} が適用される。
$g$ が中央であることも $C_X,C_Y$ がスカラーであることも使用していない。

\begin{remark}[交差積の場合にも使える部分]
\label{rem:w-crossed-products}
二つの行列 $P,Q$ の支持が異なる三次元空間 $U_1,U_2$ で、
それぞれの係数行列空間の trace-annihilator が可逆行列で生成され、
$PQ=QP=0$ が $B$ で成立する場合も同じ結論が成り立つ。
この場合、四次元の関係空間は
$R_{12}\subset U_1\otimes U_2$ と $R_{21}\subset U_2\otimes U_1$ にあり、
それらの交わりが $W^{\otimes2}$ に含まれることを使う。
中央切断の通常の行列分解 $PQ=QP=gI$ はこの状況の候補となるが、
一般の正規切断では一方の積に正規化自己同型が入る。
後者への適用には、twist された支持の一致を別に確定する必要がある。
\end{remark}

\subsection{全次数での量子平面・Jordan 平面の埋込み}

\begin{proposition}
\label{prop:w-embedded-plane}
$W\subset B_1$ が二次元で $K\cap W^{\otimes2}=kr\ne0$ とする。
このとき $W$ の基底 $u,v$ を適切に選ぶと、$W$ が生成する $B$ の部分環は
次のいずれかと同型である。
\[
 Q_q=k\langle u,v\rangle/(vu-quv),\quad q\in k^\times,
 \qquad
 Q_J=k\langle u,v\rangle/(vu-uv-u^2).
\]
これは二次関係だけの同定でなく、自然な写像 $Q_q\to B$ または $Q_J\to B$
の全次数での単射を主張する。
\end{proposition}
\begin{proof}
補題~\ref{lem:w-binary-relations} により $r$ の係数行列は可逆である。
$r$ を対称部分と反対称部分に分ける。
対称部分の階数が二なら、代数閉体上の二変数非退化対称形式を
$uv+vu$ に正規化できる。従って $r=\alpha uv+\beta vu$ と書け、
係数行列の可逆性から $\alpha\beta\ne0$ である。これは $Q_q$ の関係である。
対称部分が零なら $r$ は $vu-uv$ の定数倍である。
対称部分の階数が一なら、$r=\alpha u^2+\beta(vu-uv)$ と書ける。
可逆性から $\beta\ne0$ なので、基底を定数倍して Jordan 平面の関係となる。

次に単射性を証明する。いずれの環 $Q$ も $k[u]$ の Ore 拡大で、
単項式 $u^iv^j$ ($i,j\ge0$) が基底をなし、domain である。
$u$ は正規なので可逆化できる。
$z=u^{-1}v\in Q[u^{-1}]$ と置けば、
\[
 zu=u\tau(z),\qquad
 \tau(z)=qz\ \ (Q=Q_q),\qquad
 \tau(z)=z+1\ \ (Q=Q_J)
\]
となる。従って次数 $n$ の任意の非零斉次元 $f$ は一意に
\[
 f=u^n p(z),\qquad 0\ne p\in k[z],\qquad \deg p\le n
\]
と書ける。実際、$u^{n-j}v^j$ をこの形に直すと、
$p$ は次数 $j$ の非零先頭係数を持つ多項式になり、三角的に基底を与える。
代数閉性から $p(z)=a\prod_{i=1}^d(z-\lambda_i)$ と書く。
恒等式
\[
 \prod_{i=1}^{d}\bigl(u f_i(z)\bigr)
 =u^d\prod_{i=1}^{d}\tau^{d-i}\bigl(f_i(z)\bigr)
\]
に $f_i(z)=\tau^{-(d-i)}(z-\lambda_i)$ を入れれば、
$u^dp(z)$ は $Q_1$ の非零元 $d$ 個の積の定数倍となる。
各 $f_i$ が一次式なので $u f_i(z)\in ku+kv$ である。
残る $u^{n-d}$ も合わせると、$f$ は非零一次元 $n$ 個の積である。

自然な全射 $Q\to k\langle W\rangle\subset B$ は次数一で単射である。
その核が非零なら非零斉次元 $f$ を取り、上の一次因子分解を写すと、
$B$ で非零一次元の積が零になる。これは domain 性に反する。
\end{proof}

\begin{remark}[ここからはまだ正規化できない]
$Q_q$ の $u,v$ は $Q_q$ の中で正規であり、$Q_J$ の $u$ は $Q_J$ の中で正規である。
しかし $B$ 内での正規性は、この埋込みの命題からは従わない。
また $q\ne1$ の量子平面や Jordan 平面の場合、$W$ の線形基底変更だけで
二つの生成元を $B$ の中心に入れることはできない。
そのような基底があれば $W$ の全ての元が互いに交換するが、
命題の全次数単射に反するからである。
Zhang twist や pointed 同値によって別の表示へ移る可能性は残る。
\end{remark}

\subsection{同一支持と中央化についての障害}

\begin{lemma}[同一支持の場合に確定すること]
\label{lem:w-equal-support}
三次元部分空間 $U\subset B_1$ について
$\dim(K\cap U^{\otimes2})\le4$ である。
従って定理~\ref{thm:w-common-plane} の平方零行列で $U_1=U_2=U$ なら、
\[
 R_1=R_2=K\cap U^{\otimes2}.
\]
\end{lemma}
\begin{proof}
$\mathbb P(U)\times\mathbb P(U)$ の Segre 像は
$\mathbb P(U^{\otimes2})=\mathbb P^8$ 内の四次元射影多様体である。
$K\cap U^{\otimes2}$ が五次元以上なら、その射影化の次元は四以上であり、
射影次元定理によりこの Segre 像と交わる。
この交点は非零の階数一テンソルを与え、domain 性に反する。
後半は各 $R_i$ が四次元であることから従う。
\end{proof}

\paragraph{残る点.}
$R_1=R_2$ から二つの余核が同型になるかは、この次元論だけでは分からない。
従って二つの異なる MCM 型で支持が異なるという仮定を外してはいけない。
次の具体的課題は、三次元の同じ成分空間と同じ四次元の平方関係空間を持つ
二つの行列について、定数行列による左右同値を証明するか、反例を得ることである。

\paragraph{中央化は適合表示の中ですら自動でない.}
命題~\ref{prop:nonstandard-rank-three-example} で使う
\[
 B_0=k[u,v]\langle c,d\rangle/(c^2-uv,d^2-uv)
\]
を考える。この表示では $u,v$ は中央である。
$t\in k^\times$ に対して
\[
 \sigma_t(u)=t^2u,\quad \sigma_t(v)=v,\quad
 \sigma_t(c)=tc,\quad\sigma_t(d)=td
\]
は次数付き自己同型である。
$a\star b=a\sigma_t^{\deg a}(b)$ の規約の Zhang twist を $B_t$ と書くと、
関係は
\[
\begin{gathered}
 vu=t^2uv,\qquad cu=tuc,\qquad cv=t^{-1}vc,\\
 du=tud,\qquad dv=t^{-1}vd,\qquad
 c^2=tuv,\qquad d^2=tuv.
\end{gathered}
\]
ここでは積記号を再び通常の並置で書いた。
行列
\[
 M_r=\begin{pmatrix}r&u\\-tv&-tr\end{pmatrix},\qquad r=c,d,
\]
はそれぞれ $M_r^2=0$ を満たす。
この二つの支持は $\langle u,v,c\rangle$, $\langle u,v,d\rangle$ であり、
交わりはパラメータ $q=t^2$ の量子平面を生成する。
$t^2\ne1$ なら、同じ $B_t$ の中で $W$ の基底を変えて中央化することは不可能である。

特に $t=\mathrm i$ は命題~\ref{prop:nonstandard-rank-three-example} の反 Serre 平方根に適合した表示そのものであり、
このとき $vu=-uv$ である。
従って「$H_2$ の適合表示なら $W$ を線形基底変更だけで中央化できる」
という主張は、既存例によって既に否定される。
一方 $u,v$ はこの例では $B_t$ 全体でも正規なので、
一般の $W$ の正規化という課題の反例にはなっていない。
任意の $t$ について次数シフト自体が反 Serre の平方根になるとは主張していない。

% !TeX root = ../nc-quadric-surface-classification.tex
\section{standard 型の帰着と非 cubic 部分の完全分類}
\label{sec:standard-rank-three}

\begin{proposition}[standard 二型の両側零積]
\label{prop:standard-mutual-zero-matrices}
$B=A/(g)$ を $\HH$ の適合表示とし、二つの非自由 MCM 型を $X,Y$ とする。
$\Omega X(1)\simeq Y$, $\Omega Y(1)\simeq X$ の場合、
両者の極小線形表示行列 $M,N\in M_2(B_1)$ は
\[
 MN=NM=0
\]
を同時に満たすように取れる。
両合成階数が三なら、両行列の成分間の唯一の一次関係の係数行列は可逆である。
\end{proposition}
\begin{proof}
評価列を
\[
 0\longrightarrow Y(-1)\xrightarrow{i_Y}F_X
 \xrightarrow{q_X}X\longrightarrow0,
 \qquad
 0\longrightarrow X(-1)\xrightarrow{i_X}F_Y
 \xrightarrow{q_Y}Y\longrightarrow0
\]
と取る。ここで $F_X,F_Y$ はそれぞれ二元の次数零自由加群である。
$M=i_Yq_Y(-1):F_Y(-1)\to F_X$、
$N=i_Xq_X(-1):F_X(-1)\to F_Y$ と定めれば、
$q_Xi_Y=q_Yi_X=0$ より、適切なシフトを付した二つの積は零である。
また $q_X,q_Y$ は全射なので、これらの余核はそれぞれ $X,Y$ である。
従ってこれは極小線形表示を与える。

仮に $M$ の唯一の一次関係の係数行列の階数が一なら、
$M'=PMQ$、$N'=Q^{-1}NP^{-1}$ と同時に基底変更し、
$M'_{11}=0$ とできる。この変更後も $M'N'=N'M'=0$ である。
$M'=\left(\begin{smallmatrix}0&a\\b&c\end{smallmatrix}\right)$
と書くと、三次元の成分空間より $a\ne0$ である。
$M'N'=0$ の第一行と domain 性から $N'$ の第二行が零となり、
$N'$ の成分空間が三次元であることに反する。
$N$ の関係も同様である。
\end{proof}

\begin{remark}[正規切断での持上げの範囲]
$A_1=B_1$ による一意な持上げを同じ $M,N$ で表す。
各行列は線形行列分解の最初の行列であるから $A$ 上で単射であり、
\[
 MN=gC_X,\qquad NM=gC_Y,\qquad C_X,C_Y\in\mathrm{GL}_2(k)
\]
となる。定数行列が特異なら積が非零の定数列を殺すので単射性に反する。
$ag=g\nu(a)$ の規約で結合則を使うと
\[
 \nu(M)=C_XMC_Y^{-1},\qquad
 \nu(N)=C_YNC_X^{-1}
\]
を得る。一般の正規切断で二つの定数行列を同時に恒等へ直したとは主張しない。
以下の二次関係の議論は $B$ 上の両側零積のみを使うため、この障害を受けない。
\end{remark}

\begin{proposition}[standard 型の支持が異なる場合]
\label{prop:standard-crossed-common-plane}
命題~\ref{prop:standard-mutual-zero-matrices} の両階数三の場合で、
$U=U(M)$、$U'=U(N)$ が異なるとする。
$R_{12}$ を形式的な $MN$ の四成分の張る二次関係空間、
$R_{21}$ を $NM$ の同様の空間と置く。
このとき $W=U\cap U'$ は二次元であり、
\[
 \dim R_{12}=\dim R_{21}=4,\qquad
 R_{12}\cap R_{21}=K\cap W^{\otimes2}=kr,\qquad
 \operatorname{rank}r=2,\qquad K=R_{12}+R_{21}
\]
が成り立つ。$W$ が生成する部分環は全次数で量子平面または Jordan 平面である。
\end{proposition}
\begin{proof}
前命題の非退化な一次関係と行列乗法の全射性
（補題~\ref{lem:w-matrix-products}）により、各 $R$ は四次元である。
また $R_{12}\subset U\otimes U'$、$R_{21}\subset U'\otimes U$ であり、
両テンソル空間の交わりは $W^{\otimes2}$ である。
$\dim K=7$ と、domain の二次元部分空間には高々一つの非退化二次関係しか
存在しないことから、主張する次元と等式が従う。
全次数の単射は命題~\ref{prop:w-embedded-plane} による。
\end{proof}

\begin{lemma}[一つの支持を持つ零積行列は conic の twist を与える]
\label{lem:standard-one-support-conic}
$B$ を連結次数付き domain とし、$M,N\in M_2(B_1)$ が $MN=0$ を満たすとする。
$U(M)=U(N)=U$ が三次元であると仮定する。
このとき $U$ が生成する部分環は可換滑らか conic の
次数付き自己同型による Zhang twist である。
また $M,N$ は定数行列による左右変換で互いに移り、
$B$ 上の次数付き余核は同型である。
この命題には $NM=0$ も $\dim B_1=4$ も必要ない。
\end{lemma}
\begin{proof}
まず、各行列の四成分の唯一の一次関係は可逆な係数行列を持つ。
実際、$M$ の関係の係数行列の階数が一なら、
$M\mapsto PMQ$, $N\mapsto Q^{-1}NP^{-1}$ の同時基底変更で
$M_{11}=0$ にできる。
$M=\left(\begin{smallmatrix}0&a\\b&c\end{smallmatrix}\right)$ と書くと
$a\ne0$ であり、$MN=0$ の第一行と domain 性から $N$ の第二行が零となる。
これは $N$ の成分空間が三次元であることに反する。
$N$ の関係の場合は $N_{11}=0$ とし、積の第一列を使えば同様である。

それぞれの係数行列空間の trace-annihilator を用いて
\[
 M=M_0G,\qquad N=N_0H,\qquad G,H\in\mathrm{GL}_2(k),
 \qquad \operatorname{tr}M_0=\operatorname{tr}N_0=0
\]
と書く。
$M_0=\left(\begin{smallmatrix}a&b\\c&-a\end{smallmatrix}\right)$ とすると
$a,b,c$ は $U$ の基底である。
$\psi:U\to U$ を
\[
 \psi(M_0)=G N_0G^{-1}
\]
で定める。右辺も三次元の traceless 成分空間を持つので、$\psi$ は可逆である。
形式的テンソル積として
\[
 MN=M_0\psi(M_0)GH
\]
であるから、$MN$ の四成分の張る関係空間は
\[
 R=(1\otimes\psi)K_0,\qquad
 K_0=\bigwedge^2U+kq,\qquad q=a^2+bc
\]
である。ここでは $q$ を対称テンソルとして代表させている。
$K_0$ は可換 conic $C_0=k[a,b,c]/(a^2+bc)$ の四次元二次関係空間である。

$\sigma=\psi^{-1}\in\mathrm{GL}(U)$ を三変数可換多項式環
$P=\operatorname{Sym}(U)$ の次数付き自己同型へ延長する。
この段階では $\sigma$ が $q$ を保つとは仮定しない。
$A=P^\sigma$ と置けば、$A$ の三つの二次交換関係は
$(1\otimes\psi)\bigwedge^2U$ であるから、$A\to B$ の準同型がある。
underlying vector space $A_2=P_2$ の中で、残る関係はちょうど元 $q$ である。
domain の三次元部分空間の二次関係は高々四次元である
（補題~\ref{lem:w-equal-support} の Segre 次元論）。
従ってこの準同型の次数二の核はちょうど $kq$ である。

$\sigma$ が conic を保つことを、domain 性から導く。
任意の非零 $u\in U$ に対し、$A=P^\sigma$ の積の規約から
\[
 u\star q=u\,\sigma(q)
 =\sigma(q)\,u
 =\sigma(q)\star\sigma^{-2}(u).
\]
中央の積は underlying 可換多項式環 $P$ の積である。
$B$ では $q=0$ であり、$\sigma^{-2}(u)$ の像は非零の次数一元である。
従って domain 性により $\sigma(q)$ も零に写る。
次数二の核が $kq$ なので
\[
 \sigma(q)\in k^\times q
\]
となる。よって $\sigma$ は $C_0$ の自己同型に降り、
$A/(q)=C_0^\sigma$ から $B$ への写像を得る。

この写像は次数一で単射であり、全次数でも単射である。
実際、$C_0=k[s,t]^{(2)}$ の全斉次元は代数閉性により次数一元の積に分解でき、
$C_0^\sigma$ でも因子に $\sigma$ の逆冪を掛ければ同様である。
核の非零斉次元は $B$ における非零一次元の零積を与えてしまう。
従って $k\langle U\rangle=C_0^\sigma$ である。

最後に余核を比較する。
$MN=0$ を可換環 $C_0$ の積へ戻すと $M\sigma(N)=0$ である。
$M$ も $\sigma(N)$ も非零なので、商体上では両方の階数が一であり、
従って $\det M=\det N=0$ となる。
$a=st,b=s^2,c=-t^2$ と書くと、両行列は二変数二次式からなる階数一行列であり、
成分は全ての二次式の三次元空間を張る。
一意分解性により各行列は一次式の二成分列と一次式の二成分行の積に分解する。
それぞれの成分は独立である。
具体的には、成分の共通因子は $s^2,t^2$ を張ることから定数であり、
列または行の次数が零なら成分空間が高々二次元になるので、両方の次数は一である。
従って各行列は
\[
 \begin{pmatrix}st&s^2\\-t^2&-st\end{pmatrix}
 =\begin{pmatrix}s\\-t\end{pmatrix}\begin{pmatrix}t&s\end{pmatrix}
\]
と定数行列による左右変換で同値である。
この一次元の等式は $B$ 上でも成立し、余核の同型を与える。
\end{proof}

\begin{corollary}[standard 二型の支持も異なる]
\label{cor:standard-supports-distinct}
$\HH$ の standard 適合表示で、二つの MCM 型の合成階数がともに三なら、
それらの支持超平面は異なる。
\end{corollary}
\begin{proof}
評価列から得た行列 $M,N$ は $MN=NM=0$ を満たし、
余核は異なる二つの MCM 型である。
同じ支持を持てば補題~\ref{lem:standard-one-support-conic} が余核の同型を与えて矛盾する。
\end{proof}

\begin{theorem}[standard 型の異なる支持の同時 untwist]
\label{thm:standard-distinct-untwist}
$B=T(V)/(K)$ を noetherian 連結二次次数付き domain とし、
$\dim_kB_n=(n+1)^2$ とする。
三次元の異なる成分空間 $U,U'\subset V$ を持つ行列
$M,N\in M_2(B_1)$ が $MN=NM=0$ を満たすと仮定する。
このとき $B$ の次数付き自己同型 $\tau$ が存在して、
$B^\tau$ は中央二重平方環
\[
 k[u,v]\langle z_1,z_2\rangle/
 (z_1^2-f_1(u,v),\ z_2^2-f_2(u,v))
\]
と同型になる。ここで $u,v$ は中央、$f_1,f_2$ は非退化二次形式である。
従って standard な両階数三表示で支持が異なる場合も、
pointed exact 同値の意味で定理~\ref{thm:ore-anticommuting-central} の族に帰着する。
\end{theorem}
\begin{proof}
命題~\ref{prop:standard-mutual-zero-matrices} と同じ議論で、
両係数行列空間の trace-annihilator は可逆である。
それらを $kD_M,kD_N$ と書く。
$H^2=D_MD_N$ となる可逆行列 $H$ を取り、$Q=HD_N^{-1}$ と置けば
$QD_NQ=D_M$ である。
$M\mapsto MQ$, $N\mapsto Q^{-1}N$ により、二つの trace-annihilator を一致させる。
この同時変更は両側零積も各支持も保つ。

共通の一次関係から、成分の対応が線形同型
$\varphi:U\to U'$、$\varphi(M)=N$ を定める。
共通の係数 hyperplane の trace-annihilator を $kD$ とし、
$G=D^{-1}$ と置くと
\[
 M=M_0G,\qquad
 M_0=\begin{pmatrix}a&b\\c&-a\end{pmatrix},\qquad
 C_0=k[a,b,c]/(a^2+bc)
\]
と書ける。
$\sigma_0(M_0)=G^{-1}M_0G$ は $C_0$ の次数付き自己同型を定める。
$K_0\subset U^{\otimes2}$ を $C_0$ の二次関係空間とする。
形式的な $M^2$ の四成分は
$R_0=(1\otimes\sigma_0^{-1})K_0$ を張る。
これは $C_0^{\sigma_0}$ における
$M\star M=M\sigma_0(M)=M_0^2G^2=0$ と、
四成分の独立性から従う。
ここで $M^2=0$ を $B$ で仮定しているわけではない。

従って両側零積の二次関係は
\begin{equation}
 R_{12}=(1\otimes\varphi\sigma_0^{-1})K_0,
 \qquad
 R_{21}=(\varphi\otimes\sigma_0^{-1})K_0.
 \label{eq:standard-cross-conic-kernels}
\end{equation}
$W=U\cap U'$ と置くと、命題~\ref{prop:standard-crossed-common-plane} の次元論より
$K=R_{12}+R_{21}$ および
$R_{12}\cap R_{21}=K\cap W^{\otimes2}=kr$ であり、$r$ は階数二である。
$Q_W=k\langle W\rangle$ は全次数で量子平面または Jordan 平面である。

可換 conic を $C_0=k[s,t]^{(2)}$ と書く。
まず $W\subset U=k[s,t]_2$ が基点を持たないことを示す。
仮に $W=h\,k[s,t]_1$ とする。
$r$ を \eqref{eq:standard-cross-conic-kernels} の第一の kernel で解釈すると、
\[
 \ker\left(W\otimes\sigma_0\varphi^{-1}(W)
       \longrightarrow k[s,t]_4\right)\ne0
\]
である。$r$ の階数二より第二因子の空間も
$\sigma_0\varphi^{-1}(W)=h'k[s,t]_1$ と書ける。
実際、$hs\,A+ht\,B=0$ は $A=t\ell$, $B=-s\ell$ を与え、
$A,B$ は独立なので第二空間全体が $\ell k[s,t]_1$ となる。

$\sigma_0\varphi^{-1}(q)=hh'$ となる非零元 $q\in W$ を取る。
任意の $z\in U=k[s,t]_2$ について
\[
 z\cdot\sigma_0\varphi^{-1}(q)=hh'z
 \in W\cdot\sigma_0\varphi^{-1}(W)
\]
である。
第一の $R_{12}$ は写像
$U\otimes U'\to k[s,t]_4$, $z\otimes z'\mapsto z\sigma_0\varphi^{-1}(z')$
の全 kernel なので、$B$ において $zq\in(Q_W)_2$ となる。
従って $U\subset\operatorname{Frac}Q_W$ である。

第二の $R_{21}$ において、$W$ の二つの像は
$\varphi^{-1}(W)$ と $\sigma_0(W)$ であり、両方とも共通一次因子を持つ pencil である。
同じ議論を、写像
$U'\otimes U\to k[s,t]_4$, $z'\otimes z\mapsto\varphi^{-1}(z')\sigma_0(z)$
へ適用すると、適切な非零 $q'\in W$ に対し
$z'q'\in(Q_W)_2$ ($z'\in U'$) を得る。
よって $U'\subset\operatorname{Frac}Q_W$ でもある。
これらの分数は全て次数一である。

以上の分数を $B$ の商斜体の中で取ると、$B$ は $Q_W$ の斉次商環の
次数付き部分環になる。
量子平面および Jordan 平面の斉次商環は
$k(\mathbb P^1)[T,T^{-1};\psi]$、$\psi\in\operatorname{Aut}\mathbb P^1$
の形である。
補題~\ref{lem:w-rational-growth} より $\dim_kB_n=O(n)$ となり、
仮定した $(n+1)^2$ と矛盾する。
従って $W$ は基点を持たない。

$W$ の互いに素な二元を $w_1,w_2$ とする。
$r$ が与える $w_1A+w_2B=0$ に一意分解性を使うと、
$A=\lambda w_2$, $B=-\lambda w_1$ となる。
従って
\[
 \sigma_0\varphi^{-1}(W)=W.
\]
特に $\varphi^{-1}(W)$ も基点を持たず、
第二の $R_{21}$ に同じ議論を適用して
$\sigma_0(W)=\varphi^{-1}(W)$ を得る。
そこで
\[
 T_1=\varphi\sigma_0^{-1}:U\longrightarrow U',\qquad
 T_2=\sigma_0^{-1}\varphi^{-1}:U'\longrightarrow U
\]
はいずれも $W$ を保つ。
$K_1=(\varphi\otimes\varphi)K_0$ と書けば
$R_{12}=(1\otimes T_1)K_0$、$R_{21}=(1\otimes T_2)K_1$ である。
共通の関係は
\[
 kr=(1\otimes T_1)(\bigwedge^2W)
    =(1\otimes T_2)(\bigwedge^2W)
\]
なので $T_1|_W$ と $T_2|_W$ は比例する。
$N$ を定数倍すると $\varphi$ が同じ倍率で変わり、
$T_1,T_2$ は互いに逆の倍率で変わる。
基礎体上で平方根を取ることにより、両制限を一致させられる。

$\tau|_U=T_1$, $\tau|_{U'}=T_2$ として $V$ 上へ貼り合わせる。
これは $U/W$ と $U'/W$ を交換する可逆線形写像である。
$\sigma_0$ が $C_0$ の自己同型であることを使うと
\begin{align*}
 (\tau\otimes\tau)R_{12}
  &=(\varphi\sigma_0^{-1}\otimes\sigma_0^{-2})K_0
    =(\varphi\otimes\sigma_0^{-1})K_0=R_{21},\\
 (\tau\otimes\tau)R_{21}
  &=(\sigma_0^{-1}\otimes\varphi\sigma_0^{-2})K_0
    =(1\otimes\varphi\sigma_0^{-1})K_0=R_{12}.
\end{align*}
従って $\tau$ は $B$ の次数付き自己同型である。

$B^\tau$ の二次関係は $(1\otimes\tau^{-1})K$ なので、
\[
 (1\otimes\tau^{-1})R_{12}=K_0,\qquad
 (1\otimes\tau^{-1})R_{21}=K_1
\]
となる。よって二つの支持の部分環は同時に可換 conic になり、$W$ は全環で中央となる。
定理~\ref{thm:w-simultaneous-untwist} の証明末尾の平方完成の議論により、
非退化な $f_1,f_2$ を持つ中央二重平方表示を得る。
Zhang 同値は構造対象を保つ。
\end{proof}

\begin{theorem}[非 cubic 部分の完全分類]
\label{thm:noncubic-classification}
$\HH$ のうち pointed cubic 表示を持たない対象の pointed exact 同値類は、
$\PP^1$ 上の二つの被約な二点集合の非順序対を、
$\PGL_2$ の同時作用で割った集合と一対一に対応する。
具体的には、二対が一致する一点、一点を共有する一点、
および四点が異なる場合の
\[
 \lambda\in k\setminus\{0,1\},\qquad \lambda\sim\lambda^{-1}
\]
からなる。各対象は、二点集合を零点とする非退化二次形式 $f_1,f_2$ を用いた
中央二重平方環の pointed 圏として実現される。
non-standard な適合表示を持つものは、二対が一致する一点だけである。
\end{theorem}
\begin{proof}
定理~\ref{thm:cubic-coverage} により、cubic 表示を持たないことは
二つの MCM 型の合成階数がともに三であることと同値である。
その二型に対する $p=\Omega(1)$ は恒等か交換かのいずれかである。
恒等の場合は定理~\ref{thm:w-simultaneous-untwist}、交換の場合は
命題~\ref{prop:standard-mutual-zero-matrices}、
系~\ref{cor:standard-supports-distinct}、
定理~\ref{thm:standard-distinct-untwist} により、中央二重平方環へ帰着する。
これらの Zhang 同値は構造対象を保つ。
逆の構成、全ての $\HH$ 条件と cubic 被覆不能性、
二つの根対による必要十分な同値判定は
定理~\ref{thm:ore-anticommuting-central} とその直後の系による。
non-standard な適合表示についての最後の主張は
定理~\ref{thm:w-compatible-singleton} である。
\end{proof}

% !TEX root = ../nc-quadric-surface.tex
\section{重み付き四周期表示の偶数環}
\label{sec:weighted-inverse}

以下、$D$ を局所有限で正に次数付けされた連結 $\mathbb Z$-algebra とし、
$P_j=e_jD$ と書く。$D$ は左右 noetherian とし、右単純加群 $S_j$ は
\[
 0\longrightarrow P_{j+4}\longrightarrow
 P_{j+3}^2\oplus P_{j+2}^{d_j}\longrightarrow
 P_{j+1}^2\oplus P_{j+2}^{d_j}\longrightarrow
 P_j\longrightarrow S_j\longrightarrow0
 \tag{*}\label{eq:weighted-inverse-rank-resolution}
\]
という極小分解を持つとする。ここで $d_j=0$（$j$ 奇数）、
$d_j\in\{0,1\}$（$j$ 偶数）、$d_{j+2}=d_j$ とする。
この節ではこの完全性を仮定して逆向きに何が従うかを調べる。

\begin{proposition}[周期偶数環の Koszul 性]
\label{prop:weighted-even-koszul}
$D$ の偶数部分を二添字の周期移動で次数付けして得られる環を $B$ とする。
すると $B$ は左右 noetherian、一次生成、二次 Koszul であり、
$H_B(z)=(1+z)/(1-z)^3$ である。
\end{proposition}

\begin{proof}
極小分解の第一項から、$D$ は隣接する二つの生成元と
偶数添字から始まる距離二の生成元で生成される。
偶数添字間の道は長さ一か二で次の偶数添字を通るから、
$B$ は一次生成である。偶数部分の斉次右イデアルを
$P_j$（$j$ 偶数）の部分加群として $D$ へ拡張し、
偶数成分に制限すると元のイデアルに戻る。
実際、拡張で最終添字が偶数になる項に使われる積は
両端とも偶数の成分だけである。よって斉次右イデアルの昇鎖は停止する。
左側も同様であり、一般のイデアルには最高次の初期項を用いる。

$h_{jk}=\dim D_{jk}$ と書く。極小分解の各成分で Euler 標数を取り、
同じ添字に現れる距離二の項を消去すると、
\[
 h_{jk}-2h_{j+1,k}+2h_{j+3,k}-h_{j+4,k}=\delta_{jk}.
\]
$h_{jk}=0$（$j>k$）から順に解けば、$h_{jk}$ は $k-j$ のみの関数となり、
その母関数は $(1-2z+2z^3-z^4)^{-1}$ である。
偶数差 $2n$ の係数は $(n+1)^2$ になる。

奇数 $i$ に対し、$N_i=P_i|_{\mathrm{even}}$ と置く。
式 \eqref{eq:weighted-inverse-rank-resolution} を偶数成分に制限すると
$S_i$ は零になるため、
\[
 0\to N_{i+4}\to (P_{i+3}|_{\mathrm{even}})^2
   \to(P_{i+1}|_{\mathrm{even}})^2\to N_i\to0
\]
を得る。$i$ を四ずつ増やしてこの列をつなぐと、
$N_i$ の自由分解の各項は、最初の生成次数から一つずつ増加する次数で生成される。
特に $N_1$ の左端の生成次数は $1$、$N_3$ では $2$ である
（偶数添字 $2j$ を次数 $j$ と数える）。

次に $j=0$ の式 \eqref{eq:weighted-inverse-rank-resolution} を
偶数成分に制限し、$k$ の分解を得る。
複体の第 $p=0,1,2,3$ 項は順に
$B$、$N_1^2\oplus B(-1)^{d_0}$、
$N_3^2\oplus B(-1)^{d_0}$、$B(-2)$ である。
この複体に $-\otimes_B^{\mathbf L} k$ を適用した
スペクトル系列の第一頁では、各第 $p$ 項の
$\operatorname{Tor}_q^B(-,k)$ は、上で得た線形分解により
内部次数が $p+q$ 以下にある。
従って $\operatorname{Tor}_n^B(k,k)$ は内部次数 $n$ 以下にある。
一方、$B$ は連結一次生成なので、極小分解から同じ Tor は内部次数 $n$ 以上にある。
よって $B$ は Koszul であり、特に二次である。
\end{proof}

\begin{proposition}[残る許容条件への還元]
\label{prop:weighted-inverse-reduction}
命題~\ref{prop:weighted-even-koszul} の仮定に加え、
$D$ を四周期、次元三・parameter 四の AS 正則環とする。
偶数部分の次数付けを与える周期移動の二乗は、
$D$ の逆 Nakayama 同型を偶数部分に制限したものと
対角型自己同型を除いて一致するとする。
このとき、次の条件が成立すれば $(\qgr D,\pi e_0D)\in\HH$ である。
\begin{enumerate}
\item $B$ は domain である。
\item ある正次数の非零斉次元 $f\in B$ に対して
$\pi(B/fB)$ は $\qgr B$ で有限長でない。
\item 定理~\ref{thm:finite-ambient-test} の有限階数条件を満たす
超平面がある。
\end{enumerate}
これらは、この周期移動を適合根とする逆構成に必要でもある。
\end{proposition}

\begin{proof}
偶数部分への制限により $\qgr D\simeq\qgr B$ となる。
ここでは制限は exact であり、左随伴は偶数頂点の表現可能加群を
対応する $P_{2j}$ に送る加群の延長である。
この随伴の単位は偶数成分で同型であり、
余単位の核と余核は偶数成分が零となる。
偶数成分が零の有限生成加群は、奇数始点の任意の正距離の道が
最初に偶数頂点を通るため、有限次元である。
また有限生成加群の偶数部分も有限生成であり、
奇数頂点に置かれた生成元については、次の偶数頂点への二本の矢でその制限を生成できる。
このためこれらの操作は tails において互いに逆である。
添字零の表現可能加群は $\pi B$ に移る。

\cite[Theorems 4.19, 6.8 and 6.11, Proposition 4.11]{MN2025} から、
$D$ 側の標識は exceptional で、四添字の Nakayama 移動が
反 Serre 関手となる。仮定した対角型の差は加群の各成分を定数倍すれば
自然同型になるので、偶数部分のシフトはその平方根である。
$\qgr D$ の豊富な表現可能加群列は、奇数項を次の偶数項から
生成できるため、偶数部分列でも豊富である。
表現可能加群間の Hom が $D$ の成分と一致することから、
この根の切断環は全次数で $B$ となる。
従って条件 (3) が所要の正則正規被覆を与える。

条件 (1) を仮定する。条件 (3) の被覆からすべての有限生成加群の
Hilbert 多項式が得られ、次数が高々一の対象は $\cC_1$ に入る
（命題~\ref{prop:sklyanin-h2} と同じ長さの議論）。
$f$ は正規である必要はなく、
$f^jB/f^{j+1}B\simeq(B/fB)(-j\deg f)$ が domain 性から従う。
条件 (2) のとき、この列は $\cC/\cC_0$ でも真に下降するから
$\pi B\notin\cC_1$ である。
$\pi B$ の任意の非零部分対象が $\pi B$ のシフトを含むことと、
それによる余商の Hilbert 多項式の次数が高々一であることから、
rank one torsion-free 性が従う。

逆に、適合根を得たなら $B$ の domain 性と被覆の有限条件は既に示した。
もしすべての非零斉次 $f$ に対して $\pi(B/fB)$ が有限長なら、
$\pi B$ の任意の非零部分対象による商は有限長となる。
$\cC/\cC_0$ の部分対象を持ち上げると、
そこにおける $\pi B$ の像は零または単純となり、
$\pi B\in\cC_1$ となって矛盾する。
\end{proof}

この還元では $D$ の四周期最小分解の完全性と noetherian 性を仮定した。
残る分類課題は、前節の距離二・距離三の関係からこれらを保証する条件、
および周期偶数環の domain 性と条件 (2) の具体的な判定である。

% !TeX root = ../nc-quadric-classification.tex

\section{整な楕円四つ組からの構成}
\label{sec:elliptic-inverse}

$D$ を Van den Bergh の三次三次元正則 $\mathbb Z$-algebra とする。
$P_i=e_iD$ を右加群、$S_i$ をその単純商と書き、
$V_i=D_{i,i+1}$、$R_i$ を三次関係空間、
$W_i=k w_i=(V_i\otimes R_{i+1})\cap(R_i\otimes V_{i+3})$ とする。
\cite[arXiv v3, \S 5.2]{VdB2011} に従い、$w_i$ の非零の代表を選び、
\[
 w_i=\sum_j v_j\otimes r_j,\qquad
 w_{i+1}=\sum_j r_j\otimes\theta_i(v_j)
\]
で次数 $+4$ の周期移動 $\theta$ を定める。各 $w_i$ のスカラー変更により、
$\theta$ は $D_{ij}$ に $c_i/c_j$ を掛ける対角型の自己同型まで定まる。

\begin{lemma}[四つの添字の移動と Serre 関手]
\label{lem:cubic-serre-period}
$\theta$ は、Mori--Nyman の Nakayama 同型
$\nu:D\to D(-4)$ の逆と上記の対角型自己同型を除いて一致する。
従って $\qgr D$ 上で $(4)_\theta\simeq S^{-1}[2]$ である。
\end{lemma}

\begin{proof}
右単純加群の極小分解を $F^i$ と書く。その項は
\[
 F^i_0=P_i,\quad F^i_1=V_i\otimes P_{i+1},\quad
 F^i_2=R_i\otimes P_{i+3},\quad F^i_3=W_i\otimes P_{i+4}
\]
であり、微分はテンソルの包含と積で与えられる
\cite[arXiv v3, Theorem 5.5.7]{VdB2011}。
特に $\Ext_D^q(S_i,P_j)$ は $(q,j)=(3,i+4)$ 以外で零である。
局所コホモロジー $R^3\tau(D)_{j,i}$ は
$\Ext^3_D(P_i/(P_i)_{\ge N},P_j)$ の直極限である。
これは torsion 部分の第 $i$ 成分を計算する
$\varinjlim_N\Hom_D(P_i/(P_i)_{\ge N},-)$ を導来すればよい。
右からの作用は始域における射での引き戻し、左からの作用は終域における射との合成である。

$M_i=P_i/(P_i)_{\ge i+2}$ とする。
$j=i+4,i+5$ について $\Ext^3_D(M_i,P_j)\to R^3\tau(D)_{j,i}$ は同型である。
さらに切り詰め次数を上げたときの核の単純因子は $S_l$ ($l\ge i+2$) なので、
上記の消滅と長完全列を用いて確かめられる。
短完全列 $0\to V_i\otimes S_{i+1}\to M_i\to S_i\to0$ に対し、
その拡大類を表す $F^i_1\to V_i\otimes S_{i+1}$ は自然な商写像である。
この写像を $V_i\otimes F^{i+1}$ に持ち上げると、
第一段階は $V_i\otimes P_{i+1}$ の恒等写像、第二段階は
$R_i\subset V_i\otimes V_{i+1}\otimes V_{i+2}$ の最後の因子を
$P_{i+3}$ への左乗法に用いる写像、第三段階は
$W_i\subset V_i\otimes R_{i+1}$ である。
これらが微分と可換であることは、それぞれテンソルの最後の因子を
二通りに取り出して同じ積になることで従う。

従って $F^i$ と $V_i\otimes F^{i+1}$ を項ごとに直和した
$M_i$ の分解では、第三微分の交差項は符号を除いて
$W_i\subset V_i\otimes R_{i+1}$ である。
$\Hom_D(-,P_{i+5})$ を取ると、
第三項は $V_i^*\otimes W_{i+1}^*$ のみで、第二項は零である。
一方、$\Hom_D(-,P_{i+4})$ を取った第三項は
\[
 W_i^*\oplus(V_i^*\otimes W_{i+1}^*\otimes V_{i+4})
\]
であり、第二項 $V_i^*\otimes R_{i+1}^*$ の像で割る。
その第二成分への写像は $w_{i+1}\in R_{i+1}\otimes V_{i+4}$ の縮約なので
同型である。第一成分への写像は $w_i\in V_i\otimes R_{i+1}$ の縮約である。
よって、$a\in V_i$ を用いる $S_{i+1}\to M_i$ での引き戻しと
$\theta_i(a)$ を用いる $P_{i+5}\to P_{i+4}$ との合成は、
双対基底 $w_i^*,w_{i+1}^*$ を使えば、いずれも
$V_i^*$ 成分を $a$ で評価する写像である。交差項の符号は基底変更で吸収できる。

より具体的には $R^3\tau(D)$ の成分ごとの $k$-双対を $L$ とすると、
適切に基底を取り直した $m_i\in L_{i,i+4}$ について
\[
 a m_{i+1}=m_i\theta_i(a)\qquad(a\in V_i)
\]
が成り立つ。右 AS 正則性より $L$ の第 $j$ 列は左加群 $D e_{j-4}$ である
\cite[arXiv v1, Theorem 4.8 and Remark 4.10]{MN2025}。
$D$ は $V_i$ で生成されるため、上の関係が双加群構造を決め、
$L\simeq D(0,-4)_{\theta^{-1}}$ となる。
従って \cite[arXiv v1, Definition 4.15 and Theorem 6.8]{MN2025} の
Nakayama 同型と Serre 関手の式が得られる。
$D_{ij}$ に $c_i/c_j$ を掛ける対角型自己同型による引き戻しは、
第 $j$ 成分を $c_j$ 倍して恒等と自然同型になる。
\end{proof}

\begin{lemma}[整な楕円四つ組の二周期の平方]
\label{lem:elliptic-period-square}
幾何的データが整な算術種数一の曲線 $E$ 上の許容四つ組の場合、
\cite[arXiv v3, Proposition 5.6.1]{VdB2011} の平行移動から生じる
次数 $+2$ の周期移動 $\beta$ は、対角型自己同型を除いて
$\beta^2=\theta$ を満たす。
\end{lemma}

\begin{proof}
滑らかな場合は $E$ に原点を選び、因子和を楕円曲線の群法則で取る。
特異な場合、正規化は種数零で、全 $\delta$ 不変量が一である。
従って正規化は $\PP^1$ で、特異点は一つである。
二枝なら完備局所環は正規化環のうち二つの剰余値が等しい部分環である。
実際、剰余値の一致だけで余次元一となる。
一枝なら一様化変数を $u$ として完備局所環は $k+u^2 k[[u]]$ である。
一次の項が非零の元を含めば、形式的逆写像により正規化環全体を含むことになり、
そのような元を含まない部分環は余次元一である。

二枝の場合は正規化上の二点を $0,\infty$、一枝の場合はその点を $\infty$ とする。
滑らかな部分にそれぞれ乗法群 $k^\times$ と加法群 $k$ の群構造を入れ、
以下では両方を加法記号で書く。この群の平行移動は正規化上の
拡大写像または平行移動で実現され、上記の剰余値または一次項の条件を保つので
$E$ の自己同型となる。滑らかな部分に台を持つ次数零の因子
$\sum_j m_j (p_j)$ が主因子となる条件は、群での和 $\sum_j m_j p_j=0$ である。
実際、正規化での有理関数 $\prod_j(t-p_j)^{m_j}$ は $\infty$ で $1$ となる。
二枝の場合は $0$ での値も $1$ であること、一枝の場合は $1/t$ の一次項が
零であることが特異点で可逆となる条件であり、いずれも上記の群の等式である。
従って滑らかな部分の因子和によって、これらの因子から得た同次数の
直線束の同型を判定できる。

各 $L_i$ を与える線型系の一般のファイバーは、
特異点を避ける相異なる二点からなる。実際、$L_i$ は二重被覆を与え、
標数零なので生成する関数体の拡大は分離的である。
この二点の和を $a_i$ とすると、同じファイバーの相手は
$\iota_i(p)=a_i-p$ と書ける。楕円 helix の関係から
$a_{i+2}-a_i$ は全ての $i$ で一定である。
$2v=a_0-a_2$ を満たす点 $v$ を選べる。群が楕円曲線、
$k^\times$ または $k$ のいずれでも、二倍写像は全射だからである。
この $v$ による平行移動で $L_i$ を引き戻すと $L_{i+2}$ となり、
切断への作用が $\beta$ を与える。

$\theta_i:V_i\to V_{i+4}$ の作用を評価写像で比較する。
一般の $p\in E$ に対し $p'=\iota_{i+2}(p)$ と置き、
$V_j=H^0(E,L_j)$ における $q$ での評価の核を
$k\ell_j(q)$ と書く。
$w_i$ を $L_{i+2}$ の $p$ における評価で縮約する。
$p$ と $p'$ は同じ射影評価を与える。
さらに $L_{i+1}$ の $p$ における評価で縮約すると、
$w_i$ が二つの関係空間に属することから
$\ell_i(p)\otimes\ell_{i+3}(p)$ の非零スカラー倍を得る。
$p'$ でも同様である。非零性は、隣接する二因子の縮約が非零という
幾何的五つ組の条件による \cite[arXiv v3, Definition 5.3.1]{VdB2011}。
$L_{i+1}$ の $p,p'$ での評価は独立なので、
$L_{i+2}$ のみを評価したテンソルは、適当な $\lambda,\gamma\ne0$ に対し
\[
 \lambda\ell_i(p)\otimes\ell_{i+1}(p')\otimes\ell_{i+3}(p)
 +\gamma\ell_i(p')\otimes\ell_{i+1}(p)\otimes\ell_{i+3}(p')
\]
と書ける。$L_{i+3}$ の $p$ における評価で縮約すると、
$L_{i+3}\not\simeq L_{i+2}$ より第二項だけが残る。
$w_{i+2}$ は $w_i$ の二回の回転により与えられ、
回転で末尾へ移る二因子は $\theta_i,\theta_{i+1}$ で移す。
従って $w_{i+2}$ を各 $L_{i+2},L_{i+3},L_{i+4}$ の $p$ における評価で
縮約して零になることから、
$\theta_i(\ell_i(p'))$ は $p$ で消える。ゆえに $\theta_i$ による
$L_{i+4}$ の $p$ での射影評価の引き戻しは $L_i$ の $p'$ での評価である。
$L_i$ に関して $p'$ は
$\iota_i(p')=p+a_i-a_{i+2}=p+2v$ と同じ射影評価を与える。
これは $\beta^2$ の作用である。
全ての $V_i$ で二つの写像はスカラー倍で一致し、
$V_i$ の生成性と添字直線上でのスカラーの逐次選択により
その差は対角型自己同型である。
\end{proof}

\begin{theorem}[整な楕円四つ組からの対象]
\label{thm:integral-elliptic-inverse}
標数零の代数閉体上で、$U=(E,L_0,L_1,L_2)$ を整な算術種数一の曲線上の
Van den Bergh の許容四つ組、$D=D(U)$ とする。
このとき $(\qgr D,\pi P_0)\in\HH$ である。
平方根は VdB の平行移動に由来する偶数添字の周期で与えられ、
二次 AS 正則被覆はその周期で次数付けした標準 hull から得られる。
\end{theorem}

\begin{proof}
$\beta$ を補題~\ref{lem:elliptic-period-square} の周期移動とする。
\cite[arXiv v3, Theorem 6.1, Lemmas 6.4--6.5 and Proposition 6.6]{VdB2011} の
標準 hull はこの周期移動を持ち上げる。
偶数部分を $\beta$ で次数付けした環 $B$ は、
型 $(1,4,6,4,1)$ の四次元二次 AS 正則環 $A$ の
正則正規二次元 $g$ による商である。ここで被覆の標準分解は線型なので
$A$ は Koszul である。左右の AS 条件については
\cite[arXiv v1, Proposition 4.4 and Theorem 4.5]{MN2025} を参照できる。
右 noetherian 性は同じ hull 定理による。
$D$ の反対環は添字を反転すると三次正則型であるため
\cite[arXiv v1, Theorem 4.5]{MN2025}、同じ VdB の分類と
\cite[arXiv v3, Corollary 5.5.9 and Lemma 3.5]{VdB2011} から
$D$ と $B$ は左 noetherian でもある。
$g$ によるフィルトレーションの次数環は $B$ 上の skew polynomial 環なので
$A$ にも左 noetherian 性が昇る。実際、斉次左イデアルについては固定した
内部次数でこのフィルトレーションは有限であり、
初期項の生成元から元のイデアルの生成元を持ち上げられる。
一般の左イデアルについては内部次数による最高次成分を使う。

偶数部分への制限は、指定対象を保つ $\qgr D\simeq\qgr B$ を与える
\cite[arXiv v3, Lemmas 3.4--3.5]{VdB2011}。
$\pi P_{-j}$ の列は ample であり
\cite[arXiv v1, Theorem 6.11]{MN2025}、奇数添字の $\pi P_i$ は
次の偶数添字の $\pi P_{i+1}$ の有限直和の商になる。
これは $D$ の隣接成分による生成性から従う。
従って偶数部分列も ample である。シフト $(1)$ は
$\qgr D$ 側の $(2)_\beta$ であり、
補題~\ref{lem:cubic-serre-period} と補題~\ref{lem:elliptic-period-square}
からその平方は $S^{-1}[2]$ である。
\cite[arXiv v1, Proposition 4.11]{MN2025} により
$\Hom(\pi P_j,\pi P_i)=D_{ij}$ であり、
$\pi P_0$ は exceptional である。この Hom の同一視は合成を保つため、
$\pi P_0$ の切断環は全非負次数で $B$ と同型となる。
実際、次数 $n$ の切断の空間 $D_{-2n,0}$ を $\beta^n$ で $D_{0,2n}$ に移すと、
次数 $m,n$ の積は $a\beta^m(b)$ となり、周期による $B$ の積に一致する。
$\qgr D$ の Ext 有限性は \cite[arXiv v3, Corollary 5.5.9]{VdB2011} による。

最後に rank 条件を確認する。
楕円 helix の切断 $\mathbb Z$-algebra $D_{\mathrm{geom}}$ への商の核は、
次数四の正則な一元列 $(h_i)$ で左右から生成される
\cite[arXiv v3, Theorem 5.5.6]{VdB2011}。
$\beta$ は幾何データ上の平行移動から来るためこの核を保つ。
従って偶数部分の次数付き商では、列 $(h_{2i})$ が一つの
正則正規元 $h\in B_2$ に降りる。
$B/(h)$ は $E$ の次数四の直線束 $L_0\otimes L_1$ と
平行移動による捩れ斉次座標環である。
従ってこれは domain であり、
$h$ による分離フィルトレーションの初期項から $B$ も domain である。
また $\qgr(B/(h))\simeq\coh E$ であり、
その構造対象のシフトは可逆層に対応する
\cite[arXiv v3, Corollary 5.5.9 and Lemma 3.5]{VdB2011}。
$h$ の冪が与える $\pi B$ の下降列の隣接商は、従って有限長でない。
命題~\ref{prop:sklyanin-h2} の証明と同じ Hilbert 多項式の議論を適用すると、
$\pi B$ は Gabriel--Krull の $\cC_1$ に属さず、
任意の非零部分対象は $\pi B$ のシフトを含み、その余商は $\cC_1$ に属する。
ここで使うのは $B$ の Hilbert 級数、domain 性と、
有限生成 $B$ 加群が $A$ 上で有限自由分解を持つことである。
よって $\pi B$ は rank one torsion-free であり、全ての条件が成立する。
\end{proof}

% !TEX root = ../nc-quadric-surface.tex
\section{Sklyanin pencil の一次近傍と同値判定}
\label{sec:pencil-deformation}

命題~\ref{prop:sklyanin-antiserre-center} の無限位数の範囲で、
同一の四次元 Sklyanin 環 $A$ の pencil を比較する。
$P=A^{(2)}$、$Z\subset P_1$ を元の中心二次 pencil、
$D=P/(Z)$ を楕円商と置く。直線 $\ell\subset Z$ に対して $R_\ell=P/(\ell)$ であり、
滑らかな対象に対応する場合にはこれが反 Serre 環である。
$H_\ell=Z/\ell$ の像は $R_\ell$ の中心一次部分である。

\begin{proposition}[中心方向の一次変形]
\label{prop:pencil-tangent}
$U=D_1$、$K=\ker(U\otimes_k U\to D_2)$ とする。
一次補空間を選ぶと、次数二の積により線形写像
\[
 \Delta:Z^*\longrightarrow
 \operatorname{Hom}_k(K,U)/\operatorname{im}\partial,\qquad
 (\partial f)(r)=(f\otimes1+1\otimes f)(r),\quad f\in U^*
\]
が得られ、右辺の商において補空間の選択に依存しない。
$\ell$ を消す非零汎関数 $\lambda$ に対して
$\Delta(\lambda)$ は $R_\ell/(H_\ell^2)$ の変形を与える。
この商の同型は $D$ の次数付き自己同型と $H_\ell$ の基底の変更の下で
$\Delta$ の値を保つ。
\end{proposition}

\begin{proof}
$g,h$ を $Z$ の任意の基底とする。$g$ は $P$ 上で正則であり、
$h$ は $P/(g)$ 上で正則である。これは Sklyanin の中心正則列から従う。
実際、適合する滑らかな一切断 $g$ については domain 性が保証し、
任意の基底については、中心正則列の Koszul 複体が基底変更で同型であることを使える。
後者から任意の非零 $g$ に対する商上の正則性が従う
（$P$ 上の $g$ 自身の正則性は domain 性による）。
この正則列により
\[
 (Z)/(Z)^2 \simeq Z\otimes_k D
\]
となる。例えば $g x+h y\in (Z)^2$ なら右辺の二次式を
$g,h$ の係数に吸収でき、残る関係は $g x'+h y'=0$ となる。
これを $(g)$ で割れば $y'\in(g)$ なので $x'\in(h)$ となる。

$U$ の線形持上げ $U\to P_1$ を選ぶ。
$r\in K$ の積を $P/(Z)^2$ で取ると $Z\otimes U$ に入る。
$\lambda$ で $Z$ を評価して $U$ への写像を得る。
持上げの差 $U\to Z$ は積に $\partial f$ を加え、二つの差の積は零となる。
従って商での類は内在的である。
さらに $(Z)^2+(\ell)$ で割ると $R_\ell/(H_\ell^2)$ となり、
その一次変形が評価した写像で与えられる。
商の同型は $H_\ell$ の同定と $U$ の持上げを移すので、
基底変更と $\partial$ の範囲で作用する。
\end{proof}

\begin{lemma}[持上げ変更の単射性]
\label{lem:pencil-lift-injective}
命題~\ref{prop:sklyanin-antiserre-center} の無限位数の範囲では、
$\partial:U^*\to\operatorname{Hom}_k(K,U)$ は常に単射である。
\end{lemma}
\begin{proof}
$D$ は二次環である。実際、二次環 $A$ の二次 Veronese $P$ の二次性は
命題~\ref{prop:antiserre-quadratic} と同じ関係の挿入による議論から従い、
一次空間 $Z$ で割っても二次性は保たれる。
$\partial f=0$ となる非零 $f\in U^*$ が存在すると仮定する。
各 $r\in K$ に対して
\[
 0=f((\partial f)(r))=2(f\otimes f)(r)
\]
である。標数零と二次性から、$u\mapsto f(u)z$ は次数付き全射
$D\twoheadrightarrow k[z]$ を定める。

この全射による右 $D$-加群を $M=k[z]$ と書く。
非零次数付き部分加群は $z^n k[z]$ なので、$\pi M$ は $\qgr D$ の単純対象である。
また $z$ 倍写像 $M\to M(1)$ の余核は有限次元であり、
$\pi M\simeq(\pi M)(1)$ となる。
命題~\ref{prop:sklyanin-antiserre-center} の楕円商の同定
$\qgr D\simeq\coh E$ の下で、単純対象は閉点 $p\in E$ の skyscraper 層である。
次数シフトは台を $2\tau$ または $-2\tau$ だけ平行移動する。
従って上の同型は $p$ がこの平行移動の固定点であることを意味するが、
$2\tau\ne0$ なので矛盾する。
\end{proof}

\begin{proposition}[固定 pencil の比較を有限化する条件]
\label{prop:pencil-finite-candidates}
$\Delta$ が単射であるとする。
固定 $P$ の滑らかな二切断の pointed exact 同値を判定するには、
有限個の一次線形写像だけを二次関係に照合すればよい。
さらに、各切断と同値な固定 pencil 内の切断は有限個である。
\end{proposition}

\begin{proof}
補題~\ref{lem:pencil-lift-injective} により $\partial$ も単射である。
まず $D$ の次数付き自己同型群を次数スカラーで割ると有限になる。
$D$ は次数八の可逆層を持つ楕円曲線の twisted homogeneous
coordinate ring（またはその反対環）である。
次数付き自己同型は次数零の分数体 $k(E)$ に作用する。
次数一切断の比で与えられる完全線形系の埋込みにより、
この作用は次数八の偏極を保存する $E$ の自己同型である。
詳しくは、基底切断 $u_0$ で割った切断 $u_i/u_0$ を座標比とすると、
代数自己同型が送る比は新基底の座標比であり、埋込みは射影線形変換を介して保たれる。
同じ分数体作用を持つ二自己同型の差は、比を固定するので
$U$ に共通の有理関数を掛ける。
この有理関数は非零有限次元空間を保つため、固有ベクトルに作用させれば
スカラーと分かる。$D$ は一次生成なので、差は次数スカラーである。
次数一の斉次可逆元による分数体上の作用は二倍平行移動であるため、
分数体上の自己同型はこれと可換である。
平行移動を除いた楕円群の自己同型が無限位数の移動点を固定すれば恒等になる。
実際、その自己同型から恒等を引いた非零準同型の核は有限である。
従って残る作用は次数八の偏極を保つ平行移動であり、$E[8]$ に属する。
特に候補は高々64個である。

比較する二切断の $H_\ell$ に基底を選ぶ。
$D$ 上の自己同型の代表 $\varphi$ ごとに
命題~\ref{prop:pencil-tangent} の必要条件を調べる。
二つの変形類は非零なので、それらの比較に必要な
$H_\ell$ 上の倍率は存在するなら一意である。
代表写像の差が $\partial f$ であることを確認すれば、
$\partial$ の単射性により $f$ も一意に定まる。
この $f$ は次数一における持上げのずれである。
従って商 $R_\ell/(H_\ell^2)$ 上の候補から、
$R_\ell$ の次数一の線形写像が一意に決まる。
次数スカラーは常に全環の自己同型なので、$\varphi$ は各剰余類で一つ取ればよい。
最後にその線形写像が全二次関係を保つか検査する。
命題~\ref{prop:antiserre-quadratic} と定理~\ref{thm:antiserre-invariant}
により、この最後の条件は十分でもある。

固定した切断の非零変形類の射影化に $\varphi$ を作用させると、
$\Delta$ の単射性により各 $\varphi$ につき一致する切断は高々一つである。
\end{proof}

\begin{proposition}[通常の Sklyanin パラメータでの開条件]
\label{prop:pencil-tangent-generic}
通常の Sklyanin のパラメータ平面のチャート
$\gamma=-(\alpha+\beta)/(1+\alpha\beta)$ において、
$\partial$ と $\Delta$ がともに単射となる条件は非空 Zariski 開条件を与える。
無限位数の範囲では $\partial$ の単射性は補題~\ref{lem:pencil-lift-injective}
により自動的であり、有限比較に残る仮定は $\Delta$ の単射性だけである。
従ってこの開条件を満たし、かつ命題~\ref{prop:sklyanin-antiserre-center} の仮定を
満たす固定 pencil について、滑らかな切断の同値判定は命題~\ref{prop:pencil-finite-candidates} の有限手順で行える。
\end{proposition}

\begin{proof}
$\dim U=8$、$\dim D_2=16$ なので $K$ は48次元である。
一次持上げと関係空間の基底を替えても、
$\partial$ の階数と、$\partial$ の像に中心二方向の列を加えた行列の階数は不変である。
二つの単射性はこれらがそれぞれ $8,10$ であることと同値であり、
定ランクの座標チャートでは極大小行列式の非消滅で書ける。

非空性を有理数上で確認した。
$(\alpha,\beta,\gamma)=(2,3,-5/7)$ において、
$\sum a_i x_i^2$ の四生成元との交換子を次数三の関係空間で割って
二次元の中心 pencil の基底を求める。
次数四では、二次関係を四文字の三箇所へ挿入した空間で割る。
得られる35次元空間に、中心の三つの二乗積と
二つの中心元と八つの $U$ の持上げとの積を並べると、階数は19となる。
さらにその像で割ると16次元になり、
$U\otimes U$ からの積の核が48次元と求まる。
この核に対し、中心方向の積の係数と $\partial$ の列を並べると
有理数上の階数がそれぞれ $8,10$ となる。
計算の入力と行操作は
\texttt{computations/verify\_pencil\_tangent.py} に収めた。
この計算は開条件の非空性の証明に用い、当該有理パラメータの
平行移動の無限位数性は用いていない。
\end{proof}


% !TEX root = ../nc-quadric-classification.tex
% 2026-11-21 / v1.1
\section{有限 foundation と反 Serre 表示のスタック}
\label{sec:foundation-stack}

この節の有限代数スタックは直ちに構成できる。
それを $\HH$ の全相対モジュライと同一視するためには、末尾に記す
相対存在・開性・降下の証明がさらに必要である。

\subsection{二つの線形写像によるパラメータ}

$P,Q$ を二次元、$W_1,W_2,W_3$ をそれぞれ二、四、六次元のベクトル空間とする。
箙 $1\overset{2}{\longrightarrow}2\overset{2}{\longrightarrow}3$ の道代数を $H$ と書く。
隣接矢空間は $P,Q$ であり、$H_{13}=P\otimes Q$ である。
次元ベクトル $(2,4,6)$ の $H$-表現は
\[
 F:P\otimes W_1\longrightarrow W_2,\qquad
 G:Q\otimes W_2\longrightarrow W_3
\]
の組で与えられる。従ってそのパラメータ空間は
\[
 \mathcal R_{\mathrm{fnd}}
 =\Hom(P\otimes W_1,W_2)\oplus\Hom(Q\otimes W_2,W_3)
 \simeq\mathbb A^{64}.
\]
表現 $M=(W_1,W_2,W_3;F,G)$ に対応する一頂点拡大を $\Lambda_M$ とする。
具体的には頂点を $0,1,2,3$ とし、
\[
 (\Lambda_M)_{0i}=W_i\ (1\le i\le3),\qquad
 (\Lambda_M)_{ij}=H_{ij}\ (1\le i\le j\le3),\qquad
 (\Lambda_M)_{00}=k
\]
と置く。積は $H$ の道の積と表現 $M$ の作用による。
これは二十四次元の有向代数である。
箙の関係式の有限型スタックの一般的構成は
Abdelgadir--Okawa--Ueda \cite[Definition 3.5, Proposition 3.6]{AOU2014} にある。
ここでは全ての合成則を明示し、結合律による方程式も消去できる。

群
\[
 \mathcal G=\GL(W_1)\times\GL(W_2)\times\GL(W_3)
             \times\GL(P)\times\GL(Q)
\]
を、$g=(g_1,g_2,g_3,a,b)$ に対し
\[
 (F,G)\longmapsto
 \bigl(g_2F(a^{-1}\otimes g_1^{-1}),\
       g_3G(b^{-1}\otimes g_2^{-1})\bigr)
\]
で作用させる。中心トーラス
\[
 K=\left\{
 \left(\lambda_1 I,\lambda_2 I,\lambda_3 I,
       \frac{\lambda_2}{\lambda_1}I,
       \frac{\lambda_3}{\lambda_2}I\right)
 \right\}\simeq\mathbb G_m^3
\]
は全空間で自明に作用する。

\begin{proposition}[有限 foundation の商スタック]
\label{prop:foundation-stack}
商
\[
 \mathfrak F^{\mathrm{lin}}=[\mathcal R_{\mathrm{fnd}}/\mathcal G],\qquad
 \mathfrak F=[\mathcal R_{\mathrm{fnd}}/(\mathcal G/K)]
\]
は有限型の滑らかな Artin スタックである。
前者は上の形の有向代数を頂点とともに記述し、後者は中心 $K$ の rigidification である。
$\dim\mathfrak F=3$ である。
全ての $\HH$ 対象の定理~\ref{thm:normal-helix} の foundation は、
その End 代数を通じて $\mathfrak F^{\mathrm{lin}}$ の幾何点を与える。
\end{proposition}

\begin{proof}
$\mathcal R_{\mathrm{fnd}}$ は滑らかな六十四次元アフィン空間であり、
$\dim\mathcal G=4+16+36+4+4=64$ である。
$K$ は滑らかな中心トーラスなので $\mathcal G/K$ も滑らかな代数群である。
従って二つの商は滑らかな Artin スタックであり\cite[Tag 0DLS]{StacksModuli}、
$\dim\mathfrak F=64-(64-3)=3$ である。

頂点を保つ $H$ の自己同型は $P,Q$ 上の基底変更で尽くされる。
したがって $\mathcal G$ の作用は、頂点を保つ $\Lambda_M$ の同型を正確に与える。
対角単元 $\sum_{i=0}^3\lambda_i e_i$ による内部自己同型は、共通倍率で割り
$\lambda_0=1$ とすれば上の $K$ となる。
これは各 foundation 対象の比較同型のスカラーによる変更を、全体の自然スカラーで割ったものに対応する。
従って $K$ の除去は、$E_0$ の標識だけを残して、
$E_1,E_2,E_3$ を保つ比較同型の選択を忘れる rigidification である。
幾何点では、foundation の各項を同型類として保つ Morita 同値の自己同型群が
$\operatorname{Aut}_{\{e_i\}}(\Lambda_M)/K$ となる。
任意の底上では、これらの対象の存在と線束 twist の降下を区別する必要がある。

$\HH$ から得た四項 $(E_0,E_1,E_2,E_3)$ では
\[
 \dim\Hom(E_i,E_j)=
 \begin{pmatrix}
 1&2&4&6\\0&1&2&4\\0&0&1&2\\0&0&0&1
 \end{pmatrix}_{ij}
\]
であり、$\Hom(E_1,E_2)\otimes\Hom(E_2,E_3)\to\Hom(E_1,E_3)$ は同型である。
後者は系~\ref{cor:normal-cubic-defect} の奇数頂点間の積の計算による。
よって $E_1,E_2,E_3$ の End 代数が $H$ となり、
$\Hom(E_0,E_i)$ が上の表現 $M$ を与える。
\end{proof}

\begin{remark}
$\mathfrak F$ は foundation の線形化を除いた有限代数のパラメータスタックである。
任意の点の $\operatorname{Perf}(\Lambda_M)$ に、$\HH$ の heart、
相対 Serre 根、豊富性、AS 正則被覆が存在するとは主張していない。
また foundation を保たない全ての pointed 同値を、この群 $\mathcal G/K$ だけで割ったとも主張しない。
一般に有限代数の導来圏の同値と元の heart の exact 同値は区別する。
本稿の体上では構造対象を保つ導来同値の exact 性を
系~\ref{cor:pointed-derived-exact} で証明しているが、族での底変更には別途確認を要する。
\end{remark}

\subsection{Hochschild 障害の消滅}

\begin{proposition}
\label{prop:foundation-hochschild}
任意の $M\in\mathcal R_{\mathrm{fnd}}$ に対し
\[
 \operatorname{HH}^n(\Lambda_M)=0\quad(n\ge3),\qquad
 \dim\operatorname{HH}^2(\Lambda_M)
 -\dim\operatorname{HH}^1(\Lambda_M)=3.
\]
従ってこれらの代数の通常の結合代数変形について、三次 Hochschild コホモロジーにある障害は消える。
これは相対 $\HH$ 条件の無障害性とは別の主張である。
\end{proposition}

\begin{proof}
$A=\Lambda_M$、$E=\bigoplus_{i=0}^3ke_i\simeq k^4$ とする。
$E$ は分離代数なので、$E$ に相対的な正規化 bar 複体が $\operatorname{HH}^*(A)$ を計算する。
有向性からそのコチェイン項は
\[
 C^n=\bigoplus_{i_0<\cdots<i_n}
 \Hom_k\bigl(A_{i_0i_1}\otimes\cdots\otimes A_{i_{n-1}i_n},
             A_{i_0i_n}\bigr)
\]
であり、$C^0=k^4$、$C^n=0$ ($n>3$) である。
特に
\[
 C^3=\Hom(A_{01}\otimes A_{12}\otimes A_{23},A_{03}).
\]
$C^2$ の $(0,1,3)$ 成分
$\Hom(A_{01}\otimes A_{13},A_{03})$ から $C^3$ への微分を取る。
他の成分を零とすれば、bar 微分の四項のうち残るものは
\[
 (a,b,c)\longmapsto\varphi(a,bc)
\]
だけである。$A_{12}\otimes A_{23}\xrightarrow\sim A_{13}$ なので
この写像は二つの四十八次元空間の同型である。
従って $C^2\to C^3$ は全射で、$\operatorname{HH}^{\ge3}(A)=0$ となる。

各コチェイン項の次元は
\[
 (\dim C^0,\dim C^1,\dim C^2,\dim C^3)
 =(4,80,128,48).
\]
ここで $\dim C^1=2^2+4^2+6^2+2^2+4^2+2^2=80$、
$\dim C^2=16+48+48+16=128$ である。
中心の元は全 $e_i$ と可換なので対角であり、
各 $A_{0i}=W_i\ne0$ と可換であるため四つの対角係数は全て等しい。
従って $\operatorname{HH}^0(A)=Z(A)=k$ である。
Euler 標数を取ると
\[
 1-\dim\operatorname{HH}^1(A)+\dim\operatorname{HH}^2(A)
 =4-80+128-48=4
\]
となり、第二式が従う。
\end{proof}



\subsection{二つの ruling と foundation の内在性}

\begin{proposition}[ruling 対象の内在的認識]
\label{prop:intrinsic-ruling-pair}
$(\mathcal C,\mathcal O)\in\HH$、$t=S^{-1}[2]$ とする。
Gabriel--Krull rank one torsion-free exceptional 対象 $F$ で
\[
 \chi(\mathcal O,t^nF)=(2n+1)(2n+2)\qquad(n\gg0)
\]
を満たすものの同型類は、正確に二つである。
任意の適合表示の二つの次数零生成 MCM 型の tails 像が、この二類を与える。
従ってその非順序対は $(\mathcal C,\mathcal O)$ から内在的に定まる。
\end{proposition}

\begin{proof}
適合表示を一つ選び、二つの MCM 型の tails 像を $X,Y$ と書く。
定理~\ref{thm:normal-helix} の評価列により、次数付き加群 $Z=X,Y$ と
$p(Z)$ は自由加群に埋め込まれる。
$B$ の商斜体へ移せば両者の正の rank の和は二なので、各 rank は一である。
従って $Z(-1)\hookrightarrow B^2$ の非零成分 $Z(-1)\to B$ は単射である。
実際、商斜体へ移した核は零であり、$Z$ は自由加群の部分加群なので
この核はもともと零である。
また $Z$ の非零斉次元は単射 $B(-d)\hookrightarrow Z$ を与える。
tails と Gabriel--Krull 商 $q:\mathcal C\to\mathcal C/\mathcal C_1$ に移すと
\[
 q\pi B(-d)\hookrightarrow q\pi Z\hookrightarrow q\pi B(1)
\]
となり、外側は非零単純対象なので $q\pi Z$ も非零単純である。
$\pi Z\hookrightarrow\pi B(1)$ によって torsion-free 性も従う。
exceptionality は同定理の geometric helix から従い、
Hilbert 級数 $2/(1-z)^3$ から指定の Euler 多項式を得る。

命題~\ref{prop:first-orbit-object} と同じ順序で
$K_0(\mathcal C)$ を $K_0(\mathbb P^1\times\mathbb P^1)$ と同定する。
$[F]$ の Chern 指標を $r+aH_1+bH_2+cH_1H_2$ と書くと
\[
 \chi(\mathcal O,t^nF)
 =r(2n+1)^2+(a+b)(2n+1)+c,\qquad
 \chi(F,F)=r^2+2rc-2ab.
\]
係数比較と例外性から $r=1$、$a+b=1$、$c=0$、$ab=0$ となる。
従って可能な類は $(a,b)=(1,0),(0,1)$ の二つだけである。

同じ $K_0$ 類を持つ二つの rank one torsion-free exceptional 対象 $L,F$ が
この多項式を満たすなら $\chi(L,F)=1$ である。
もし $\Ext^2(L,F)\ne0$ なら Serre 双対性から $F\to t^{-1}L$ が非零であり、
rank one torsion-free 性により単射となる。
十分大きい $n$ で Hom 次元は Euler 特性と一致する。
この点は命題~\ref{prop:first-orbit-object} の証明の自由被覆の議論による。
具体的には任意の $T$ を $\mathcal O$ のシフトの有限直和 $P$ で覆い、
核も同様に $P'$ で覆う。自由シフト間の $\Ext^1$ の消滅から
$\Ext^1(t^{-n}\mathcal O,T)$ は $\Ext^2(t^{-n}\mathcal O,P')$ の部分商であり、
後者は Serre 双対性と負次数切断の消滅で零となる。
$\Ext^2(t^{-n}\mathcal O,T)$ についても $P$ で同じ議論が使える。
従って上の単射は
\[
 (2n+1)(2n+2)\le(2n-1)(2n)
\]
を与えて矛盾する。従って $\Hom(L,F)\ne0$ であり、同じ理由で単射である。
その余核 $T$ は $K_0$ 類が零なので Hilbert 多項式が零である。
ここで有限生成 $B$-加群を $A$-加群と見ると有限自由分解を持つため
Hilbert 関数は最終的に多項式であり、有限次元加群を除いてもその多項式は変わらない。
従って零 Hilbert 多項式を持つ加群は有限次元であり、$T=0$ となる。
よって各 $K_0$ 類に高々一つの同型類しかない。
$X,Y$ は互いに非同型なので二つの類を尽くす。
\end{proof}

\begin{corollary}[指定構造対象から始まる二つの foundation]
\label{cor:two-intrinsic-foundations}
$M=s\mathcal O$ を命題~\ref{prop:first-orbit-object} の内在的対象とする。
定理~\ref{thm:normal-helix} で作る、添字零が $\mathcal O$ である foundation は、
各項の同型まで正確に二つである。
それぞれ最初の三項が $(\mathcal O,X,M)$ または $(\mathcal O,Y,M)$ となり、
第四項はこれらの三項から内在的に決まる。
pointed 自己同値群はこの二元集合に作用する。
\end{corollary}

\begin{proof}
初めの三項の内在性は直前の命題と第一軌道対象の認識による。
一つの選択 $E_1=X$ または $Y$ を固定する。
既存の full foundation の完全性から
\[
 {}^{\perp}\langle\mathcal O,E_1,M\rangle\simeq D^b(\operatorname{mod}k).
\]
この左直交の exceptional 生成元は、第四項の cohomological shift に限られる。
条件 $\operatorname{RHom}(M,E_3)\simeq k^2[0]$ はその shift を一意に固定する。
従って第四項も一意である。
\end{proof}

\begin{remark}[相対的な二重被覆へ残る点]
上の二択は代数閉体上の対象についての結論である。
相対的な有限 étale 二重被覆、あるいは単一の $S_2$ 商による
$\HH$ スタックの構成には、次を別途示す必要がある。
\begin{enumerate}
\item 相対 Serre と Euler 多項式によって定義した $M,X,Y$ の
      fppf 局所的存在、底変更、線束 twist を除く降下。
\item $\HH$ に対応する $\mathcal R_{\mathrm{fnd}}$ の部分関手が
      有限表示の開または局所閉条件となること。
\item ruling の交換で得る新しい foundation が族で代数的に構成され、
      二回の交換と恒等の整合性を持つこと。
\item 相対 AZ 再構成と heart の内在的認識により、
      このgroupoidの全矢が元の $\mathcal O$-marked exact 同値を与えること。
\end{enumerate}
exceptional 対象の $\Ext^1$ と $\Ext^2$ の消滅は、微小拡大上の一意延長の
根拠になるが、上記の有限表示・開性・代数化を単独で証明するものではない。
$\mathfrak F$ の滑らかさや $\operatorname{HH}^3$ の消滅だけから、
$\HH$ の相対条件が全ての変形で保たれるとは結論しない。
\end{remark}

\subsection{反 Serre 環の商と標識に由来する慣性}
\label{subsec:antiserre-stack-candidate}

ここでは $k$ 上の pointed 圏の射を
$(F,u):(\mathcal C,\mathcal O)\to(\mathcal C',\mathcal O')$、
$u:F\mathcal O\xrightarrow\sim\mathcal O'$ とし、
二つの射の間の二射を $u'\eta_{\mathcal O}=u$ を満たす自然同型
$\eta:F\Rightarrow F'$ とする。
以下の自己同型群は、この二groupoidの一射を二射で割った群を意味する。

\begin{lemma}[標識付き圏の二射と次数スカラー]
\label{lem:pointed-degree-scalar}
$(\mathcal C,\mathcal O)\in\HH$ とし、
$R=\bigoplus_{n\ge0}\Hom(\mathcal O,t^n\mathcal O)$ とする。
$\operatorname{Aut}(\operatorname{Id}_{\mathcal C})=k^\times$ であり、
pointed 一射の二自己同型は自明である。
また完全列
\[
 1\longrightarrow\mathbb G_m(k)
 \xrightarrow{\lambda\mapsto\xi_\lambda}
 \operatorname{Aut}_{\mathrm{gr}}(R)
 \longrightarrow
 \operatorname{Aut}_{\mathrm{pt}}(\mathcal C,\mathcal O)
 \longrightarrow1,
 \qquad \xi_\lambda|_{R_n}=\lambda^n\operatorname{id}
\]
がある。従って graded algebra の商をそのまま用いると、
pointed 圏の自己同型群にない次数スカラーが慣性として一つ余る。
一方、標識対象自身のスカラーをさらに慣性へ加える必要はない。
\end{lemma}

\begin{proof}
自然自己同型 $\eta$ の各 $t^{-n}\mathcal O$ 上の成分は、例外性により
スカラー $\lambda_n$ である。
$\Hom(t^{-n}\mathcal O,\mathcal O)=R_n\ne0$ と自然性から
$\lambda_n=\lambda_0$ となる。
反 Serre 軌道の豊富性により全ての対象はこれらの有限直和の商なので、
$\eta$ は全ての対象で同じスカラーである。
pointed 二射は標識上で恒等であるため、このスカラーは一に限られる。
特に $(\operatorname{Id},\lambda\operatorname{id}_{\mathcal O})$ と
$(\operatorname{Id},\operatorname{id}_{\mathcal O})$ は互いに二同型である。

graded 自己同型は AZ 再構成を通じて pointed exact 自己同値を与える。
逆に pointed exact 自己同値は Serre 関手と可換するので、
その可換同型を選ぶことで $R$ の graded 自己同型を与える。
可換同型の異なる選択は、前段の自然自己同型の計算により一つのスカラーだけ異なり、
$R_n$ 上ではその $n$ 乗だけ異なる。これで全射性を得る。

核を計算する。graded 自己同型 $\phi$ の誘導する pointed 関手が恒等と二同型なら、
各自由シフト $\pi R(i)$ 上のその二同型は、
基準同型 $b\mapsto\phi^{-1}(b)$ のスカラー $c_i$ 倍である。
隣接シフト間の Hom に対する自然性は
$c_{i+1}\phi_1=c_i\operatorname{id}_{R_1}$ を与える。
従って $\phi_1$ はスカラーで、$R$ の一次生成性により $\phi=\xi_\lambda$ である。
逆に次数スカラーでの加群 twist は、次数 $n$ 上で $\lambda^{-n}$ を掛ける写像により
恒等関手と自然同型であり、次数零の標識も保つ。
\end{proof}

九次元空間 $U$ を固定し、
\[
 \operatorname{Gr}\bigl(56,U\otimes U\bigr)
\]
上の関係部分束 $J$ から二次環 $R_J=T(U)/(J)$ を作る。
反 Serre 環として実現される部分を記号的に $X_{\mathrm{as}}$ と書く。
ここで $X_{\mathrm{as}}$ はまだ単なる幾何点集合以上の意味には構成していない。
$\GL(U)$ の中心は Grassmannian に自明に作用するため、
実現部分が有限表示の不変部分スキームとして構成できれば
\[
 [X_{\mathrm{as}}/\GL(U)]\longrightarrow
 [X_{\mathrm{as}}/\PGL(U)]
\]
は次数スカラーを除く $\mathbb G_m$-rigidification となる。
補題の幾何点慣性の計算から、pointed 圏のモジュライに対応すべきものは右辺である。
ただしこの式だけでは相対 $\HH$ モジュライの構成を完了しない。

\begin{proposition}[九生成元・二十五次元二次成分の gerbe は中立]
\label{prop:antiserre-gerbe-neutral}
$Y=\operatorname{Gr}(56,U\otimes U)$ とし、
$\mathcal J\subset U\otimes U$ を普遍部分束とする。
$[Y/\GL(U)]$ 上の普遍環の初めの二成分を
\[
 \mathcal R_1=U,\qquad
 \mathcal R_2=(U\otimes U)/\mathcal J
\]
と書く。その階数はそれぞれ九、二十五である。
線束
\[
 \mathcal L=(\det\mathcal R_1)^{-11}
                 \otimes(\det\mathcal R_2)^2
\]
は中心 $\mathbb G_m\subset\GL(U)$ に関して重み一である。
従って $\mathbb G_m$-gerbe
\[
 [Y/\GL_9]\longrightarrow[Y/\PGL_9]
\]
は中立であり、選んだ $\mathcal L$ によって
$[Y/\GL_9]\simeq[Y/\PGL_9]\times B\mathbb G_m$ と同値になる。
$\mathcal R_n\otimes\mathcal L^{-n}$ は $[Y/\PGL_9]$ 上へ降下し、
その全体は通常の準連接次数付き代数をなす。
特にこの商に残る独立の Brauer 障害はない。
同じ結論は任意の $\PGL_9$-不変部分への制限に成り立つ。
\end{proposition}

\begin{proof}
中心の元 $\lambda I$ は $\mathcal R_1$ に $\lambda$、
$\mathcal R_2$ に $\lambda^2$ で作用する。
従って $\det\mathcal R_1$ と $\det\mathcal R_2$ の重みは九と五十であり、
\[
 -11\cdot9+2\cdot50=1
\]
から $\mathcal L$ の重みは一である。
$\mathcal V=\mathcal R_1\otimes\mathcal L^{-1}$ の中心作用は自明なので、
$\GL_9$-同変降下により $[Y/\PGL_9]$ 上の階数九のベクトル束となる。
その射影枠束は普遍 $\PGL_9$-torsor と一致するので、
通常の枠束はこの torsor の $\GL_9$ への持ち上げを与える。
これが gerbe の中立化である。

同じ事実を局所的に見ると、$\PGL_9$-遷移関数の持ち上げ
$g_{ij}$ による三重重なりの障害が $\lambda_{ijk}$ であるとき、
$\mathcal L$ の遷移係数 $\ell_{ij}$ は重み一のため
$\ell_{ij}\ell_{jk}\ell_{ki}=\lambda_{ijk}$ を満たす。
従って $\ell_{ij}^{-1}g_{ij}$ は真の $\GL_9$-cocycle となる。
一般にこのような重み一の線束を持つ $\mathbb G_m$-gerbe は、
線束の枠束を取ることにより底の $B\mathbb G_m$ との積に同値となる。

$\mathcal R_n$ に対する中心の重みは $n$ なので
$\mathcal R_n\otimes\mathcal L^{-n}$ の重みは零である。
積もこの twist と両立する。
具体的には $\mathcal J\otimes\mathcal L^{-2}\subset
\mathcal V\otimes\mathcal V$ が部分束として降下し、
\[
 T(\mathcal V)/(\mathcal J\otimes\mathcal L^{-2})
\]
が求める代数である。
三次以上の成分の平坦性は Grassmannian 全体では主張せず、
別に定める相対 Hilbert 条件の locus で課す。
\end{proof}

\begin{remark}[有限型 locus と相対再構成の不足]
まず、幾何点の実現集合やその被約スキーム像を取るだけでは、
非被約底上の変形を指定したことにならない。
相対 Hilbert 関数・full section algebra 条件、AZ 豊富性、相対 Serre 同定、
平方根の存在と二次 AS 正則被覆の条件を、
有限表示の部分関手として記述する必要がある。
体上で成立する有限回の同値判定も、全底上で一様な次数の
Isom 関手や降下の表現可能性をそのまま与えるものではない。

直前の命題により、$\PGL_9$-torsor の Brauer 障害そのものは解消される。
しかし線束で次数ごとに twist した環の shift は、元の shift に
底からの線束をテンソルしたものに対応する。
従って相対 tails 圏の構成、相対 Serre 関手との同定、
その $\Hom(\mathcal O,t^n\mathcal O)$ の底変更を確認する仕事はなお残る。
gerbe の中立性は、この相対再構成を代用する主張ではない。

従って現段階で証明した内容は、有限 foundation の明示商スタック、
その Hochschild 計算、幾何点上の内在的二択、および反 Serre 表示で除くべき慣性の同定である。
これらを全 $\HH$ の有限型 Artin スタックと同一視する相対定理は残る。
\end{remark}

% !TeX root = ../nc-quadric-classification.tex
% 2026-11-21 / v1.1
\section{中央二重平方族の表示スタック}
\label{sec:noncubic-stack}

$W=k u\oplus k v$ とし、非退化二次形式の開集合を
$Q^{\mathrm{nd}}=\operatorname{Sym}^2W\setminus V(\operatorname{disc})$ と書く。
\[
 P=Q^{\mathrm{nd}}\times Q^{\mathrm{nd}},\qquad
 \widetilde G=(\operatorname{GL}(W)\times\mathbb G_m^2)\rtimes S_2
\]
と置く。$S_2$ は二つの形式と二つのスカラーを交換し、恒等成分の作用は
\[
 (A,a_1,a_2)\cdot(f_1,f_2)
 =\bigl(a_1^{-2}\operatorname{Sym}^2(A)f_1,
        a_2^{-2}\operatorname{Sym}^2(A)f_2\bigr)
\]
である。これは環
\[
 B_f=k[u,v]\langle c,d\rangle/(c^2-f_1,d^2-f_2)
\]
の $W$ の基底変更、$c,d$ のスカラー変更と交換に対応する。
中央部分群
\[
 D=\{(\lambda I_W,\lambda,\lambda)\mid\lambda\in\mathbb G_m\}
\]
は $P$ に自明に作用する。$G=\widetilde G/D$ とする。
$G$ は $W\oplus kc\oplus kd$ の分解を保ち、後二因子を交換してよい
$\operatorname{PGL}_4$ の閉部分群であり、滑らかな五次元アフィン代数群である。

\begin{proposition}[表示から構成するスタック]
\label{prop:noncubic-presentation-stack}
\[
 \mathfrak N^{\mathrm{pres}}=[P/G]
\]
は滑らかな有限型 Artin スタックであり、次元は一である。
$P$ 上の普遍次数付き環は全ての次数で有限自由である。
その幾何点の同型類は全 noncubic $\mathcal H_2$ 対象の pointed exact 同値類と一致する。
\end{proposition}

\begin{proof}
$P$ は六次元の滑らかな開集合であり、$G$ は五次元の滑らかな代数群なので、
商スタックの一般論により最初の主張を得る
\cite[Tag 0DLS]{StacksModuli}。

普遍環は係数環上で $u^a v^b w$ を基底とする。
ここで $w$ は $c,d$ の交互語であり、空語も許す。
$c^2\mapsto f_1$, $d^2\mapsto f_2$ の書換えは長さを下げ、
$c^3,d^3$ の重なりは係数の中心性によって解消する。
従って基底は任意の底変更と整合し、各次数は有限自由である。

幾何点の同型類については定理~\ref{thm:noncubic-classification} と
定理~\ref{thm:ore-anticommuting-central} を使う。非退化二次形式は
被約二点集合をスカラーを除いて一意に定め、代数閉体上ではそのスカラーを
$c,d$ の平方によって吸収できる。$D$ は軌道を変えない。
\end{proof}

\begin{proposition}[根配置スタックへの射]
\label{prop:noncubic-sign-gerbe}
$\mathcal D=\mathbb P(\operatorname{Sym}^2W)\setminus V(\operatorname{disc})$ とし、
\[
 \mathfrak R=[(\mathcal D\times\mathcal D)/
       (\operatorname{PGL}(W)\times S_2)]
\]
と置く。自然な射 $\mathfrak N^{\mathrm{pres}}\to\mathfrak R$ は gerbe であり、
二つの形式を順序付けた局所表示上では相対慣性が $\mu_2^2$ である。
$S_2$ はこの二因子を交換する。
従って根配置の商スタック自体は pointed 慣性を保持しない。
\end{proposition}

\begin{proof}
$G\to\operatorname{PGL}(W)\times S_2$ の核は二次元トーラスであり、
形式の二つのスカラーに対して逆二乗により作用する。
その作用の核は $\mu_2^2$ で、$c,d$ の独立の符号変更を表す。
$P\to\mathcal D^2$ は二つの形式の代表を選ぶトーラス torsor である。
任意の根配置の同型は、射影変換の線型持上げと単元の平方根を選べば
局所的に $G$ の元へ持ち上がる。
標数零ではこれらの選択は fppf、実際には適当な\'etale 局所化で可能である。
従って対象と同型の両方が局所的に持ち上がり、gerbe の条件を満たす
\cite[Tag 06PD]{StacksModuli}。
\end{proof}

\begin{remark}[GIT の同一視]
$u,v$ の基底を固定し、$f_1=uv$, $f_2=u(v-t u)$ とする。
$t\ne0$ では根対が一根共有、$t=0$ では二対が一致する。
$t\ne0$ の全点は一つの軌道に属するが、その閉包は $t=0$ の軌道を含む。
従って通常の GIT 商は、これら二つの異なる pointed 類を分離しない。
同型類集合の二つの特別点を、一つの既存 GIT 商から読み取ることはできない。
また $\operatorname{Sym}^2\mathcal D$ を単なるスキームとして先に作ると、
一致した二対の $S_2$ 慣性も失うため、上の $S_2$ 商はスタックとして取る必要がある。
\end{remark}

\begin{remark}[相対版として未証明の部分]
上で無条件に構成したのは、普遍中央二重平方環とその射影化した座標変換からなる
表示スタックである。普遍環は $P$ 上に存在し、$D$ は次数 $n$ に重み $n$ で
作用するため、剛体化後のスタック上への通常の環束としての降下はここでは主張しない。
幾何点の同型類は一致するが、任意の基底上の
$\mathcal O$-marked 圏の族との同値を主張するには、相対再構成と相対正規形が必要である。
特に「全ての幾何ファイバーが noncubic」という条件だけでは、
階数四を検出する小行列式の冪零成分が消えるとは限らない。
例えば双対数上の行列は全幾何ファイバーで階数三でも、四次小行列式が
非零の冪零元となり得る。階数三の閉部分スタックを定めるなら、
その小行列式をスキーム論的に零とする条件が必要である。
このため、点の分類をそのまま全 family の分類へ読み替えない。
\end{remark}

% !TEX root = ../nc-quadric-classification.tex
% 2026-11-21 / v1.8
\section{三次平坦性から全次数平坦性への相対的な認識}
\label{sec:relative-koszul}

ここでは三次までの関係式の平坦性と、全次数の平坦性との間を閉じる。
用いるのは Polishchuk--Positselski の局所環上の Koszul 変形原理
\cite[Chapter 6, Theorem 7.3]{PolishchukPositselski2005} である。
この定理は被約基底や一径数の形式基底に限らず、
任意の Noetherian 局所環に適用できる。
従って本節の結論は双対数上の一次変形だけでなく、
任意の Artin 基底と非被約 Noetherian 基底を含む。
一方、得られるのは次数付き代数とその線型解消の相対的な結論であり、
相対 Abel heart、Serre 関手、平方根、AS 被覆の認識は別に必要である。

\begin{proposition}[全 $\HH$ の切断環と反 Serre 環の Koszul 性]
\label{prop:all-h2-antiserre-koszul}
$k$ を本稿の基礎体とし、$(\mathcal C,\mathcal O)\in\HH$ とする。
任意の適合平方根 $s$ に対して
\[
 B=B(\mathcal C,\mathcal O,s),\qquad
 R=R^{\mathrm{as}}(\mathcal C,\mathcal O)
\]
はともに Koszul である。
従ってこの性質に cubic/noncubic の型別仮定は不要である。
\end{proposition}
\begin{proof}
定義より $B\simeq A/(g)$ であり、$A$ は連結 Koszul 環、
$g\in A_2$ は左右正則正規元である。
正則正規二次商が Koszul 性を保つことは
\cite[Chapter 2, Corollary 6.3]{PolishchukPositselski2005}
に含まれる。中心性を必要としないことを明確にするため、
この場合の環変換スペクトル系列による証明を記す。

右 $A$-加群としての分解
\[
 0\longrightarrow A(-2)\xrightarrow{\,g\cdot\,}A
 \longrightarrow B\longrightarrow0
\]
から、$\operatorname{Tor}^A_q(B,k)$ は
$q=0$ で $k$、$q=1$ で $k(-2)$、それ以外で零である。
それぞれ一つの内部次数に集中するので、
正次数の $B$ の作用は零である。
環変換スペクトル系列は
\[
 E^2_{p,q}
 =\operatorname{Tor}^B_p
       \bigl(k,\operatorname{Tor}^A_q(B,k)\bigr)
 \ \Longrightarrow\ \operatorname{Tor}^A_{p+q}(k,k)
\]
となり、$q=0,1$ の二行だけを持つ。
$B$ が Koszul でないと仮定し、
$\operatorname{Tor}^B_p(k,k)$ が内部次数 $p$ に集中しない
最小の $p$ を取る。
連結次数付き環の極小解消では内部次数はホモロジー次数以上なので、
内部次数 $j>p$ の非零元が存在する。
この $E^2_{p,0}$ の元に入る微分はなく、
出る可能性のある微分は
\[
 d_2:E^2_{p,0}\longrightarrow E^2_{p-2,1}
\]
だけである。
しかし最小性により右辺は内部次数 $(p-2)+2=p$ に集中する。
従って内部次数 $j$ の元は終端まで残り、
$A$ の Koszul 性に反する。これで $B$ は Koszul となる。
反対環にも同じ議論が適用できる。

$s^2\simeq S^{-1}[2]$ により $R\simeq B^{(2)}$ である。
Koszul 環の Veronese 部分環は再び Koszul である
\cite[Chapter 3, Proposition 2.2(i)]{PolishchukPositselski2005}。
別の具体表示で Zhang twist を用いる場合にも、
次数付き加群圏の同値が線型自由解消を移すので結論は変わらない。
\end{proof}

\begin{lemma}[Noetherian 基底上の三次平坦性判定]
\label{lem:relative-koszul-degree-three}
$S$ を Noetherian スキーム、$\mathcal V$ を有限階数の局所自由層、
$\mathcal J\subset\mathcal V^{\otimes2}$ を局所直和因子とし、
\[
 \mathcal R=T_{\mathcal O_S}(\mathcal V)/(\mathcal J)
\]
と置く。次を仮定する。
\begin{enumerate}
\item $\mathcal R_3$ は有限局所自由である。
\item 全ての幾何ファイバー $\mathcal R_{\bar s}$ は Koszul である。
\end{enumerate}
このとき全ての $\mathcal R_n$ は有限局所自由である。
さらに、$n\ge2$ に対して
\[
 \mathcal K_n=
 \bigcap_{i=1}^{n-1}
 \mathcal V^{\otimes(i-1)}\otimes\mathcal J
       \otimes\mathcal V^{\otimes(n-i-1)}
 \ \subset\mathcal V^{\otimes n},
 \qquad
 \mathcal K_1=\mathcal V,\quad \mathcal K_0=\mathcal O_S
\]
は局所直和因子であり、任意の底変更と両立する。
左右の Koszul 複体は $\mathcal O_S$ の線型解消となり、
各内部次数で局所分裂完全である。
従ってこれらの解消は任意の底変更後にも完全である。
\end{lemma}
\begin{proof}
まず任意の点 $s\in S$ の局所環 $D=\mathcal O_{S,s}$ で考える。
$V=\mathcal V_s$、$J=\mathcal J_s$ とする。
各次数の $R=T_D(V)/(J)$ は有限生成 $D$-加群であり、
$R_1,R_2,R_3$ は自由である。
Koszul 性は体の拡大に関して忠実平坦に検出されるため、
幾何ファイバーの仮定から $R\otimes_D\kappa(s)$ も Koszul である。
従って
\cite[Chapter 6, Theorem 7.3]{PolishchukPositselski2005}
を適用して全ての $R_n$ が自由となる。
各 $\mathcal R_n$ は有限表示なので、全ての茎での自由性から
有限局所自由性が従う。

解消と底変更についても同定理の証明を用いる。
$W=V^{\otimes n}$ の部分加群
\[
 J_i=V^{\otimes(i-1)}\otimes J
       \otimes V^{\otimes(n-i-1)}\quad(1\le i<n)
\]
を考える。
$W/J_i$ は $R_2$ と $V$ のテンソル積である。
$|i-j|=1$ ならば $W/(J_i+J_j)$ は $R_3$ と $V$ の
テンソル積であり、$|i-j|>1$ ならば二つの $R_2$ と $V$ の
テンソル積である。従ってこれらは全て自由である。
剰余体上では Koszul 性によって関係部分空間の格子が分配的である。
\cite[Chapter 6, Proposition 7.2]{PolishchukPositselski2005}
より、$J_i$ が生成する格子は $D$ 上でも分配的であり、
その全ての商は自由である。
同章 Proposition 7.1 により $W$ は有限個の自由直和因子へ分解でき、
各 $J_i$ はそのうちのいくつかの和となる。
この分解において和と交叉は任意のテンソル積で保存される。
これにより $K_n=\bigcap_iJ_i$ は直和因子となり、
任意の底変更と両立する。
構成は有限表示の層の間のものなので、
各茎の分解を近傍へ延長して同じ結論を $S$ 上で得る。

例えば左 Koszul 複体は
\[
 \cdots\longrightarrow
 \mathcal R\otimes\mathcal K_2
 \longrightarrow \mathcal R\otimes\mathcal K_1
 \longrightarrow\mathcal R\longrightarrow\mathcal O_S
 \longrightarrow0
\]
であり、$\mathcal K_i$ の生成次数を $i$ とする。
各内部次数 $m$ では項
$\mathcal R_{m-i}\otimes\mathcal K_i$ からなる有限な局所自由複体である。
その幾何ファイバーは Koszul 環の解消の対応する次数成分なので完全である。
有限な局所自由複体の完全性は幾何ファイバーで検出されるから、
この複体は完全である。
アフィン開上の有界な完全射影複体は分裂するため、
これは局所分裂完全であり、任意の底変更後も完全である。
右 Koszul 複体についても同じ証明を適用する。
\end{proof}

\begin{corollary}[全幾何ファイバーが $\HH$ 型である関係式族]
\label{cor:h2-geometric-family-koszul}
$S$ を Noetherian $k$-スキームとし、
$S\to X_3$ を命題~\ref{prop:truncated-antiserre-moduli-stack}
の関係式空間への射とする。
その二次環 $\mathcal R$ の全幾何ファイバーが
$\HH$ 対象の反 Serre 環に次数付き同型であると仮定する。
このとき全ての次数で
\[
 \operatorname{rank}\mathcal R_n=(2n+1)^2
\]
であり、$\mathcal R$ は全次数平坦で相対的に Koszul である。
同じ結論は商スタック $[X_3/\operatorname{PGL}_9]$ への射と、
同命題の正規化された普遍次数付き環にも成り立つ。

特に $\HH$ 点を閉ファイバーとする任意の局所 Artin 基底上の
$X_3$ の変形は、全次数平坦な二次代数の変形を与える。
この結論は Artin 基底の冪零次数を制限しない。
\end{corollary}
\begin{proof}
$X_3$ の定義によって次数一、二、三の成分はそれぞれ階数
$9,25,49$ の局所自由層である。
命題~\ref{prop:all-h2-antiserre-koszul} と
補題~\ref{lem:relative-koszul-degree-three} を適用すると、
全次数の局所自由性と相対 Koszul 性を得る。
階数は幾何ファイバー上の Hilbert 関数から決まる。
商スタックの場合は fppf 局所に基底を選び、
同じ議論を適用すればよい。
次数 $n$ の成分への線束の $n$ 乗による twist は
局所自由性と Koszul 複体の完全性を保つ。
これらの性質は忠実平坦に降下するので結論が従う。
局所 Artin 基底には一つの位相的点しかなく、
その幾何ファイバーに命題を適用すれば最後の主張を得る。
\end{proof}

\begin{corollary}[既存の二次 jet の全次数平坦性]
\label{cor:noncubic-second-order-all-degrees}
命題~\ref{prop:noncubic-transverse-second-order} で得られた
$k[t]/(t^3)$ 上の各関係式族 $R_t$ は、
三次成分だけでなく全ての次数で自由であり、
\[
 \operatorname{rank}_{k[t]/(t^3)}(R_t)_n=(2n+1)^2
\]
を満たす。左右の相対 Koszul 解消も存在する。
\end{corollary}
\begin{proof}
同命題の関係部分束は二次成分の自由性を保ち、
Schur complement の等式により三次成分は自由階数 $49$ である。
閉ファイバーは既存の noncubic 反 Serre 環であるから、
直前の系を $S=\operatorname{Spec}k[t]/(t^3)$ に適用する。
\end{proof}

\begin{remark}[全次数平坦性と内在的相対モジュライの区別]
本節により、既存の二次 jet に対して
「四次以上の平坦性も別に検証する」という課題は不要になる。
また、任意の固定次数の Artin 持上げが $X_3$ 内で得られれば、
その持上げの全次数平坦性は自動的である。
一方、$X_3$ 内で次の Artin 段階へ持ち上がること自体は証明していない。
まして、得た環から相対 noetherian Abel heart と相対 Serre 関手を作り、
適合平方根、Artin--Zhang 豊富性、AS 正則被覆を得ることは
本節の結論に含めない。
各幾何ファイバーが $\HH$ 型である条件の有限表示性も、
局所環ごとの Koszul 変形原理からは従わない。
\end{remark}

% !TEX root = ../nc-quadric-classification.tex
% 2026-11-22 / v1.11（相対認識の四命題の証明を論文Bへ移管）（相対flatness剛性補題をBに移管）
\section{相対的な対象の定義と反 Serre による認識}
\label{sec:relative-h2-definition}

本節では $k$ を標数零の代数閉体とし、まず $k$ 上局所有限型の
スキームの big fppf site を用いる。この範囲は非被約スキームと
有限型の局所 Artin 環を含む。任意の noetherian $k$-代数に対する
連続拡張や極限保存性を、本節の構成だけから結論しない。
とくに finite type の基底での局所 noetherian 性と、任意の
noetherian 基底変更での strong noetherian 性を区別する。

\subsection{Abel 圏の基底変更を指定する}

局所有限型の affine スキーム $U=\operatorname{Spec}A$ 上で、
$A$-線形 noetherian Abel 圏を $\mathcal C_U$ とし、
$\mathcal A_U=\operatorname{Ind}\mathcal C_U$ と置く。
$A\to A'$ に対し $(\mathcal A_U)_{A'}$ は
$\mathcal A_U$ 内の $A'$-加群対象の圏、すなわち
$A'$ の作用を備えた対象とその作用を保つ射の圏である。
この構成において引戻しは
$A'\otimes_A-$ であり、一般には右完全、$A'$ が $A$ 上平坦なら完全である。
本稿でいう圏の基底変更は、Hom 加群を単に tensor して対象集合を固定する
線形圏の基底変更ではない。
この区別と基本的な随伴は
\cite[\S4]{LowenVanDenBergh2006}、任意の有限型 affine 基底に必要な
加群対象の定式化と flat descent は
\cite[\S7.1, Definition 7.3, Remarks 7.4--7.5]{AHLHExistence}
を参照する。後者は Artin--Zhang の abstract Hilbert schemes の
B2.2, B3.16, B6.3, C8.6 を明記して用いている。

\begin{definition}[基底変更を備える Abel 圏の族]
\label{def:relative-abelian-system}
$T$ 上の Abel 圏の族とは、$T$ 上の locally finite type affine
$U$ に $\mathcal C_U$ を割り当てる pseudofunctor であって、次を満たすものとする。
\begin{enumerate}
\item 各 $\mathcal C_U$ は $A$-線形で本質的に小さい noetherian Abel 圏であり、
  その Ind 圏は局所 noetherian Grothendieck 圏である。
\item affine の間の射 $U'=\operatorname{Spec}A'\to U$ について
  \[
    \mathcal A_{U'}\simeq(\mathcal A_U)_{A'}
  \]
  という基底変更同値があり、引戻しは $A'\otimes_A-$ に対応する。
  これらの比較同値は合成と恒等射に対して coherent である。
\item 対象、射、完全列について fppf descent が成立する。
  とくに $\mathcal C$ は対象と全ての射を保持する stack である。
\item $\operatorname{Hom}_{\mathcal C_U}(E,F)$ は有限 $A$-加群であり、
  affine flat base change と可換する。
\end{enumerate}
射はこの系の間の、基底変更と可換する線形完全同値とする。
\end{definition}

この定義は圏の局所表示を指定するものではない。
基底変更同値は、族という語の意味を固定するためのデータであり、
射はその coherent な同型を許す。
有限型 affine 基底変更での局所 noetherian 性は
Artin--Zhang B6.3 の範囲内である。
任意の noetherian 基底変更へ、この結論を無条件に拡張しない。

$\mathcal C_U$ の対象 $E$ が $A$ 上平坦とは、
$-\otimes_AE:\operatorname{Mod}A\to\mathcal A_U$ が完全であることをいう。
本稿では
\[
 \mathcal D_U=D(\mathcal A_U)^c
\]
を標準の導来 enhancement における compact 対象の圏とする。
この圏は $\mathcal C$ から決まり、別の dg enhancement を標識に加えない。
以下の導来条件は $T$ の fppf 局所で課す。
すなわち fppf 被覆 $T_i\to T$ が存在して、その各 $T_i$ への
全ての本節の affine 基底変更 $U$ 上で、$D(\mathcal A_U)$ は
compactly generated であり、$\mathcal D_U$ は $D^b(\mathcal C_U)$ に含まれ、
相対 smooth proper な $\operatorname{Perf}A$-線形圏となり、導来基底変更
\[
 \mathcal D_U\otimes_{\operatorname{Perf}A}\operatorname{Perf}A'
       \simeq\mathcal D_{U'}
 \tag{RP-BC}\label{eq:relative-perfect-basechange}
\]
と整合することを要求する。
$\mathcal O$ はこれらの局所 chart 上でこの compact 部分に属するものとする。
chart の選択は標識に含めない。
本節の stack 性の証明には、貼合せ後の大域的な
$D(\operatorname{Ind}\mathcal C_U)$ が局所導来圏の descent と一致するという
追加主張を用いない。定義そのものが fppf 局所条件であることを用いる。
特異な基底上で $\mathcal D_U=D^b(\mathcal C_U)$ としたり、
compact 部分だけに bounded heart を課したりしない。
また全ての底上平坦 coherent 対象が compact であることを
今回の定義には追加しない。この一致は別の相対認識問題である。
体上では、以下の foundation による議論で
$\mathcal D_k=D^b(\mathcal C_k)$ が従う。

以下の Hom、双対、導来 tensor はこの標準 enhancement で解釈する。
(\ref{eq:relative-perfect-basechange}) は導来圏の基底変更であり、
heart の通常の引戻しが全ての基底変更で完全であるという主張ではない。

\begin{definition}[相対 Serre 条件]
\label{def:relative-serre-condition}
上の族が相対 Serre 条件を満たすとは、前記の局所 chart の各 affine $U$ 上で
次が成立し、任意の本節の基底変更と整合することをいう。
\begin{enumerate}
\item compact 対象 $E,F$ の間の $\mathbf R\operatorname{Hom}(E,F)$
  は $A$ 上 perfect で、導来基底変更と可換する。
\item $\mathcal D_U$ に自己同値 $S_U$ が存在し、
  functorial な同型
  \[
   \mathbf R\operatorname{Hom}(E,F)^\vee_A
       \simeq\mathbf R\operatorname{Hom}(F,S_UE)
  \]
  を与える。
\end{enumerate}
Serre 関手および duality の選択自体は標識に含めない。
二つの選択の比較は、表現する双対 functor の Yoneda 一意性による。
\end{definition}

\subsection{相対的な強い\texorpdfstring{$\mathcal H_2$}{H2}条件}

相対 quadratic AS4 被覆とは、局所 affine 基底 $A$ 上の次数付き環
$\mathcal B=T_A(V)/(I)$ であって、$\operatorname{rk}V=4$、
$I\subset V^{\otimes2}$ が階数六の部分束、
各次数が有限局所自由、左右 noetherian、
全幾何 fiber が型 $(1,4,6,4,1)$ の AS 正則環となるものをいう。
加えて、標準の左右 Koszul 複体は全次数で完全であり、
基底変更後にも同じ完全列を与えるものとする。
ここで幾何 fiber の定義には、記述された有限個の係数を含む体からの
体拡大で作る次数付き環を用いる。
被覆の二次の正則正規元 $g$ は、左右乗法の単射性と正規性に加え、
$\mathcal B/(g)$ の次数別局所自由性を要求する。
従って切断の完全列は任意の本節の基底変更と両立する。
固定した相対切断環への被覆の有限判定は定理~\ref{thm:relative-finite-as-cover} に与える。
平方根を選ばず、その存在だけを要求する部分関手の代数性は別問題である。

\begin{definition}[有限型基底上の内在的 $\mathfrak H_2$]
\label{def:relative-h2}
$\mathfrak H_2(T)$ の対象は、定義~\ref{def:relative-abelian-system} の族
$\mathcal C$ と対象 $\mathcal O\in\mathcal C_T$ の対で、以下を満たすものとする。
\begin{enumerate}[label=(RQ\arabic*)]
\item $\mathcal O$ は基底上平坦で、fppf 局所 chart の compact 導来部分は
  (\ref{eq:relative-perfect-basechange}) と相対 Serre 条件を満たす。
  これらの chart の全ての affine 基底変更上で
  $\mathbf R\operatorname{Hom}(\mathcal O,\mathcal O)\simeq A[0]$。
\item 各幾何 fiber の対 $(\mathcal C_{\bar t},\mathcal O_{\bar t})$ は、
  体上の定義~\ref{def:h2} の $\mathcal H_2$ に属する。
  特に標識は定義~\ref{def:gk-rank} の意味で rank one torsion-free であり、
  Serre 条件は fiber の全 $D^b(\mathcal C_{\bar t})$ に成立する。
  fiber 圏は module-object 基底変更、または下記の全軌道環の
  $\qgr$ によって計算し、その一致を要求する。
\item $T$ の fppf 局所で、系全体の線形完全自己同値 $s$ が存在する。
  前記の chart と共通細分を取った局所 compact 導来部分で
  $s^2\simeq S^{-1}[2]$ であり、全ての affine 基底変更上で
  $(\mathcal O,s)$ は Artin--Zhang ample である。
  ここで生成条件は有限直和からの全射、Hom の条件は各全射に対する
  十分大きなシフトでの全射性として課す。
\item 同じ局所的な $s$ について、前記の相対 AS4 環 $\mathcal B$ と
  二次の正則正規元 $g$ が存在し、全次数で
  \[
   \bigoplus_{n\ge0}\operatorname{Hom}(\mathcal O,s^n\mathcal O)
     \simeq\mathcal B/(g)
  \]
  となる。この同型は任意の本節の基底変更と可換する。
\end{enumerate}
射は pointed exact 同値 $(F,u)$、二射は
$u'\eta_{\mathcal O}=u$ を満たす自然同型とする。
$s$、根同型、被覆、切断元、軌道環の表示のいずれも対象の標識ではない。
\end{definition}

この定義は「全幾何 fiber が $\mathcal H_2$ に属する」ことだけを課す定義より強い。
fppf 局所に存在すべき平方根と AS 被覆は、非被約底の変形方向を含めて要求される。
体上の分類と fiberwise な判定だけで、この相対条件を確認したとは扱わない。

\begin{lemma}[Abel flatness と標識付き二射の剛性]
\label{lem:relative-h2-flat-center}
定義~\ref{def:relative-h2} の局所族について、次が成り立つ。
\begin{enumerate}
\item $\mathcal C_U$ は Lowen--Van den Bergh の意味で $A$ 上 flat な Abel 圏である。
\item $\operatorname{Aut}(\operatorname{Id}_{\mathcal C_U})=A^\times$。
  特に pointed 一射の二自己同型は自明である。
\end{enumerate}
\end{lemma}
\begin{proof}
相対 Abel flatness と pointed 2-自己同型の剛性の完全な証明は、
相対モジュライ別稿 \cite[相対的な対象と基底変更]{QuadricModuli2026} に置く。
本稿では後続の相対比較のために主張と仮定を保持する。
\end{proof}

\begin{proposition}[相対的な反 Serre 再構成]
\label{prop:relative-h2-antiserre-reconstruction}
定義~\ref{def:relative-h2} の族について、相対反 Serre 関手の heart 上の代表
$t=S^{-1}[2]$ は局所的な $s^2$ と一致し、選択を除いて降下する。
局所的な軌道環
\[
 R=\bigoplus_{n\ge0}\operatorname{Hom}(\mathcal O,t^n\mathcal O)
\]
は、次数スカラーの比較を許して、基底変更と整合する。
各 $R_n$ は階数 $(2n+1)^2$ の有限局所自由加群であり、
\[
 R\simeq T(R_1)/(J),\qquad
 J=\ker(R_1^{\otimes2}\to R_2),\qquad\operatorname{rk}J=56.
\]
さらに局所 affine 基底で
\[
 (\mathcal C_U,\mathcal O,t)\simeq(\qgr R,\pi R,(1))
\]
が成立する。
\end{proposition}
\begin{proof}
完全な証明は相対モジュライ別稿
\cite[相対反 Serre 再構成と三次表示の認識]{QuadricModuli2026}
に置く。同稿では本稿の体上分類・normal helix・Veronese・foundation の結果を
入力として、相対基底変更と pointed 同型の比較を証明する。
\end{proof}

\begin{proposition}[fppf stack 性と三次表示への全忠実性]
\label{prop:relative-h2-stack-recognition}
定義~\ref{def:relative-h2} の二 groupoid は、二射を除くと
fppf stack in groupoids となる。
相対反 Serre 関係式は全忠実な射
\[
 \Phi:\mathfrak H_2\longrightarrow[X_3/\operatorname{PGL}_9]
\]
を与える。
従ってその像は、同型で飽和した全部分 stack として定義される。
この主張は像の有限表示性・局所閉性・代数性を含まない。
\end{proposition}
\begin{proof}
完全な証明は相対モジュライ別稿
\cite[相対反 Serre 再構成と三次表示の認識]{QuadricModuli2026}
に置く。同稿では本稿の体上分類・normal helix・Veronese・foundation の結果を
入力として、相対基底変更と pointed 同型の比較を証明する。
\end{proof}

\begin{remark}[普遍環の次数別線束と認識条件]
\label{rem:relative-h2-line-twist}
$[X_3/\operatorname{PGL}_9]$ 上の普遍環は
$L=(\det R_1)^{-11}(\det R_2)^2$ を用いて
$R_n^{\mathrm{norm}}=R_n\otimes L^{-n}$ と正規化している。
この正規化は pointed $\qgr$ 圏を変えないが、その次数移動を
元の移動と同じにするわけではない。
自然な同値の下で
\[
 (1)_{R^{\mathrm{norm}}}\simeq t\otimes L^{-1}.
\]
従って普遍環に $(1)=S^{-1}[2]$ を大域に課して像を定義してはならない。
相対反 Serre 環との次数別線束 twist を許すか、fppf 局所でこの条件を課す。
平方根の存在条件自体は、常に真の反 Serre 関手 $t$ に対して課す。
\end{remark}

\begin{proposition}[体上の対象の回復]
\label{prop:relative-h2-field-recovery}
$k$ 上の体としての値を取る評価 functor は
\[
 \mathfrak H_2(k)\simeq\mathcal H_2(k)
\]
という pointed groupoid の同値を与える。
右辺では一射を標識に適合する自然同型で割る。
この主張は任意 noetherian 基底への延長を含まない。
\end{proposition}
\begin{proof}
完全な証明は相対モジュライ別稿
\cite[相対反 Serre 再構成と三次表示の認識]{QuadricModuli2026}
に置く。同稿では本稿の体上分類・normal helix・Veronese・foundation の結果を
入力として、相対基底変更と pointed 同型の比較を証明する。
\end{proof}


\subsection{対角と代数性の判定}

\begin{proposition}[相対認識から対角と代数性判定への還元]
\label{prop:relative-h2-algebraicity-reduction}
$\Phi:\mathfrak H_2\to\mathfrak P=[X_3/\operatorname{PGL}_9]$
が全忠実な fppf スタックの射であるとする。
$X_{\mathrm{as}}=X_3\times_{\mathfrak P}\mathfrak H_2$ と置く。
\begin{enumerate}
\item $\mathfrak H_2$ の対角は affine かつ finite presentation である。
\item $X_{\mathrm{as}}$ は $X_3$ の $\operatorname{PGL}_9$-不変 fppf 部分層であり、
  fppf スタックとして $\mathfrak H_2\simeq
  [X_{\mathrm{as}}/\operatorname{PGL}_9]$ である。
\item 本節の有限型基底の site 上で、$\mathfrak H_2$ が
  $k$ 上有限型の Artin スタックにより表現されるための必要十分条件は、
  $X_{\mathrm{as}}$ が $k$ 上有限型の代数空間により表現されることである。
\end{enumerate}
ここでは $X_{\mathrm{as}}\to X_3$ が開または局所閉 immersion であるとは結論しない。
\end{proposition}
\begin{proof}
完全な証明は相対モジュライ別稿
\cite[相対反 Serre 再構成と三次表示の認識]{QuadricModuli2026}
に置く。同稿では本稿の体上分類・normal helix・Veronese・foundation の結果を
入力として、相対基底変更と pointed 同型の比較を証明する。
\end{proof}


本節は内在的な fppf スタックと代数性の判定を構成した。
この判定に必要な $X_{\mathrm{as}}$ の表現可能性は、
定理~\ref{thm:global-h2-smooth-finite-type} と系~\ref{cor:global-xas-smooth-scheme} で証明する。
同空間は有限型 smooth スキームとなるが、
命題~\ref{prop:global-h2-not-ambient-open} により一致根対では $X_3$ の開部分ではない。
既約包絡の幾何点集合や incidence の像との同一視は用いない。
任意 Noetherian 基底上の素朴な Abel 公理との一致、および既存の foundation 表示との
相対同値は、この代数性判定とは別に検討する。

% !TEX root = ../nc-quadric-classification.tex
% 2026-11-21 / v1.11
\section{Artin 基底での foundation と反 Serre 偏極}
\label{sec:artin-foundation}

この節の基底は、剰余体 $k$ を持つ有限次元可換局所 Artin
$k$-代数 $(A,\mathfrak m,k)$ に限る。
非被約な一般有限型基底への代数化は、この節の結論には含めない。
$(\mathcal C,\mathcal O)\in\HH$ とし、
$\mathcal G=\operatorname{Ind}\mathcal C$ と置く。
flat Abel 変形は Lowen--Van den Bergh の意味で取り、
Grothendieck 圏で記すときは $\mathcal G_A$ とする。

\begin{theorem}[Artin 変形への foundation の持上げ]
\label{thm:artin-foundation-lift}
定理~\ref{thm:normal-helix} の一方の helix $(E_i)_{i\in\mathbb Z}$ を選び、
$E_0=\mathcal O$、$T=\bigoplus_{i=0}^{3}E_i$、
$\Lambda=\operatorname{End}(T)$ と置く。
任意の flat Grothendieck 変形 $\mathcal G_A$ において、次が成立する。
\begin{enumerate}
\item 各 $E_i$ は flat な noetherian 対象 $E_{i,A}$ に持ち上がり、
  その持上げの同型類は一意である。同型そのものの一意性は主張しない。
\item 持上げた列は geometric helix であり、
  $T_A=\bigoplus_{i=0}^{3}E_{i,A}$ は $D(\mathcal G_A)$ の compact tilting generator である。
\item $\Lambda_A=\operatorname{End}(T_A)$ は階数二十四の自由 $A$-代数で、
  $\Lambda_A\otimes_A k\simeq\Lambda$。
  これは補題~\ref{lem:foundation-artin-deformation} の有限 foundation 表示に入る。
\item $D(\mathcal G_A)\simeq D(\operatorname{Mod}\Lambda_A)$、
  $D(\mathcal G_A)^c\simeq\operatorname{Perf}\Lambda_A$ である。
  compact 部分は $A$ 上 smooth proper であり、相対 Serre 関手を持つ。
\end{enumerate}
さらに、閉ファイバーを同定した Abel 変形と有限 foundation 代数の線形変形は、
De Deken--Lowen の変形比較により対応する。
この対応は適合平方根または AS4 被覆の持上げを主張しない。
\end{theorem}
\begin{proof}
完全な証明は相対モジュライ論文 \cite[\texttt{artin_foundation}]{QuadricModuli2026} に置く。
\end{proof}

\begin{lemma}[Artin 基底での exactness の認識]
\label{lem:artin-morita-heart-exact}
定理~\ref{thm:artin-foundation-lift} の表示の下で、$P$ を
$A$-線形の可逆 perfect $\Lambda_A$-双加群とし、
$F_A=-\otimes_{\Lambda_A}^{\mathbf L}P$ を $D(\mathcal G_A)$ の自己同値とする。
その特殊ファイバー $F_0$ が $\mathcal G$ の完全自己同値から誘導されるならば、
$F_A$ は $\mathcal G_A$ の完全自己同値から誘導される。
この自己同値は noetherian 対象および $A$ 上平坦な対象を保つ。
\end{lemma}
\begin{proof}
完全な証明は相対モジュライ論文 \cite[\texttt{artin_foundation}]{QuadricModuli2026} に置く。
\end{proof}

\begin{lemma}[Artin 基底での豊富性の持上げ]
\label{lem:artin-ample-lift}
$F_A$ を前補題の完全自己同値とし、$\mathcal O_A$ を $\mathcal O$ の flat な持上げとする。
$(\mathcal O,F_0)$ が AZ ample であり、全ての $N\in\mathcal C$ について
\[
 \operatorname{Ext}^{j}_{\mathcal C}(F_0^{-n}\mathcal O,N)=0
 \quad(j>0,\ n\gg0)
\]
と仮定する。このとき $(\mathcal O_A,F_A)$ は
$\mathcal C_A=(\mathcal G_A)^{\mathrm{noeth}}$ 上 AZ ample である。
\end{lemma}
\begin{proof}
完全な証明は相対モジュライ論文 \cite[\texttt{artin_foundation}]{QuadricModuli2026} に置く。
\end{proof}

\begin{proposition}[全 Abel 変形の canonical 反 Serre 環]
\label{prop:artin-canonical-antiserre}
任意の flat pointed Abel 変形 $(\mathcal C_A,\mathcal O_A)$ について、
relative compact 圏の反 Serre 関手 $t_A=S_A^{-1}[2]$ は
$\mathcal C_A$ の完全自己同値に延長し、$(\mathcal O_A,t_A)$ は AZ ample である。
さらに
\[
 R_A^{\mathrm{can}}
 =\bigoplus_{n\ge0}\operatorname{Hom}(\mathcal O_A,t_A^n\mathcal O_A)
\]
は全次数自由、階数 $(2n+1)^2$ の二次環である。
生成元の枠と次数線束の標準的な扱いによって $[X_3/\operatorname{PGL}_9]$ の
Artin 変形を与え、
\[
 (\mathcal C_A,\mathcal O_A)
 \simeq(\operatorname{qgr}R_A^{\mathrm{can}},\pi R_A^{\mathrm{can}}).
\]
この構成は Artin 基底変更と両立する。
\end{proposition}
\begin{proof}
完全な証明は相対モジュライ論文 \cite[\texttt{artin_foundation}]{QuadricModuli2026} に置く。
\end{proof}

\begin{proposition}[Artin 関係式族は表示された全軌道環である]
\label{prop:artin-presentation-full-orbit}
命題~\ref{prop:artin-relations-flat-abelian} の $R_A$ について、
$\mathcal O_A=\pi R_A$、$F_A=(1)$ と置くと、
$(\mathcal O_A,F_A)$ は AZ ample であり、自然な次数付き射
\[
 R_A\longrightarrow
 \bigoplus_{n\ge0}\operatorname{Hom}(\mathcal O_A,F_A^n\mathcal O_A)
\]
は全次数で同型である。
\end{proposition}
\begin{proof}
完全な証明は相対モジュライ論文 \cite[\texttt{artin_foundation}]{QuadricModuli2026} に置く。
\end{proof}

\begin{corollary}[形式的な偏極忘却の分裂]
\label{cor:artin-polarization-retraction}
$x\in[X_3/\operatorname{PGL}_9](k)$ を $\HH$ の反 Serre 環の点とする。
Artin 基底上で、関係式から pointed Abel 圏へ移る変形写像
\[
 \mathrm{Def}_{[X_3/\operatorname{PGL}_9],x}
 \longrightarrow\mathrm{Def}_{\mathrm{ab},(\mathcal C,\mathcal O)}
\]
は、canonical 反 Serre 環を取る右逆を持つ。
従って一次変形空間上の写像は分裂全射である。さらに標準的な分裂完全列
\[
 0\longrightarrow\operatorname{HH}^{1}(\Lambda)
 \longrightarrow T_x[X_3/\operatorname{PGL}_9]
 \longrightarrow\operatorname{HH}^{2}(\Lambda)
 \longrightarrow0
 \tag{Artin-pol}\label{eq:artin-polarization-tangent-split}
\]
を得る。ここで中央項は接複体の $H^0$、すなわち閉ファイバーの同定を保つ
一次変形の同型類のベクトル空間である。
二根対一致点ではその核は一次元で、
命題~\ref{prop:antiserre-polarization-tangent} の Zhang 方向が生成する。
\end{corollary}
\begin{proof}
完全な証明は相対モジュライ論文 \cite[\texttt{artin_foundation}]{QuadricModuli2026} に置く。
\end{proof}

\begin{corollary}[Abel 変形の無障害性と三基準点の接次元]
\label{cor:artin-abelian-unobstructed}
$\HH$ の任意の対象の pointed Abel 変形は、Artin 基底の小拡大に対して無障害である。
また
\[
 \dim T_x[X_3/\operatorname{PGL}_9]
 =3+2\dim\operatorname{HH}^{1}(\Lambda).
\]
既存の noncubic 三基準点では、Abel 変形の接次元は $3,3,4$、
偏極を保持した関係式変形の接次元は $3,3,5$ となる。
\end{corollary}
\begin{proof}
完全な証明は相対モジュライ論文 \cite[\texttt{artin_foundation}]{QuadricModuli2026} に置く。
\end{proof}

\begin{corollary}[$\operatorname{HH}^{1}=0$ の点での全 Artin 偏極認識]
\label{cor:artin-rigid-polarization}
$\operatorname{HH}^{1}(\Lambda)=0$ と仮定する。
$x$ を閉ファイバーとする任意の Artin $X_3$-変形について、
表示された shift $F_A=(1)$ は、その pointed Abel 圏の canonical な反 Serre 関手
$t_A=S_A^{-1}[2]$ と自然同型である。
従って偏極忘却は、Artin 変形の同型類の関手上で canonical 反 Serre 環の構成と互いに逆になる。
特に noncubic の二つの根対が異なる全ての点で、この結論が成立する。
根を一つ共有する場合も含む。
\end{corollary}
\begin{proof}
完全な証明は相対モジュライ論文 \cite[\texttt{artin_foundation}]{QuadricModuli2026} に置く。
\end{proof}

\begin{corollary}[非一致根対での関係式空間の滑らかさ]
\label{cor:artin-noncubic-relations-smooth}
$x\in[X_3/\operatorname{PGL}_9](k)$ が、二つの根対の異なる noncubic 対象の
反 Serre 環を表すとする。根を一つ共有する場合も含める。
このとき $[X_3/\operatorname{PGL}_9]$ は $x$ で滑らか、局所次元は三である。
生成元の枠を選んだ対応点 $J\in X_3(k)$ では、$X_3$ は滑らかで局所次元は八十三である。
\end{corollary}
\begin{proof}
完全な証明は相対モジュライ論文 \cite[\texttt{artin_foundation}]{QuadricModuli2026} に置く。
\end{proof}

この滑らかさは $X_3$ とその商についての結論である。
平方根と AS4 被覆の相対存在条件を課す $X_{\mathrm{as}}$ の開性や滑らかさを
結論してはいない。

\begin{remark}
この分裂は Artin 変形の比較である。
canonical な反 Serre 偏極の存在は、適合平方根と AS4 被覆の存在とは異なる。
従って全 Abel 変形が $\mathfrak H_2$ の変形であることや、
$X_{\mathrm{as}}$ の代数性・開性を結論していない。
また、四つの foundation 対象の像だけを比べても偏極は認識できない。
それらは例外対象として一意に持ち上がるため、Zhang 方向でも各像は同型になる。
関手を比較するには、有限 foundation 代数の射への作用との整合も必要である。
\end{remark}

\subsection{Artin 上の平方根・被覆の相対族への昇格}
\label{subsec:artin-relative-promotion}

Artin 基底で対象を作っただけでは、定義~\ref{def:relative-h2} が要求する
全有限型基底変更上の圏の系を確認したことにはならない。
次の命題では、既に平方根と AS 被覆が与えられた場合について、その比較を閉じる。
任意の Abel 圏に対する非平坦基底変更の導来降下定理は用いない。

\begin{proposition}[Artin 上の適合表示から相対 $\mathfrak H_2$ へ]
\label{prop:artin-witness-relative-promotion}
$(A,\mathfrak m,k)$ を有限次元可換局所 Artin $k$-代数とする。
$(\mathcal C_A,\mathcal O_A)$ を $(\mathcal C,\mathcal O)\in\HH$ の
flat pointed Abel 変形とし、次のデータが与えられたと仮定する。
\begin{enumerate}
\item 適合する体上の平方根 $s_0$ の $A$-線形 Morita 持上げ $s_A$。
  有限 foundation の Morita 表示において
  $s_A^2\simeq S_A^{-1}[2]$ という双加群の同型を持つ。
\item 全軌道環の同型
  \[
   B_A:=\bigoplus_{n\ge0}
     \operatorname{Hom}(\mathcal O_A,s_A^n\mathcal O_A)
   \simeq\widetilde B_A/(g_A).
  \]
  ここで $\widetilde B_A$ は各次数有限自由な相対 quadratic AS4 環であり、
  $g_A$ は二次の左右正則正規元、商も全次数 $A$-平坦である。
\end{enumerate}
このとき $s_A$ は heart 上完全で $(\mathcal O_A,s_A)$ は AZ ample である。
さらに、有限型 $A$-代数 $A'$ に対して
\[
 \mathcal C_{A'}=
 \bigl((\operatorname{Ind}\mathcal C_A)_{A'}\bigr)^{\mathrm{noeth}}
 \simeq \operatorname{qgr}(B_A\otimes_A A')
\]
と置き、$\mathcal O_{A'}=A'\otimes_A\mathcal O_A$ とすれば、
これらは定義~\ref{def:relative-h2} の相対族を与える。
従って $(\mathcal C_A,\mathcal O_A)$ は $\mathfrak H_2(A)$ の対象となる。
平方根と被覆はこの相対族の存在を証明するための補助データであり、標識に加えない。
\end{proposition}
\begin{proof}
完全な証明は相対モジュライ論文 \cite[\texttt{artin_foundation}]{QuadricModuli2026} に置く。
\end{proof}

\begin{remark}
本命題は平方根・AS 被覆の\emph{存在}を新たに作っていない。
Artin 基底上でそれらを構成したとき、任意の有限型基底変更を含む相対族の条件を
追加の未証明仮定として残さずに確認するための命題である。
一般の $X_3$-変形に対して AS 被覆が存在するという主張は含まない。
\end{remark}

% !TEX root = ../nc-quadric-classification.tex
% 2026-11-21 / v1.11
\section{有限 foundation 上の反 Serre 偏極の相対認識}
\label{sec:finite-polarization-recognition}

第~\ref{sec:antiserre-polarization}節の Zhang 族が示すように、平坦な pointed Abel 圏を構成することと、
表示された次数移動を内在的反 Serre 関手と同定することは別問題である。
本節では有限 foundation と次数移動の双加群表示が得られている場合に、
後者を非被約基底も含む有限個の方程式と開条件で判定する。
任意の $X_3$ 族からこの有限表示を構成することは、本節の仮定に含めない。

\begin{lemma}[最低次 Hom を線束にする scheme-theoretic 条件]
\label{lem:lowest-hom-line-locus}
$S$ を局所 noetherian スキーム、$K$ を $S$ 上 perfect 複体とする。
各幾何点で $H^i(K_{\bar s})=0$（$i<0$）、
$\dim H^0(K_{\bar s})\le1$ と仮定する。
このとき有限表示閉部分スキーム $j:Z\hookrightarrow S$ が存在して、
次を満たす。
\begin{enumerate}
\item $\mathcal L=H^0(\mathbf L j^*K)$ は線束であり、
その形成は $Z$ 上の任意の基底変更と可換する。
\item $f:T\to S$ が $Z$ を経由することと、全ての
$g:T'\to T$ に対し $H^0(\mathbf L(fg)^*K)$ が線束となることは同値である。
\end{enumerate}
\end{lemma}
\begin{proof}
perfect 複体の局所的な単位成分の消去により、
$K$ は局所的に有限自由加群の複体
$[K^0\xrightarrow{d^0}K^1\to\cdots\to K^b]$ で表せる。
特に負次数の項は不要である。
$\operatorname{rk}K^0=r$ とする。
仮定より、各幾何点で $d^0$ の階数は $r-1$ 以上である。
$r=0$ の領域には求める $Z$ の点がない。
$r\ge1$ の場合は、$d^0$ のある $(r-1)$ 次 minor が可逆となる
開集合で覆える（$r=1$ では空の minor を $1$ とする）。
その開集合上で行列の初等変形を行えば、
\[
 d^0=\begin{pmatrix}I_{r-1}&0\\0&v\end{pmatrix}
\]
と書ける。$v$ は一本の列であり、行数零の場合も許す。
この表示で $Z$ を $v$ の全成分が零となる閉部分スキームとする。
これは幾何点上の階数条件だけでなく、基底環での実際の消滅条件である。
その上では $\ker d^0$ は最後の一座標が張る直和因子であり、
任意の基底変更後も同じ線束を与える。

逆に任意の環への基底変更後、$H^0(K)$ が線束であるとする。
上の表示では、それは $\operatorname{Ann}(v)$ という基底環のイデアルである。
線束は基底環上忠実であり、$v$ の各成分はこのイデアルを零化するから、
各成分自体が零でなければならない。
従ってこの表示では、$H^0$ が線束であるという条件だけでも
上の閉部分スキームを経由する。普遍的特徴付けによって局所構成は貼り合う。
これは \cite[Tag 0BDH, Lemma 36.31.4]{StacksModuli} の
最低次数 $a=0$、階数 $r=1$ の場合の具体的な行列表示でもある。
\end{proof}

\begin{theorem}[scalar 自己同型を持つ双加群の比較部分スキーム]
\label{thm:bimodule-line-comparison}
$S$ を局所 noetherian スキーム、$\Lambda$ を有限局所自由
$\mathcal O_S$-代数とする。任意の $T\to S$ に対して
\[
 Z(\Lambda_T)=\mathcal O_T
\]
と仮定する。$P,Q$ を perfect な $\Lambda$-双加群とし、
それぞれ導来 tensor 積について可逆とする。
このとき $S$ 上の fppf 部分層
\[
 \mathscr E(P,Q)(T)=
 \begin{cases}
 \{*\},&P_T\simeq Q_T\text{ が }T\text{ の fppf 局所で成立する},\\
 \varnothing,&\text{それ以外}
 \end{cases}
\]
は有限表示の局所閉部分スキーム $S_{P\sim Q}\hookrightarrow S$ によって表される。
その上には標準的な線束
\[
 \mathcal L=H^0\mathbf R\mathcal Hom_{\Lambda^e}(P,Q)
\]
と評価同型 $P\otimes\mathcal L\simeq Q$ が存在する。
この線束と同型は、線束の同型を許して一意である。
全ての構成は任意の基底変更と可換する。
\end{theorem}
\begin{proof}
$\Lambda^e=\Lambda\otimes\Lambda^{\mathrm{op}}$ と置く。
$\Lambda^e$ は基底上有限局所自由であり、$P,Q$ はその perfect 加群なので、
\[
 K=\mathbf R\mathcal Hom_{\Lambda^e}(P,Q)
\]
は基底上 perfect で、任意の導来基底変更と可換する。
$P$ の可逆性により
\[
 \mathbf R\mathcal Hom_{\Lambda_T^e}(P_T,P_T)
 \simeq
 \mathbf R\mathcal Hom_{\Lambda_T^e}(\Lambda_T,\Lambda_T).
\]
右辺は負次数 cohomology を持たず、次数零は仮定より $\mathcal O_T$ である。
従って $\operatorname{Aut}(P_T)=\mathbb G_m(T)$ であり、
この scalar 記述は任意の基底変更後にも成立する。

perfect 複体の幾何ファイバーの cohomology 次元の上半連続性により、
\[
 U=\{s\in S\mid H^i(K_{\bar s})=0\ (i<0),\ 
                      \dim H^0(K_{\bar s})\le1\}
\]
は開集合である。比較同型が局所的に存在する任意の $T\to S$ は
$U$ を経由する。実際、その全幾何点で $K$ は $P$ の自己 Hom と同型となる。
上半連続性と局所的な非負次数表示は、\cite[Tag 0BDH, Lemma 36.31.1]{StacksModuli}
の有限自由複体の消去による。

補題~\ref{lem:lowest-hom-line-locus} を $K|_U$ に適用し、
閉部分スキーム $Z\hookrightarrow U$ と線束
$\mathcal L=H^0(K_Z)$ を得る。
負次数がないため、標準写像 $\mathcal L[0]\to K_Z$ があり、
tensor--Hom 随伴により評価射
\[
 e:P_Z\otimes\mathcal L\longrightarrow Q_Z
\]
を定める。その cone は基底上 perfect なので、それが零となる最大の
開部分 $W\subset Z$ が存在する。
この開条件は、有限自由表示の複体が完全になる条件、または幾何ファイバーの
全 Betti 数が零になる条件である。perfect 複体の Nakayama により、
任意の $T\to Z$ に対し $e_T$ が同型であることと、$T\to Z$ が $W$ を経由する
ことが同値となる。従って $W$ は $S$ の有限表示局所閉部分スキームである。

$T\to S$ 上 fppf 局所に $P_T\simeq Q_T$ とする。
この局所同型の下で $K_T$ の次数零 Hom は $\mathcal O_T$ となり、
任意の更なる基底変更でもそうである。線束性と評価同型は忠実平坦降下するので、
$T\to S$ はまず $Z$、次いで $W$ を経由する。
逆に $T\to W$ なら、線束 $\mathcal L_T$ を局所自明化することで
$P_T\simeq Q_T$ を得る。これで $W$ が所要の部分層を表すことが分かった。

任意の線束 $L$ と同型 $P_T\otimes L\simeq Q_T$ について、
相対 Hom を取ると $L\simeq H^0(K_T)$ となり、評価同型へ移る。
scalar な自己同型の記述から、この比較は一意である。
構成の特徴付けは任意の $T\to S$ を用いており、基底変更との整合性も従う。
\end{proof}

\begin{remark}[fppf 像の一般的障害との関係]
ここでは単に有限表示の比較パラメータ空間の像を取ったのではない。
自己 Hom の次数零が基底環そのものであるため、比較が存在する場合の
Hom は線束になり、その普遍性を有限行列の方程式で課すことができた。
同型の空間はこの部分スキーム上の $\mathbb G_m$-torsor となる。
有限型射の任意の fppf 像が代数的になるという主張ではない。
\end{remark}

\begin{corollary}[有限 foundation 上の反 Serre 比較]
\label{cor:foundation-antiserre-comparison-locus}
$\Lambda$ を本稿の相対 foundation 代数とし、
$P$ を表示された次数移動を表す可逆 perfect 双加群とする。
相対反 Serre 関手 $S^{-1}[2]$ の双加群を $Q$ とする。
次数移動が反 Serre 関手と基底線束 twist を除いて一致する条件は、
基底の有限表示局所閉部分スキームで表される。
\end{corollary}
\begin{proof}
foundation 代数の対角頂点成分は $\mathcal O_S$ である。
中心元は冪等元と可換なので頂点ごとの scalar であり、
頂点間の正階数の局所自由成分と可換することから全 scalar は等しい。
この議論は任意の基底変更後にも成立するので $Z(\Lambda_T)=\mathcal O_T$。
有限な相対 bar 解消による smooth proper 性と、仮定した双加群表示により
前定理を適用できる。
\end{proof}

\begin{remark}[次数の正規化と適用範囲]
反 Serre 軌道環の正規化で
$L_{\mathrm{deg}}=(\det R_1)^{-11}(\det R_2)^2$ とすると、
正規化した次数移動は $t\otimes L_{\mathrm{deg}}^{-1}$ を表す。
従って前定理で $P$ をこの次数移動、$Q=t$ とすれば、
比較線束は $L_{\mathrm{deg}}$ である。
大域的な $P\simeq Q$ だけを課してこの線束を落としてはならない。

本節は有限 foundation と次数移動の可逆 perfect 双加群表示を与えた後の
偏極認識を解決する。任意の $X_3$ 族に相対 foundation が存在すること、
heart の次数移動がその表示へ移ること、平方根・AS 被覆の局所存在が
同じ部分スキーム上で成立することは別に確認しなければならない。
従って本節だけから $X_{\mathrm{as}}$ 全体の表現可能性は結論しない。
\end{remark}

% !TEX root = ../nc-quadric-classification.tex
% 2026-11-22 / v1.9（全証明はBへ移管、主張のみ保持）
\section{固定した相対切断環の二次 AS 被覆}
\label{sec:relative-as-cover}

平方根を固定した相対軌道環が既に得られている場合には、
その AS 被覆の存在を有限表示の問題へ帰着できる。
以下では被覆を選ぶ空間と、被覆の存在だけを要求する部分関手を区別する。
特に前者の代数性だけから後者の代数性を結論しない。

\begin{theorem}[二次被覆の相対有限判定]
\label{thm:relative-finite-as-cover}
$S$ を Noetherian スキーム、$\mathcal V$ を階数四の局所自由層、
$\mathcal K\subset\mathcal V^{\otimes2}$ を階数七の局所直和因子とする。
二次環
\[
 \mathcal B=T_{\mathcal O_S}(\mathcal V)/(\mathcal K)
\]
について、各 $\mathcal B_n$ が有限局所自由で階数 $(n+1)^2$ を持ち、
全ての代数閉体からの射 $\bar s\to S$ に対して
$\mathcal B_{\bar s}$ が左右 noetherian な Koszul 環であると仮定する。
さらに $S$ の各アフィン開上で $\mathcal B$ は左右 noetherian とする。
この条件の任意の底変更への保存は仮定しない。

$G=\operatorname{Gr}_S(6,\mathcal K)$ とし、
普遍部分束を $\mathcal J$、商直線束を
$\mathcal L=\mathcal K/\mathcal J$ とする。
$G$ 上の二次環を
\[
 \mathcal A=T(\mathcal V)/\mathcal J,\qquad
 \mathcal D=\mathcal A^!=T(\mathcal V^*)/(\mathcal J^\perp)
\]
と書く。
次の条件を満たす局所閉部分スキーム $E_{\mathcal B}\subset G$ が存在し、
$S$ 上有限表示である。
\begin{enumerate}
\item $\mathcal A_3,\mathcal A_4$ が階数 $20,35$ の局所自由層である。
\item $\mathcal D_3,\mathcal D_4,\mathcal D_5$ がそれぞれ
階数 $4,1,0$ の局所自由層である。
次数零、一、二の階数は既に $1,4,6$ である。
\item 次数三の左右の積の像が、層として
\[
 \operatorname{im}(\mathcal L\otimes\mathcal V\to\mathcal A_3)
 =\operatorname{im}(\mathcal V\otimes\mathcal L\to\mathcal A_3)
\]
である。
\item $i=1,2$ に対する左右の積
\[
 \mathcal L\otimes\mathcal A_i\longrightarrow\mathcal A_{i+2},
 \qquad
 \mathcal A_i\otimes\mathcal L\longrightarrow\mathcal A_{i+2}
\]
が局所直和因子への単射である。
\item $i=1,2,3$ に対する積
$\mathcal D_i\otimes\mathcal D_{4-i}\to\mathcal D_4$
が完全な双線型組合せである。
\end{enumerate}

任意の Noetherian $S$-スキーム $T$ と
$\mathcal J_T\subset\mathcal K_T$ について、
$T\to G$ が $E_{\mathcal B}$ を経由することは、次と同値である。
\begin{enumerate}[label=(\alph*)]
\item $\mathcal A_T=T(\mathcal V_T)/(\mathcal J_T)$ の全次数が
有限局所自由であり、全ての幾何ファイバーが型 $(1,4,6,4,1)$ の
AS 正則環である。
\item $\mathcal L_T\subset(\mathcal A_T)_2$ は、局所生成元を選ぶと
左右正則正規元で生成される。
その正則性は任意の底変更後にも保たれ、
$\mathcal A_T/(\mathcal L_T)\simeq\mathcal B_T$ である。
\end{enumerate}
このとき $\operatorname{rank}(\mathcal A_T)_n=\binom{n+3}{3}$ であり、
左右の相対 Koszul 解消および双加群 Koszul 解消は長さ四で、
任意の底変更と両立する。
また、左右それぞれの相対 AS Gorenstein 条件は
\[
 \underline{\operatorname{Ext}}_{\mathcal A_T}^i
       (\mathcal O_T,\mathcal A_T)
 \simeq
 \begin{cases}
  (\mathcal D_T)_4\otimes\mathcal O_T(4),&i=4,\\
  0,&i\ne4
 \end{cases}
\]
となる。最高次数の直線束は局所的に自明化できる。
ここでの長さ四は相対解消の長さであり、
特異な基底上の環全体の大域次元を四と主張するものではない。

各アフィン開上で $\mathcal A_T$ が左右 noetherian であることは、
$\mathcal B_T$ が左右 noetherian であることと同値である。
従って $\mathcal B$ が全ての Noetherian 底変更の後にも左右 noetherian
であるという追加仮定の下では、$E_{\mathcal B}$ は
固定された $\mathcal B$ への全ての相対 noetherian AS 被覆を表現する。
この追加仮定を置かない場合は、$\mathcal B_T$ が左右 noetherian となる
テスト基底に限って同じ結論を用いる。
\end{theorem}
\begin{proof}
完全な証明は相対モジュライ論文 \cite[相対 AS 被覆の有限判定と fppf 像の反例]{QuadricModuli2026} に置く。
\end{proof}

\begin{remark}[相対被覆の存在として残る問題]
固定された平方根の相対軌道環 $\mathcal B$ に上の定理を適用できるとき、
被覆が fppf 局所に存在することは、有限表示の射
$E_{\mathcal B}\to S$ が fppf 局所に切断を持つことに等しい。
しかし、その射が滑らか・平坦・全射であるとは証明していない。
平方根そのものの相対存在と、得た被覆が相対 Serre 公式に適合することも
この定理の入力の外にある。
従ってこの有限表示への帰着だけを $\HHstack$ の代数性の証明としてはならない。
\end{remark}

\begin{example}[有限型射の fppf 像は代数的とは限らない]
\label{ex:fppf-image-nonalgebraic}
体 $k$ 上の有限型射
\[
 f:\mathbb A^2_{a,c}\longrightarrow\mathbb A^2_{a,b},
 \qquad (a,c)\longmapsto(a,ac)
\]
を考え、その fppf 層としての像を $I$ とする。
$I$ は $\mathbb A^2_{a,b}$ の部分層であるが、
代数空間では表現できない。
\end{example}
\begin{proof}
完全な証明は相対モジュライ論文 \cite[相対 AS 被覆の有限判定と fppf 像の反例]{QuadricModuli2026} に置く。
\end{proof}

この例は、幾何点の構成可能な像や scheme-theoretic image を求めることと、
非被約基底上の fppf 局所存在条件を表すこととの相違を示す。
\emph{この例から $\HHstack$ 自身の非代数性は結論しない。}
本研究では、平方根と被覆を忘れた部分関手の代数性を得るために、
さらに局所切断、相対認識、または同値関係の表示を証明する必要がある。

% !TEX root = ../nc-quadric-classification.tex
% 2026-11-21 / v1.8
\section{反 Serre 偏極の認識と次数付き表示の変形}
\label{sec:antiserre-polarization}

全次数平坦な二次環から相対 Abel 圏を作れることと、
その環の次数移動が圏の内在的反 Serre 関手であることは異なる。
この相違は noncubic の二根対一致点で実際に現れる。
以下では全次数平坦な代数的な族を用い、
pointed Abel 圏を固定しても関係式空間の像が動くことを示す。

\begin{proposition}[pointed Abel 圏が一定となる非自明な二次環の族]
\label{prop:antiserre-polarization-family}
$k$ を標数零の代数閉体とし、
\[
 B=k[u,v]\langle c,d\rangle/(c^2-uv,d^2-uv),\qquad
 \iota(u)=u,\quad\iota(v)=v,\quad\iota(c)=-c,\quad\iota(d)=-d
\]
と置く。
$R=(B^{(2)})^{\iota}$ を
命題~\ref{prop:noncubic-antiserre-descent} の noncubic 反 Serre 環とする。
この表示では、$a\in B_{2m}$、$b\in B_{2n}$ の積は
\[
 a\mathbin{\circ} b=a\,\iota^m(b)
\]
である。
基底 $S=\operatorname{Spec}k[q,q^{-1}]$ 上で
\[
 \sigma_q(u)=qu,\qquad \sigma_q(v)=q^{-1}v,\qquad
 \sigma_q(c)=c,\qquad \sigma_q(d)=d
\]
と定める。
$\sigma_q$ は $\iota$ と可換な次数付き自己同型であり、
$R\otimes k[q,q^{-1}]$ の自己同型でもある。
その\emph{右} Zhang twist を $\mathcal R^{(q)}$ と書き、積を
\[
 a*_q b=\sigma_q^{\deg_R b}(a)\mathbin{\circ}b
 \tag{ZT}\label{eq:polarization-right-zhang}
\]
と定める。
この族は次を満たす。
\begin{enumerate}
\item 全次数有限自由で
$\operatorname{rank}\mathcal R^{(q)}_n=(2n+1)^2$ となり、
左右 noetherian な二次環の族である。
従って $S\to X_3$ および
$S\to[X_3/\operatorname{PGL}_9]$ を与える。
\item Zhang twist の同値により
\[
 (\operatorname{qgr}\mathcal R^{(q)},\pi\mathcal R^{(q)})
 \simeq
 (\operatorname{qgr}(R\otimes k[q,q^{-1}]),\pi(R\otimes k[q,q^{-1}]))
\]
である。この同値は基底上線型、pointed exact であり、
任意の底変更と両立する。
\item $q=1$ のファイバーは $R$ であるが、$q=2$ のファイバーは
$R$ に次数付き同型ではない。
しかも $q=2$ のファイバーは、いかなる $\HH$ 対象の
内在的反 Serre 環にも次数付き同型ではない。
\end{enumerate}
従って $X_3$ には、$\operatorname{qgr}$ と構造対象からは
$\HH$ 対象を与えながら、内在的反 Serre 環の表示ではない幾何点が存在する。
\end{proposition}
\begin{proof}
$uv$ は $\sigma_q$ で固定されるので、二つの二乗関係は保存される。
$\sigma_q$ と $\iota$ の可換性も表示から直接従う。
Zhang twist は次数別の加群を変えないため、平坦性と Hilbert 関数は
一定の族 $R\otimes k[q,q^{-1}]$ から従う。
$R=T(R_1)/(J)$ とすれば、二次関係は
\[
 J_q=(\sigma_q^{-1}\otimes\operatorname{id})(J)
\]
である。全次数の Zhang 同値から、この二次関係が全関係を生成する。
元の $B$ は中心 $k[u,v,h]$（$h=cd+dc$）上で有限である。
従って基底拡大後も左右 noetherian であり、
Veronese および Zhang twist の同値により $\mathcal R^{(q)}$ もそうである。
この議論は Noetherian 基底への底変更後にも適用できる。

次数付き加群の同値は Zhang の構成である
\cite[Theorem 3.1]{Zhang1996}。
係数環と左右の規約を明確にするため、ここでは具体式も記す。
次数 $n$ で $\phi_n=\sigma_q^{-n}$ とすると、
式~(\ref{eq:polarization-right-zhang}) の環は
\[
 a\diamond b=a\mathbin{\circ}\sigma_q^{-\deg_R a}(b)
\]
という通常の左次数による twist に $\phi$ で次数付き同型となる。
この後者の環に対して、次数付き右加群 $M$ の作用を
\[
 m\diamond a=m\,\sigma_q^{-\deg m}(a)
\]
と置けば、同じ次数付き加群上に新しい右作用を得る。
逆自己同型による同じ構成が逆関手となる。
従ってこれは係数環上で exact な同値であり、
自由加群を自由加群へ、torsion 加群を torsion 加群へ送る。
従って表示した pointed exact な $\operatorname{qgr}$ の同値を得る。
全構成は $\sigma_q$ と次数による有限の積から定義されるので、
底変更と両立する。

次数一の中心を比較して最後の主張を示す。
$R_1$ の基底を
\[
 (u^2,uv,v^2,uc,vc,ud,vd,cd,dc)
 \tag{P-basis}\label{eq:polarization-generator-basis}
\]
とする。
元の $R$ では $e=cd-dc$ が非零の中心一次元である。
実際、$e$ は $u,v$ と可換、$c,d$ と反可換であり、
$\iota(e)=e$ である。
従って上の基底の各元に対して、積 $\circ$ に入る $\iota$ の符号が
反可換性の符号を打ち消す。
$cd$ と $dc$ の線型独立性は PBW 基底から従う。

一方、上の九基底の $\sigma_q$-weight は
\[
 (2,0,-2,1,-1,1,-1,0,0)
\]
である。
基底元 $x$ の $\iota$-符号を $\eta_x\in\{1,-1\}$ とすると、
\[
 x*_q(uv)-(uv)*_q x
 =(q^{\operatorname{wt}(x)}-\eta_x)\,uvx.
\]
$q=2$ ではこの係数が零となる基底元は $uv,cd,dc$ だけである。
$B$ の PBW 基底から $uv$ の乗法は単射なので、
中心一次元があれば $z\in\langle uv,cd,dc\rangle$ である。
これらは $\sigma_q$ と $\iota$ で固定され、元の $B$ で $u^2$ と可換だから、
\[
 z*_q u^2-u^2*_q z=(1-q^2)u^2z.
\]
$u^2$ の乗法も単射なので、$q=2$ では $z=0$ となる。
従って
\[
 Z(R)_1\ne0,\qquad Z(\mathcal R^{(2)})_1=0
\]
であり、二つの環は次数付き同型ではない。

もし $\mathcal R^{(2)}$ が別の $\HH$ 対象
$(\mathcal C',\mathcal O')$ の反 Serre 環なら、
反 Serre 豊富性と再構成により
$(\mathcal C',\mathcal O')\simeq
(\operatorname{qgr}\mathcal R^{(2)},\pi\mathcal R^{(2)})$ である。
これは既に示した Zhang 同値によって元の pointed 圏と同値である。
内在的反 Serre 環は pointed exact 同値で不変なので
（定理~\ref{thm:antiserre-invariant}）、
$\mathcal R^{(2)}\simeq R$ となって矛盾する。
\end{proof}

\begin{proposition}[一定の Abel 変形へ移る非零接方向]
\label{prop:antiserre-polarization-tangent}
前命題の族を $q=1+\varepsilon$、$\varepsilon^2=0$ に制限する。
次数零の導分
\[
 \delta(u)=u,\quad\delta(v)=-v,\quad\delta(c)=\delta(d)=0
\]
を用いると、積の一次項は
\[
 \mu_1(a,b)=\deg_R(b)\,\delta(a)\mathbin{\circ}b.
\]
これは $\operatorname{HH}^2(R)_0$ の非零類を与える。
従って対応する $X_3$ の接方向は $\operatorname{PGL}_9$ 軌道の接方向ではない。
一方、対応する pointed Abel 圏の一次変形は自明である。
\end{proposition}
\begin{proof}
族全体が結合的なので $\mu_1$ は Hochschild $2$-cocycle である。
その非自明性は中心元上の交換子で検出できる。
\[
 U=u^4,\qquad V=v^4
\]
はともに $R$-次数二であり、元の $R$ の中心に属する。
実際、元の $B$ で中心、$\iota$ で固定され、
次数二の左因子に現れる $\iota^2$ も恒等写像である。
$\delta(U)=4U$、$\delta(V)=-4V$ なので、
変形した積で
\[
 U*V-V*U
 =\varepsilon\bigl(\mu_1(U,V)-\mu_1(V,U)\bigr)
 =16\varepsilon u^4v^4\ne0.
 \tag{P-comm}\label{eq:polarization-central-commutator}
\]
最後の非零性は中央の $u,v$ と交互語からなる PBW 基底による。

もし $\mu_1$ が Hochschild 境界なら、
$\operatorname{id}+\varepsilon h$ による一次の基底変更で
一定の積へ移せる。
しかし元の中心元 $U,V$ について、任意の一次基底変更が交換子へ与える補正は
\[
 [h(U),V]+[U,h(V)]=0
\]
である。
これは式~(\ref{eq:polarization-central-commutator}) に反する。
従って非零の Hochschild 類であり、特に次数一生成元の基底変更に由来する
$\operatorname{PGL}_9$ 軌道方向でもない。
pointed Abel 変形の自明性は前命題の Zhang 同値を双対数へ制限すれば従う。
\end{proof}

\begin{remark}[反 Serre 認識に対する意味]
上の一次変形では、構造対象付き Abel 圏が一定なので、
その内在的反 Serre 関手と反 Serre 軌道環も一定のものに同定される。
従って $\mathcal R^{(q)}$ の次数一移動を、
この内在的反 Serre 関手と同定することはできない。
全次数平坦性と $\operatorname{qgr}$ の再構成だけでは、
関係式空間から内在的 $\HHstack$ への全忠実性は得られない。
必要なのは、表示された次数移動が相対 Serre 双対から得られる
$S^{-1}[2]$ と一致するという追加の認識である。

中心元の交換子による上の証明は標数零で直接成立し、
接空間の数値比較を仮定しない。
独立な有限線型代数の確認では、既存の Grassmann 座標で
$\operatorname{PGL}_9$ 軌道行列の階数は $79$、
この方向を加えた階数は $80$ となる。
さらに六係数方向も加えた行列の階数は $81$、
そこへこの方向を加えると $82$ となる。
これらは
\texttt{computations/verify\_antiserre\_polarization.py}
による法 $101$ の厳密計算で確認した。
軌道行列の次数零導分 $\delta$ による非零核ベクトルを
整数上でも確認するので、最初の二つの階数 $79,80$ は
標数零でも正確な値である。
\end{remark}

% !TEX root = ../nc-quadric-classification.tex
% 2026-11-21 / v1.11
\section{Serre 平方根の Artin 持上げ}
\label{sec:artin-serre-roots}
\subsection{Artin 基底での Serre 根の平方関係}

以下で $k$ は標数零の体、$(A,\mathfrak m,k)$ は可換局所 Artin
$k$-代数とする。$\Lambda_A$ は有限自由 $A$-代数で、相対的に smooth、
$\Lambda=\Lambda_A\otimes_Ak$ の中心は $k$ と仮定する。
$\operatorname{DPic}_A(\Lambda_A)$ は可逆で、$\Lambda_A^e$ 上 perfect な双加群の同型類の群であり、
積は導来 tensor とする。相対 Serre 関手の逆に $[2]$ を合成した元を
$t_A$ と書く。これは任意の $A$-線形 Morita 自己同値と可換である。
ここでは平方根との同型そのものを標識にしない。

\begin{lemma}[恒等に還元される Morita 自己同値]
\label{lem:artin-infinitesimal-dpic}
還元準同型の核
\[
 K_A=\ker\bigl(\operatorname{DPic}_A(\Lambda_A)
           \longrightarrow\operatorname{DPic}_k(\Lambda)\bigr)
\]
は、次の nilpotent Lie 代数の Baker--Campbell--Hausdorff 群である：
\[
 \frac{\{\delta\in\operatorname{Der}_A(\Lambda_A):
                    \delta(\Lambda_A)\subset\mathfrak m\Lambda_A\}}
      {\operatorname{ad}(\mathfrak m\Lambda_A)}.
\]
特に $K_A$ の任意の元は一意な平方根を持つ。
\end{lemma}
\begin{proof}
まず $P\in K_A$ を右 $\Lambda_A$-加群として考える。
$\Lambda_A$ は有限次元 $k$-代数なので、$P$ は有限生成 projective
右加群の有界 minimal 複体で表せる。
$\mathfrak m\Lambda_A$ は Jacobson radical に含まれ、
還元後もこの複体は minimal である。
還元は次数零の $\Lambda$ と同型なので、minimal 複体の一意性と
nilpotent Nakayama により、元の複体も次数零の projective 加群となる。
projective 加群の nilpotent ideal に沿う持上げの一意性から、これは
右 $\Lambda_A$-加群 $\Lambda_A$ と同型である。
従って双加群 $P$ は環の自己同型 $\alpha$ による twist である。
還元 $P_k\simeq\Lambda$ は $\alpha_k$ が inner であることを意味するので、
単元を持ち上げて共役することで $\alpha\equiv1\pmod{\mathfrak m}$ とできる。

このような自己同型全体を $\operatorname{Aut}_1(\Lambda_A)$ と書く。
二つの同型が同じ双加群を与える差は inner である。
その共役単元 $u$ の還元は $Z(\Lambda)^\times=k^\times$ に属するから、
スカラーで割れば $u\in1+\mathfrak m\Lambda_A$ とできる。従って
\[
 K_A=\operatorname{Aut}_1(\Lambda_A)
           /\operatorname{Ad}(1+\mathfrak m\Lambda_A).
\]

近恒等自己同型について $\alpha-1$ は nilpotent な $A$-線形作用素である。
有限和で定義される $\alpha^x=(1+(\alpha-1))^x$ は、$x$ の多項式として
任意の非負整数 $x$ で積を保つ。標数零なので多項式恒等式により
任意の $x\in k$ で積を保つ。その $x=0$ での微分
$\log\alpha$ は $A$-導分となる。
逆に、像が $\mathfrak m\Lambda_A$ に含まれる導分の有限指数和は
自己同型である。これらは互いに逆で、積は BCH 積に対応する。
また
\[
 \log(\operatorname{Ad}(u))=\operatorname{ad}(\log u),
 \qquad u\in1+\mathfrak m\Lambda_A.
\]
inner 導分は Lie ideal なので、上の商も nilpotent Lie 代数の
BCH 群である。$\exp(v)^2=\exp(2v)$ と $2\in k^\times$ より最後の主張を得る。
\end{proof}

\begin{proposition}[平方関係の補正]
\label{prop:artin-serre-square-correction}
$p_0\in\operatorname{DPic}_k(\Lambda)$ が $p_0^2=t_k$ を満たすとする。
$p_0$ の Morita 自己同値としての持上げ $p\in\operatorname{DPic}_A(\Lambda_A)$
が存在すれば、同じ還元を持つ $s$ で $s^2=t_A$ となるものが存在する。
さらに商 $A\twoheadrightarrow A'$ 上で $p_{A'}^2=t_{A'}$ が既に成立する場合、
$s_{A'}=p_{A'}$ となるように補正できる。
\end{proposition}
\begin{proof}
$d=p^2t_A^{-1}\in K_A$ とする。
Serre 関手の自然性より $p$ と $t_A$ は可換なので、$d$ も $p$ と可換する。
補題の一意な平方根を用いて $e=d^{-1/2}\in K_A$ と置く。
$pep^{-1}$ も $d^{-1}$ の平方根なので、一意性より $pe=ep$ である。
従って $s=ep$ は
\[
 s^2=e^2p^2=d^{-1}p^2=t_A
\]
を満たす。追加条件の下では $e_{A'}^2=1$ であり、$e_{A'}$ は
$K_{A'}$ に属する。同補題の一意性により $e_{A'}=1$ である。
\end{proof}

この命題は underlying Morita 自己同値の持上げが存在することを仮定する。
それ自体の障害を消したという主張ではない。また一般の基底上で
根忘却射が滑らかであるとは主張していない。

\begin{lemma}[合成と持上げ障害]
\label{lem:root-obstruction-composition}
$A\twoheadrightarrow A'$ を核 $I$ が $\mathfrak m I=0$ を満たす
small extension とする。$p'\in\operatorname{DPic}_{A'}(\Lambda_{A'})$
の閉 fiber を $p_0$ とし、$p'^2=t_{A'}$ と仮定する。
その持上げ障害を、一方の Morita 同定によって
\[
 b\in\operatorname{HH}^2(\Lambda)\otimes_kI
\]
と書くと、
\[
 b+(p_0)_*b=0.
\]
特に $(p_0)_*$ が $\operatorname{HH}^2(\Lambda)$ 上恒等なら、
$p'$ は $A$ 上の Serre 平方根へ持ち上がる。
\end{lemma}
\begin{proof}
$p'$ を表す有限 projective $\Lambda_{A'}$-双加群の有界複体を選ぶ。
相対 smoothness により可逆双加群は双加群として perfect である。
nilpotent ideal に沿う projective 加群の持上げによって、その graded
複体と微分を $\Lambda_A^e$ 上へ持ち上げる。
持ち上げた微分 $d$ の $d^2$ は $I$ に値を取り、そのコホモロジー類が
通常の perfect 複体の持上げ障害である。
$\mathfrak m I=0$ により、この群は
$\operatorname{Ext}^2_{\Lambda^e}(p_0,p_0)\otimes_kI$ と同定できる。
類が零なら次数一の項で $d$ を補正して $d^2=0$ とできる。

合成を表す tensor 複体の微分の平方は
\[
 (d\otimes1+1\otimes d)^2=d^2\otimes1+1\otimes d^2
\]
である（複体の tensor の符号規約を含む）。
具体的には、曲率類を $c_{p_0}:p_0\to p_0[2]$ とし、
$b=c_{p_0}\otimes p_0^{-1}$、
$(p_0)_*a=p_0\otimes a\otimes p_0^{-1}$ と同定する。
左右の Morita 同定の比較がこの $(p_0)_*$ なので、
合成 $p'^2$ の障害は $b+(p_0)_*b$ に対応する。
ところが $p'^2\simeq t_{A'}$ には canonical な持上げ $t_A$ があるため、
この障害は零である。

最後の仮定の下では $2b=0$ だから $b=0$ である。
得られる perfect 双加群の持上げは、閉 fiber で可逆なので可逆である。
実際、その左右随伴との評価・余評価の cone は閉 fiber で零であり、
perfect 複体の nilpotent derived Nakayama により零である。
命題~\ref{prop:artin-serre-square-correction} を適用して平方関係を補正する。
この補正は $A'$ 上の既存の根を保つ。
\end{proof}

\begin{corollary}[Hochschild 二次に自明に作用する根]
\label{cor:hhtrivial-root-artin}
$p_0^2=t_k$ とし、$(p_0)_*$ が $\operatorname{HH}^2(\Lambda)$ 上恒等とする。
任意の有限自由かつ相対 smooth な局所 Artin 変形 $\Lambda_A$ 上で、
$p_0$ は $s_A^2=t_A$ を満たす Morita 自己同値へ持ち上がる。
\end{corollary}
\begin{proof}
$A$ を small extension の有限列に分解し、直前の補題を帰納的に適用する。
各段の障害に作用するのは閉 fiber の $(p_0)_*$ であり、
途中の持上げの選択によってこの恒等作用は変わらない。
\end{proof}

\subsection{可換 quadric における根の存在と四つの一次条件}

以下は $Q=\mathbb P^1\times\mathbb P^1$ の full strong foundation を用いた
有限自由相対 smooth Artin Morita 変形についての結論である。
その閉 fiber の中心は $H^0(Q,\mathcal O_Q)=k$ なので、上の議論が適用できる。
Abel 圏の変形への適用には、その foundation の相対持上げと
heart との比較を別途用いる。

\begin{proposition}[恒等幾何型の根は全 Artin Morita 変形に延長する]
\label{prop:quadric-root-all-artin}
$s_0=(-)\otimes\mathcal O_Q(1,1)$ は、$Q$ の任意の上記 Artin Morita 変形へ
相対 Serre 平方根として延長する。
\end{proposition}
\begin{proof}
$s_0^2=(-)\otimes\mathcal O_Q(2,2)=S_Q^{-1}[2]$ である。
また $H^1(T_Q)=H^2(\mathcal O_Q)=0$ なので
\[
 \operatorname{HH}^2(Q)=H^0(Q,\textstyle\bigwedge^2T_Q)
                     =H^0(Q,\mathcal O_Q(2,2)).
\]
線束 tensor は局所的に恒等関手であり、この次数では他の HKR 成分がない。
従って局所 bivector への制限の単射性により、$s_0$ はこの九次元空間に
恒等に作用する。系~\ref{cor:hhtrivial-root-artin} を適用する。
\end{proof}

\begin{corollary}[可換点の全 Abel 変形における exact ample な根]
\label{cor:quadric-exact-root-all-artin}
$(\operatorname{coh}Q,\mathcal O_Q)$ の任意の flat pointed Abel 変形
$(\mathcal C_A,\mathcal O_A)$ には、$s_0=(-)\otimes\mathcal O_Q(1,1)$ を還元とし、
$s_A^2\simeq S_A^{-1}[2]$ を満たす完全自己同値 $s_A$ が存在する。
$(\mathcal O_A,s_A)$ は AZ ample で、その全軌道環は
各次数の階数 $(n+1)^2$ の自由 $A$-加群からなる二次 noetherian 環である。
この主張だけでは AS4 被覆の存在を結論しない。
\end{corollary}
\begin{proof}
定理~\ref{thm:artin-foundation-lift} により有限 foundation の変形に移し、
前命題で Morita 平方根を持ち上げる。
補題~\ref{lem:artin-morita-heart-exact} が heart 上の完全性を与える。
特殊ファイバーの $s_0$ は豊富な線束による tensor なので、
Serre 消滅と補題~\ref{lem:artin-ample-lift} から相対 AZ 豊富性が従う。
$s_A^n\mathcal O_A$ は normal helix の $E_{2n}$ の一意な flat lift であり、
De Deken--Lowen の geometric Hom の平坦性から全軌道の階数と基底変更を得る。
二次性は補題~\ref{lem:relative-quadratic-from-fibers}、
noetherian 性は補題~\ref{lem:artin-noetherian-lift} による。
\end{proof}

\begin{proposition}[四つの根の一次持上げ条件]
\label{prop:quadric-four-root-tangents}
$k=\mathbb C$ とし、$Q$ の一次変形類
$u\in\operatorname{HH}^2(Q)$ を取る。
定理~\ref{thm:commutative-four-roots} の四代表 $s_\varphi$ の underlying
Morita 自己同値が同じ一次変形へ延長するための必要十分条件は
\[
 (s_\varphi)_*u=u.
\]
この条件が成立すれば、平方関係も満たすように持上げを選べる。
固定部分空間の次元は、同定理の順に $9,3,5,3$ である。
\end{proposition}
\begin{proof}
必要性は補題~\ref{lem:autoequivalence-deformation-obstruction}、十分性は
Toda~\cite[Theorem 4.7, arXiv v3]{Toda2009} による。
得られた持上げの平方関係は
命題~\ref{prop:artin-serre-square-correction} で補正できる。

$V=H^0(\mathbb P^1,T_{\mathbb P^1})$ と置くと
$\operatorname{HH}^2(Q)=V\otimes V$ である。
対合 $z\mapsto-z$ は $\partial_z,z\partial_z,z^2\partial_z$ に対し
固有値 $-1,1,-1$ を持つので、$V$ の正負固有空間の次元は $1,2$ である。
従って一因子対合の固定空間は次元 $1\cdot3=3$、
二因子対合では $1\cdot1+2\cdot2=5$ となる。
因子交換は bivector の wedge の順序を逆転するので、$V\otimes V$ には
単なる転置ではなく $v\otimes w\mapsto-w\otimes v$ と作用する。
その固定空間は $\bigwedge^2V$、次元 $3$ である。
恒等幾何型では全空間が固定され、次元は $9$ である。
\end{proof}

従って、可換 quadric で選んだ非自明な根の持上げ障害から、
「いずれの根も存在しない」一次方向を得ることはできない。
恒等幾何型の根は全九次元の一次方向を覆い、さらに上述の有限 Artin
Morita 変形の全段階へ持ち上がる。
この結論自体は二次 AS 正則被覆の存在を含まず、
相対 $\mathcal H_2$ 条件をすべて確認したという主張ではない。


% !TeX root = ../nc-quadric-classification.tex
% 2026-11-21 / v1.32（証明本文をBへ移管）
\section{相対 cubic tails の導来認識と効果性}
\label{sec:relative-cubic-tails}

\subsection{三次解消による相対 tails の compact 認識}

以下では $R$ を任意の可換環とする。$R$ の正則性も
Noetherian 性も、まず仮定しない。
$C$ は正の向きの $R$-線形 $\mathbb Z$-代数、すなわち
$C_{ij}=0$（$j<i$）、$C_{ii}=R$ であり、積を
$C_{ij}\otimes_R C_{jk}\to C_{ik}$ と書くものとする。
右加群 $P_i=e_iC$ の成分は $(P_i)_j=C_{ij}$ である。
$S_i$ は頂点 $i$ にだけ $R$ を置く右加群とする。
次数 torsion は、各元が十分遠い頂点への全ての射で
消されるという意味で用いる。

\begin{theorem}[相対 cubic tails の四項 tilting]
\label{thm:relative-cubic-tails-tilting}
$C$ は隣接成分で生成され、全ての $C_{ij}$ は
有限 projective $R$-加群であるとする。
有限 projective $R$-加群 $V_i,R_i$ と可逆 $R$-加群 $W_i$ により
\[
 0\longrightarrow W_i\otimes_RP_{i+4}
 \longrightarrow R_i\otimes_RP_{i+3}
 \longrightarrow V_i\otimes_RP_{i+1}
 \longrightarrow P_i\longrightarrow S_i\longrightarrow0
 \tag{CT-res}\label{eq:cubic-tails-simple-resolution}
\]
が全ての $i$ で完全であると仮定する。さらに
\[
 \operatorname{Ext}^q_{\operatorname{Gr}C}(S_l,P_j)=0
 \quad\text{unless }q=3\text{ and }j=l+4
 \tag{CT-dual}\label{eq:cubic-tails-relative-dual}
\]
とし、$\mathcal L_l=\operatorname{Ext}^3(S_l,P_{l+4})$ は
可逆 $R$-加群であるとする。
$\mathcal O_i=\pi P_i\in\operatorname{QGr}C$ と置く。
このとき次が成立する。
\begin{enumerate}
\item $\mathbf R\omega$ の cohomological dimension は高々二であり、
  $D(\operatorname{QGr}C)$ は $\mathcal O_i$（$i\in\mathbb Z$）で
  compactly generated される。
\item 任意の $i,j$ について
  \[
  \operatorname{Ext}^q_{\operatorname{QGr}C}
       (\mathcal O_i,\mathcal O_j)
   \simeq
   \begin{cases}
     C_{ji},&q=0,\\
     \operatorname{Hom}_R(C_{i,j-4},\mathcal L_{j-4}),&q=2,\ j\ge i+4,\\
     0,&\text{otherwise}.
   \end{cases}
  \tag{CT-coh}\label{eq:cubic-tails-line-cohomology}
  \]
  とくに任意の連続四項の直和
  $T_a=\bigoplus_{i=a}^{a+3}\mathcal O_i$ は compact tilting generator
  である。その通常の自己準同型環を $\Gamma_a$ とすると
  \[
    D(\operatorname{QGr}C)\simeq D(\operatorname{Mod}\Gamma_a),\qquad
    D(\operatorname{QGr}C)^c\simeq\operatorname{Perf}\Gamma_a.
    \tag{CT-Morita}\label{eq:cubic-tails-canonical-morita}
  \]
\end{enumerate}
全ての仮定と有限 projective 係数の解消が任意の $R$-基底変更で
保たれる場合、これらの compact 比較もその基底変更と整合する。
本定理は $C$ の Noetherian 性や coherent heart の存在を結論しない。
\end{theorem}
\begin{proof}
証明は相対モジュライ論文 \cite[\texttt{relative_cubic_tails}]{QuadricModuli2026} に置く。
\end{proof}

\begin{lemma}[有限型 heart の一様な cohomology]
\label{lem:relative-cubic-coherent-cohomology}
前定理で $R$ を Noetherian とする。
さらに $\operatorname{QGr}C$ が locally Noetherian であり、
$\mathcal O_i$ が Noetherian な生成系をなすと仮定する。
任意の Noetherian 対象 $M$ について
\[
 \operatorname{Ext}^q(\mathcal O_i,M)
 \text{ は有限 }R\text{-加群},\qquad
 \operatorname{Ext}^{q}(\mathcal O_i,M)=0
 \quad(q>0,\ i\gg0).
 \tag{CT-finite}\label{eq:cubic-tails-coherent-finiteness}
\]
任意の Noetherian 対象 $M,N$ の Hom も有限 $R$-加群である。
同じ仮定を満たす平坦基底変更の下では、これらの Hom は基底変更と可換する。
cofinal な偶数添字部分列による生成条件も成立するなら、
その二項移動は Artin--Zhang ample である。
\end{lemma}
\begin{proof}
証明は相対モジュライ論文 \cite[\texttt{relative_cubic_tails}]{QuadricModuli2026} に置く。
\end{proof}

\subsection{右軌道規約と既存 foundation への適用}

ここで既存の $D_{ij}=\operatorname{Hom}(L_i,L_j)$ と右軌道環
$B$ の規約を明示的に合わせる。正の向きの右加群に使う
$\mathbb Z$-代数は
\[
 C_{ij}=D_{-j,-i}=\operatorname{Hom}(L_{-j},L_{-i})\quad(i\le j)
 \tag{CT-right}\label{eq:cubic-tails-right-convention}
\]
であり、積は射の合成 $f\otimes g\mapsto f\circ g$ である。
これは $D$ の反対代数を添字反転したもの。
もとの $D$ の右 tails と、添字を変えず同一視してはならない。
この規約では $\pi P_i$ が $L_{-i}$ に対応し、
$i=-3,-2,-1,0$ の四項窓が元の foundation を与える。

\begin{proposition}[幾何的 cubic 開の canonical Grothendieck 導来比較]
\label{prop:quadric-canonical-cubic-tails}
命題~\ref{prop:quadric-uniform-cubic-hom-open} の有限型開を用いる。
任意の affine 基底変更 $T=\operatorname{Spec}R\to U^\circ$ について
自然な導来 Morita 比較
\[
 D(\operatorname{QGr}B_T)\simeq D(\operatorname{Mod}\Lambda_T),\qquad
 D(\operatorname{QGr}B_T)^c\simeq\operatorname{Perf}\Lambda_T
 \tag{QT-compact}\label{eq:quadric-tails-compact-comparison}
\]
を得る。この比較は構造対象を $F_0$ に、次数移動 $(1)$ を
既存の $\mathsf s$ に対応させ、任意の導来基底変更と整合する。
特に compact 部分は相対 smooth proper であり、
その canonical Serre 関手について
$(1)^2\simeq S^{-1}[2]$ が成立する。
これは Noetherian 性を仮定しない Grothendieck tails の結論である。
Noetherian 基底では、後の定理~\ref{thm:quadric-finite-type-effectivity} が
$B_T$ の左右 Noetherian 性を証明するので、
$\operatorname{QGr}B_T=\operatorname{Ind}(\operatorname{qgr}B_T)$
も成立する。
\end{proposition}
\begin{proof}
証明は相対モジュライ論文 \cite[\texttt{relative_cubic_tails}]{QuadricModuli2026} に置く。
\end{proof}

\begin{corollary}[Noetherian な心上の Hom と豊富性]
\label{cor:quadric-noetherian-base-recognition}
上の affine 基底 $R$ を Noetherian とし、$B_T$ を左右 Noetherian
と仮定する。このとき
$\operatorname{QGr}B_T=\operatorname{Ind}(\operatorname{qgr}B_T)$ であり、
前命題はこの heart に固有の canonical compact 導来圏の比較となる。
一般の coherent 対象間の Hom は有限で、平坦基底変更と可換する。
$\pi B_T$ は基底上平坦かつ exceptional、$(\pi B_T,(1))$ は AZ ample である。
これらの主張は左右 Noetherian 性が保たれる全有限型基底変更で成立する。
特に系~\ref{cor:quadric-point-local-noetherian-as4} の
任意の Noetherian 局所基底 $R$ と、その任意の有限型可換 $R$-代数で
成立する。後の定理~\ref{thm:quadric-finite-type-effectivity} により、
単一の有限型開全体と任意の Noetherian 基底変更にも適用できる。
\end{corollary}
\begin{proof}
証明は相対モジュライ論文 \cite[\texttt{relative_cubic_tails}]{QuadricModuli2026} に置く。
\end{proof}

\begin{lemma}[相対 tails の一様 cohomology から全ての $\chi_q$ へ]
\label{lem:relative-cubic-global-chi}
系~\ref{cor:quadric-noetherian-base-recognition} の Noetherian 基底と
右 Noetherian 環 $B=B_T$ において、任意の有限生成次数付き右加群 $M$ と
任意の $q\ge0$ に対して
\[
   \underline{\operatorname{Ext}}^q_B(R,M)
   \quad\text{は有限 }R\text{-加群である}.
\]
従って右相対 $\chi_q$ が全ての $q$ で成立する。
反対環にも同じ cubic 議論を適用すれば左側も成立する。
この結論は、左右 Noetherian 性がある同じ affine 基底全体での主張である。
\end{lemma}
\begin{proof}
証明は相対モジュライ論文 \cite[\texttt{relative_cubic_tails}]{QuadricModuli2026} に置く。
\end{proof}

\subsection{相対内在的条件への昇格}

\begin{theorem}[Noetherian な基底上での相対条件の全確認]
\label{thm:quadric-noetherian-relative-promotion}
命題~\ref{prop:quadric-canonical-cubic-tails} の幾何的 cubic 開
$U^\circ$ への affine 射 $\operatorname{Spec}R\to U^\circ$ を取り、
$R$ は Noetherian 可換環、$B_R$ は左右 Noetherian 環であると仮定する。
任意の可換有限型 $R$-代数 $R'$ に対して
\[
 \mathcal C_{R'}=\operatorname{qgr}(B_R\otimes_RR'),
 \qquad\mathcal O_{R'}=\pi(B_R\otimes_RR')
\]
と置く。この系は、有限型 $R$-基底の site 上で
定義~\ref{def:relative-abelian-system} および
定義~\ref{def:relative-h2} の全ての条件を満たす。
その局所的な平方根は次数移動 $(1)$ であり、
canonical compact 導来圏は
\[
 D(\operatorname{Ind}\mathcal C_{R'})^c
     \simeq\operatorname{Perf}(\Lambda_R\otimes_RR')
\]
である。この比較は module-object 引戻しの導来関手による
任意の有限型基底変更と整合する。

従って $\operatorname{Spec}R$ 自身が元の局所有限型 $k$-site に属する場合、
$(\mathcal C_R,\mathcal O_R)$ は実際に
$\mathfrak H_2(R)$ の対象を与える。
また系~\ref{cor:quadric-point-local-noetherian-as4} の任意の
Noetherian 局所基底 $R$ は本定理の仮定を満たす。
その $R$ が元の site に属さない場合にも、有限型 $R$-基底の系に関する
上の全条件は成立する。ただしこの記述だけで
$\mathfrak H_2$ の一般基底への極限保存的な延長を定義したとはしない。
\end{theorem}
\begin{proof}
証明は相対モジュライ論文 \cite[\texttt{relative_cubic_tails}]{QuadricModuli2026} に置く。
\end{proof}

% !TeX root = ../nc-quadric-classification.tex
% 2026-11-21 / v1.34
\section{次数別完備化による Noetherian 性の貼合せ}
\label{sec:formal-noetherian-patching}

\paragraph{主張の範囲.}
以下は可換 Noetherian 基底上の非可換次数付き代数と
正の向きの $\mathbb Z$ 代数に対する、証明済みの環論的補題である。
閉部分上とその開補集合上で実際の代数族が Noetherian であることを仮定する。
幾何学的ファイバーの Noetherian 性だけを仮定する主張ではない。
各次数成分の基底上の有限性を用いるが、平坦性、射影性、二次表示、
成長次数、導来同値は仮定しない。新規性は未確認である。

\subsection{中央変数と次数付き有限生成}

\begin{lemma}[中央変数に対する Hilbert 基底論法]
\label{lem:patch-central-variable}
$A$ を環、$M$ を右 Noetherian $A$ 加群とする。
$t$ を中央変数とすれば $M[t]$ は右 Noetherian $A[t]$ 加群である。
内部次数を持つ $A,M$ について、Noetherian 性を次数付き部分加群だけに対して
仮定・結論する形も成立する。この場合 $t$ の内部次数を $0$ とする。
同じ結論は、冪等元 $e_i$ を持つ $\mathbb Z$ 代数 $C$ と
右 $C$ 加群の圏にも成立する。この場合、$t$ は頂点を変えず、
$C[t]_{ij}=C_{ij}[t]$ と定める。
\end{lemma}

\begin{proof}
$N\subset M[t]$ を部分加群とする。$n\geq0$ に対して
$L_n\subset M$ を、$N$ に属し $t$ 次数が $n$ 以下である多項式の
$t^n$ の係数全体とする。これは $A$ 部分加群であり、
$t$ 倍により $L_n\subset L_{n+1}$ である。
$M$ が Noetherian なので、この列はある $n_0$ から安定する。
$0\leq n\leq n_0$ に対して $L_n$ の有限生成元を取り、
それぞれを $N$ の $t$ 次数 $n$ 以下の多項式へ持ち上げる。
この有限個の多項式は $N$ を生成する。
実際、$v\in N$ の最高次係数を消すため、$\deg_t v\leq n_0$ のときは
その次数の持上げを用い、$\deg_t v>n_0$ のときは $n_0$ 次の持上げの
適当な $t$ 冪倍を用いる。差の次数が下がるので、次数に関する帰納法が使える。

内部次数を持つ場合、次数付き $N$ から作る $L_n$ は次数付き部分加群なので、
同じ証明が使える。
$\mathbb Z$ 代数の場合も $L_n$ は $M$ の右部分加群である。
有限生成元を各頂点に同次に取り、多項式の持上げも同じ頂点内で選べば、
まったく同じ帰納法が成り立つ。変数を一つずつ付加することで、
有限個の中央変数についても結論を得る。
\end{proof}

\begin{lemma}[次数付き Noetherian 性から通常の Noetherian 性へ]
\label{lem:patch-graded-to-ungraded}
$A=\bigoplus_{d\geq0}A_d$ を非負次数付き代数とし、
$M$ を下に有界な次数付き右 $A$ 加群とする。
$M$ の次数付き部分加群がすべて有限生成なら、$M$ は
次数を忘れても Noetherian である。
\end{lemma}

\begin{proof}
$N\subset M$ を次数付きとは限らない部分加群とする。
内部次数による増大フィルトレーションを入れると、
$\operatorname{gr} N\subset\operatorname{gr} M=M$ は次数付き部分加群である。
その有限個の同次生成元を $u_1,\ldots,u_s\in N$ に持ち上げる。
任意の $x\in N$ の最高内部次数項を $u_\ell$ の右倍で消せる。
差の最高内部次数は真に下がる。$M$ が下に有界であるから
この操作は有限回で終了し、$u_1,\ldots,u_s$ が $N$ を生成する。
\end{proof}

\subsection{次数ごとの完備性による持上げ}

\begin{lemma}[完備基底上の持上げ]
\label{lem:patch-complete-lifting}
$R$ を可換 Noetherian 環、$I=(f_1,\ldots,f_r)\subset R$ をイデアルとし、
$R$ は $I$ 進完備かつ分離的であると仮定する。
$A=\bigoplus_{d\geq0}A_d$ を中央 $R$ 代数とし、各 $A_d$ は有限 $R$ 加群とする。
このとき $A/IA$ が右 Noetherian なら $A$ も右 Noetherian である。

より一般に、$M$ を各成分が有限 $R$ 加群である次数付き右 $A$ 加群とする。
$M/IM$ が次数付き右 $A/IA$ 加群として Noetherian なら、
$M$ も次数付き右 $A$ 加群として Noetherian である。
$M$ が下に有界なら、次数を忘れても Noetherian である。
\end{lemma}

\begin{proof}
各有限 $R$ 加群 $A_d,M_d$ は $I$ 進完備かつ分離的である。
これは有限加群の完備化が $-\otimes_R\widehat R_I$ と一致することと
$R=\widehat R_I$ から従う
\cite[Tag 00MA]{StacksCompletionPatching}。

$\overline A=A/IA$、$\overline M=M/IM$ と置く。
内部次数と $I$ 進次数を持つ対象に対し、自然な全射
\begin{equation}\label{eq:patch-gr-surjections}
 \overline A[t_1,\ldots,t_r]\twoheadrightarrow\operatorname{gr}_I A,
 \qquad
 \overline M[t_1,\ldots,t_r]\twoheadrightarrow\operatorname{gr}_I M
\end{equation}
がある。$t_\alpha$ は内部次数 $0$、$I$ 進次数 $1$ であり、
$f_\alpha$ の初項へ写す。これらが全射である理由は、
$I^p A_d$ と $I^p M_d$ が、次数 $p$ の $f_\alpha$ の単項式を
係数とする和で生成されることである。
ここでは全射だけを用い、平坦性を仮定しないので同型とは主張しない。
補題~\ref{lem:patch-central-variable} より、
$\operatorname{gr}_I M$ は右 $\operatorname{gr}_I A$ 加群として Noetherian である。

$N\subset M$ を次数付き部分加群とする。
$F^pN=N\cap I^pM$ による誘導フィルトレーションを入れれば
$\operatorname{gr}_F N\subset\operatorname{gr}_I M$ は二重次数付き部分加群である。
有限個の二重同次生成元を取り、それぞれを
\[
 u_\ell\in N_{d_\ell}\cap I^{a_\ell}M_{d_\ell}
 \qquad (1\leq\ell\leq s)
\]
に持ち上げる。これらが $N$ を生成することを示す。

$x\in N_d$ を固定する。$I$ 進次数 $p=0,1,\ldots$ ごとに初項を消せるので、
$a_\ell\leq p$ のとき
\[
 c_{\ell,p}\in I^{p-a_\ell}A_{d-d_\ell},
\]
$a_\ell>p$ のとき $c_{\ell,p}=0$ を選んで、
\begin{equation}\label{eq:patch-approximation}
 x-\sum_{p=0}^{q}\sum_{\ell=1}^s u_\ell c_{\ell,p}
 \in N_d\cap I^{q+1}M_d
 \qquad(q\geq0)
\end{equation}
とできる。存在しない負次数成分は $0$ とする。
固定した $\ell$ について $p-a_\ell\to\infty$ であるから、
$c_\ell=\sum_{p\geq0}c_{\ell,p}$ は有限 $R$ 加群
$A_{d-d_\ell}$ 内で収束する。$M_d$ の分離性と
積の $I$ 進連続性から、\eqref{eq:patch-approximation} の極限は
\[
 x=\sum_{\ell=1}^s u_\ell c_\ell
\]
である。したがって $N$ は有限生成である。
この証明は $N$ 全体の完備性を仮定せず、
固定した内部次数の係数加群の完備性だけを使っている。

$M=A$ とすれば、$A$ のすべての次数付き右イデアルが有限生成である。
補題~\ref{lem:patch-graded-to-ungraded} により通常の右 Noetherian 性を得る。
一般の下に有界な $M$ も同様である。
\end{proof}

\begin{remark}[全体の完備化との区別]
一般の $R$ に対して以下で使う代数は
\[
 A^{\mathrm{deg}\wedge}
 :=A\otimes_R\widehat R_I
 =\bigoplus_{d\geq0}(A_d\otimes_R\widehat R_I)
\]
である。これは各内部次数を個別に完備化した次数付き代数であり、
無限個の内部次数にまたがる元を含み得る環
$\varprojlim_n A/I^nA$ と同一視しない。
補題~\ref{lem:patch-complete-lifting} の収束は一つの内部次数内で行われるため、
$A^{\mathrm{deg}\wedge}$ 全体が $I$ 進完備である必要はない。
\end{remark}

\subsection{閉部分と開補集合からの貼合せ}

\begin{theorem}[Noetherian 性の形式的貼合せ]
\label{thm:patch-noetherian}
$R$ を可換 Noetherian 環、$I=(f_1,\ldots,f_r)\subset R$ とする。
$A=\bigoplus_{d\geq0}A_d$ を各 $A_d$ が有限 $R$ 加群である中央 $R$ 代数とする。
次を仮定する。
\begin{enumerate}
 \item $A/IA$ は右 Noetherian である。
 \item 各 $\alpha$ について $A\otimes_R R_{f_\alpha}$ は右 Noetherian である。
\end{enumerate}
このとき $A$ は右 Noetherian である。左についても同じ結論が成り立つ。

より一般に、各成分が有限 $R$ 加群である次数付き右 $A$ 加群 $M$ について、
$M/IM$ とすべての $M_{f_\alpha}$ がそれぞれ次数付き Noetherian なら、
$M$ も次数付き Noetherian である。
$M$ が下に有界なら通常の右加群としても Noetherian である。
\end{theorem}

\begin{proof}
$\widehat R=\widehat R_I$ と置く。
$\widehat R$ は Noetherian かつ $I\widehat R$ 進完備であり、
$\widehat R/I\widehat R=R/I$ である
\cite[Tags 031C, 05GH]{StacksCompletionPatching}。
$A^{\mathrm{deg}\wedge}/IA^{\mathrm{deg}\wedge}=A/IA$ なので、
補題~\ref{lem:patch-complete-lifting} により
$A^{\mathrm{deg}\wedge}$ は右 Noetherian である。

有限直積
\begin{equation}\label{eq:patch-faithful-cover}
 R^{\ast}=\widehat R_I\times\prod_{\alpha=1}^r R_{f_\alpha}
\end{equation}
を取る。完備化の平坦性 \cite[Tag 00MB]{StacksCompletionPatching} と局所化の平坦性により
$R\to R^{\ast}$ は平坦である。
$\operatorname{Spec}\widehat R_I\to\operatorname{Spec} R$ の像は、
$\widehat R_I/I\widehat R_I=R/I$ より $V(I)$ を含む。
残りは $D(f_\alpha)$ で覆われるから、スペクトル上で全射である。
したがって $R\to R^{\ast}$ は忠実平坦である
\cite[Tag 00HQ]{StacksCompletionPatching}。
なお $I=R$ の場合は開被覆だけを用い、$I=0$ の場合は仮定 (1) で結論が出る。

$A\otimes_R R^{\ast}$ は右 Noetherian 環の有限直積である。
$J\subset A$ を任意の右イデアルとすると、平坦性により
$J\otimes_R R^{\ast}$ は $A\otimes_R R^{\ast}$ の右イデアルであり、有限生成である。
その有限個の生成元をテンソルの有限和として書き、
現れる有限個の $J$ の元で生成される右イデアルを $J_0$ とする。
すると
\[
 (J/J_0)\otimes_R R^{\ast}=0.
\]
忠実平坦性より $J=J_0$ である。したがって $A$ は右 Noetherian である。
左の場合は反対代数に適用する。

一般の $M$ については、補題~\ref{lem:patch-complete-lifting} の加群版で
$M^{\mathrm{deg}\wedge}$ の次数付き Noetherian 性を得て、
任意の次数付き部分加群 $N\subset M$ に上と同じ有限生成の降下を適用する。
最後の主張は補題~\ref{lem:patch-graded-to-ungraded} による。
\end{proof}

\begin{corollary}[閉集合だけを指定した形]
$Z\subset\operatorname{Spec} R$ を閉集合とし、$I$ をその被約閉部分の定義イデアルとする。
$A/IA$ が右 Noetherian であり、$\operatorname{Spec} R\setminus Z$ を覆う有限個の
主開集合上で $A$ が右 Noetherian なら、$A$ は右 Noetherian である。
\end{corollary}

\begin{proof}
有限個の主開集合を $D(g_1),\ldots,D(g_s)$ と書く。
$J=(g_1,\ldots,g_s)$ とすれば $\sqrt J=I$ である。
$R$ が Noetherian なので、ある $m$ について $I^m\subset J$ である。
$A/JA$ は中央の冪零イデアル $IA/JA$ を持ち、その商は $A/IA$ である。
各 $(I^pA+JA)/(I^{p+1}A+JA)$ は有限個の中央係数を用いて
$A/IA$ の有限直和の商となるので、右 Noetherian である。
この有限フィルトレーションにより $A/JA$ は右 Noetherian である。
定理~\ref{thm:patch-noetherian} を $J$ に適用する。
\end{proof}

\subsection{Z 代数の各 row への適用}

\begin{definition}
$C$ を、$C_{ij}=0$ for $j<i$、恒等元 $e_i\in C_{ii}$、
積 $C_{ij}\otimes_R C_{jk}\to C_{ik}$ を持つ中央 $R$ 線形 $\mathbb Z$ 代数とする。
右加群は $M=\bigoplus_jM_j$ と作用 $M_j\otimes_R C_{jk}\to M_k$ からなり、
部分加群はこの頂点分解を保つものとする。
$P_i=e_iC$、すなわち $(P_i)_j=C_{ij}$ と置く。
$C$ が右 Noetherian であるとは、すべての $P_i$ が
この右加群圏で Noetherian であることをいう。
これは有限生成右 $C$ 加群がすべて Noetherian であることと同値である。
\end{definition}

\begin{theorem}[$\mathbb Z$ 代数の形式的貼合せ]
\label{thm:patch-z-algebra}
$R$ と $I=(f_1,\ldots,f_r)$ は定理~\ref{thm:patch-noetherian} の通りとし、
$C$ を各 $C_{ij}$ が有限 $R$ 加群である中央 $R$ 線形 $\mathbb Z$ 代数とする。
$C/IC$ と各 $C_{f_\alpha}$ が右 Noetherian なら、$C$ は右 Noetherian である。
左 Noetherian 性、すなわち各 column $Ce_j$ の Noetherian 性についても同様である。

各 $M_j$ が有限 $R$ 加群である右 $C$ 加群 $M$ に対しても、
$M/IM$ と各 $M_{f_\alpha}$ の Noetherian 性から $M$ の Noetherian 性が従う。
\end{theorem}

\begin{proof}
まず $R$ が $I$ 進完備な場合を考える。
補題~\ref{lem:patch-central-variable} の $\mathbb Z$ 代数版により、各 row
\[
 \overline P_i[t_1,\ldots,t_r]
 \quad\text{は}\quad
 \overline C[t_1,\ldots,t_r]
 \text{ 上 Noetherian である},
\]
ただし $\overline C=C/IC$、$\overline P_i=e_i\overline C$ である。
全射
\[
 \overline C[t_1,\ldots,t_r]\twoheadrightarrow\operatorname{gr}_I C,
 \qquad
 \overline P_i[t_1,\ldots,t_r]\twoheadrightarrow\operatorname{gr}_I P_i
\]
より、$\operatorname{gr}_I P_i$ は右 Noetherian 加群である。

$N\subset P_i$ を右部分加群とし、
$\operatorname{gr}_F N$ の有限生成元を $I$ 進次数と頂点について同次に選ぶ。
それらを
\[
 u_\ell\in N_{j_\ell}\cap I^{a_\ell}C_{i,j_\ell}
\]
に持ち上げる。固定した $x\in N_j$ の $I$ 進初項を逐次消す係数は
\[
 c_{\ell,p}\in I^{p-a_\ell}C_{j_\ell,j}
 \quad(p\geq a_\ell)
\]
に取れる。各 $C_{j_\ell,j}$ は有限 $R$ 加群なので完備であり、
係数の級数が収束する。$C_{ij}$ の分離性から
$x=\sum_\ell u_\ell c_\ell$ を得る。
よって $P_i$ のすべての右部分加群が有限生成である。
ここで $i$ ごとに有限生成元を取っており、$i$ に関する一様な生成次数の上界は不要である。

一般の $R$ では、各 $C_{ij}$ を $\widehat R_I$ に基底変換して
上の完備基底上の結論を得る。
次に \eqref{eq:patch-faithful-cover} を使い、各 $N\subset P_i$ の有限生成を
忠実平坦降下する。テンソルの生成元を展開する際に各元を頂点ごとに分ければ、
有限個の頂点同次生成元を降下できる。
したがってすべての row は Noetherian である。

一般の $M$ も、$\overline M[t_1,\ldots,t_r]\twoheadrightarrow\operatorname{gr}_I M$ を用いて同様に証明できる。
左の場合は各 column に同じ議論を適用すればよい。
\end{proof}

\begin{remark}[普通の環の Noetherian 性との区別]
無限個の直交冪等元を持つ $\bigoplus_{i,j}C_{ij}$ 全体を
単位的な Noetherian 環と呼ぶ主張ではない。
必要なのは各 row、各 column、およびその有限直和の商の Noetherian 性である。
正規族による商 $C/(g)$ を用いる場合も、閉部分上・各主開集合上で得る入力は
この row ごとの Noetherian 性でよい。
正規族の交換関係と正則性から $C/(g)$ の結果を $C$ へ持ち上げるための補題を、
次の定理~\ref{thm:patch-normal-family} で与える。
\end{remark}

\subsection{正則正規線族からの row Noetherian 性の持上げ}

\begin{theorem}[正則正規族による $\mathbb Z$ 代数の拡大]
\label{thm:patch-normal-family}
$R$ を可換環とし、$C$ を正の向きの中央 $R$ 線形 $\mathbb Z$ 代数とし、$C_{ii}=R$ とする。
整数 $\delta>0$ と可逆 $R$ 加群
$G_i\subset C_{i,i+\delta}$ $(i\in\mathbb Z)$ が与えられ、ある両側イデアル $K\subset C$ に対して
乗法が全ての $i,j$ で同型
\begin{equation}\label{eq:patch-normal-family-input}
 G_i\otimes_R C_{i+\delta,j}
       \xrightarrow{\sim}K_{ij}
       \xleftarrow{\sim}C_{i,j-\delta}\otimes_R G_{j-\delta}
\end{equation}
を与えると仮定する。
$\overline C=C/K$ が右 Noetherian なら $C$ も右 Noetherian である。
左についても同様である。
特に有限生成右加群、有限生成左加群はそれぞれ Noetherian になる。

この定理では $G_i$ の周期性、無限個の $G_i$ を同時に自明化する基底の開被覆、
各 $C_{ij}$ の $R$ 上の有限性を仮定しない。
\end{theorem}

\begin{proof}
\emph{(a) 頂点の線束 twist による正規族の自明化.}
各 $\delta$ を法とする剰余類で一つの $L_i$ を $R$ と定め、
全ての整数 $i$ について再帰的に
\begin{equation}\label{eq:patch-vertex-lines}
 L_{i+\delta}=G_i^{-1}\otimes_R L_i
\end{equation}
と定める。各固定した $i$ の $L_i$ は有限個の可逆加群の tensor 積なので、
可逆 $R$ 加群として定義される。
次の $\mathbb Z$ 代数を取る：
\[
 C'_{ij}=L_i^{-1}\otimes_R C_{ij}\otimes_R L_j.
\]
右加群の圏の完全同値は、各頂点で
$M_j\mapsto M_j\otimes_R L_j$ とすることで得られる。
この同値は $P_i$ を $L_i\otimes_R P'_i$ へ送るため、
各 row の Noetherian 性を保つ。
また
\[
 G'_i=L_i^{-1}\otimes_R G_i\otimes_R L_{i+\delta}\simeq R.
\]
これは基底環を変更する操作ではなく、頂点ごとの線束 twist である。
無限個の線束の同時 Zariski 自明化を仮定していない。

従って $C'$ を $C$ と書き直し、全ての $G_i$ に生成元
$g_i\in C_{i,i+\delta}$ がある場合を証明すればよい。
\eqref{eq:patch-normal-family-input} により、左からの $g_i$ 倍と
右からの $g_j$ 倍はそれぞれ単射で、同じ正規イデアルを与える。

\emph{(b) 添字を $\delta$ 移す自己同型.}
$a\in C_{ij}$ に対して、正規性と左からの正則性により一意な
$\sigma(a)\in C_{i+\delta,j+\delta}$ があり、
\begin{equation}\label{eq:patch-normal-shift}
 g_i\sigma(a)=a g_j
\end{equation}
を満たす。
左右の同型 \eqref{eq:patch-normal-family-input} から各 $\sigma:C_{ij}\to C_{i+\delta,j+\delta}$ は同型である。
左から $g_i$ を消去すると
$\sigma(ab)=\sigma(a)\sigma(b)$ が従い、$\sigma(e_i)=e_{i+\delta}$ である。
従って $\sigma$ は頂点を $\delta$ だけ移す $R$ 線形自己同型である。
さらに $a=g_i$ を \eqref{eq:patch-normal-shift} に代入して左から消去すれば
$\sigma(g_i)=g_{i+\delta}$ となる。
従って $\sigma$ は $K$ を保ち、$\overline C$ の自己同型も誘導する。

\emph{(c) 正の skew polynomial $\mathbb Z$ 代数の row Hilbert 基底論法.}
$D=\overline C[t;\sigma]$ を
\[
 D_{ij}=\bigoplus_{p\geq0}\overline C_{i,j-p\delta}\,t^p
\]
と定める。ここで $a t^p$ は、係数 $a\in\overline C_{i,j-p\delta}$ に
頂点 $j-p\delta$ から $j$ への形式的な $p$ 個の $t$ を付した記号である。
積は
\[
 (a t^p)(b t^q)=a\sigma^{-p}(b)t^{p+q}
\]
で定める。$\overline C$ が正の向きなので、固定した $i,j$ について
和に現れる $p$ は有限個である。

$N\subset e_iD$ を右部分加群とする。
各 $n\geq0$ に対して $L_n\subset e_i\overline C$ を、
頂点 $\ell$ の成分が
\[
 (L_n)_\ell=
 \left\{a\in\overline C_{i\ell}:
 \begin{array}{l}
 N_{\ell+n\delta}\text{ のある元が }t\text{ 次数 }n\text{ 以下であり、}\\
 \text{その }t^n\text{ の係数が }a\text{ である}
 \end{array}\right\}
\]
となる部分加群として定める。
これは実際に $\overline C$ 部分加群である。
$a\in(L_n)_\ell$ と $c\in\overline C_{\ell m}$ に対し、
$a$ の持上げに右から $\sigma^n(c)$ を掛ければ、最高係数は $ac$ となるためである。
また右から $t$ を掛けることで $L_n\subset L_{n+1}$ を得る。
$e_i\overline C$ は Noetherian なので、この昇鎖はある $n_0$ から安定する。

$0\leq n\leq n_0$ の各 $L_n$ の有限個の頂点同次生成元を取り、
$N$ の $t$ 次数 $n$ 以下の元へ持ち上げる。
これら有限個の元が $N$ を生成する。
実際、$x\in N_j$ の最高 $t$ 次数を $d$ とする。
$d\leq n_0$ なら $L_d$ の生成元を用い、
$d>n_0$ なら $L_{n_0}$ の持上げを右から $t^{d-n_0}$ 倍して用いる。
最高係数の表示に現れる $c\in\overline C_{\ell,j-d\delta}$ に対し
右から掛ける係数を $\sigma^d(c)$ とすれば、最高係数が正しく消える。
差の $t$ 次数が下がるので、有限回の次数帰納法が完了する。
よって $D$ は各 row について右 Noetherian である。

\emph{(d) $K$ 進次数付き代数.}
$g_i^{(p)}=g_i g_{i+\delta}\cdots g_{i+(p-1)\delta}$ と置く。
正規性から
\[
 K^p_{ij}=g_i^{(p)}C_{i+p\delta,j}
          =C_{i,j-p\delta}g_{j-p\delta}^{(p)}.
\]
両側の正則性により、固定した $p$ の右辺を $K^{p+1}_{ij}$ で割ることは
係数を $K_{i,j-p\delta}$ で割ることと一致する。
従って
\[
 \overline C_{i,j-p\delta}\xrightarrow{\sim}
 K^p_{ij}/K^{p+1}_{ij},
 \qquad
 \overline a\longmapsto
 \overline{a g_{j-p\delta}^{(p)}}
\]
を得る。式 \eqref{eq:patch-normal-shift} が積の整合性を与えるので、
\[
 \operatorname{gr}_K C\simeq\overline C[t;\sigma]=D
\]
である。

\emph{(e) 元の row への有限回の持上げ.}
$N\subset P_i$ に $F^pN=N\cap e_iK^p$ を入れる。
$\operatorname{gr}_F N\subset e_i\operatorname{gr}_K C$ は右有限生成である。
その有限個の $K$ 進次数と頂点についての同次生成元を
$u_1,\ldots,u_s\in N$ に持ち上げる。
固定した $x\in N_j$ の初項をこれらの右倍で逐次消去する。
正の向きと $\delta>0$ から
\[
 K^p_{ij}=0\qquad(p\delta>j-i)
\]
なので、この操作は有限回で終了し、$x\in\sum_\ell u_\ell C$ となる。
従って $P_i$ の全ての部分加群が有限生成であり、$C$ は右 Noetherian である。

左の主張は、添字を反転した反対 $\mathbb Z$ 代数
$C^{\mathrm{op,ref}}_{ij}=C_{-j,-i}$ に右の結果を適用して得る。
この変換は正の向きと、左右両方の
\eqref{eq:patch-normal-family-input} を保つ。
\end{proof}

\begin{remark}[偶数環による証明との関係]
二周期の族では、偶数部分から得る通常の次数付き環で
正則正規な次数二の線束を扱う証明も可能である。
定理~\ref{thm:patch-normal-family} はその周期性を仮定しないので、
根を選んでいない geometric cubic foundation の族にも適用できる。
必要な幾何的入力は、相対曲線商 $\overline C$ の row Noetherian 性と、
正則正規族の普遍的同型 \eqref{eq:patch-normal-family-input} である。
この二つをファイバーだけでなく実際の族として証明する必要は残る。
\end{remark}

\subsection{適用時に確認すべき入力}

閉部分 $Z\subset U=\operatorname{Spec} R$ と開補集合 $U\setminus Z$ に分ける適用では、
必要な入力は次の有限個の主張である。
\begin{enumerate}
 \item $R$ は Noetherian であり、すべての構造成分 $A_d$ または $C_{ij}$ は有限 $R$ 加群である。
 \item 閉部分を定義する実際の商環 $R/I$ 上で、$A/IA$ または $C/IC$ が Noetherian である。
 \item 開補集合を覆う有限個の主開集合 $D(f_\alpha)$ 上で、代数全体または各 row が Noetherian である。
\end{enumerate}
この三条件から結論は上の定理で従う。
Noetherian 性を各点の局所環で証明した後、任意のイデアルの生成元が
自動的に一点の近傍で生成を続ける、と推論する必要はない。
形式的貼合せは、閉部分に沿う完備化という一つの環上で全ての次数を扱い、
有限個の開局所化と合わせた忠実平坦被覆で有限生成を降下する。


% !TeX root = ../nc-quadric-classification.tex
% 2026-11-21 / v1.34（証明をBへ移管）
\section{相対 cubic 環の普遍 Noetherian 性と有限型効果性}
\label{sec:relative-cubic-noetherian}

\paragraph{本稿の役割.}
全次数の局所自由性と標準 cubic 解消が既に構成された族に対し、
体上の Noetherian 性を各局所環へ持ち上げるだけでなく、
一つの Noetherian 基底上の全 projective row と column の Noetherian 性を導く。
周期性や Serre 根の選択は仮定せず、二周期の場合には通常の偶数環の左右 Noetherian 性も従う。
相対 point curve と正規族の構成を明記し、楕円部分の Noetherian 性には
Keeler の相対 Fujita 消滅を用いる。
linear 閉部分はテンソルの rank-one flattening により別に扱う。
最後に第\ref{sec:formal-noetherian-patching}節の形式的貼合せを適用する。
本稿の相対化および貼合せとの組合せの新規性は未確認である。

\subsection{仮定と結論}

$k$ を研究で固定した体、$R$ を可換 Noetherian $k$-代数、$S=\operatorname{Spec} R$ とする。
幾何ファイバーでは常に代数閉包へ移る。
正の $R$-線形 $\mathbb Z$-代数 $D$ について次を仮定する。
\begin{enumerate}
\item $D_{ii}=R$, $D_{ij}=0$ $(j<i)$ であり、全ての $D_{ij}$ は有限 projective $R$-加群である。
\item $V_i=D_{i,i+1}$ は階数二で、$D$ は $V_i$ により生成される。
  定義関係は階数二の直和因子
  $R_i\subset V_i\otimes V_{i+1}\otimes V_{i+2}$ により生成される。
\item 四次 syzygy は可逆加群 $W_i$ であり、普遍的に
  \[
  W_i\subset (R_i\otimes V_{i+3})\cap(V_i\otimes R_{i+1})
  \subset V_i\otimes\cdots\otimes V_{i+3}
  \]
  となる。縮約は同型
  \begin{equation}\label{noeth-eq:contraction}
  W_i\otimes V_{i+3}^{\vee}\simeq R_i,
  \qquad W_i\otimes V_i^{\vee}\simeq R_{i+1}
  \end{equation}
  を与える。
\item 全幾何ファイバーの cubic 代数は幾何的 quintuple に対応する AS 正則代数である。
  特に $W_i$ の任意の非零隣接二縮約は零でない。
\end{enumerate}
これらは既に得られた有限型 geometric cubic 開の相対 cubic 解消から供給する入力である。
\begin{theorem}[周期性を仮定しない普遍 Noetherian 性]
\label{thm:relative-geometric-cubic-noetherian}
上の仮定の下で、$D$ の全ての projective row と column は Noetherian である。
偶数添字に制限した $D^{\mathrm{ev}}_{ij}=D_{2i,2j}$ も同じ性質を持つ。
これらの結論は任意の可換 Noetherian $R$-代数 $R'$ への基底変更後にも成立する。
\end{theorem}

\begin{corollary}[二周期の偶数環]
\label{cor:relative-periodic-even-noetherian}
上の $D$ に積と整合する二周期同型 $\theta:D\xrightarrow{\sim}D(2)$ が指定されているとする。
\begin{equation}\label{noeth-eq:even-ring}
 B_n=D_{0,2n},\qquad
 a\star b=a\,\theta^m(b)\quad(a\in B_m, b\in B_n)
\end{equation}
と置くと、$B$ は左右 Noetherian である。
さらに任意の可換 Noetherian $R$-代数 $R'$ に対して
$B\otimes_RR'$ も左右 Noetherian である。
\end{corollary}

\begin{proof}
全証明は相対モジュライ論文 \cite[\texttt{relative_cubic_noetherian}]{QuadricModuli2026} に置く。
\end{proof}

以下ではまず nonlinear 部分で相対曲線と普遍正規族を構成し、
linear 閉部分と合わせて定理を証明する。
stalk ごとの Noetherian 性から開近傍上の Noetherian 性を推論する箇所はない。

\subsection{linear 閉部分と相対 point curve}

\subsubsection{閉部分の定義}

局所的に $W_0$ の生成元を $w$ とし、
\[
 w\in (V_0\otimes V_2)\otimes(V_1\otimes V_3)
\]
を $02\mid13$ flattening の $4\times4$ 行列として読む。
その $2\times2$ 小行列式が生成するイデアルを $I\subset R$ とする。
基底と $W_0$ の生成元の変更はこのイデアルを変えない。
\[
 Z=V(I),\qquad S_{\mathrm e}=S\setminus Z
\]
と置く。Van den Bergh の linear tensor の標準形
\cite[Lemma 5.2.4, (5.3)]{VdB2011} により、
$Z$ の幾何点はちょうど linear quintuple の点である。
実際、その標準形は
\[
 (x_0y_2-y_0x_2)(x_1y_3-y_1x_3)
\]
である。

これは reduced locus の定義だけではない。
$02\mid13$ の rank-one determinantal scheme は Segre scheme そのものであり、
$W_0$ が至る所非零であるため、$Z$ 上では可逆部分束
$A_0\subset V_0\otimes V_2$ と $A_1\subset V_1\otimes V_3$ が存在して
$W_0=A_0\otimes A_1$ となる。
局所的には一つの可逆な成分と小行列式の方程式を用いて
\begin{equation}\label{noeth-eq:linear-factorization}
 w=a_{02}\otimes b_{13}
\end{equation}
と因数分解でき、その二つの像線は一意なので重なりで貼り合わさる。
両双線形形式は非退化である。
もし一方がある幾何点で退化すれば、その零方向と隣接因子の任意の非零方向による
縮約が零となり、geometric 条件に反する。
従ってそれらの行列式は可逆である。
この議論は $Z$ が非被約の場合も適用できる。

\subsubsection{楕円部分の行列式}

以下しばらく $S_{\mathrm e}$ の任意の affine open を $S$ と書き直す。
$P_i=\PP(V_i^{\vee})$ とし、$P_i\times_S P_{i+1}$ 上で $W_i$ を最初の二因子により評価すると
\begin{equation}\label{noeth-eq:relation-matrix}
 \mu_i:W_i\otimes V_{i+3}^{\vee}\otimes\OO(-1,-1)
 \longrightarrow V_{i+2}\otimes\OO
\end{equation}
を得る。
geometric 条件により、これは全ての幾何点で零行列でない。
また nonlinear 条件と \cite[\S5.4, 特に Lemma 5.4.10]{VdB2011} により、
その行列式は各基底ファイバー上で恒等的に零ではない。
従って
\[
 C_i=V(\det\mu_i)\subset P_i\times_S P_{i+1}
\]
は相対有効 Cartier 因子であり、$S$ 上平坦な $(2,2)$ 曲線族である。
各幾何ファイバーは算術種数一の Gorenstein 曲線である。
ここで非被約性や可約性を除外しない。

\begin{lemma}\label{noeth-lem:scheme-shift}
$R_i$ の二つの multilinear relation により定義される
\[
 \Gamma_i\subset P_i\times_S P_{i+1}\times_S P_{i+2}
\]
の二つの隣接射影は同型
\[
 \Gamma_i\xrightarrow{\sim}C_i,
 \qquad \Gamma_i\xrightarrow{\sim}C_{i+1}
\]
を与える。従って canonical な同型 $\psi_i:C_i\xrightarrow{\sim}C_{i+1}$ がある。
これらは任意の基底変更と可換する。
\end{lemma}

\begin{proof}
全証明は相対モジュライ論文 \cite[\texttt{relative_cubic_noetherian}]{QuadricModuli2026} に置く。
\end{proof}

\begin{remark}[二種類の $(2,2)$ 行列式を区別する]
ここで使う $C_i$ は\emph{隣接因子} $(i,i+1)$ 上の point curve である。
その各 $\PP^1$ への射影が有限であることは仮定せず、一般に成分を縮約し得る。
必要なのは隣接二射影を合わせた閉埋め込みと、$\OO(1,1)|_{C_i}$ の相対豊富性である。
これに対し、層双加群を作る際の\emph{反対因子} $(0,2)$ 上の行列式は
geometric 条件により両射影有限であるが、上の point curve と一様には同じでない。
特に linear 点では point scheme は曲面 $\PP^1\times\PP^1$ となり、
隣接行列式は恒等的に零となる。
従って本節の flat curve を linear 点までそのまま延長してはいけない。
\end{remark}

\subsection{相対 helix 座標代数と普遍正規族}

補題 \ref{noeth-lem:scheme-shift} により全 $C_i$ を $C=C_0$ と同定する。
その座標射を $\rho_i:C\to P_i$ とし、
\[
 \pi:C\to S,\qquad \mathcal L_i=\rho_i^*\OO_{P_i}(1),\qquad
 \mathcal L_{ij}=\mathcal L_i\otimes\cdots\otimes\mathcal L_{j-1}\quad(j>i)
\]
と置く。
幾何ファイバーではこれは Van den Bergh の admissible quadruple の elliptic helix そのものである
\cite[\S5.4, Lemma 5.5.1, Theorem 5.5.10]{VdB2011}。
従って各ファイバーの各既約成分 $E$ について
\begin{equation}\label{noeth-eq:parity-degrees}
 \deg(\mathcal L_{i+2}|_E)=\deg(\mathcal L_i|_E)
\end{equation}
である。
この数値等式は非被約ファイバーでもその reduced 成分上で読む。
本稿は相対線束としての helix 関係を fiberwise 同型から推論しない。
必要なのは \eqref{noeth-eq:parity-degrees} のみである。

\begin{lemma}\label{noeth-lem:cohomology-bc}
$j>i$ に対し
\[
 H_{ij}:=\pi_*\mathcal L_{ij}
\]
は階数 $2(j-i)$ の有限 locally free 層であり、
$R^{q}\pi_*\mathcal L_{ij}=0$ $(q>0)$ である。
この pushforward は任意の基底変更と可換する。
また $V_i\to\pi_*\mathcal L_i$ は同型で、$\mathcal L_i$ は相対大域生成される。
\end{lemma}

\begin{proof}
全証明は相対モジュライ論文 \cite[\texttt{relative_cubic_noetherian}]{QuadricModuli2026} に置く。
\end{proof}

$H_{ii}=R$, $H_{ij}=0$ $(j<i)$ とし、線束の積により $H$ を $\mathbb Z$-代数とする。
$\pi_*\OO_C=R$ は $(2,2)$ 因子の標準完全列から直接従う。
各 $R_i$ は $C$ 上で消えるので、canonical な準同型
\[
 D\longrightarrow H
\]
が得られる。

\begin{proposition}\label{noeth-prop:normal-family}
上の準同型は全成分で全射であり、核を $K$ とすると全 $K_{ij}$ は有限 projective である。
$K_{ij}=0$ $(j-i<4)$、$G_i:=K_{i,i+4}$ は可逆加群である。
さらに積は普遍同型
\begin{equation}\label{noeth-eq:normal-family}
 G_i\otimes_R D_{i+4,j}\xrightarrow{\sim}K_{ij}
 \xleftarrow{\sim}D_{i,j-4}\otimes_R G_{j-4}
 \qquad(j\ge i+4)
\end{equation}
を与える。従ってこれは左右 regular な正規族である。
構成、核、同型の全てが任意の基底変更と可換する。
\end{proposition}

\begin{proof}
全証明は相対モジュライ論文 \cite[\texttt{relative_cubic_noetherian}]{QuadricModuli2026} に置く。
\end{proof}

特に両側の表示を比較すると canonical な同型
\begin{equation}\label{noeth-eq:normal-commutation}
 G_i\otimes D_{i+4,j+4}\simeq D_{ij}\otimes G_j
\end{equation}
を得る。これが基底上の線束を保った正規交換式である。

\subsubsection{二周期性と通常の次数付き正規線}

この小節だけ、系~\ref{cor:relative-periodic-even-noetherian} の二周期同型 $\theta$ を仮定する。
相対 point scheme の構成は $D$ とその添字の同型に関して functorial である。
従って $\theta$ は $H$, $K$, $G_i$ にも二周期同型を誘導し、
$D\to H$ および \eqref{noeth-eq:normal-family} と整合する。
具体的には、$\theta$ が与える $C_2\to C_0$ と
$\psi_1\psi_0:C_0\to C_2$ の合成を $\sigma:C\to C$ とすれば、
\[
 \sigma^*\mathcal L_i\simeq\mathcal L_{i+2}
\]
となる。これは座標の tautological line bundle と $\theta$ の可換図式から得られる。

$\overline B_n=H_{0,2n}$ に \eqref{noeth-eq:even-ring} と同じ二周期積を入れると、
$\overline B$ は $\mathcal L_0\mathcal L_1$ と $\sigma$ による twisted homogeneous coordinate ring である。
また $G=G_0\subset B_2$ とすると
\begin{equation}\label{noeth-eq:graded-normal-line}
 \ker(B\to\overline B)=GB=BG,
 \qquad G\otimes B(-2)\simeq\ker(B\to\overline B)\simeq B(-2)\otimes G.
\end{equation}
これは $G_i$ の二周期同一視を \eqref{noeth-eq:normal-family} に適用したものである。
$G$ を自明化する開集合では、その生成元 $g\in B_2$ は左右 regular な正規元である。
全ての主張は任意の基底変更後も成立する。

\subsection{相対豊富性と座標代数の Noetherian 性}

\subsubsection{Keeler の消滅の適用}

用いる一次文献の正確な形は次である。
Keeler \cite[Theorem 1.5]{Keeler2003} は、可換 Noetherian 環上の射影スキーム、
豊富な線束 $H$、任意の coherent 層 $F$ に対し、ある $m_0$ が存在して
\begin{equation}\label{noeth-eq:keeler}
 H^q(F\otimes H^m\otimes N)=0
 \quad(q>0, m\ge m_0)
\end{equation}
が全ての relatively nef 線束 $N$ について成立することを与える。
確認した arXiv v5 には erratum が添付され、そこで Lemma 3.3 の修正後にも
Theorem 1.5 が成立することが再確認されている
\cite[付属 Erratum, \S2]{Keeler2003}。

固定した $i$ に対し $H_i=\mathcal L_i\mathcal L_{i+1}$ は相対豊富である。
実際、これは隣接二座標による閉埋め込みの $\OO(1,1)$ である。
式 \eqref{noeth-eq:parity-degrees} から
\begin{equation}\label{noeth-eq:nef-prefix}
 \mathcal L_{i,i+2m}\otimes H_i^{-m}
\end{equation}
は各幾何ファイバーの各既約成分上で次数零であり、relatively numerically trivial、
従って relatively nef である。
奇数長の積ではさらに大域生成線束一つを tensor するので、残余因子はやはり nef である。
逆方向の積 $\mathcal L_{i-2m,i}$ についても同じ議論が働く。
従って \eqref{noeth-eq:keeler} により、前向き・後向き双方の積は ample sequence を成す。
これは全ての affine Noetherian 基底変更後にも適用できる。

\subsubsection{row ideal の有限生成を直接証明する}

\begin{lemma}\label{noeth-lem:row-noetherian}
$H$ は左右 Noetherian な $R$-線形 $\mathbb Z$-代数である。
偶数添字に制限した $H^{\mathrm{ev}}$ も左右 Noetherian である。
二周期同型が指定されている場合、その二周期環 $\overline B$ も左右 Noetherian である。
\end{lemma}

\begin{proof}
全証明は相対モジュライ論文 \cite[\texttt{relative_cubic_noetherian}]{QuadricModuli2026} に置く。
\end{proof}

\subsubsection{正規線からの持上げ}

\begin{lemma}\label{noeth-lem:normal-lift}
$A$ を非負次数付き $R$-代数とし、正次数にある可逆 $R$-部分加群 $G\subset A_d$ が
左右 regular な正規線を成すとする。
すなわち $GA=AG$ であり、積が
$G\otimes A(-d)\simeq GA\simeq A(-d)\otimes G$ を与えるとする。
$A/GA$ が左右 Noetherian なら $A$ も左右 Noetherian である。
\end{lemma}

\begin{proof}
全証明は相対モジュライ論文 \cite[\texttt{relative_cubic_noetherian}]{QuadricModuli2026} に置く。
\end{proof}

命題~\ref{noeth-prop:normal-family} の正則正規族と
補題~\ref{noeth-lem:row-noetherian} に、周期性を仮定しない
定理~\ref{thm:patch-normal-family} を $\delta=4$ として適用する。
これにより $S_{\mathrm e}$ 上の $D$ の全ての row と column は Noetherian である。
無限個の正規線 $G_i$ は同定や自明化を仮定せず、同定理の頂点線束 twist で扱う。
同じ構成と Keeler の定理は任意の可換 Noetherian 基底変更後に適用できるので、
楕円 open 上で普遍 Noetherian 性が得られる。
補題~\ref{noeth-lem:normal-lift} の通常の次数付き環としての形は、後で AS4 被覆に用いる。

\subsection{linear 閉部分}

ここでは再び元の $S$ に戻り、$Z=V(I)$ 上の代数を考える。
この小節では $R/I$ を $R$ と書き、全ての cubic データを $Z$ 上へ制限する。
式~\eqref{noeth-eq:linear-factorization} の二つの非退化線束を $A_0,A_1$ と書く。
$V_3^\vee$ による縮約から $R_0=A_0\otimes V_1$ を得る。
以下では二周期同型を用いず、普遍的な縮約同型~\eqref{noeth-eq:contraction} により
この表示を全整数へ伝播し、線束の twist を明示する。
添字順と異なる順に書いた tensor は、元の添字順への shuffle を介して読む。

\subsubsection{linear relation の両方向への伝播}

\begin{lemma}[一つの線型関係からの伝播]
\label{lem:linear-propagation-without-period}
ある $i$ において非退化な可逆部分束
$A_i\subset V_i\otimes_R V_{i+2}$ があり、
添字順への shuffle を介して
\[
 R_i=A_i\otimes_R V_{i+1}
\]
であるとする。このとき、一意な可逆部分束
$A_j\subset V_j\otimes_R V_{j+2}$ $(j\in\mathbb Z)$ が存在し、全て非退化で、
\begin{equation}\label{eq:linear-all-relations}
 R_j=A_j\otimes_R V_{j+1},\qquad
 W_j=A_j\otimes_R A_{j+1}
\end{equation}
となる。この構成は任意の基底変更と可換する。
ここで非退化とは、縮約
$A_j\otimes_R V_j^\vee\to V_{j+2}$ が同型であることをいう。
\end{lemma}

\begin{proof}
全証明は相対モジュライ論文 \cite[\texttt{relative_cubic_noetherian}]{QuadricModuli2026} に置く。
\end{proof}

\subsubsection{係数線束を吸収する頂点 twist}

\begin{theorem}[周期を用いない線型標準形]
\label{thm:linear-normalization-without-period}
補題~\ref{lem:linear-propagation-without-period} の設定で、$D$ を $V_i$ により生成され、三次関係が $R_i$ である
正の向きの $R$ 線形 $\mathbb Z$ 代数とする。
$E=V_0$、$F=V_1$ と置く。
このとき頂点ごとの可逆加群 $N_i$ が存在し、
\[
 D'_{ij}=N_i^{-1}\otimes_R D_{ij}\otimes_R N_j
\]
は、偶数位置の矢空間が $E$、奇数位置の矢空間が $F$、
三次関係が同じ parity の二文字の交換関係である標準 linear $\mathbb Z$ 代数となる。
特に $n\geq0$ に対し
\begin{equation}\label{eq:linear-symmetric-components}
 D'_{i,i+n}\simeq
 \Sym^{\lceil n/2\rceil}E_{\bar i}
       \otimes_R\Sym^{\lfloor n/2\rfloor}E_{1-\bar i},
 \qquad E_0=E,\ E_1=F.
\end{equation}
全ての同定は積と整合する。
\end{theorem}

\begin{proof}
全証明は相対モジュライ論文 \cite[\texttt{relative_cubic_noetherian}]{QuadricModuli2026} に置く。
\end{proof}

\subsubsection{各 row と column の普遍 Noetherian 性}

\begin{corollary}\label{cor:linear-universal-row-noetherian}
前定理で $R$ が Noetherian なら、$D$ の全ての row と column は Noetherian である。
同じ結論は任意の可換 Noetherian $R$ 代数への基底変更後にも成り立つ。
\end{corollary}

\begin{proof}
全証明は相対モジュライ論文 \cite[\texttt{relative_cubic_noetherian}]{QuadricModuli2026} に置く。
\end{proof}

\begin{remark}[二周期の場合の Zhang twist]
二周期同型 $\theta$ が指定される場合、$V_{2r}=E$, $V_{2r+1}=F$ と同定し、
$E,F$ と最初の二つの係数線束を有限開被覆上で自明化する。
非退化形式を $a=(\mathrm{id}\otimes P)\varepsilon_E$、
$b=(\mathrm{id}\otimes Q)\varepsilon_F$ と書く。$P,Q$ は可逆行列である。
添字 $2r,2r+1$ でそれぞれ $P^r,Q^r$ を使う隣接写像は、
\[
 (P^r\otimes P^{r+1})\varepsilon_E=\det(P)^r a
\]
と $Q$ に対する同式により、標準 linear 代数から元の $D$ への同型を与える。
この座標で $\theta$ は $P^{-1},Q^{-1}$ となるので、偶数環 $B_Z$ は
Segre 環に $(P^{-1},Q^{-1})$ が誘導する自己同型 $\tau$ による
$a*b=a\tau^{\deg a}(b)$ という Zhang twist として具体的に書ける。
主定理の row と column の証明は、この二周期性を用いない。
\end{remark}

\subsection{形式貼合せと普遍性}

ここから再び元の基底環 $R$ とそのイデアル $I$ を用いる。
第\ref{sec:formal-noetherian-patching}節の $\mathbb Z$ 代数に対する定理~\ref{thm:patch-z-algebra}を用いる。

\begin{proof}
全証明は相対モジュライ論文 \cite[\texttt{relative_cubic_noetherian}]{QuadricModuli2026} に置く。
\end{proof}

\begin{corollary}[AS4 被覆への帰結]
系~\ref{cor:relative-periodic-even-noetherian} の $B$ が、同じ基底上の次数付き代数 $A$ の regular normal な二次線による商であり、
その regular 性と商の同定が任意の基底変更と可換するとする。
このとき $A$ も任意の可換 Noetherian 基底変更後に左右 Noetherian である。
\end{corollary}

\begin{proof}
全証明は相対モジュライ論文 \cite[\texttt{relative_cubic_noetherian}]{QuadricModuli2026} に置く。
\end{proof}

\paragraph{ここで証明していない事項.}
この結果は Noetherian 性に関するものであり、任意の導来 Morita 同値が
指定された相対 heart や偏極を保つことまでは主張しない。
AS4 被覆の全次数平坦性、正則正規元、相対 Gorenstein 解消は前提となる別の構成結果である。


\subsection{可換二次曲面の有限型近傍への適用}

\begin{theorem}[可換二次曲面の有限型相対実現]
\label{thm:quadric-finite-type-effectivity}
命題~\ref{prop:quadric-zariski-homological-as-cover} と
命題~\ref{prop:quadric-canonical-cubic-tails} の共通の有限型 affine 開
$U^\circ$ を取る。この開は可換二次曲面 $Q$ の点 $u$ を含む。
このとき以下が成立する。
\begin{enumerate}
\item 任意の可換 Noetherian affine 基底変更
$T=\operatorname{Spec}R\to U^\circ$ に対し、軌道環 $B_T$ と同じ基底上の被覆 $A_T$ は
左右 Noetherian である。
\item $T$ が局所有限型 $k$-site に属する場合、
$(\qgr B_T,\pi B_T)$ とその全有限型基底変更は
定義~\ref{def:relative-h2} の全公理を満たす。
これらは基底変更・fppf 降下と整合し、射
\[
 f:U^\circ\longrightarrow\mathfrak H_2
\]
を定める。$u$ の像は $(\coh Q,\OO_Q)$ である。
\item この族の canonical compact 圏は $\operatorname{Perf}\Lambda_T$ であり、
次数移動は既存の Morita 根 $\mathsf s_T$ に対応する。
その全軌道環は $B_T$、同じ根の AS4 被覆は $A_T$ である。
\end{enumerate}
元の site に属さない Noetherian $R$ に対しては、
有限型 $R$-test 上で同じ全公理を満たす系という意味で (2) を読む。
本定理だけから $\mathfrak H_2$ 全体の代数性は結論しない。
\end{theorem}

\begin{proof}
全証明は相対モジュライ論文 \cite[\texttt{relative_cubic_noetherian}]{QuadricModuli2026} に置く。
\end{proof}

% !TEX root = ../nc-quadric-classification.tex
% 2026-11-21 / v1.34（証明本文をBへ移管）
\section{幾何的二周期 cubic 族の普遍 canonical hull}
\label{sec:universal-cubic-hull}
\paragraph{主張の範囲.}
任意の幾何的二周期 cubic 族について、全幾何ファイバーの canonical hull を
有限 projective 加群間の射で相対化する。
可換二次曲面の特殊性、Artin 厚み、Krull 交叉定理を用いない。
体上の存在定理は Van den Bergh の Theorem~6.1 に依存し、
本節ではそこで存在する hull の六次元関係空間を、逆 coevaluation の核として
正確に同定する。相対化の新規性は未確認である。
二周期の存在と、その二乗が反 Serre 関手であることは別条件である。

\subsection{入力と定理}

$k$ を標数零の体、$R$ を可換 Noetherian $k$-代数とする。
$D$ を正の $R$-線形 $\mathbb Z$-代数とし、次を仮定する。
\begin{enumerate}
\item $D_{ii}=R$、$D_{ij}=0$ $(j<i)$、全ての $D_{ij}$ は有限 projective である。
  $V_i=D_{i,i+1}$ の階数は二であり、$D$ は三次関係の直和因子
  $R_i\subset V_i\otimes V_{i+1}\otimes V_{i+2}$、$\rank R_i=2$ により提示される。
\item 四次の交わり
  \[
  W_i=(R_i\otimes V_{i+3})\cap(V_i\otimes R_{i+1})
  \]
  は可逆 $R$-加群である。交わりは任意の基底変更と可換し、縮約により
  \[
  W_i\otimes V_{i+3}^{\vee}\xrightarrow{\sim}R_i,
  \qquad W_i\otimes V_i^{\vee}\xrightarrow{\sim}R_{i+1}
  \tag{UH-1}\label{hull-eq:coevaluation}
  \]
  が得られる。
\item 全幾何ファイバーは geometric quintuple に対応する標準 cubic AS 正則
  $\mathbb Z$-代数である。標準 cubic 解消は各内部次数で有限 projective の
  分裂完全複体となり、任意の基底変更に整合する。
\item 積と整合する実際の二周期同型 $\beta:D\xrightarrow{\sim}D(2)$ がある。
\end{enumerate}
これらは補題~\ref{lem:relative-geometric-cubic-flatness} の相対 cubic 解消から得られる入力である。
最後の条件は、幾何ファイバーごとの二周期同型の存在より強い。
本節の積の規約は
\[
 B_n=D_{0,2n},\qquad a\star b=a\,\beta^m(b)
 \quad(a\in B_m,\ b\in B_n)
 \tag{UH-2}\label{hull-eq:product}
\]
である。右軌道積 $a*b=s^n(a)\circ b$ と比較する際は、
式 (\ref{eq:cubic-tails-right-convention}) の
$C_{ij}=D_{-j,-i}$、$\beta_C=s^{-1}$ を用いて規約を合わせる。
偶数環と右軌道環の積の一致は定理~\ref{thm:quadric-finite-type-effectivity} の証明と同じである。

$U_i=V_i\otimes V_{i+1}$ および
\[
 S_i=R_i\otimes V_{i+3}+V_i\otimes R_{i+1}
       \subset U_i\otimes U_{i+2}
 \tag{UH-3}\label{hull-eq:relations}
\]
と置く。$S_i$ は階数七の直和因子である。
$U=U_0$ とし、$U\otimes U\xrightarrow{1\otimes\beta}U_0\otimes U_2$
に沿って $S_0$ を $U^{\otimes2}$ の部分加群とも考える。
この規約で $B=T_R(U)/(S_0)$ である。

\begin{theorem}[一般二周期 cubic 族の canonical AS4 被覆]
\label{thm:universal-cubic-hull}
上の仮定の下で、基底変更に自然な全射
\[
 \Phi_i:S_i\longrightarrow W_i,\qquad
 \Phi_i|_{W_i}=2\,\mathrm{id}_{W_i}
 \tag{UH-4}\label{hull-eq:Phi}
\]
があり、$H_i=\ker\Phi_i$ は階数六の直和因子である。
$\beta(H_i)=H_{i+2}$ であり、
\[
 A=T_R(U)/(H_0)\longrightarrow B
 \tag{UH-5}\label{hull-eq:hull}
\]
は次を満たす。
\begin{enumerate}
\item 各 $A_n$ は階数 $\binom{n+3}{3}$ の有限 projective 加群であり、
  任意の基底変更後にも同じ提示と次数別比較が成立する。
\item $P=A^!$ は階数列 $(1,4,6,4,1)$ の有限 projective な二次代数で、
  $P_n=0$ $(n\ge5)$ である。$P$ は相対 Koszul Frobenius であり、
  $A$ の左右および双加群 Koszul 解消は長さ四で、任意の基底変更後にも完全である。
\item $G=S_0/H_0\cong W_0$ は $A_2$ の直和因子である。
  $J=\ker(A\to B)$ に対し、全次数で普遍的な同型
  \[
   G\otimes_R A_{n-2}\xrightarrow{\sim}J_n
     \xleftarrow{\sim}A_{n-2}\otimes_R G
   \tag{UH-6}\label{hull-eq:regular-normal}
  \]
  がある。従って $A/(G)=B$、$GA=AG$ であり、$G$ の任意の局所生成元は
  左右正則正規な二次の元である。
\item 任意の可換 Noetherian $R$-代数 $R'$ に対して、$A_{R'}$ は左右 Noetherian である。
\end{enumerate}
各幾何ファイバーの被覆は Van den Bergh の canonical hull を
指定された二周期に沿って graded 化したものに一致する。
ここで Serre 平方根に関する仮定は用いない。
\end{theorem}

Noetherian 性には定理~\ref{thm:relative-geometric-cubic-noetherian} を用いる。
それ以外の結論は本節の有限 projective と tensor の議論による。
相対 AS4 という語は基底の解消を指し、特異な $R$ 上で環 $A$ 自身の
絶対大域次元が四であると主張しない。

\subsection{逆 coevaluation と次数三の正規性}

\subsubsection{評価の規約と線束値の写像}

\eqref{hull-eq:coevaluation} の逆を使うと
\[
 R_i\otimes V_{i+3}\cong
 W_i\otimes V_{i+3}^{\vee}\otimes V_{i+3}\longrightarrow W_i
\]
および
\[
 V_i\otimes R_{i+1}\cong
 V_i\otimes W_i\otimes V_i^{\vee}\longrightarrow W_i
\]
という二つの評価写像が得られる。
両者の $W_i$ への制限は階数二の恒等写像の trace、すなわち $2\,\mathrm{id}$ である。
従って交わりで一致し、\eqref{hull-eq:Phi} を定める。
この全射は $W_i\hookrightarrow S_i$ の半分を切断に持ち、
$S_i=H_i\oplus W_i$ が成り立つ。
構成は有限 projective の評価と普遍的な交わりのみを用いるため、
任意の基底変更および $D$ の同型に自然である。

局所生成元 $w_i$ を選び $\Phi_i(x)=\phi_i(x)w_i$ と書くと、
\[
 w_i=\sum_{a=1}^2r_a\otimes v_a
 \quad\Longrightarrow\quad
 \phi_i(r_a\otimes v_b)=\delta_{ab}.
 \tag{UH-7}\label{hull-eq:inverse-coeval}
\]
他方の分解でも同じ規約を用いる。
特に $\phi_i(w_i)=2$ であり、$w_i$ を $u w_i$ に変えると
$\phi_i$ は $u^{-1}\phi_i$ になる。
従って本来の大域対象は単一の scalar-valued functional ではなく、
\eqref{hull-eq:Phi} とその双対線 $W_i^{\vee}\subset S_i^{\vee}$ である。

\subsubsection{四周期と二周期の区別}

$W_0,W_1$ を局所的に自明化し、$w_{i+2}=\beta(w_i)$ と選ぶ。
$w_i\in V_i\otimes R_{i+1}$ と
$w_{i+1}\in R_{i+1}\otimes V_{i+4}$ の二つの非退化 tensor の比較により、
一意な同型 $\theta_i:V_i\xrightarrow{\sim}V_{i+4}$ が定まる。
具体的には $w_i=\sum_a v_a\otimes r_a$ と書いたとき
$w_{i+1}=\sum_a r_a\otimes\theta_i(v_a)$ とする。
四回の巡回から $\theta_i\otimes\theta_{i+1}\otimes\theta_{i+2}$ は
$R_i$ を $R_{i+4}$ に移す。
従って $\theta:D\to D(4)$ は代数同型である。
自然性から
\[
 \theta_{i+2}\beta_i=\beta_{i+4}\theta_i
 \tag{UH-8}\label{hull-eq:commute}
\]
であり、選んだ生成元の下で $\theta$ と $\beta$ は可換する。
この $\theta$ は \cite[\S5.2, (5.2)]{VdB2011} の canonical 四周期である。
一般に $\theta=\beta^2$ は成立しない。

\begin{lemma}[偶数環の相対 Koszul 性]
\label{hull-lem:even-koszul}
$B$ は相対 Koszul であり、$E=B^!$ の全次数は有限 projective で、
任意の基底変更と両立する。階数列は
\[
 (\rank E_n)_{n\ge0}=(1,4,7,8,8,\ldots).
 \tag{UH-9}\label{hull-eq:E-ranks}
\]
\end{lemma}
\begin{proof}
証明は相対モジュライ論文 \cite[\texttt{universal_cubic_hull}]{QuadricModuli2026} に置く。
\end{proof}

\begin{lemma}[普遍的な次数三分解]
\label{hull-lem:K3}
$E_3^{\vee}$ の対応する添字付き成分は
\begin{align*}
 K_{3,i}
 &= (S_i\otimes V_{i+4}\otimes V_{i+5})
       \cap(V_i\otimes V_{i+1}\otimes S_{i+2})\\
 &= (R_i\otimes R_{i+3})
       \oplus(V_i\otimes W_{i+1}\otimes V_{i+5}).
 \tag{UH-10}\label{hull-eq:K3}
\end{align*}
この同定は任意の基底変更と両立する。
\end{lemma}
\begin{proof}
証明は相対モジュライ論文 \cite[\texttt{universal_cubic_hull}]{QuadricModuli2026} に置く。
\end{proof}

\begin{lemma}[functional 線の普遍正規性]
\label{hull-lem:normal}
\eqref{hull-eq:Phi} の双対は可逆加群 $N=W_0^{\vee}\subset E_2$ を与え、
$NE=EN$ が任意の基底変更後に成立する。
\end{lemma}
\begin{proof}
証明は相対モジュライ論文 \cite[\texttt{universal_cubic_hull}]{QuadricModuli2026} に置く。
\end{proof}

\subsection{全幾何ファイバーでの canonical hull との同定}

本節のみ基礎環を代数閉体 $K$ とする。
\cite[Theorem 6.1]{VdB2011} により、canonical 四周期 $\theta$ の逆の
正規化自己同型の類を持つ、偶数 $\mathbb Z$-代数 $D^{\mathrm{ev}}$ の hull が存在する。
以下では、そこに現れる関係空間が実際に $H_i=\ker\phi_i$ であることを証明する。
これは、既知の hull の存在だけから候補の同定を省略しないための一段である。

\begin{lemma}[正規化式が決める評価の一意性]
\label{hull-lem:fibre-identification}
canonical hull の六次元関係空間を $H'_i\subset S_i$ とする。
このとき $H'_i=H_i$ である。
\end{lemma}
\begin{proof}
証明は相対モジュライ論文 \cite[\texttt{universal_cubic_hull}]{QuadricModuli2026} に置く。
\end{proof}

\begin{remark}[Bauwens--Van den Bergh の式との直接照合]
\cite[Proposition 2.7, (2.11), Remark 3.3]{BauwensVanDenBergh1998} は graded cubic の場合、
\[
 w=x^tf=(Qf)^tx,\qquad
 \psi(f_i x_j)=\eta(\Lambda^{-1})_{ji},\qquad
 \psi(x_i f_j)=\eta(\Lambda Q^{-t})_{ij}
\]
と書く。canonical choice は $\Lambda=Q^t$ であるから
\[
 \psi(x_i f_j)=\eta\delta_{ij},\qquad
 \psi(f_i x_j)=\eta(Q^{-1})_{ij}.
 \tag{UH-15}\label{hull-eq:BV-convention}
\]
これは $w$ の二つの分解に対する \eqref{hull-eq:inverse-coeval} と正確に一致する。
原典では偶数環の元の次数を元の cubic 環の次数で記している箇所があるため、
そこに現れる次数四の functional は本節の $E_2$ である。
一般 $\mathbb Z$-代数での同定は補題~\ref{hull-lem:fibre-identification} が直接与える。
\end{remark}

\begin{corollary}[体上の双対と正則性]
\label{hull-cor:fibre-dual}
任意の実際の二周期 $\beta$ に対して、$T(U)/(H_0)$ は canonical hull の
graded 化である。$\phi\in E_2$ は左右正則正規で、
$E/(\phi)$ は階数列 $(1,4,6,4,1)$ の Koszul Frobenius 代数である。
\end{corollary}
\begin{proof}
証明は相対モジュライ論文 \cite[\texttt{universal_cubic_hull}]{QuadricModuli2026} に置く。
\end{proof}

\subsection{非被約基底上での相対化}

再び一般の Noetherian $R$ に戻る。
定理~\ref{thm:universal-cubic-hull} の証明を完結する。

\subsubsection{まず双対の商を平坦にする}

補題~\ref{hull-lem:even-koszul} より $E_n$ は有限 projective である。
補題~\ref{hull-lem:normal} の正規線 $N=W_0^{\vee}$ による左右乗法
\[
 N\otimes E_{n-2}\longrightarrow E_n,
 \qquad E_{n-2}\otimes N\longrightarrow E_n
 \tag{UH-16}\label{hull-eq:E-inject}
\]
は、系~\ref{hull-cor:fibre-dual} により全幾何ファイバーで単射である。
有限 projective 加群間の写像のファイバー単射性から、
これらは局所直和因子への単射になる。
従って
\[
 P=E/(N)
 \tag{UH-17}\label{hull-eq:P}
\]
の全次数は有限 projective であり、商は任意の基底変更と可換する。
\eqref{hull-eq:E-ranks} の差から階数列は $(1,4,6,4,1,0,\ldots)$ となる。
ここでは商のファイバー次元が一定であることだけから平坦性を推論していない。
既に projective である加群間の写像 \eqref{hull-eq:E-inject} を用いている。

$P$ は二次代数であり、$P_1,P_2,P_3$ は projective、全幾何ファイバーは Koszul である。
従って再び \cite[Chapter 6, Theorem 7.3]{PolishchukPositselski2005} を \emph{$P$ に}適用する。
全ての関係交わりが直和因子となるため、二次双対
\[
 P^!=T_R(U)/(H_0)=A
\]
の全次数も有限 projective で、任意の基底変更に整合する。
全ファイバーでの canonical hull の Hilbert 級数により
$\rank A_n=\binom{n+3}{3}$ となる。
この順序により、$A_3$ の局所自由性を証明前に仮定する循環を避けている。

\subsubsection{Frobenius、AS 解消、二次切断}

$L=P_4$ は可逆加群である。
各積の adjoint
\[
 P_i\longrightarrow\operatorname{Hom}_R(P_{4-i},L)
 \qquad(0\le i\le4)
 \tag{UH-18}\label{hull-eq:Frobenius}
\]
は有限 projective 加群間の射で、全幾何ファイバーで同型である。
従って同型であり、任意の基底変更後もそうである。
これが相対 Frobenius 性である。
Koszul 複体の各内部次数は有限な projective 複体で、
そのファイバーは体上の Koszul 解消だから分裂完全となる。
$P_i=0$ $(i>4)$ より、左右および対角双加群の解消は長さ四である。
Frobenius pairing による双対化から
\[
 \underline{\operatorname{Ext}}_A^j(R,A)
 \cong
 \begin{cases}
 L\otimes_RR(4),&j=4,\\
 0,&j\ne4
 \end{cases}
 \tag{UH-19}\label{hull-eq:AS}
\]
を左右に得る。ここで表示は対応する側の加群としてのものとし、
双加群同定には通常の Frobenius 自己同型を含める。
解消は任意の基底変更後にも完全なので、これらは相対 AS4 条件である。

$A\to B$ は提示から全射である。
$A_n,B_n$ が有限 projective だから $J_n=\ker(A_n\to B_n)$ も projective であり、
核は任意の基底変更と可換する。
$G=S_0/H_0\cong W_0$ は $J_2$ と同一視される。
左右乗法 \eqref{hull-eq:regular-normal} は有限 projective 加群間の射で、
各幾何ファイバーでは canonical hull の正則二次切断の同型である。
従って全次数で同型であり、任意の基底変更後も同型である。
これにより正則性、正規性、商 $A/(G)=B$ を同時に得る。
特に局所生成元 $g$ について $gA=Ag$ と両側正則性があるので、
$ga=\nu(a)g$ によって定まる写像 $\nu$ は全単射の代数自己同型である。
規約 \eqref{hull-eq:graded-normality} では、$\nu|_U=T^{-1}$ となる。
実際、$\phi(c)=1$ となる $c\in S_0$ で $g$ を表し、
$c\otimes u-\nu(u)\otimes c=h_L+h_R$、
$h_L\in H_0\otimes U$、$h_R\in U\otimes H_0$ と書く。
すると $z=c\otimes u-h_L=\nu(u)\otimes c+h_R\in K_3$ であり、
二つの縮約はそれぞれ $u$ と $\nu(u)$ になる。
\eqref{hull-eq:normal-contraction} を graded 化した等式から $u=T\nu(u)$ を得る。
この確認も幾何ファイバーへの還元を用いない。

\subsubsection{有限判定と Noetherian 性}

定理~\ref{thm:relative-finite-as-cover} の五条件は、本構成では全て基底上で成立する。
具体的には
\[
 \rank A_3=20,\quad\rank A_4=35,\qquad
 \rank P_3=4,\quad\rank P_4=1,\quad P_5=0,
\]
左右の $G$ 乗法の像の一致、次数一・二に対する局所直和因子への単射、
および \eqref{hull-eq:Frobenius} である。
従って canonical な階数六部分束 $H_0\subset S_0$ は、
被覆の有限判定部分スキームへの実際の切断を与える。
全ファイバーで有限条件が成り立つという主張のみで代用していない。

最後に任意の Noetherian $R'$ に移る。
系~\ref{cor:relative-periodic-even-noetherian} より $B_{R'}$ は左右 Noetherian である。
$G_{R'}$ を有限アフィン開被覆上で自明化する。
補題~\ref{noeth-lem:normal-lift} の議論を適用する。
正則正規な正次数元 $g$ に対して、$\operatorname{gr}_{(g)}A_{R'}$ は
$B_{R'}$ 上の skew polynomial 環となるため左右 Noetherian である。
固定内部次数では $(g)$ 進フィルトレーションは有限なので、
初期斉次生成元を有限個持ち上げる還元は有限回で止まる。
これにより斉次イデアルは有限生成、最高内部次数による還元で一般のイデアルも有限生成となる。
基底の有限開被覆から忠実平坦に降下すれば $A_{R'}$ は左右 Noetherian である。
これで定理~\ref{thm:universal-cubic-hull} の全ての主張を得た。

\subsection{Serre 平方根の条件は独立である}

\begin{remark}[二周期と反 Serre 関手]
任意の実際の二周期 $\beta$ が上の AS4 被覆を持つことと、
そこから来る exact 自己同値 $s$ が
\[
 s^2\simeq S^{-1}[2]
 \tag{UH-20}\label{hull-eq:Serre-root}
\]
を満たすことは別である。
後者は本定理の仮定にも結論にも含まれない。
有限 foundation $\Lambda$ 上で $s$ が可逆 perfect 双加群 $P_s$ により表され、
$\mathcal S_\Lambda=\mathbf R\operatorname{Hom}_R(\Lambda,R)$ が Serre 双加群である場合、
正確な追加条件は
\[
 P_s\mathbin{\otimes^{\mathbf L}}_\Lambda P_s
    \simeq\mathcal S_\Lambda^{-1}[2]
 \quad\text{in }\operatorname{Perf}(\Lambda^e).
 \tag{UH-21}\label{hull-eq:Morita-Serre}
\]
右辺の逆は可逆双加群の Morita 逆である。
根・根同型を fppf 局所の存在条件として扱う定義では、
基底線束 $\ell$ を許した
$P_s\mathbin{\otimes^{\mathbf L}} P_s\simeq\mathcal S_\Lambda^{-1}[2]\otimes_R\ell$ を
局所自明化して同じ条件を表せる。
この比較は定理~\ref{thm:bimodule-line-comparison} と
系~\ref{cor:foundation-antiserre-comparison-locus} で認識する対象である。
単に対象ごとに $s(F_0)\cong F_2$, $s(F_1)\cong F_3$ であることだけでは
\eqref{hull-eq:Morita-Serre} は従わない。
\end{remark}

\begin{example}[period two は自動的な Serre 根ではない]
$Q=\mathbb P^1\times\mathbb P^1$、$\sigma\in\operatorname{Aut}(Q)$ を
$\sigma^2\ne1$ かつ両 ruling を保つ自己同型とする。
例えば標数零で
$\sigma([x_0:x_1],[y_0:y_1])=([2x_0:x_1],[y_0:y_1])$ とできる。
$s=(-\otimes\OO_Q(1,1))\circ\sigma^*$ は標準 cubic helix の各対象を二項先へ送り、
実際の二周期を与える。
$s^{-n}\OO_Q\cong\OO_Q(-n,-n)$ なので $(\OO_Q,s)$ は AZ ample でもある。
従ってその偶数軌道環には本定理の canonical AS4 被覆がある。
しかし
\[
 s^2=(-\otimes\OO_Q(2,2))\circ\sigma^{2*},\qquad
 S^{-1}[2]=(-\otimes\OO_Q(2,2))
\]
であり、動かされる点の skyscraper sheaf を比べれば同型でない。
従って同じ指定二周期を H2 の witness とするには、
\eqref{hull-eq:Serre-root} を別途確かめる必要がある。
この例は underlying $(\coh Q,\OO_Q)$ が H2 に属することを否定しない。
\end{example}

\paragraph{一般 cubic 点への適用で残る入力.}
本節により、実際の geometric cubic 二周期族が構成された後の AS4 被覆は
Q 固有の障害ではなくなる。
一般の内在的 H2 族に適用するには、構造対象を含む適合する四項 foundation、
その上で代数化された実際の Morita 根、
および \eqref{hull-eq:Morita-Serre} と全軌道環の比較を供給する必要がある。
これらの存在範囲および根を忘れた相対実現部分層の代数性は、本節では主張しない。


% !TEX root = ../nc-quadric-classification.tex
% 2026-11-21 / v1.35
\section{一般の幾何的 cubic 族の相対内在的実現}
\label{sec:relative-cubic-families}

本節では、可換二次曲面を中心とするという制限を外し、
普遍 Noetherian 性と canonical hull を相対 tails の認識に結び付ける。
Serre 平方関係は、まず指定された周期の独立な入力として明記する。
次節の Picard 構成が非線型の普遍族上でこの入力を供給する。

\subsection{周期を選ぶ前の canonical compact 圏}

$R$ を可換 Noetherian $k$-代数とし、
$C$ を定理~\ref{thm:relative-geometric-cubic-noetherian} の仮定を満たす
正の向きの $R$-線形 $\mathbb Z$-代数とする。
右表現可能加群を $P_i=e_iC$、$\mathcal O_i=\pi P_i$ と書く。
三次関係を $R_i$、四次 overlap を $W_i$ とする。

\begin{proposition}[一般 geometric cubic 族の canonical tails]
\label{prop:general-cubic-canonical-tails}
任意の可換 Noetherian $R$-代数 $R'$ 上で
$\operatorname{QGr}(C_{R'})$ は locally Noetherian であり、
任意の連続四項の直和はその非有界導来圏の compact tilting generator である。
特に
\[
 T=\bigoplus_{i=0}^3\mathcal O_{-i},\qquad
 \Gamma=\operatorname{End}(T)
\]
と置くと、$\Gamma$ は階数 $24$ の有限 projective $R$-代数であり、
\[
 D(\operatorname{QGr}C)\simeq D(\operatorname{Mod}\Gamma),\qquad
 D(\operatorname{QGr}C)^c\simeq\operatorname{Perf}\Gamma.
 \tag{GC-Morita}\label{eq:general-cubic-canonical-morita}
\]
この比較は実際の module-object 引戻しの導来関手と整合する。
compact 圏は相対 smooth proper で、相対 Serre 関手とその任意基底変更を持つ。
この命題では二周期の存在を仮定しない。
\end{proposition}

\begin{proof}
完全な証明は相対モジュライ論文 \cite[\texttt{relative_cubic_families}]{QuadricModuli2026} に置く。
\end{proof}

\subsection{実際の二周期からの全相対条件}

$C$ に二周期同型 $\beta:C\to C(2)$ を指定し、
\[
 B_n=C_{0,2n},\qquad a\star b=a\,\beta^m(b)
 \quad(a\in B_m,\ b\in B_n)
 \tag{GC-prod}\label{eq:general-cubic-period-product}
\]
と定める。$\check B\simeq C^{\mathrm{ev}}$ の下で、
$B$ の次数移動 $(1)$ は $C$ の right tails 上で
$\mathcal O_i\mapsto\mathcal O_{i-2}$ という完全自己同値 $s$ を与える。
矢も $\beta^{-1}$ で移す。単に対象の置換を指定しているのではない。
式~\eqref{eq:general-cubic-canonical-morita} によって
$s$ を表す可逆 perfect $\Gamma$ 双加群を $P_s$ と書く。

\begin{theorem}[Serre 適合二周期族の相対効果性]
\label{thm:geometric-cubic-relative-effectivity}
上の $C,\beta$ に対し、基底の fppf 局所で
\[
 P_s\otimes_\Gamma^{\mathbf L}P_s
       \simeq\mathcal S_\Gamma^{-1}[2]
 \tag{GC-root}\label{eq:general-cubic-serre-input}
\]
が与えられ、任意の基底変更ではこの同型を引き戻すとする。
このとき有限型 $R$-基底上の系
\[
 \mathcal C_{R'}=\operatorname{qgr}(B\otimes_RR'),\qquad
 \mathcal O_{R'}=\pi(B\otimes_RR')
 \tag{GC-family}\label{eq:general-cubic-relative-family}
\]
は定義~\ref{def:relative-abelian-system}、
\ref{def:relative-h2} の全条件を満たす。
特に $\operatorname{Spec}R$ が元の局所有限型 $k$-site に属すれば、
実際の $\mathfrak H_2(R)$ の対象を与える。
平方根、平方同型、二周期、cubic helix および AS4 被覆は、
構成に用いる witness であり、出力の標識には加えない。

各 affine Noetherian 基底変更上で全相対 $\chi_q$ も成立する。
$R$ が元の site に属さない場合は、
有限型 $R$-test 上で同じ全公理が成立するという意味に限る。
\end{theorem}

\begin{proof}
完全な証明は相対モジュライ論文 \cite[\texttt{relative_cubic_families}]{QuadricModuli2026} に置く。
\end{proof}

\begin{remark}[Hom の向きで表示する場合]
\label{rem:general-cubic-right-orientation}
framed 五つ組を $D_{ij}=\operatorname{Hom}(E_i,E_j)$ として
回復する場合には、右加群のために
\[
 C_{ij}=D_{-j,-i},\qquad \beta_C=\beta_D^{-1}
\]
を用いる。このとき $E_i$ は $\pi e_{-i}C$ に対応し、
四項 foundation の頂点・全 Hom・積が一致する。
canonical 四周期も反転し、$\beta_D^2=\theta_D$ が
頂点スカラーを除いて成立すれば、$C$ 側でも対応する平方関係が成立する。
相対 Serre の比較は $C$ の canonical compact 圏で行う。
反対の積を暗黙に同じ環と同一視しない。
\end{remark}

% !TEX root = ../nc-quadric-classification.tex
% 2026-11-21 / v1.35
\section{相対 Picard 作用による cubic Serre 平方根の構成}
\label{sec:relative-picard-roots}

\paragraph{主張の範囲と依存.}
本稿は標数零における\emph{非線型} geometric cubic 五つ組の普遍族上に、
étale 局所に実際の二周期を構成する。
その平方は四次 tensor の巡回から定まる canonical 四周期に
頂点スカラーを除いて一致する。
全ての恒等式を被約な普遍基底上で証明してから引き戻すので、
非被約な Artin 基底も結論に含む。
原典の既知入力は、Artin--Tate--Van den Bergh の
\emph{相対} Picard 作用と、Van den Bergh の geometric 五つ組の構成である。
相対 tails と有限 foundation の比較、および Serre の双加群比較には、
定理~\ref{thm:relative-cubic-tails-tilting} と定理~\ref{thm:bimodule-line-comparison} を用いる。
本稿は任意の二周期が Serre 平方根であるとは主張しない。
相対化の新規性は未確認である。

\subsection{一次資料の正確な内容}

Van den Bergh \cite[Proposition 5.6.1]{VdB2011} は、
代数閉体 $k$、$\operatorname{char}k\ne2$ 上の全ての quadric
$\mathbb Z$ 代数 $D$ に対して $D\simeq D(2)$ を示す。
全 $\mathbb Z$ 代数の二周期性が主張されている。
その後の Corollary 5.6.2 が偶数部分の通常の次数付き環による表示を与える。

非線型の場合、楕円 helix $(C,(\mathcal L_i)_i)$ に対し
\[
 N=\mathcal L_2\mathcal L_0^{-1}\in\Pic^0(C),\qquad
 \mathcal L_{i+2}\simeq\mathcal L_iN
\]
である。$a\in\Pic^0(C)$、$a^{-2}=N$ を選び、
\cite[Theorem 4.2.3]{VdB2011} の作用
\[
 \eta:\Pic^0(C)\longrightarrow\Aut(C),\qquad
 \eta(a)^*L=L\otimes a^{-\deg L}
 \tag{PR-1}\label{eq:field-eta}
\]
を用いる。$\sigma=\eta(a)$ は各次数二の束 $\mathcal L_i$ を $\mathcal L_{i+2}$ へ移す。
原典の Proposition 6.6 は二周期を標準 hull へ持ち上げるが、
任意の二周期 $\beta$ に対する $\beta^2=\theta$ を主張してはいない。

\begin{proposition}[ATV の相対 Picard 作用]\label{prop:ATV-relative-action}
$C\to S$ を $\PP^2_S$ の三次因子、または
$\PP^1_S\times_S\PP^1_S$ の bidegree $(2,2)$ の因子の平坦族とする。
$P=\Pic^0_{C/S}$ を、各幾何ファイバーの各被約既約成分で次数零の
直線束の類を表す相対 Picard 部分群とする。
$P$ の $C$ への作用があり、その構成は任意の底変更および族の同型に自然である。
作用の符号を \eqref{eq:field-eta} と合わせて $\eta$ と書く。
\end{proposition}

\begin{proof}
完全な証明は相対モジュライ論文 \cite[\texttt{relative_picard_roots}]{QuadricModuli2026} に置く。
\end{proof}

\subsection{被約な普遍基底で証明してから底変更する}

以下、$k$ は代数閉体、$\operatorname{char}k=0$ とする。
四つの二次元空間 $V_0,\ldots,V_3$ を固定し、
\[
 \mathscr G\subset\PP(V_0\otimes V_1\otimes V_2\otimes V_3)
\]
を geometric 五つ組の開部分とする。
$02\mid13$ flattening の階数一の閉部分を除き、
\[
 S=\mathscr G\setminus\mathscr G_{\mathrm{lin}}
\]
と置く。$S$ は滑らかで被約、既約な有限型 $k$ スキームである。
以下の構成は $S$ の affine open 上で行って貼り合わせてもよい。

普遍直線 $W\subset(V_0\otimes\cdots\otimes V_3)\otimes\OO_S$ を巡回し、
$V_{i+4}=V_i$ として $W_i$ と三次関係 $R_i$ を定める。
これにより普遍 cubic $\mathbb Z$ 代数 $D$ を得る。
$W_i$ は全て同じ基底線 $W$ の巡回像であり、この共通の同定の下で
canonical 四周期
\[
 \theta:D\longrightarrow D(4)
 \tag{PR-2}\label{eq:theta}
\]
は各 $V_i=V_{i+4}$ 上の恒等である。
この記述に $W$ の大域的な生成元は必要ない。
局所生成元 $w_i$ を使えば
\[
 w_i=\sum_a v_a\otimes r_a,
 \qquad w_{i+1}=\sum_a r_a\otimes\theta_i(v_a)
\]
という \cite[\S5.2]{VdB2011} の定義そのものである。

\begin{lemma}[普遍族の次数別平坦性]\label{lem:universal-flatness}
各 $D_{ij}$ は有限局所自由で、標準 cubic 解消は各内部次数で分裂完全である。
その双対複体も AS 双対の所定の次数以外で完全で、残る cohomology は可逆層である。
従ってこれらの提示・解消・双対・関係の交わりは任意の底変更と両立する。
\end{lemma}

\begin{proof}
完全な証明は相対モジュライ論文 \cite[\texttt{relative_picard_roots}]{QuadricModuli2026} に置く。
\end{proof}

普遍 point curve を $\pi:C\to S$ とする。
隣接二因子の評価による $2\times2$ 行列の determinant が
$\PP(V_0^\vee)\times\PP(V_1^\vee)$ 内に $C$ を定める。
geometric 条件は、determinant が零である点でもこの行列の階数が一であることを保証する。
三点 correspondence の隣接二点への射影は、単元小行列式による掃き出しから
スキームとして同型になる。
従って全 $\mathcal L_i$ を同じ $C$ 上の直線束として定義できる。
これは第~\ref{sec:relative-cubic-noetherian}節の相対 point-curve 構成である。

$S$ は線型 locus を除いているので、determinant は各幾何ファイバーで非零である。
したがって $C/S$ は平坦な relative Cartier divisor である。
その標準 exact sequence
\[
 0\longrightarrow\OO_{\PP^1_S\times\PP^1_S}(-2,-2)\otimes\pi^*\ell
 \longrightarrow\OO_{\PP^1_S\times\PP^1_S}
 \longrightarrow\OO_C\longrightarrow0
 \tag{PR-3}\label{eq:curve-sequence}
\]
（$\ell$ は基底の可逆層）から、$\pi_*\OO_C=\OO_S$ が普遍的に成り立ち、
$R^1\pi_*\OO_C$ は可逆層、$R^{>1}\pi_*\OO_C=0$ となる。

\begin{lemma}[被約基底での相対 Picard の一致判定]\label{lem:reduced-equalizer}
$X$ を被約スキーム、$Y\to X$ を相対 diagonal が immersion である代数空間とする。
二つの section $f,g:X\to Y$ が全幾何点上で一致すれば、$f=g$ である。
この判定は $Y=\Pic_{C/X}$ に適用できる。
\end{lemma}

\begin{proof}
完全な証明は相対モジュライ論文 \cite[\texttt{relative_picard_roots}]{QuadricModuli2026} に置く。
\end{proof}

\begin{lemma}[普遍 helix の相対的な差]\label{lem:relative-helix-difference}
$N=[\mathcal L_2\mathcal L_0^{-1}]$ は $P=\Pic^0_{C/S}$ の section であり、相対 Picard 群で
\[
 [\mathcal L_{i+2}\mathcal L_i^{-1}]=N
 \tag{PR-4}\label{eq:helix-difference}
\]
が全ての $i$ について成り立つ。
\end{lemma}

\begin{proof}
完全な証明は相対モジュライ論文 \cite[\texttt{relative_picard_roots}]{QuadricModuli2026} に置く。
\end{proof}

\subsection{平方根の étale パラメータ空間}

\begin{lemma}[相対 Picard の二等分]\label{lem:picard-halves}
$P=\Pic^0_{C/S}$ は滑らかな相対次元一の可換群代数空間であり、
\[
 [-2]:P\longrightarrow P
\]
は étale かつ全射である。
従って
\[
 T=S\mathbin{\times}_{N,P,[-2]}P
 \tag{PR-5}\label{eq:T}
\]
は $S$ 上 étale 全射であり、普遍 section $a\in P(T)$ は
$a^{-2}=N$、すなわち $a^2=N^{-1}$ を満たす。
\end{lemma}

\begin{proof}
完全な証明は相対モジュライ論文 \cite[\texttt{relative_picard_roots}]{QuadricModuli2026} に置く。
\end{proof}

$T$ は $k$ 上滑らかで被約である。
必要なら $T$ の étale scheme chart を取り、以下では $T$ 自体をスキームとして扱う。
命題~\ref{prop:ATV-relative-action} により、
普遍 Picard class $a$ から実際の曲線の自己同型
\[
 \sigma=\eta(a):C_T\longrightarrow C_T
 \tag{PR-6}\label{eq:sigma}
\]
が得られる。
ここでは $a$ を表す大域的な直線束を選ぶ必要はない。
群作用そのものが相対 Picard 空間上に定義されている。

\begin{lemma}[線束の比較と二周期]\label{lem:relative-beta}
$T$ の Zariski 開被覆上で同型
\[
 \varepsilon_i:\sigma^*\mathcal L_i\xrightarrow{\sim}\mathcal L_{i+2}
 \qquad(i=0,1,2,3)
 \tag{PR-7}\label{eq:epsilons}
\]
を選べ、それらが実際の代数同型
\[
 \beta:D_T\xrightarrow{\sim}D_T(2)
 \tag{PR-8}\label{eq:beta}
\]
を与える。
\end{lemma}

\begin{proof}
完全な証明は相対モジュライ論文 \cite[\texttt{relative_picard_roots}]{QuadricModuli2026} に置く。
\end{proof}

\subsection{平方と canonical 四周期の比較}

\begin{theorem}[相対的な適合 Serre 根]\label{thm:relative-compatible-root}
上の $T$ をさらに Zariski 局所化した基底上で、
\eqref{eq:beta} は頂点対角型の自己同型 $\delta$ に対し
\[
 \beta^2=\delta\theta
 \tag{PR-9}\label{eq:beta-square}
\]
を満たす。
従って tails 上では $(2)_\beta^2\simeq(4)_\theta$ である。
この同型は基底上に構成されており、任意の底変更後にも成り立つ。
特に非被約基底へ引き戻しても、相対 Serre 関手を持つ有限 foundation の比較の下で
\[
 s=(2)_\beta,\qquad s^2\simeq S^{-1}[2]
 \tag{PR-10}\label{eq:Serre-square}
\]
を得る。
\end{theorem}

\begin{proof}
完全な証明は相対モジュライ論文 \cite[\texttt{relative_picard_roots}]{QuadricModuli2026} に置く。
\end{proof}



\begin{remark}[比較線束の局所自明化]
相対 Picard の等式からは、まず二つの直線束の差が基底から来ることが分かる。
\eqref{eq:epsilons} の前に四つの比較直線束を自明化したことが必要である。
その選択を変えると $\beta_i$ は基底単元倍されるが、
projective な平方比較は変わらず、(d) の頂点スカラーがその差を吸収する。
さらに (e) では Serre 双加群の比較線束 $\ell$ を自明化している。
従って幾何ファイバー上の平方根の存在だけから、
基底線束を省略した大域的な平方同型を主張してはいない。
必要なのは有限個の基底直線束の局所自明化である。
\end{remark}

\begin{remark}[四項 foundation と右加群の向き]
Abelian helix から $D_{ij}=\Hom(E_i,E_j)$ を回復する場合は、
右 tails の構成に添字を反転した反対代数
\[
 C_{ij}=D_{-j,-i},\qquad \beta_C=\beta_D^{-1}
\]
を用いる。
すると $E_i$ は $\pi e_{-i}C$ に対応し、四項窓の自己準同型代数は
元の foundation の積と一致する。
$\beta_D^2=\delta\theta_D$ を反転すれば、$C$ 上でも
$\beta_C^2=\theta_C$ が頂点スカラーを除いて成り立つ。
ここで $\theta_C$ は $\theta_D^{-1}$ に対応する canonical 四周期である。
(e) の双加群 $P,Q$ とその比較 locus は、この $C$ の右 tails 上で直接取れる。
従って Serre の比較と右軌道積の規約は同じ向きで成立する。
\end{remark}

\subsection{Artin 持上げと atlas 構成への入力}

\begin{corollary}[有限型の étale 根 chart]\label{cor:finite-type-root-chart}
$S$ を覆う有限個の affine scheme $T_\alpha$ と étale 射 $T_\alpha\to S$ を選び、
その各々で上の実際の二周期および Serre 平方同型を構成できる。
従って $\coprod_\alpha T_\alpha$ は有限型の滑らかな $k$ スキームであり、
全ての非線型 cubic geometric 五つ組の点を覆う。
\end{corollary}

\begin{proof}
完全な証明は相対モジュライ論文 \cite[\texttt{relative_picard_roots}]{QuadricModuli2026} に置く。
\end{proof}

\begin{corollary}[指定した Picard 根の全 Artin 持上げ]\label{cor:all-artin-root}
$x\in S(k)$ と $a_x\in\Pic^0(C_x)$、$a_x^{-2}=N_x$ を固定する。
\eqref{eq:T} の中の $a_x$ に対応する点を含む affine étale scheme chart を選び、
その点を $t$ とする。必要なら前の有限被覆にこの chart を一枚追加する。
任意の局所 Artin $k$-代数 $A$ と、その閉点が $x$ となる framed geometric
五つ組 $\operatorname{Spec} A\to S$ に対し、$t$ を指定した $\operatorname{Spec} A\to T$ の持上げは一意に存在する。
対応する二周期は Serre 適合であり、任意の Artin 商・小拡大に整合する。
\end{corollary}

\begin{proof}
完全な証明は相対モジュライ論文 \cite[\texttt{relative_picard_roots}]{QuadricModuli2026} に置く。
\end{proof}

\paragraph{Abelian 変形と atlas への適用.}
相対根の構成自体は、非線型 locus 全体で特異・可約・非被約な曲線ファイバーを含む。
平坦 Abelian Artin 変形からの適用では、
補題~\ref{lem:cubic-artin-overlap-recovery} と命題~\ref{prop:framed-cubic-artin-lift} が
DDL による helix の持上げから、可逆部分束 $W_0$ とその overlap・縮約を含む
指定基底付き五つ組を回復する。
補題~\ref{lem:artin-two-heart-comparison} は、同じ四項 foundation を持つ二つの圏の
標準 dg Morita 同値が、指定された exact な還元を保って heart-exact となることを示す。
従って系~\ref{cor:all-artin-root} の根を元の Abelian 変形へ移せる。
この指定同型の持上げと有限型効果性を合わせた smooth atlas は
定理~\ref{thm:quintuple-chart-artin-lifting} と定理~\ref{thm:cubic-smooth-open} で構成する。
線型 cubic 点は全て pointed $(\operatorname{coh}(\PP^1\times\PP^1),\OO)$ に対応し、
既存の可換二次曲面の étale Morita-root chart を用いればよい。


% !TeX root = ../nc-quadric-classification.tex
% 2026-11-22 / v1.34（4証明を論文Bへ移管）
\section{可換二次曲面を含む滑らかな Artin 開部分}
\label{sec:quadric-smooth-atlas}

\paragraph{範囲と依存関係.}
定理~\ref{thm:quadric-finite-type-effectivity} が与える実際の
$\mathfrak H_2$ 族と canonical compact 圏の比較から、
可換 quadric を含む有限型で滑らかな Artin 開部分を構成する。
最初に一般の効果的 foundation 族について、必要な入力を明記して表示定理を証明する。
$\mathfrak H_2$ 全体の代数性は結論しない。新規性は未確認である。

\subsection{入力の正確な形}

$k$ は標数零の代数閉体とする。
$\mathfrak H_2$ は局所有限型 $k$-スキームの big fppf site 上で定義された
pointed Abel 圏の族のスタックであり、既存の相対認識定理により
全忠実な射
\[
 \Phi:\mathfrak H_2\longrightarrow\mathfrak P
 :=[X_3/\operatorname{PGL}_9]
\]
を持つとする。$X_3$ は有限型分離 $k$ スキームであり、
$\mathfrak P$ の対角は affine かつ有限表示である。
したがって $\mathfrak H_2$ の Isom 層も同じ性質を持つ。
ここまでは既存の
\ref{prop:relative-h2-stack-recognition} と
\ref{prop:relative-h2-algebraicity-reduction} の結論である。

有限 foundation の標準表示空間を
\[
 S_0=\Hom_k(P\otimes W_1,W_2)
       \oplus\Hom_k(Q\otimes W_2,W_3)
       \simeq\mathbb A_k^{64}
\]
とする。各空間の次元は
$\dim P=\dim Q=\dim W_1=2$、$\dim W_2=4$、$\dim W_3=6$ である。
点 $(F,G)$ は四頂点の有向代数 $\Lambda$ を定め、
頂点 $1,2,3$ の部分代数では積
$\Lambda_{12}\otimes\Lambda_{23}\to\Lambda_{13}$ が同型である。
これらは既存の \ref{prop:foundation-stack} の表示である。
以下では $S\subset S_0$ を開部分とする。
一般の滑らかな foundation chart を用いる場合も、後述の基底の持上げを
その smooth chart の無限小持上げに置き換えれば同じ証明になる。

\begin{theorem}[効果的 foundation 族からの局所 atlas]
\label{thm:atlas-from-effective-foundation}
有限型スキーム $U$、\'etale 射 $a:U\to S$、および射
$q:U\to\mathfrak H_2$ が与えられているとする。
さらに $q$ の族の中に、元の site 内の任意の affine 基底変更と整合する対象の列
$(E_i)_{i\in\mathbb Z}$ があり、次の条件を満たすと仮定する。
\begin{enumerate}
 \item $E_0=\OO$ であり、各 $E_i$ は基底上平坦な noetherian 対象かつ
 canonical 導来圏の compact 対象である。
 \item 元の site 内の任意の affine 基底変更 $\operatorname{Spec} A\to U$ において
 $(E_{i,A})_{i\in\mathbb Z}$ は De Deken--Lowen の意味で geometric $(4,3)$-helix である。
 すなわち各四連続項は $D(\operatorname{Ind}\mathcal C_A)$ の compact generators である
 strong exceptional collection となり、前向きの導来 Hom は次数零の
 有限射影 $A$ 加群であり、helix の mutation 条件を満たす。
 \item $T_A=\bigoplus_{i=0}^3E_{i,A}$ に対して
 \[
  \End(T_A)\simeq a^*\Lambda\otimes_{\OO_U}A
 \]
 が、頂点と積を保ち基底変更と整合して与えられる。
 特に Hom の階数行列は
 \[
 \begin{pmatrix}
 1&2&4&6\\0&1&2&4\\0&0&1&2\\0&0&0&1
 \end{pmatrix}
 \]
 であり、$\Hom(E_1,E_2)\otimes\Hom(E_2,E_3)\to\Hom(E_1,E_3)$ は同型である。
\end{enumerate}
このとき $q$ はスキームで表現可能な smooth 射である。
その fppf 像 $\mathfrak V\subset\mathfrak H_2$ は開部分スタックであり、
有限型で滑らかな Artin スタックにより表現される。
具体的には
\[
 R=U\times_{\mathfrak H_2}U
   \simeq U\times_{\mathfrak P}U,
 \qquad
 \mathfrak V\simeq[U/R].
\]
$R$ は有限型スキームで、$R\to U$ の二つの射は smooth である。
$U$ に可換 quadric $Q$ を表す点があれば、$\mathfrak V$ は $Q$ の開近傍である。
\end{theorem}

この主張で $a$ が \emph{$U$ の全点で} \'etale であることは重要である。
可換点一つでの形式的変形関手の同型だけを仮定していない。

\subsection{指定された Artin 同型を延長する}

\begin{lemma}[foundation の持上げと指定同型]
\label{lem:atlas-artin-lift}
前定理の仮定の下で、$A'\twoheadrightarrow A$ を剰余体 $k$ を持つ
局所 Artin $k$ 代数の small extension とする。
$u_A:\operatorname{Spec} A\to U$、$c_{A'}\in\mathfrak H_2(A')$、
pointed exact 同型
\[
 \alpha_A:q(u_A)\xrightarrow{\sim}c_{A'}|_A
\]
が与えられれば、$u_A$ を延長する $u_{A'}:\operatorname{Spec} A'\to U$ と
$\alpha_A$ を延長する pointed exact 同型
\[
 \alpha_{A'}:q(u_{A'})\xrightarrow{\sim}c_{A'}
\]
が存在する。
\end{lemma}

\begin{proof}
完全な証明は相対モジュライ別稿 \cite[効果的foundationの滑らかな局所atlas]{QuadricModuli2026} に置く。
\end{proof}

\subsection{任意の有限型基底上の smoothness}

\begin{proof}
完全な証明は相対モジュライ別稿 \cite[効果的foundationの滑らかな局所atlas]{QuadricModuli2026} に置く。
\end{proof}

\begin{remark}[全 site への拡張の範囲]
上で構成する $[U/R]$ は通常の全スキーム上の Artin スタックである。
その局所有限型スキームの site への制限が $\mathfrak V$ を表す。
この構成のために $\mathfrak H_2$ 全体の任意基底への連続拡張や極限保存性を
先に証明する必要はない。
一方、$\mathfrak V$ の外側にある $\mathfrak H_2$ の対象については何も結論していない。
\end{remark}

\subsection{可換点近傍での仮定の確認}

\paragraph{\'Etale 性.}
既存の \ref{prop:quadric-etale-morita-root} は、平方根空間の写像が
可換点で \'etale であることを示した後、その \'etale locus に制限し、
可逆双加群の gerbe の \'etale 細分と比較線束の Zariski 自明化を行っている。
従ってそこで取る $U\to S$ は全点で \'etale である。
以後の $U^\circ\subset U$ への開制限もこの性質を保つ。

\paragraph{任意の閉ファイバーでの geometric helix.}
今回の cubic tails 比較は、$E_i=L_i$ と構造 row の対応を全整数 $i$ について与える。
すべての四項窓が compact tilting をなし、全ての $i\leq j$ に対し
\[
 \mathbf R\Hom(E_i,E_j)=\Hom(E_i,E_j)[0]
\]
となる。Serre と次数移動の比較は
\[
 E_{i+4}\simeq S^{-1}E_i[2]
\]
を与え、full exceptional 窓の Serre rotation は
\cite[Definition 5.12]{DeDekenLowen2011} の mutation 条件と一致する。
これらは $Q$ だけでなく、$U^\circ$ の全ての閉ファイバーで成立する。
元の site 内の任意の affine 基底変更でも同じ四項窓・Hom・完全列が存在するので、
補題~\ref{lem:atlas-artin-lift} は $A'\to A$ に直接適用できる。

\paragraph{平坦性と noetherian 性.}
各 row は次数成分ごとに基底上有限射影であり、
その像 $E_i$ への商関手は完全で tensor と可換する。
従って $E_i$ は基底上平坦である。
定理~\ref{thm:quadric-finite-type-effectivity} により
全有限型基底変更後の Noetherian 性が成立し、
これらは coherent heart の noetherian 対象である。
compact 性と unbounded 導来圏の生成は relative cubic tilting 定理による。
特異な基底において全ての coherent 対象が compact であるという主張は使わない。

\paragraph{豊富性と Z-generating 条件.}
全構造 Ext の公式と Noetherian heart の有限表示から、任意の coherent $M$ に対し
\[
 \Ext^p(E_i,M)=0\qquad(p>0,\ i\ll0)
\]
を得る。偶数部分 $E_{2n}=s^n\OO$ の AZ 豊富性は、列全体の
十分負の項の有限直和からの生成を与える。
Grothendieck 圏の任意の全射についても、$E_i$ からの一つの射の像は noetherian なので、
有限な noetherian 部分対象の全射へ帰着できる。
その核への上の Ext 一次消滅が tails-projectivity を与える。
付随する体上の cubic $\mathbb Z$ 代数は Noetherian であり、
$L_{\mathrm{tails}}=T_{\mathrm{tails}}$ である。
従って \cite[Corollary 3.16]{DeDekenLowen2011} の $\mathbb Z$-generating 条件を満たす。
前向き Hom の消滅は、有限次元 $k$ ベクトル空間 $V$ に対する
$\Ext^{1,2}(E_i,V\otimes_k E_j)=0$ $(i\leq j)$ も与える。
これで \cite[Theorems 4.5 and 5.15]{DeDekenLowen2011} による Abel 変形と
finite foundation 変形の比較も、全閉ファイバーで適用できる。

\paragraph{比較の対象.}
relation scheme $R=U^\circ\times_{\mathfrak P}U^\circ$ は
全ての pointed exact 同値を記録する。
foundation の選ばれた四対象を保つ線形群による商だけに置き換えていない。
とくに ruling の交換など、選ばれた窓を別の窓へ送る同値も
$\mathfrak P$ の Isom 層を通じて含まれる。

\begin{corollary}[可換 quadric の有限型 smooth 開近傍]
\label{cor:quadric-effective-smooth-neighborhood}
定理~\ref{thm:quadric-finite-type-effectivity} の
$q:U^\circ\to\mathfrak H_2$ に対し、$Q$ は $\mathfrak H_2$ の有限型で滑らかな Artin 開部分
$\mathfrak V_Q$ に属し、
\[
 U^\circ\longrightarrow\mathfrak V_Q,
 \qquad
 \mathfrak V_Q\simeq
 [U^\circ/(U^\circ\times_{\mathfrak P}U^\circ)]
\]
は smooth atlas とその groupoid 表示である。
\end{corollary}
\begin{proof}
完全な証明は相対モジュライ別稿 \cite[効果的foundationの滑らかな局所atlas]{QuadricModuli2026} に置く。
\end{proof}

\begin{remark}[代数性について残る範囲]
結論は $\mathfrak H_2$ 内部の開部分についてのものであり、
$\mathfrak P$ の開部分と同一視する主張ではない。
また全 $\mathfrak H_2$ の framed 実現部分層 $X_{\mathrm{as}}$ の表現可能性は、
$\mathfrak V_Q$ の外側を覆う同様の効果的 chart が得られるまで残る。
可換点だけの形式的比較、あるいは有限型の一族の存在だけから
この開性を主張してはならない。
本稿では任意の有限型 $T\to\mathfrak H_2$ に対する射 $Z_T\to T$ を実際に構成し、
その全閉点での Artin 持上げを確認したことが追加の根拠である。
\end{remark}


\begin{corollary}[可換点の局所次元]
\label{cor:quadric-open-stack-dimension}
上の $\mathfrak V_Q$ の可換点における接空間の同型類は九次元、
無限小自己同型は六次元であり、スタックの局所次元は三である。
必要なら $Q$ の連結開近傍に制限して、$\mathfrak V_Q$ を純次元三に取れる。
\end{corollary}
\begin{proof}
完全な証明は相対モジュライ別稿 \cite[効果的foundationの滑らかな局所atlas]{QuadricModuli2026} に置く。
\end{proof}

% !TeX root = ../nc-quadric-classification.tex
% 2026-11-21（証明をBへ移管）
\section{一般 cubic helix の持上げと cubic locus の滑らかな開性}
\label{sec:cubic-smooth-atlas}

\paragraph{主張の範囲と依存関係.}
本節では、一般 geometric cubic 点について、De Deken--Lowen の helix 持上げから
四つの隣接 Hom 空間と四次 syzygy の直線を回復する。
可換二次曲面上の Künneth 計算は用いず、指定された pointed 同型と Hom の基底を保って
五つ組の有限パラメータ空間へ持ち上げる。
相対 Picard 根と定理~\ref{thm:geometric-cubic-relative-effectivity} による実際の族を入力とし、
幾何点がちょうど cubic 対象である有限型で滑らかな Artin 開部分を得る。
noncubic 境界の相対認識と $\mathfrak H_2$ 全体の代数性は別の問題である。
本節の相対回復と atlas 構成の新規性は未確認である。

\subsection{五つ組と四項 foundation}

$k$ は標数零の代数閉体とする。
四つの二次元 $k$ 空間 $U_0,U_1,U_2,U_3$ を固定し、
\[
 S\subset\mathscr G\subset
 \mathbb P(U_0\otimes_kU_1\otimes_kU_2\otimes_kU_3)
\]
を geometric 五つ組の開部分 $\mathscr G$ の任意の開部分とする。
ここで $\mathbb P$ は直線をパラメータ化する規約であり、
$S$ 上の tautological 直線を $W$ と書く。
非線型 locus に制限する場合も、以下の持上げ補題はそのまま適用できる。

$S$ 上で縮約により
\[
 R_0=\operatorname{im}
 \bigl(W\otimes U_3^\vee\longrightarrow U_0\otimes U_1\otimes U_2\bigr)
\]
を定める。geometric 条件によりこれは階数二の部分束である。
五つ組の四項 foundation 代数 $\Lambda(W)$ は、頂点 $0,1,2,3$、
隣接矢空間 $U_0,U_1,U_2$ を持ち、三次関係が $R_0$ である代数である。
従ってその前向き成分は
\[
 \Lambda(W)_{ii}=\OO_S,\quad
 \Lambda(W)_{i,i+1}=U_i,\quad
 \Lambda(W)_{i,i+2}=U_i\otimes U_{i+1},\quad
 \Lambda(W)_{03}=(U_0\otimes U_1\otimes U_2)/R_0.
\]
全階数は $24$ である。積は path の連結であり、圏の Hom と比較するときは合成による積を用いる。
同じ右加群規約を両辺で固定する。

普遍 cubic 代数 $D(w)$ を前向きの $\Hom(E_i,E_j)$ の順で読む場合、
右 tails の入力は、添字を反転した反対代数
\begin{equation}\label{eq:quintuple-right-tails-convention}
 C_{ij}=D(w)_{-j,-i},\qquad E_i=\pi e_{-i}C
\end{equation}
である。定理~\ref{thm:relative-cubic-tails-tilting} の全構造 Hom の計算により
\[
 \Hom(E_i,E_j)=C_{-j,-i}=D(w)_{ij}\qquad(i\le j)
\]
となり、上の foundation と一致する。
実際の二周期 $\beta_D:D\to D(2)$ に対しては
$\beta_C=\beta_D^{-1}:C\to C(2)$ を用いる。
この変換は反転した canonical 四周期との平方比較を保つ。
相対族への昇格では、定理~\ref{thm:geometric-cubic-relative-effectivity} をこの $C$ に適用する。
偶数環と右軌道環の比較は定理~\ref{thm:quadric-finite-type-effectivity} の証明と同じで、
右軌道積は $a*b=s^{\deg b}(a)\circ b$ である。
$D$ の通常の右 tails と $C$ の右 tails を同一視することはしない。

\subsection{五つの対象からの有限次数回復}

\begin{lemma}[四次 overlap の普遍的回復]
\label{lem:cubic-artin-overlap-recovery}
$B$ を剰余体 $k$ の可換局所 Artin $k$ 代数とし、
$\mathcal D$ を $B$ 線形 Grothendieck 圏とする。
$E_0,\ldots,E_4$ は compact な $B$-flat 対象であり、
その導来 Hom は前向きに次数零へ集中すると仮定する。
また $(E_0,E_1,E_2,E_3)$ と $(E_1,E_2,E_3,E_4)$ は strong exceptional collections とする。
\[
 D_{ij}=\Hom_{\mathcal D}(E_i,E_j)\quad(0\le i\le j\le4)
\]
は有限自由で、剰余体上への基底変更と可換し、
その剰余体上の積は標準 geometric cubic 代数の対応する窓と一致するとする。
このとき、次の構成が成立する。
\begin{enumerate}
\item $V_i=D_{i,i+1}$ は階数二である。二重積は同型、三重積と四重積は全射であり、
\[
 \operatorname{rk}D_{i,i+r}=2,4,6,9\qquad(r=1,2,3,4)
\]
である（添字が $0,\ldots,4$ に入る範囲で読む）。
\item 三重積の核
\[
 R_i=\ker\bigl(V_i\otimes V_{i+1}\otimes V_{i+2}\longrightarrow D_{i,i+3}\bigr),
 \qquad i=0,1,
\]
は階数二の直和因子である。
\item $T_4=V_0\otimes V_1\otimes V_2\otimes V_3$ と置く。
\[
 W_0=(R_0\otimes V_3)\cap(V_0\otimes R_1)\subset T_4
\]
は階数一の直和因子であり、積による列
\begin{equation}\label{eq:artin-overlap-complex}
0\longrightarrow W_0\longrightarrow
 (R_0\otimes V_3)\oplus(V_0\otimes R_1)
 \longrightarrow T_4\longrightarrow D_{04}\longrightarrow0
\end{equation}
は分裂完全な有限自由複体である。
最初の射は $w\mapsto(w,-w)$、中央の射は二つの包含の和である。
\item 縮約は同型
\begin{equation}\label{eq:artin-contraction-isos}
 W_0\otimes V_3^\vee\xrightarrow{\sim}R_0,
 \qquad
 W_0\otimes V_0^\vee\xrightarrow{\sim}R_1
\end{equation}
を与える。従って $(V_0,V_1,V_2,V_3,W_0)$ は $B$ 上の geometric 五つ組である。
\end{enumerate}
これらの有限加群、核、交わり、縮約の構成は、任意の可換 $B$ 代数への
基底変更と可換する。
圏の Hom の基底変更と比較するときは、与えられた相対 compact 圏の基底変更同型を用いる。
\end{lemma}

\begin{proof}
全証明は相対モジュライ論文 \cite[\texttt{cubic_smooth_atlas}]{QuadricModuli2026} に置く。
\end{proof}

\begin{corollary}[四項 foundation の完全な同定]
\label{cor:quintuple-four-window}
前補題で四つの基底同型 $U_i\otimes B\xrightarrow{\sim}V_i$ を選び、
それによる五つ組の $B$ 点を $w_B$ とする。
このとき合成は、頂点と全ての Hom の作用を保つ同型
\[
 \Lambda(w_B)\xrightarrow{\sim}
 \End_{\mathcal D}\left(\bigoplus_{i=0}^3E_i\right)
\]
を与える。
\end{corollary}

\begin{proof}
全証明は相対モジュライ論文 \cite[\texttt{cubic_smooth_atlas}]{QuadricModuli2026} に置く。
\end{proof}

\subsection{指定された還元同型と基底を持つ持上げ}

\begin{proposition}[一般 cubic 点での framed Artin 持上げ]
\label{prop:framed-cubic-artin-lift}
$A'\twoheadrightarrow A$ を剰余体 $k$ を持つ局所 Artin $k$ 代数の冪零拡大とする。
$\mathcal G_A$ は flat な $A$ 線形 Grothendieck 圏であり、
$A$-flat な対象の geometric $(4,3)$-helix $(E_{i,A})_{i\in\mathbb Z}$ を持つとする。
全ての前向き Hom は有限自由で基底変更と可換し、閉ファイバーは標準 geometric cubic helix であるとする。
特に各四項窓は $D(\mathcal G_A)$ の compact tilting generator である。

$\mathcal G_{A'}$ を $\mathcal G_A$ の flat Grothendieck 変形とする。
同時に指定された還元同型を持つ flat 構造対象 $\OO_{A'}$ があり、
$E_{0,A}=\OO_A$ と同定されているとする。
四つの基底同型
\[
 \varphi_{i,A}:U_i\otimes A\xrightarrow{\sim}
 \Hom(E_{i,A},E_{i+1,A}),\qquad i=0,1,2,3,
\]
を指定し、回復される五つ組の点を $w_A\in S(A)$ とする。
このとき次が存在する。
\begin{enumerate}
\item flat な geometric helix $(E'_{i})_{i\in\mathbb Z}$ と指定還元同型
  $A\otimes_{A'}E'_i\xrightarrow{\sim}E_{i,A}$。
  $E'_0=\OO_{A'}$ とでき、その還元同型も指定したものに取れる。
\item $\varphi_{i,A}$ を延長する基底同型
\[
 \varphi_{i,A'}:U_i\otimes A'\xrightarrow{\sim}\Hom(E'_i,E'_{i+1})
\]
と、これらから回復される $w_{A'}\in S(A')$。
その $A$ への還元は\emph{指定された点} $w_A$ に等しい。
\item 頂点と全ての積を保つ同型
\[
 \Lambda(w_{A'})\xrightarrow{\sim}\End\left(\bigoplus_{i=0}^3E'_i\right)
\]
であって、その還元は $\varphi_{i,A}$ と元の合成から定まる同型に一致するもの。
\end{enumerate}
$W_{0,A}$ の生成元も指定されていれば、その生成元を $W_{0,A'}$ の生成元へ持ち上げられる。
射影空間 $S$ の点の持上げはこの追加の生成元の選択を必要としない。
\end{proposition}

\begin{proof}
全証明は相対モジュライ論文 \cite[\texttt{cubic_smooth_atlas}]{QuadricModuli2026} に置く。
\end{proof}

\subsection{pointed exact 同型の還元を保つ比較}

\begin{lemma}[Artin 基底上の二つの heart の比較]
\label{lem:artin-two-heart-comparison}
$B$ は局所 Artin 環とし、$\mathcal G_1,\mathcal G_2$ を flat な $B$ 線形
Grothendieck 圏とする。それぞれ flat compact tilting generator $T_1,T_2$ を持ち、
その End 代数は $B$ 上有限自由で、Artin 商への基底変更と可換するとする。
各 Artin 商でもこれらの還元は compact tilting generators と仮定する。
頂点を保つ End 代数の同型から標準 dg Morita 同値
\[
 F:D(\mathcal G_1)\xrightarrow{\sim}D(\mathcal G_2)
\]
を定める。
その閉ファイバー $F_k$ が指定された exact 同値の導来化なら、$F$ も heart-exact である。
さらに $B\twoheadrightarrow A$ 上の還元 Morita 同値が
指定された exact 同値 $\alpha_A$ の導来化に一致すれば、
$F$ の heart への制限は $\alpha_A$ を延長する。
指定された構造対象とその比較同型を保つ場合は pointed 同値となる。
\end{lemma}

\begin{proof}
全証明は相対モジュライ論文 \cite[\texttt{cubic_smooth_atlas}]{QuadricModuli2026} に置く。
\end{proof}

\begin{theorem}[五つ組への étale chart の指定 Artin 持上げ]
\label{thm:quintuple-chart-artin-lifting}
$S\subset\mathscr G$ を上の開部分とし、$a:U\to S$ を有限型の étale 射とする。
実際の相対族 $q:U\to\mathfrak H_2$ が構成されており、任意の affine 基底変更に整合する
flat compact geometric cubic helix $(E_i)$、$E_0=\OO$ を持つとする。
さらに $a$ の普遍五つ組とこの helix の五つ組が同定され、
その四項窓の End 代数は頂点と積を保って $a^*\Lambda(W)$ と同定されるとする。
全ての幾何点で helix は標準 geometric cubic helix と仮定する。

$A'\twoheadrightarrow A$ を局所 Artin $k$ 代数の small extension とし、
$u_A\in U(A)$、$c_{A'}\in\mathfrak H_2(A')$、および\emph{指定された} pointed exact 同型
\[
 \alpha_A:q(u_A)\xrightarrow{\sim}c_{A'}|_A
\]
を取る。
このとき $u_A$ を延長する $u_{A'}\in U(A')$ と、$\alpha_A$ を延長する pointed exact 同型
\[
 \alpha_{A'}:q(u_{A'})\xrightarrow{\sim}c_{A'}
\]
が存在する。いずれの段階にも Hochschild 一次または二次上の根の作用の仮定は不要である。
\end{theorem}

\begin{proof}
全証明は相対モジュライ論文 \cite[\texttt{cubic_smooth_atlas}]{QuadricModuli2026} に置く。
\end{proof}

\subsection{smoothness の帰結と限界}

\begin{corollary}[全 closed cubic 点での smooth chart]
\label{cor:quintuple-effective-smooth-chart}
前定理の入力に加え、既存の相対認識から全忠実な射
\[
 \Phi:\mathfrak H_2\longrightarrow\mathfrak P=[X_3/\operatorname{PGL}_9]
\]
があり、$\mathfrak P$ の対角は affine かつ有限表示であるとする。
このとき $q:U\to\mathfrak H_2$ はスキームで表現可能な smooth 射である。
その fppf 像は $\mathfrak H_2$ の開部分スタックで、有限型で滑らかな Artin スタックである。
\end{corollary}

\begin{proof}
全証明は相対モジュライ論文 \cite[\texttt{cubic_smooth_atlas}]{QuadricModuli2026} に置く。
\end{proof}

\subsection{全 cubic locus の有限型 smooth 開部分}

\begin{theorem}[cubic 対象全体の smooth Artin 開部分]
\label{thm:cubic-smooth-open}
$\mathfrak H_2$ には有限型で滑らかな Artin 開部分
\[
 \mathfrak V_{\mathrm{cub}}\subset\mathfrak H_2
\]
が存在し、その幾何点はちょうど pointed cubic 表示を持つ対象である。
より正確には、任意の局所有限型基底 $T$ 上の相対 $\mathfrak H_2$ 対象が
$\mathfrak V_{\mathrm{cub}}(T)$ に属するための必要十分条件は、
全ての幾何ファイバーが pointed cubic 表示を持つことである。

有限型の滑らかなスキーム $U_{\mathrm{cub}}$ と representable smooth 全射
\[
 q_{\mathrm{cub}}:U_{\mathrm{cub}}\longrightarrow\mathfrak V_{\mathrm{cub}}
\]
を取れる。その同値関係は全ての pointed exact 同値を用いる
\[
 R_{\mathrm{cub}}=U_{\mathrm{cub}}\times_{\mathfrak H_2}U_{\mathrm{cub}}
 \simeq U_{\mathrm{cub}}\times_{\mathfrak P}U_{\mathrm{cub}},\qquad
 \mathfrak V_{\mathrm{cub}}\simeq[U_{\mathrm{cub}}/R_{\mathrm{cub}}]
\]
であり、選んだ五つ組や根を保つ同値だけに限定しない。
\end{theorem}

\begin{proof}
全証明は相対モジュライ論文 \cite[\texttt{cubic_smooth_atlas}]{QuadricModuli2026} に置く。
\end{proof}

\begin{corollary}[全 cubic 点の形式変形と次元]
\label{cor:cubic-formal-h2-abelian}
$x=(\mathcal C,\OO)$ を $\mathfrak V_{\mathrm{cub}}(k)$ の点とする。
剰余体 $k$ の局所 Artin 基底上で、閉ファイバーの同定を持つ相対 $\mathfrak H_2$ 変形から
flat pointed Abel 変形への忘却は groupoid の同値である。
従ってこの形式変形問題は無障害であり、四項 foundation 代数を $\Lambda$ とすると、
その接空間と無限小自己同型空間はそれぞれ
\[
 \operatorname{HH}^2(\Lambda),\qquad\operatorname{HH}^1(\Lambda)
\]
である。さらに $\mathfrak V_{\mathrm{cub}}$ は純次元三である。
ここでは根・根の平方同型・AS4 被覆は標識として固定しない。
\end{corollary}

\begin{proof}
全証明は相対モジュライ論文 \cite[\texttt{cubic_smooth_atlas}]{QuadricModuli2026} に置く。
\end{proof}

\begin{remark}[noncubic 境界との区別]
$\mathfrak V_{\mathrm{cub}}$ の幾何点は cubic 対象だけであり、
定理~\ref{thm:complete-h2-classification} の noncubic 対象は含まない。
従ってこの定理は $\mathfrak H_2$ 全体の代数性や、noncubic 境界に沿う相対認識を結論しない。
また $\mathfrak V_{\mathrm{cub}}$ を $[X_3/\operatorname{PGL}_9]$ の開部分と同一視する主張でもない。
結論に必要なのは、相対 Picard 作用と二周期の実際の構成、同じ根の相対 AS4 被覆、
全 Noetherian 性、canonical tails 比較、そして本節の指定 Artin 同型の持上げを合わせたことである。
\end{remark}

% !TEX root = ../nc-quadric-classification.tex
% 2026-11-21 / v1.36
\section{明示 noncubic 表示族の相対内在的実現}
\label{sec:relative-noncubic-families}

第~\ref{sec:noncubic-stack}節の中央二重平方族について、
canonical compact 圏と任意の Noetherian 基底変更を含む全相対条件を確認する。
重み付き射影スタック上の Clifford 代数として heart を表し、
その対角加群の perfectness を幾何ファイバーから持ち上げる。
これにより表示スタック $\mathfrak N^{\mathrm{pres}}$ から
内在的 $\mathfrak H_2$ への実際の射を得る。
全ての基底上平坦 coherent 対象が compact であるという追加条件は用いない。
任意の noncubic 特殊ファイバーを持つ相対族の正規形や、
cubic を横断する開近傍の効果性とは区別する。

\subsection{中央有限代数に対する fiberwise perfectness の持上げ}

\begin{lemma}\label{lem:nc-fiber-perfect}
$R\to D$ を Noetherian 可換環の準同型、$H$ を中心に $D$ の作用する
有限 $D$-代数、$M$ を有限 $D$-加群である右 $H$-加群とする。
$H,M$ は $R$ 上平坦と仮定する。
$\mathfrak q\in\operatorname{Spec} D$、$\mathfrak p=\mathfrak q\cap R$ とする。
対応する fiber の点において
\[
 M\otimes_R\kappa(\mathfrak p)
 \quad\text{が}\quad
 H\otimes_R\kappa(\mathfrak p)\text{ 上 perfect}
\]
ならば、$M$ は $\mathfrak q$ のある近傍で $H$ 上 perfect である。
\end{lemma}

\begin{proof}
完全な証明は相対モジュライ論文 \cite[\texttt{relative_noncubic_families}]{QuadricModuli2026} に置く。
\end{proof}

\begin{remark}
この補題は、$R$ 上の平坦性を省いた fiber 判定ではない。
特に有限自由解消の syzygy の $R$-平坦性と、最後の
$\operatorname{Tor}_1^R(K,\kappa(\mathfrak p))=0$ の両方を使う。
基底の被約性、正則性、有限 Krull 次元は用いない。
\end{remark}

\subsection{noncubic heart の weighted stack 表示}

$k$ を標数零の体、$R$ を可換 Noetherian $k$-代数とする。
判別式が可逆な二元二次形式 $f_1,f_2\in R[u,v]_2$ を取り、
\[
 B=R[u,v]\langle c,d\rangle/(c^2-f_1,d^2-f_2),
 \qquad h=cd+dc,
 \qquad C=R[u,v,h],\quad\deg(u,v,h)=(1,1,2)
\]
と置く。$u,v$ は中心にある。
補題~\ref{lem:noncubic-relative-clifford-center} より
\[
 B=C\oplus Cc\oplus Cd\oplus C(cd)
\]
であり、この表示は全ての基底変更と両立する。
$X=\PP_R(1,1,2)=[(\operatorname{Spec} C\setminus V(C_+))/\mathbb G_m]$ とし、
$\mathcal B$ を $B$ の対応する $X$ 上の代数層とする。
\[
 \mathcal B\simeq\OO_X\oplus\OO_X(-1)^{\oplus2}\oplus\OO_X(-2)
 \quad\text{as }\OO_X\text{-modules}.
 \tag{NCF-1}\label{eq:nc-clifford-sheaf}
\]

\begin{proposition}\label{prop:nc-qgr-weighted}
構造対象を保つ完全同値
\[
 \operatorname{QGr} B\simeq\operatorname{QCoh}(X,\mathcal B),\qquad
 \qgr B\simeq\coh(X,\mathcal B)
\]
がある。これは任意の Noetherian 基底変更における module-object 比較と整合する。
coherent 対象間の Hom は有限 $R$-加群であり、任意の平坦基底変更と可換する。
対象・射・完全列について fppf descent が成立する。
\end{proposition}

\begin{proof}
完全な証明は相対モジュライ論文 \cite[\texttt{relative_noncubic_families}]{QuadricModuli2026} に置く。
\end{proof}

\begin{lemma}[weighted cohomology]\label{lem:nc-weighted-cohomology}
任意の $j\in\mathbb Z$ について $\mathbf R\Gamma(X,\OO_X(j))$ は
有限 projective $R$-複体であり、任意の基底変更と可換する。
\[
 H^0(X,\OO_X(j))=C_j,\qquad H^1(X,\OO_X(j))=0,
 \qquad H^2(X,\OO_X(j))=C_{-j-4}^{\vee},
\]
ただし負次数の $C$ は零とする。
従って $\mathbf R\Gamma(X,\mathcal B(j))$ も同じ普遍性を持つ。
特に $n\ge0$ なら
\[
 \mathbf R\Gamma(X,\mathcal B(n))=B_n[0].
 \tag{NCF-2}\label{eq:nc-full-orbit-cohomology}
\]
\end{lemma}

\begin{proof}
完全な証明は相対モジュライ論文 \cite[\texttt{relative_noncubic_families}]{QuadricModuli2026} に置く。
\end{proof}

\subsection{canonical compact 圏と相対 smooth proper 性}

\begin{proposition}\label{prop:nc-compact-generator}
\[
 G=\bigoplus_{i=0}^3\mathcal B(i)
\]
は $D(\operatorname{QCoh}(X,\mathcal B))$ の compact generator である。
その dg 自己準同型代数を $E_R=\mathbf R\End(G)$ とすると、
$E_R$ は $R$ 上 perfect であり、任意の Noetherian $R$-代数 $R'$ に対して
\[
 E_R\otimes_R^{\mathbf L}R'\simeq E_{R'},\qquad
 D(\operatorname{QGr} B)^c\simeq\operatorname{Perf}(E_R)
\]
となる。この同定は実際の module-object 引戻しの導来関手に対して
\[
 D(\operatorname{QGr} B)^c\otimes_{\operatorname{Perf} R}\operatorname{Perf} R'
 \simeq D(\operatorname{QGr} B_{R'})^c
 \tag{NCF-3}\label{eq:nc-perfect-basechange}
\]
を与える。compact 対象は $D^b(\qgr B)$ に属する。
\end{proposition}

\begin{proof}
完全な証明は相対モジュライ論文 \cite[\texttt{relative_noncubic_families}]{QuadricModuli2026} に置く。
\end{proof}

\begin{remark}
$G$ を tilting と主張していない。
例えば負の差に \eqref{eq:nc-clifford-sheaf} を適用すると $H^2$ が現れる。
有限な通常の End 代数に置き換えず、$E_R$ は dg 代数として扱う。
\end{remark}

\begin{proposition}\label{prop:nc-relative-smooth}
上の canonical compact 圏は $R$ 上 smooth proper である。
この結論は任意の Noetherian 基底変更後に成立する。
\end{proposition}

\begin{proof}
完全な証明は相対モジュライ論文 \cite[\texttt{relative_noncubic_families}]{QuadricModuli2026} に置く。
\end{proof}

\subsection{同じ根・全軌道・相対 AS4 被覆}

\begin{proposition}\label{prop:nc-root-ample-orbit}
étale 局所に、$B$ の次数付き自己同型 $\sigma$ で
\[
 \sigma^2=\iota,\qquad
 \iota(u)=u,\quad\iota(v)=v,\quad\iota(c)=-c,\quad\iota(d)=-d
\]
となるものを選べる。
canonical compact 圏の相対 Serre 関手と、heart の完全自己同値は
\[
 S=(-)_\iota(-2)[2],\qquad s=(-)_\sigma(1),\qquad s^2=S^{-1}[2]
\]
を満たす。$(\mathcal B,s)$ は AZ ample であり、
右軌道積の全軌道環は $B^\sigma$ に等しい。
これらは任意の Noetherian 基底変更と整合する。
\end{proposition}

\begin{proof}
完全な証明は相対モジュライ論文 \cite[\texttt{relative_noncubic_families}]{QuadricModuli2026} に置く。
\end{proof}

\begin{proposition}[全相対 AS4 条件]\label{prop:nc-full-relative-as4}
前命題の選択の下で
\[
 A'=R[u,v]\langle c,d\rangle/
 \bigl((c^2-f_1)-\eta(d^2-f_2)\bigr),\qquad
 g'=(c^2-f_1)+\eta(d^2-f_2)
\]
と置く。
$A=(A')^\sigma$ と $g'$ は、定義~\ref{def:relative-h2}の相対 quadratic AS4 被覆の全条件を満たし、
\[
 A/(g')=B^\sigma
\]
である。特に全ての Noetherian 基底変更後も $A$ は左右 Noetherian であり、
左右・双加群 Koszul 解消および正則正規切断は任意の底変更後も完全である。
\end{proposition}

\begin{proof}
完全な証明は相対モジュライ論文 \cite[\texttt{relative_noncubic_families}]{QuadricModuli2026} に置く。
\end{proof}

\subsection{表示スタックからの実際の相対族と残る境界問題}

\begin{theorem}\label{thm:noncubic-relative-effectivity}
第~\ref{sec:noncubic-stack}節の表示スタック
\[
 \mathfrak N^{\mathrm{pres}}=[P/G],\qquad
 P=Q^{\mathrm{nd}}\times Q^{\mathrm{nd}}
\]
上の降下した反 Serre 環 $\mathcal R$ は、任意の有限型 affine test $T$ 上で
\[
 (\mathcal C_T,\mathcal O_T)=(\qgr\mathcal R_T,\pi\mathcal R_T)
\]
という実際の相対 $\mathfrak H_2$ 族を定める。
従って自然な射
\[
 q_{\mathrm{nc}}:\mathfrak N^{\mathrm{pres}}\longrightarrow\mathfrak H_2
\]
が存在し、その全ての幾何ファイバーは noncubic である。
ここでは平方根・AS4 被覆は局所的 witness であり、出力の標識に加えない。
\end{theorem}

\begin{proof}
完全な証明は相対モジュライ論文 \cite[\texttt{relative_noncubic_families}]{QuadricModuli2026} に置く。
\end{proof}

\begin{remark}[noncubic 境界との区別]
$q_{\mathrm{nc}}$ の存在だけでは、noncubic 点を含む開近傍の atlas は得られない。
特に任意の Artin 変形は、閉 fiber が noncubic でも
隣接乗法の四次 minor が非零の冪零元となり得る。
従って二形式表示の次元一のパラメータ族で全ての横断変形を捉えたことにはならない。
任意の相対 $\mathfrak H_2$ 族の noncubic 表示への局所還元、
この射の相対全忠実性・monomorphism 性、および cubic を横断する近傍の効果性は別途証明が必要である。
ここでの対角加群の議論は、明示的な中央有限な Clifford order を持つ族に適用したものであり、
一般の root-family が同じ中央有限表示を持つとは仮定していない。
\end{remark}

% !TEX root = ../nc-quadric-classification.tex
% 2026-11-21 / v1.39
\section{noncubic 境界の適合根と Hochschild 作用}
\label{sec:noncubic-hh-roots}


\paragraph{結論と範囲.}
中央二重平方表示の二つの根対を交換する\emph{trace 零の}線型対合を選ぶと、
対応する適合 Serre 根は有限 foundation の $\operatorname{HH}^1,\operatorname{HH}^2$ に恒等に作用する。
ただし根対が一致するときは対角の交換対合を選ぶ。
trace 零だけでは十分でないことと、選択条件の具体的な行列式を後で示す。
この計算は noncubic 表示の一次元の変形だけを扱うものではない。
foundation の全変形接空間、すなわち根対が異なる場合の三次元、
一致する場合の四次元を扱う。
従って既存の Artin 根持上げ補題を適用すると、指定した良い根は
全ての Artin Abel 変形へ持ち上がり、その Morita 類は一意である。
さらに有限 foundation の各 noncubic 点には、
この根と実際の Serre 平方同型を持つ有限型 étale 近傍が存在する。
同じ根に対する AS4 被覆の横断変形への持上げと、有限型基底上の全相対認識は
本節の結論に含めない。
本節の計算の新規性は未確認である。

\subsection{四項 foundation の二項移動}

$k$ を標数零の代数閉体とする。
full geometric helix の四項を $E_0,E_1,E_2,E_3$ とし、
\[
A=\Hom(E_0,E_1),\quad B=\Hom(E_1,E_2),\quad
C=\Hom(E_2,E_3),\quad D=\Hom(E_0,E_2)
\]
と置く。$A,B,C$ は二次元、$D$ は四次元である。
$W=\Hom(E_0,E_3)$ は六次元であり、
\[
F:A\otimes B\longrightarrow D,\qquad
G:C\otimes D\longrightarrow W
\]
が foundation の積を与える。
因子の並べ替えを省略して書く。
$G$ は全射とし、$K=\ker G$ は二次元とする。
包含 $K\subset C\otimes D$ を curry した写像を
\[
j:D^{\vee}\longrightarrow C\otimes K^{\vee}
\]
と書く。さらに
\begin{equation}\label{eq:nchh-delta}
\delta=(-F^{\vee},j):D^{\vee}\longrightarrow
H=(A^{\vee}\otimes B^{\vee})\oplus(C\otimes K^{\vee})
\end{equation}
は単射とする。これらは、以下で使う noncubic 点の近傍で開条件である。
$\mathscr L=\operatorname{im}\delta\in\Gr(4,H)$ と置く。

\begin{lemma}[二項移動の有限線型代数]\label{lem:nchh-gale}
$E_{i+4}=S^{-1}E_i[2]$ により延長し、四項窓を
$(E_2,E_3,E_4,E_5)$ へ移す。
対応するデータは
\begin{align}
A'&=C,& B'&=K^{\vee},& C'&=A,\label{eq:nchh-new-small}\\
D'&=H/\mathscr L,&
F'&:C\otimes K^{\vee}\longrightarrow D',\label{eq:nchh-new-D}\\
K'&=\operatorname{im}\bigl(B^{\vee}\longrightarrow A\otimes D'\bigr),\label{eq:nchh-new-K}
\end{align}
である。\eqref{eq:nchh-new-K} は $A$ の coevaluation と第一 summand の商写像の合成である。
$K'\simeq B^{\vee}$ となる開近傍で、この新しい四平面は
\begin{equation}\label{eq:nchh-gale}
\mathscr L'=\{(-Y,X):(X,Y)\in\mathscr L^{\perp}\}
 \subset(C^{\vee}\otimes K)\oplus(A\otimes B)
\end{equation}
である。ここで $\perp$ は $H$ と $H^{\vee}$ の標準評価による annihilator を表す。
この記述は任意の局所 Artin 変形にも同じ式で適用できる。
\end{lemma}

\begin{proof}
全証明は相対モジュライ論文 \cite[\texttt{noncubic_hh_roots}]{QuadricModuli2026} に置く。
\end{proof}

\begin{remark}[全変形空間を使うこと]
$G$ の商の基底を忘れ、代わりに $K$ の基底を選ぶと、
foundation のデータは $(F,j)$ である。
$\delta$ が単射である開部分で $D$ の基底変更を除くことは、
四平面 $\mathscr L\subset H$ を取ることに等しい。
従って全 foundation 変形の局所模型は
$\Gr(4,8)$ の開部分を
\[
\widetilde G=\GL(A)\times\GL(B)\times\GL(C)\times\GL(K)
\]
で割ったものとなる。
第一 summand への作用は $g_A^{-t}\otimes g_B^{-t}$、
第二 summand への作用は $g_C\otimes g_K^{-t}$ である。
自明に作用する三次元中心を除いた群を $G_{\mathrm{eff}}$ と書く。
その次元は十三で、Grassmannian の次元は十六である。
この模型には noncubic 表示を保たない全ての横断方向が含まれている。
\end{remark}

\subsection{中央二重平方と良い根の選択}

\begin{lemma}[trace 零の交換対合]\label{lem:nchh-exchange}
二つの非退化二元二次形式 $f_1,f_2$ に対し、座標とそれぞれの平方根変数の
定数倍を選び、
\begin{equation}\label{eq:nchh-normal-forms}
f_1=uv,\qquad f_2=\ell_1\ell_2,
\qquad
\begin{pmatrix}\ell_1\\\ell_2\end{pmatrix}
=L_0\begin{pmatrix}u\\v\end{pmatrix},\qquad
L_0=\begin{pmatrix}\alpha&\beta\\\gamma&-\alpha\end{pmatrix},\quad
\alpha^2+\beta\gamma=1
\end{equation}
とできる。二根対が一致する場合にも $L_0=\operatorname{diag}(1,-1)$ と選べる。
\end{lemma}
\begin{proof}
全証明は相対モジュライ論文 \cite[\texttt{noncubic_hh_roots}]{QuadricModuli2026} に置く。
\end{proof}

以下、
\[
\mathcal B=k[u,v]\langle c,d\rangle/(c^2-uv,d^2-\ell_1\ell_2)
\]
と置く。$u,v$ は中央である。
\begin{equation}\label{eq:nchh-sigma}
\sigma(u,v)=L_0(u,v),\qquad\sigma(c)=d,\qquad\sigma(d)=-c
\end{equation}
は次数付き自己同型で、$\sigma^2=\iota$、ただし
$\iota(u)=u,\iota(v)=v,\iota(c)=-c,\iota(d)=-d$ である。
命題~\ref{prop:noncubic-relative-serre-on-perfect} の双加群計算から
$S=(-)_\iota(-2)[2]$ であり、
\begin{equation}\label{eq:nchh-good-root}
s=(-)_\sigma(1),\qquad s^2=S^{-1}[2]
\end{equation}
は適合根である。同じ根の AS4 表示は定理~\ref{thm:ore-anticommuting-central} の Ore/Zhang 構成が与える。

\begin{proposition}[根の作用と交代 Gale 双対]\label{prop:nchh-root-gale}
\eqref{eq:nchh-good-root} の Hochschild 作用は、前節の Grassmann 模型で
\begin{equation}\label{eq:nchh-involution}
\tau(\mathscr L)=J\mathscr L^{\perp},\qquad
J(X,Y)=(D_0Y,-D_0X),\qquad D_0=\operatorname{diag}(1,-1)
\end{equation}
の微分により与えられる。$X,Y$ を二次正方行列として書いている。
$J:H^{\vee}\to H$ は交代同型であり、中央二重平方の点 $\mathscr L$ は
$\tau$ の固定点、すなわち対応する交代形式の Lagrangian 四平面である。
\end{proposition}

\begin{proof}
全証明は相対モジュライ論文 \cite[\texttt{noncubic_hh_roots}]{QuadricModuli2026} に置く。
\end{proof}

\subsection{全 Hochschild 接空間への作用}

\begin{theorem}\label{thm:noncubic-hh-trivial-root}
任意の noncubic $\mathcal H_2$ 対象の有限 foundation を $\Gamma$ とする。
補題~\ref{lem:nchh-exchange} の trace 零の交換対合を用い、
根対が一致する場合には $L_0=D_0$ と選ぶ。
この選択による適合根 $s$ は
\[
s_*|_{\operatorname{HH}^1(\Gamma)}=1,\qquad s_*|_{\operatorname{HH}^2(\Gamma)}=1
\]
を満たす。
二つの根対が異なる場合は $\dim\operatorname{HH}^1=0$, $\dim\operatorname{HH}^2=3$、
一致する場合はそれぞれ $1,4$ である。
\end{theorem}

\begin{proof}
全証明は相対モジュライ論文 \cite[\texttt{noncubic_hh_roots}]{QuadricModuli2026} に置く。
\end{proof}

\begin{corollary}[全 Artin 変形への根の持上げ]\label{cor:noncubic-root-all-artin}
前定理の根 $s$ は、閉ファイバーを固定した任意の局所 Artin Abel 変形へ、
相対 Serre 平方根として持ち上がる。
その Morita 自己同値の同型類は一意である。
この一意性は平方同型自体の一意な選択を意味しない。
\end{corollary}
\begin{proof}
全証明は相対モジュライ論文 \cite[\texttt{noncubic_hh_roots}]{QuadricModuli2026} に置く。
\end{proof}

\subsection{全 foundation 変形の Lagrangian 正規形}
\label{subsec:noncubic-lagrangian-slice}

前の接計算は、有限型の具体的な局所表示に強められる。
この表示は foundation の全横断方向を含むが、
それだけで相対 $\mathcal H_2$ の認識を与えるものではない。

四つの二次元空間をそれぞれ $k^2$ と同定し、
第\ref{sec:noncubic-hh-roots}節の $H$ と $J$ を使う。
$\omega$ を行列 $J^{-1}$ の非退化交代形式とする。
$X=\Gr(4,H)$、$L=\operatorname{LG}(4,H,\omega)$ と置けば、
$L$ は $\tau(\mathscr L)=J\mathscr L^\perp$ の固定部分スキームである。
基底変更群の作用は
\[
 \rho(g_A,g_B,g_C,g_K)
 =\operatorname{diag}(g_A^{-t}\otimes g_B^{-t},
                      g_C\otimes g_K^{-t})
\]
である。$X$ 上で自明に作用する中心を
\[
 C_0=\{(aI,bI,cI,dI):a^{-1}b^{-1}=cd^{-1}\}
       \simeq\mathbf G_m^3,
 \qquad G=\widetilde G/C_0
\]
と書く。

\begin{proposition}[Lagrangian な有限型 \texorpdfstring{\'etale}{etale} 表示]
\label{prop:noncubic-lagrangian-etale-slice}
群対合
\begin{equation}\label{eq:noncubic-group-involution}
 \widetilde\theta(g_A,g_B,g_C,g_K)
 =(D_0g_CD_0,g_K^{-t},D_0g_AD_0,g_B^{-t})
\end{equation}
は $G$ の対合 $\theta$ を誘導し、
$\tau(g\mathscr L)=\theta(g)\tau(\mathscr L)$ を満たす。
その固定群の単位成分を $G^+=(G^\theta)^\circ$ と書くと、
\[
 G^+\simeq
 (\GL_2\times\GL_2)/(\mathbf G_m\times\mu_2),
 \qquad \dim G^+=7.
\]
ここで中央部分群の埋込みは
$(z,\epsilon)\mapsto(zI,\epsilon z^{-1}I)$ である。

$G$ の Lie 代数の負固有空間を $\mathfrak g^-$ とする。
自然な軌道微分から得る束写像
\begin{equation}\label{eq:noncubic-normal-orbit-map}
 q:\mathfrak g^-\otimes\mathcal O_L
   \longrightarrow N_{L/X}
   \simeq\bigwedge^2\mathscr U^\vee
\end{equation}
を考える。$\mathscr U$ は $L$ 上の tautological 四平面束である。
その行列式の非消滅開を $L^{\mathrm{tr}}$ とすると、
\[
 [L^{\mathrm{tr}}/G^+]\longrightarrow[X/G]
 \tag{LG-slice}\label{eq:noncubic-lagrangian-stack-map}
\]
は表現可能な \texorpdfstring{\'etale}{etale} 射である。
双方のスタックの次元は三である。
$\det q$ は、定数一次元因子を除けば
Pl\"ucker 線束の三乗 $\mathcal O_L(3)$ の切断である。

中央二重平方の trace 零正規形
\eqref{eq:nchh-plane} が $L^{\mathrm{tr}}$ に属する必要十分条件は
$\alpha\ne0$ である。
根対が異なればこの条件は自動的に成立し、
根対が一致するときは補題~\ref{lem:nchh-exchange} の通り
$L_0=D_0$ と選べば成立する。
従ってその近傍の全 foundation 変形は、
\texorpdfstring{\'etale}{etale} 局所に Lagrangian 四平面で表示できる。
同じ結論は、補題~\ref{lem:nchh-gale} の階数条件を課した
foundation の開部分へ制限したときにも成立する。
\end{proposition}

\begin{proof}
全証明は相対モジュライ論文 \cite[\texttt{noncubic_hh_roots}]{QuadricModuli2026} に置く。
\end{proof}

\begin{remark}[この表示が与えるもの]
\label{rem:noncubic-lagrangian-scope}
上の射の像は foundation スタックの開部分であり、
全 noncubic 点を含む。その smooth 性は有限代数表示についての主張である。
実際の Morita 平方根との比較には
系~\ref{cor:noncubic-etale-morita-root} を使い、
必要に応じて二つの \texorpdfstring{\'etale}{etale} 近傍の積を取る。
全正軌道 Hom、相対 heart、同じ根の AS4 被覆の存在を、
この Lagrangian 正規形だけから結論してはならない。
特に \eqref{eq:noncubic-lagrangian-stack-map} の標的を
内在的 $\mathfrak H_2$ に置き換えたとは主張していない。
\end{remark}

\begin{remark}[trace 零だけでは良い根を特徴付けない]
\label{rem:noncubic-offdiagonal-root}
根対一致の場合、$L_0=\left(\begin{smallmatrix}0&1\\1&0\end{smallmatrix}\right)$
も trace 零の交換対合であるが、上の \texorpdfstring{\'etale}{etale} 開には属さない。
対応する根は連結慣性
$u\mapsto zu,v\mapsto z^{-1}v$ を逆元へ移すため、
$\operatorname{HH}^1$ に $-1$ で作用する。
\eqref{eq:nchh-tangent-sequence} の正負分解より、この場合
\[
 (\dim\operatorname{HH}^2_+,\dim\operatorname{HH}^2_-)=(3,1).
\]
従ってこの根は全ての一次 Abel 変形に延長するわけではない。
本稿で全 Artin 延長に使う良い根は、根対一致時には
$L_0=D_0$ を選ぶという条件を含む。
\end{remark}

\subsection{Hochschild 二次に自明な根の étale 代数化}
\label{subsec:hh2-etale-morita-root}

次の命題は noncubic 点に限らない。
根の存在する空間から étale な切片を取ることにより、
Hochschild 一次への作用を仮定せずに根を代数化する。

\begin{proposition}[Hochschild 二次に自明な根の étale 代数化]
\label{prop:hh2-etale-morita-root}
$S$ を標数零の代数閉体 $k$ 上の局所有限型 scheme とし、
$\Lambda$ を $S$ 上の有限局所自由代数とする。
次を仮定する。
\begin{enumerate}
\item $\operatorname{Perf}(\Lambda)$ は相対 smooth proper であり、
任意の基底変更後も
$\operatorname{HH}^{<0}(\Lambda)=0$、
$\operatorname{HH}^0(\Lambda)=\OO_S$ である。
\item 相対 inverse Serre $t=S_\Lambda^{-1}[2]$ は
可逆 perfect 双加群で表され、任意の導来基底変更と可換する。
\item $x\in S(k)$ と
$p_0\in\operatorname{DPic}_k(\Lambda_x)$ が与えられ、
$p_0^2=t_x$、
$(p_0)_*|_{\operatorname{HH}^2(\Lambda_x)}=1$ である。
\end{enumerate}
$\mathcal G\to S$ を Morita 自己同値の étale 局所同型類の代数空間、
$\tau=[t]$ とし、
\[
 \mathcal Z=
 \mathcal G\times_{\mathrm{sq},\mathcal G,\tau}S
\]
を平方根空間とする。
このとき $\mathcal Z\to S$ は $p_0$ で smooth である。
さらに有限型 affine な étale 近傍 $(U,u)\to(S,x)$ 上に、
可逆 perfect $\Lambda_U$-双加群 $P$ と実際の同型
\[
 P\otimes_{\Lambda_U}^{\mathbf L}P\simeq t_U
 \tag{N-root}\label{eq:hh2-etale-actual-root}
\]
が存在し、$[P_u]=p_0$ となる。
追加で $(p_0)_*|_{\operatorname{HH}^1(\Lambda_x)}=1$ なら、
$\mathcal Z\to S$ 自体が $p_0$ で étale であり、
同じ閉根を持つ Artin 持上げの Morita 同型類は一意である。
\end{proposition}

\begin{proof}
全証明は相対モジュライ論文 \cite[\texttt{noncubic_hh_roots}]{QuadricModuli2026} に置く。
\end{proof}

\begin{remark}[Hochschild 一次の役割]
\label{rem:hh2-etale-root-tangent}
同じ閉根を持つ根の一次的な差は
$a\in\operatorname{HH}^1(\Lambda_x)$ で
$a+(p_0)_*a=0$ を満たすものである。
従って平方根空間の相対接空間は
$\ker(1+(p_0)_*)$ である。
Hochschild 一次への恒等作用はこの相対接空間を消し、
根の空間自体を étale にするために使う。
étale 近傍上で一つの根を選ぶ存在定理には不要である。
\end{remark}

\begin{corollary}[全 noncubic 点の有限型 étale Morita 根]
\label{cor:noncubic-etale-morita-root}
$S$ を有限 foundation スタックの有限型 affine chart とし、
$x\in S(k)$ を任意の noncubic $\mathcal H_2$ 対象の
normal foundation に対応する点とする。
定理~\ref{thm:noncubic-hh-trivial-root} の良い根 $s_0$ を選ぶ。
このとき $x$ の有限型 affine étale 近傍 $(U,u)\to(S,x)$ 上に、
$s_0$ を閉 fiber とする可逆 perfect 双加群 $P$ と
実際の平方同型 \eqref{eq:hh2-etale-actual-root} が存在する。
平方根空間自体もこの点で étale である。

四つの普遍 projective 対象 $F_0,\ldots,F_3$ を、
閉 fiber で
$(\OO,\pi M_c,\OO(1),\pi M_d(1))$ に対応させる。
近傍をさらに縮めれば、$P$ による自己同値 $\mathsf s$ は
\[
 \mathsf s F_0\simeq F_2,\qquad
 \mathsf s F_1\simeq F_3
 \tag{N-window}\label{eq:noncubic-etale-root-window}
\]
を満たす。これらの同型は任意の導来基底変更に保たれる。
\end{corollary}

\begin{proof}
全証明は相対モジュライ論文 \cite[\texttt{noncubic_hh_roots}]{QuadricModuli2026} に置く。
\end{proof}

\begin{remark}[有限型近傍でまだ結論していないこと]
\label{rem:noncubic-etale-root-scope}
上の $U$ は有限代数と perfect Morita 根の実際の有限型族である。
この構成だけから相対 Abel heart、
全正軌道 Hom の同じ開近傍での次数零集中、
全軌道環の平坦性・Noetherian 性・AZ 豊富性、
同じ根に対する相対 AS4 被覆は従わない。
従って $U\to\mathfrak H_2$ や
noncubic 点を含む smooth Artin 開近傍を、ここでは構成したとは主張しない。
明示的な noncubic 二形式表示族の効果性と、
cubic を横断するこの étale 近傍の効果性も区別する。
\end{remark}

\subsection{根の有限性、計算の検証と残る問題}

根対が異なる場合、命題~\ref{prop:first-orbit-object} により任意の適合根の
$\OO$ 上の像は同じ内在的対象である。
従って二つの根の差は pointed 自己同値を与え、全根の同型類は
有限な pointed 自己同値群の torsor の部分集合に入る。
一般点でのこの群は既存原稿では $D_8\times C_2$（位数十六）であり、
二つの平方根変数の符号群 $\mu_2^2$ だけに限定してはいけない。
本節は全ての根の作用を恒等と主張せず、上記の trace 零の選択による
一つの根を特定して証明した。
他の根に対する作用は、それとの差である pointed 自己同値の作用との合成である。

補足の厳密有理計算
\texttt{computations/verify\_noncubic\_gale\_root.py} は、非交差、一根共有、一致の各配置で
$\mathscr L$、$J$、Grassmann tangent の十六次元行列、
基底変更軌道の全十六生成元を構成する。
浮動小数点や乱数を使わず、
\[
\begin{array}{c|ccc}
 &\text{非交差}&\text{一根共有}&\text{一致}\\\hline
\dim(\text{orbit}^+,\text{orbit}^-)&(7,6)&(7,6)&(6,6)\\
\dim(\operatorname{HH}^2_+,\operatorname{HH}^2_-)&(3,0)&(3,0)&(4,0)
\end{array}
\]
を確認する。一般の配置への証明はこの三例の計算から外挿せず、
定理~\ref{thm:noncubic-hh-trivial-root} の同変接複体と安定化群の議論による。

\paragraph{同じ根の AS4 被覆は別の問題である.}
持ち上げた同じ根の軌道環の $B_1$ は四次元の $D$ であり、
$B_2=\Hom(E_0,E_4)$ と積は同じ有限 foundation の bar 複体から計算できる。
根を適合するように twist する前の中央二重平方環では、二次双対に正規な二次元の
pencil が存在するが、固定した適合根の $\mathcal B^\sigma$ でも
同じ pencil が残るとは限らない。ambient の関係線の $\sigma$ 不変性が必要である。
非交差・一根共有・一致の三つの有理配置では、三次平坦性を保つ
関係空間の接次元 $18,18,19$ の全方向に対し、選んだ projective 正規線と
その normalizer が一次に一意に持ち上がることを
\texttt{computations/verify\_noncubic\_cover\_tangent.py} の有限線型代数で確認した。
これはこの特定の被覆の一次計算であり、それ自体から高次存在は従わない。
次節では別の中央 AS4 被覆を全配置で構成し、
系~\ref{cor:noncubic-all-artin-h2} で同じ良い根に対する全 Artin 延長を証明する。
有限型近傍の一様な効果性はなお別問題である。

% 2026-11-21 / v1.40
% !TEX root = ../nc-quadric-classification.tex
\section{noncubic の良い根に対する中央 AS4 被覆}
\label{sec:noncubic-central-hull}

\paragraph{範囲.}
以下の構成は、中央二重平方環を与える従来の Ore 被覆とは異なる被覆である。
同じ trace 零交換対合で選んだ Serre 根の軌道環について、切断元自体が中央となる
AS4 被覆を与える。既存の Ore ambient の中で切断元を中央化する主張ではない。
新規性は未確認である。

基礎体を標数零の代数閉体とし、\(a,b\in k\)、
\(\theta=1-ab\ne0\)、\(\delta^2=\theta\) と置く。
二つの非退化二元二次形式は座標・定数倍を選べば
\[
 f_1=uv,\qquad f_2=(u-av)(bu-v)
       =bu^2-(1+ab)uv+av^2
\]
と書ける。二根対の一致も \(a=b=0\) に含まれる。
\[
 B_0=k[u,v]\langle c,d\rangle/(c^2-f_1,d^2-f_2)
\]
で \(u,v\) は中央とする。良い交換自己同型は
\[
 \sigma(u)=(u-av)/\delta,\quad
 \sigma(v)=(bu-v)/\delta,\quad
 \sigma(c)=d/\delta,\quad \sigma(d)=-\delta c
\]
であり、\(\sigma^2=\iota\)、\(\iota(u,v,c,d)=(u,v,-c,-d)\) である。
以下では \(x*y=x\sigma^{\deg x}(y)\) の Zhang 規約を使う。

\begin{theorem}\label{thm:noncubic-central-hull}
次の六関係環を定める。
\[
 A_0=T(u,v,c,d)\Big/
 \left(
 \begin{gathered}
 [u,c],[u,d],[v,c],[v,d],\quad c^2-\frac{uv+vu}{2},\\
 d^2-bu^2+\frac{1+ab}{2}(uv+vu)-av^2
 \end{gathered}
 \right).
\]
\(A_0\) は左右 Noetherian な quadratic Koszul AS 正則環で、型は
\((1,4,6,4,1)\) である。
\(g_0=uv-vu\) は左右正則正規で、
\[
 A_0/(g_0)=B_0,\qquad
 g_0u=-ug_0,\quad g_0v=-vg_0,\quad
 g_0c=cg_0,\quad g_0d=dg_0.
\]
\(\sigma\) は \(A_0\) に延長し、\(\sigma(g_0)=-g_0\) となる。
従って \(A=A_0^\sigma\) における \(g=g_0\) は中央正則二次元であり、
\[
 B_0^\sigma=A/(g).
\]
特に、指定した良い Serre 根の全軌道環に対する中央 AS4 被覆が存在する。
さらに \(E=(B_0^\sigma)^!\) の二次中心は一次元であり、
この被覆の双対線がその全体である。
\end{theorem}

\begin{proof}
全証明は相対モジュライ論文 \cite[\texttt{noncubic_central_hull}]{QuadricModuli2026} に置く。
\end{proof}

\begin{lemma}[良い noncubic 根では二次反中心が消える]
\label{lem:noncubic-no-quadratic-anticenter}
定理~\ref{thm:noncubic-central-hull} の \(E=(B_0^\sigma)^!\) について
\[
 Z(E)_2=k\chi,\qquad
 Z^-(E)_2:=\{z\in E_2\mid zx=-xz\ (x\in E_1)\}=0
\]
である。最初の式の \(\chi\) は同定 \(E=E_0^\alpha\) による像を表す。
\end{lemma}
\begin{proof}
全証明は相対モジュライ論文 \cite[\texttt{noncubic_central_hull}]{QuadricModuli2026} に置く。
\end{proof}


\begin{remark}[相対化と横断族への適用で分ける入力]
明示 noncubic 表示族上では、\(2\theta\) が単元となる任意の可換 Noetherian 基底に
上の構成を相対化できる。双対の各次数の projectivity は既存の相対 Koszul 補題による。
\(\chi\) 乗法は各幾何 fiber で単射であるから局所直和因子への単射となり、商も
任意の基底変更と両立する有限 projective 加群となる。
上の pairing の行列式は単元なので、相対 Frobenius 条件も成立する。
商に PP の相対 Koszul 定理を適用することで双対 ambient の全次数平坦性を得る。
正規元商からの Noetherian 持上げは任意の Noetherian 基底変更後にも適用できる。

横断族では、この計算だけから全軌道の平坦性や canonical heart は従わない。
ただし、被約有限型基底 \(S\) 上にその良い根の平坦な Koszul 軌道環 \(B\) があり、
\(E=B^!\) の二次・三次が階数七・八の束であると仮定する。
稠密な cubic 開部分上で同じ根の canonical hull が中央であるなら、
commutator 写像 \(E_2\to\mathcal Hom(E_1,E_3)\) の階数は稠密開で高々六である。
被約性により七次 minors は全体で零であり、上の定理は noncubic 点で階数六を与える。
従ってその点の近傍で核は任意の基底変更と両立する可逆部分束となる。
これが延ばすべき中央線であり、従来の Ore 被覆の正規線ではない。
局所環上では、次節の定理~\ref{thm:noncubic-local-central-hull} がこの入力を確認する。
全軌道平坦性、同じ根と cubic canonical 周期との比較、有限 AS4 条件を満たす
単一近傍の構成、Noetherian heart の認識は、この線の延長と別に記載する必要がある。
\end{remark}

% !TEX root = ../nc-quadric-classification.tex
% 2026-11-21 / v1.40
\section{五頂点窓と noncubic 変形の全 Artin 認識}
\label{sec:noncubic-local-tails}

\paragraph{局所環で成立することと、一様な開で残ること.}
本節では系~\ref{cor:noncubic-etale-morita-root} の有限型近傍の
\emph{指定された noncubic 点における局所環}を取る。
その同じ局所環上で、全ての前向き Hom、全軌道環、canonical
Grothendieck 導来圏と Noetherian heart を構成する。
極小解消の Betti 数を族の中で固定せず、五頂点の非極小 bar 解消を使う。
さらに二次双対の中心・反中心を使って同じ根の中央 AS4 被覆を局所環上に構成し、
全 noncubic 点で任意の Artin 変形を相対 $\mathfrak H_2$ として認識する。
有限型の一つの開近傍上での全次数一様性と全幾何 fiber の $\HH$ 所属は別問題であり、
本節だけから noncubic 点を含む $\mathfrak H_2$ の開 chart は得られない。

\subsection{極小 Betti 数に依存しない bar 解消}

以下で $C$ は正の向きの $R$-線形 $\mathbb Z$ 代数、
$P_i=e_iC$、$S_i$ は頂点 $i$ に $R$ を置く加群とする。
積は $C_{ij}\otimes C_{jk}\to C_{ik}$ と書く。

\begin{lemma}[五頂点 bar 解消の持上げ]
\label{lem:five-window-bar-lifting}
次のいずれかを仮定する。
\begin{enumerate}
\item $R$ は可換局所環であり、全 $C_{ij}$ は有限自由である。
  以下の fiber 条件を剰余体上で仮定する。
\item $R$ は任意の可換環であり、全 $C_{ij}$ は有限 projective である。
  以下の fiber 条件を全ての幾何点上で仮定する。
\end{enumerate}
fiber 条件は、各 $S_i$ が長さ三の有限生成 projective 解消を持ち、
その項に現れる表現可能加群が $P_i,\ldots,P_{i+4}$ に限られ、
最後の項が $P_{i+4}$ 一個であり、さらに
\[
 \operatorname{Ext}_{C_{\bar x}}^q(S_i,P_j)=0
 \quad\text{unless }(q,j)=(3,i+4),\qquad
 \operatorname{Ext}_{C_{\bar x}}^3(S_i,P_{i+4})\simeq\kappa(\bar x)
 \tag{BW-fiber}\label{eq:five-window-fiber-dual}
\]
となることである。最後の項以外の極小 Betti 数は一定と仮定しない。

このとき $S_i$ は $C$ 上 perfect で、projective dimension は高々三である。
解消と AS 双対条件は任意の $R$-基底変更後にも成立し、
\[
 \mathbf R\operatorname{Hom}_{C}(S_i,P_j)\simeq
 \begin{cases}\mathcal L_i[-3],&j=i+4,\\0,&j\ne i+4\end{cases}
 \tag{BW-dual}\label{eq:five-window-relative-dual}
\]
である。ここで $\mathcal L_i$ は可逆 $R$-加群である。
\end{lemma}

\begin{proof}
全証明は相対モジュライ論文 \cite[\texttt{noncubic_local_tails}]{QuadricModuli2026} に置く。
\end{proof}

\begin{remark}[二次曲面の固定階数の五頂点複体]
\label{rem:five-window-fixed-ranks}
$\operatorname{rank} C_{i,i+r}=1,2,4,6,9$（$r=0,1,2,3,4$）の場合、
\eqref{eq:five-window-bar} は局所的に
\begin{align*}
0\longrightarrow P_{i+4}^{16}&\longrightarrow
 P_{i+3}^{8}\oplus P_{i+4}^{48}
 \longrightarrow P_{i+2}^{4}\oplus P_{i+3}^{16}\oplus P_{i+4}^{40}\\
 &\longrightarrow P_{i+1}^{2}\oplus P_{i+2}^{4}
       \oplus P_{i+3}^{6}\oplus P_{i+4}^{9}
 \longrightarrow P_i\longrightarrow S_i\longrightarrow0
\end{align*}
となる。最終添字 $i+4$ の係数複体の階数は、homological degree
$1,2,3,4$ の順に $9,40,48,16$ である。
極小形が cubic 型から noncubic 型へ変わっても、この bar 複体の階数は変わらない。
本節の方法が要求するのはこの固定階数の複体の完全性である。
\end{remark}

\begin{proposition}[五頂点解消による canonical tails]
\label{prop:five-window-canonical-tails}
前補題の仮定に加え、$C$ が距離高々二の成分で生成されるとする。
$\mathcal O_i=\pi P_i\in\operatorname{QGr}C$ と置くと、
\[
 \operatorname{Ext}^q(\mathcal O_i,\mathcal O_j)\simeq
 \begin{cases}
 C_{ji},&q=0,\\
 \operatorname{Hom}_R(C_{i,j-4},\mathcal L_{j-4}),&q=2,\ j\ge i+4,\\
 0,&\text{otherwise}
 \end{cases}
 \tag{BW-coh}\label{eq:five-window-tails-cohomology}
\]
である。$\mathbf R\omega$ の cohomological dimension は高々二であり、
任意の連続四項の直和 $T_a=\bigoplus_{i=a}^{a+3}\mathcal O_i$ は
$D(\operatorname{QGr}C)$ の compact tilting generator となる。
従って
\[
 D(\operatorname{QGr}C)\simeq D(\operatorname{Mod}\operatorname{End}(T_a)),
 \qquad D(\operatorname{QGr}C)^c\simeq\operatorname{Perf}\operatorname{End}(T_a).
 \tag{BW-Morita}\label{eq:five-window-tails-morita}
\]
全構成は任意の導来基底変更と整合する。
Noetherian 性はこの命題の仮定にも結論にも含めない。
\end{proposition}

\begin{proof}
全証明は相対モジュライ論文 \cite[\texttt{noncubic_local_tails}]{QuadricModuli2026} に置く。
\end{proof}

\subsection{良い noncubic 根の局所環上での全軌道と heart}

\begin{theorem}[noncubic 局所環における全次数 tails 認識]
\label{thm:noncubic-local-tails-effectivity}
系~\ref{cor:noncubic-etale-morita-root} の $(U,u)$ と Morita 根 $\mathsf s$ を取り、
$R=\mathcal O_{U,u}$ とする。
より一般に、閉点を $u$ へ送る Noetherian 局所 $\mathcal O_{U,u}$-代数を $R$ としてもよい。
$\Lambda=\Lambda_U\otimes R$ とし、
\[
 E_{2a}=\mathsf s^aF_0,\quad E_{2a+1}=\mathsf s^aF_1,
 \quad C_{ij}=\operatorname{Hom}_{\Lambda}(E_{-j},E_{-i})\quad(i\le j)
\]
と置く。$C_{ij}=0$（$i>j$）とし、積は射の合成とする。
次が成立する。
\begin{enumerate}
\item 全 $i\le j$ で $\mathbf R\operatorname{Hom}(E_i,E_j)$ は有限自由 $R$-加群の
  次数零成分に集中する。階数は閉 fiber の normal helix と等しく、
  偶数差 $2n$ で $(n+1)^2$、奇数差 $2n+1$ で $(n+1)(n+2)$ である。
  この同定と積は任意の導来基底変更と整合する。
\item 全軌道環
  $B=\bigoplus_{n\ge0}\operatorname{Hom}(F_0,\mathsf s^nF_0)$ は
  一次生成・二次・相対 Koszul で、全 $B_n$ は階数 $(n+1)^2$ の自由加群である。
  $B$ は左右 Noetherian であり、任意の有限型可換 $R$-代数への変更後も同様である。
\item 任意の可換 $R$-代数 $R'$ について canonical な比較
  \[
  D(\operatorname{QGr}B_{R'})\simeq D(\operatorname{Mod}\Lambda_{R'}),
  \qquad D(\operatorname{QGr}B_{R'})^c\simeq\operatorname{Perf}\Lambda_{R'}
  \]
  があり、$\pi B_{R'}$ を $F_0$ に、次数移動 $(1)$ を $\mathsf s$ に送る。
  特に compact 圏の canonical Serre 関手について $(1)^2\simeq S^{-1}[2]$ である。
\item 有限型可換 $R$-代数 $R'$ 上で $\mathcal C_{R'}=\operatorname{qgr}B_{R'}$ は
  Noetherian な相対 Abel 圏の系をなし、module-object 基底変更と fppf descent を持つ。
  coherent 対象間の Hom は有限で平坦基底変更と可換する。
  $\pi B_{R'}$ は平坦かつ exceptional で、$(\pi B_{R'},(1))$ は AZ ample である。
  左右の全 $\chi_q$ が成立し、canonical compact 圏は $\operatorname{Perf}\Lambda_{R'}$ である。
\end{enumerate}
この定理だけには同じ根の AS4 被覆を含めない。標準局所基底では、
後の定理~\ref{thm:noncubic-local-central-hull} がその被覆を与える。
また $R$ は一般に元の局所有限型 $k$-site の対象ではないため、
以上をそのまま $U$ の開近傍上の $\mathfrak H_2$ 族と呼ばない。
\end{theorem}

\begin{proof}
全証明は相対モジュライ論文 \cite[\texttt{noncubic_local_tails}]{QuadricModuli2026} に置く。
\end{proof}

\begin{remark}[局所定理と有限型一様化の役割]
有限型近傍上の各軌道 Hom は perfect であり、固定した有限個については開性を使える。
本定理は、その無限個の開条件の共通部分が開であるとは述べない。
AS4 被覆の延長が得られても、有限型近傍で使う全軌道の一様平坦性・表示比較と
全幾何 fiber の rank-one torsion-free 条件は別途確認する。
一方、同じ局所環上では全次数の bar 解消と Noetherian heart を既に得たので、
以後の局所 AS4 構成で canonical compact 圏を独立の未証明入力に戻す必要はない。
有限型一様化は定理~\ref{thm:noncubic-finite-type-effectivity} で、
全幾何 fiber の有限 foundation 認識を新しい入力として証明する。
\end{remark}


\subsection{中央被覆の局所延長と全 Artin 認識}

\begin{lemma}[cubic 根の canonical hull の scalar normalizer]
\label{lem:cubic-compatible-scalar-normalizer}
体 \(K\) 上の geometric cubic \(\mathbb Z\)-代数に実際の二周期 \(\beta\) があり、
対応する tails 自己同値 \(s\) が \(s^2\simeq S^{-1}[2]\) を満たすとする。
その偶数環 \(B\) の canonical hull の双対線 \(K\phi\subset B^!_2\) は、
ある \(\lambda\in\{1,-1\}\) に対して
\[
 \phi x=\lambda x\phi\qquad(x\in B^!_1)
\]
を満たす。ここでは最初から \(\lambda=1\) とは主張しない。
\end{lemma}
\begin{proof}
全証明は相対モジュライ論文 \cite[\texttt{noncubic_local_tails}]{QuadricModuli2026} に置く。
\end{proof}

\begin{theorem}[局所基底上の noncubic 横断中央被覆]
\label{thm:noncubic-local-central-hull}
\((R,\mathfrak m,k)\) を標数零の体上の Noetherian 正則局所環、\(K=\operatorname{Frac}R\)
とする。二次環 \(B=T_R(V)/(J)\) について次を仮定する。
\begin{enumerate}
\item \(V,J\) は階数四、七の自由加群で、\(J\subset V^{\otimes2}\) は直和因子である。
全ての \(B_n\) は階数 \((n+1)^2\) の有限自由加群であり、\(B\) は相対 Koszul である。
特に \(E=B^!\) の全次数は有限自由、任意の底変更と可換し、その階数は
\(1,4,7,8,8,\ldots\) である。
\item 閉 fiber \(B_k\) は定理~\ref{thm:noncubic-central-hull} の良い根の
noncubic 軌道環である。
\item generic fiber \(B_K\) は geometric cubic helix を二項移す実際の根の偶数環であり、
その根の平方は反 Serre 関手に同型である。
\end{enumerate}
このとき \(B\) は相対 quadratic Koszul AS4 環 \(A\) と中央正則二次元 \(g\) により
\[
 B\cong A/(g)
\]
と表示される。全次数の有限自由性、左右および双加群 Koszul 解消、中央正則切断は
任意の基底変更と両立する。\(B\) が左右 Noetherian なら \(A\) もそうである。
この定理は有限型開近傍上の一様な Noetherian 性をそれ自体では主張しない。
\end{theorem}

\begin{proof}
全証明は相対モジュライ論文 \cite[\texttt{noncubic_local_tails}]{QuadricModuli2026} に置く。
\end{proof}

\begin{remark}[Artin 結論と有限型一様性]
この定理の入力を有限 foundation の局所環上で確認できれば、
任意の Artin 厚みへの引戻しは、同じ良い根の AS4 被覆を与える。
しかし局所環上の全次数平坦性・Noetherian 性の証明だけから、
単一の有限型開近傍上で全次数の認識が成立すると推論しない。
有限型近傍の全 orbit、canonical tails と普遍 Noetherian 性はそれぞれ別の入力である。
\end{remark}

\begin{corollary}[全 noncubic 点の Artin 認識]
\label{cor:noncubic-all-artin-h2}
標数零の代数閉体上の任意の noncubic \(\mathcal H_2\) 対象
\((\mathcal C,\mathcal O)\) を取る。
その任意の flat pointed Abel Artin 変形は、指定した良い根の持上げの全軌道環に対する
中央正則二次 AS4 被覆を持ち、全有限型基底変更を備えた相対 \(\mathfrak H_2\) 族へ昇格する。
従って、局所 Artin 基底上で相対 \(\mathfrak H_2\) 変形から pointed Abel 変形への忘却は
変形 groupoid の同値である。
根対が異なる点では接次元三・無限小慣性零、根対が一致する点では接次元四・無限小慣性一次元であり、
いずれも変形は無障害である。
\end{corollary}
\begin{proof}
全証明は相対モジュライ論文 \cite[\texttt{noncubic_local_tails}]{QuadricModuli2026} に置く。
\end{proof}

\begin{remark}[形式的完了と代数性の区別]
全 cubic 点における既存の形式認識と前系を合わせると、全 \(\mathcal H_2\) 幾何点で
\(\mathfrak H_2\) と pointed Abel 変形の Artin groupoid が一致する。
この結論から内在的スタック全体の有限型代数性はまだ従わない。
局所環上で同時に成り立つ全軌道・Noetherian 条件を、単一の有限型 étale 近傍へ
一様に広げる段階と、その相対認識を別に閉じる必要がある。
\end{remark}

\subsection{中央二次元の全次数正則性の有限型一様化}

\begin{lemma}[有限な二次商から中央正則線を一様化する]
\label{lem:spread-central-regular-line}
$R$ を可換 Noetherian 環、$V$ を有限 $R$-加群とし、
$J\subset V^{\otimes2}$ とし、$E=T_R(V)/(J)$ を次数一生成の二次 $R$-代数とする。
特に $E_{<0}=0$、$E_0=R$、$E_1=V$ は有限生成である。
$\phi\in E_2$ は中央であると仮定し、$P=E/(\phi)$ と置く。
$P_0,\ldots,P_4$ は有限 projective $R$-加群、$P_5=0$ とする。
さらに、ある $\mathfrak p\in\operatorname{Spec}R$ において
$\phi$ の $E\otimes_RR_{\mathfrak p}$ 上の乗法が単射であると仮定する。
これは剰余体上だけの正則性ではなく、局所環上の正則性の仮定である。

このとき $f\in R\setminus\mathfrak p$ が存在し、単一の開 $D(f)$ 上で
\begin{enumerate}
\item $\phi$ は全次数で左右正則である。
\item 全ての $(E_f)_n$ は有限 projective $R_f$-加群であり、
  \[
  0\longrightarrow(E_f)_{n-2}\xrightarrow{\phi\cdot}(E_f)_n
       \longrightarrow(P_f)_n\longrightarrow0
  \tag{SC-exact}\label{eq:spread-central-regular-sequence}
  \]
  は $R_f$ 上分裂完全である。負次数は零とする。
\item 任意の可換 $R_f$-代数への変更後にも、中央正則性、
  上の全次数の完全列、商 $E_f/(\phi)=P_f$ が保たれる。
\end{enumerate}
さらに $\deg t=2$、$t$ の作用を $\phi$ とすれば、非標準的な次数付き同型
\[
 E_f\simeq\bigoplus_{i=0}^4(P_i)_f\otimes_{R_f}R_f[t](-i)
 \tag{SC-module}\label{eq:spread-central-projective-module}
\]
がある。従って $E_f$ は $R_f[t]$ 上有限 projective でもある。
各 $P_i$ が階数 $1,4,6,4,1$ を持つ場合、$E_n$ の階数は
$1,4,7,8,8,\ldots$ となる。
この補題は $P$ または $E$ の Koszul 性を結論しない。
したがって二次双対 $P^!$ の AS 正則性や全次数平坦性も、この補題だけからは従わない。
\end{lemma}

\begin{proof}
全証明は相対モジュライ論文 \cite[\texttt{noncubic_local_tails}]{QuadricModuli2026} に置く。
\end{proof}

\begin{remark}[中央正則性と Koszul 性の入力]
この補題の有限性は、中央二次元による商が有限個の次数にしか台を持たないことから来る。
したがって $E_n$ ごとに別々の開近傍を選ぶ必要がない。
有限型近傍上で $P_0,\ldots,P_5$ の指定階数と中央性を確立できれば、
局所環で得た $\phi$ の正則性を、全次数について同じ開へ延ばせる。
一方、指定点での $P$ の Koszul 性をその開全体へ延ばすことは別である。
既存の局所 Koszul 変形原理を用いるには、適用先の各局所環の剰余体上の
Koszul 性、またはそれを与える一様な伝播補題がさらに必要である。
第\ref{sec:noncubic-finite-type}節では、全 fiber の有限認識からこの入力を供給し、
補題~\ref{lem:spread-local-central-hull} で同じ根の AS4 被覆まで延長する。
\end{remark}

% !TeX root = ../nc-quadric-classification.tex
% 2026-11-22 / v1.43（証明をBへ移管）
\section{noncubic 境界を横断する有限型効果性}
\label{sec:noncubic-finite-type}

局所環上の構成を一つの有限型近傍へ広げる。
有限 foundation の Lagrangian 表示に全幾何点を認識する有限開条件を課し、
その同じ基底の全 fiber の geometric helix を用いて全 Hom を先に平坦化する。
二つの cubic 開部分と非被約 noncubic 閉部分から Noetherian 性を貼り合わせ、
局所中央被覆を有限型開へ延長する。
一般の Koszul 性の開性や、無限個の開集合の交叉の開性は用いない。
以下の相対結果とその組合せの新規性は未確認である。

\subsection{Lagrangian 開部分の全幾何点の認識}

以下では標数零の代数閉体上で考える。
前節の Grassmann 表示で、第二 summand の $D_L$ による基底変更を行い、
$H=U\oplus U^*$、$U=k^2\otimes k^2\simeq\operatorname{Mat}_2(k)$ と書く。
交代形式は二 summand 間の標準評価から作るものとする。
この基底変更は二次元因子上の同型なので分解テンソルの集合を保つ。
$\Sigma_U\subset\mathbb P(U)$ と $\Sigma_{U^*}\subset\mathbb P(U^*)$ を
二つの Segre 二次曲面とする。

\begin{lemma}[両階数三の Lagrangian 平面の直接正規形]
\label{lem:lagrangian-double-rank-three-recognition}
$L\subset U\oplus U^*$ を Lagrangian 四平面とする。
二つの射影の階数がともに三で、$L\cap U$ と $L\cap U^*$ の各生成元が
二次正方行列として可逆であると仮定する。
このとき $\operatorname{GL}_2\times\operatorname{GL}_2$ の $U$ 上の tensor 作用と
$U^*$ 上の反傾作用、および $L$ の基底変更により、$L$ は
\begin{equation}\label{eq:lagrangian-rank-three-normal-form}
\begin{aligned}
 L=\bigl\langle&(-E_{21},\alpha E_{12}+\gamma E_{21}),\quad
 (E_{12},\beta E_{12}-\alpha E_{21}),\\
 &(-D_0,0),\quad (0,I_2)\bigr\rangle,
 \qquad \alpha^2+\beta\gamma=1.
\end{aligned}
\end{equation}
となる。ここで $D_0=\operatorname{diag}(1,-1)$ である。
対応する有限 foundation は、中央二重平方環
\[
 k[u,v]\langle c,d\rangle/
 \bigl(c^2-uv,\ d^2-(\alpha u+\beta v)(\gamma u-\alpha v)\bigr)
\]
の既知の normal foundation と同型である。
この認識は、候補 foundation 自身の domain 性や $\mathcal H_2$ 所属を仮定しない。
\end{lemma}
\begin{proof}
全証明は相対モジュライ論文 \cite[noncubic有限型効果性]{QuadricModuli2026} に置く。
\end{proof}

\begin{proposition}[非被約閉基底上の両階数三正規形]
\label{prop:lagrangian-relative-double-rank-three}
$R$ を標数零の体上の可換環、$L\subset U_R\oplus U_R^*$ を
Lagrangian な階数四の部分束とする。
二射影の三次小行列式がそれぞれ単位イデアルを生成し、その四次行列式が
$R$ の元としてともに零であると仮定する。
さらに二つの交点線束の局所生成元は、二次正方行列として可逆とする。
このとき $\operatorname{Spec}R$ の \texorpdfstring{\'etale}{etale} 被覆上で、
前補題の式~\eqref{eq:lagrangian-rank-three-normal-form} が成立する。
従ってその有限 foundation は同じ基底上の明示中央二重平方族の
有限 foundation と同型である。基底の被約性は不要である。
\end{proposition}
\begin{proof}
全証明は相対モジュライ論文 \cite[noncubic有限型効果性]{QuadricModuli2026} に置く。
\end{proof}

\paragraph{Cross reshuffling.}
$M\in\operatorname{Hom}(U,U^*)$ を四因子テンソルと見て、その係数を
$M_{ck,ab}$ と書く。cross reshuffling を
\[
 \operatorname{Cr}(M)_{ka,bc}=M_{ck,ab}
\]
で定める。これは中間の隣接二因子を入力とする縮約を表す四次正方行列である。
$M=y\otimes\nu$ が階数一で、$y$ と $\nu$ がそれぞれ二次正方行列として可逆なら、
$\operatorname{Cr}(M)$ は行・列の並替を除いて二つの可逆二次行列の Kronecker 積であり、可逆である。

\begin{proposition}[noncubic 点を含む全幾何点認識の開部分]
\label{prop:lagrangian-geometric-recognition-open}
$L_0$ を前補題の条件を満たす点とし、有限 foundation の Grassmann chart 内にあるとする。
$\operatorname{LG}(4,H)$ における $L_0$ の開近傍 $W$ を取れて、その全幾何点の
有限 foundation は $\mathcal H_2$ 対象の normal foundation となる。
具体的に、$L$ の局所基底を列行列
\[
 \begin{pmatrix}M\\N\end{pmatrix},\qquad M,N\in\operatorname{Mat}_4
\]
で表し、次の有限開条件を課せばよい。
\begin{enumerate}
\item $\mathbb P(L\cap U)\cap\Sigma_U=\varnothing$、
      $\mathbb P(L\cap U^*)\cap\Sigma_{U^*}=\varnothing$。
\item $\operatorname{Cr}(N\operatorname{adj}M)$ と
      $\operatorname{Cr}(M\operatorname{adj}N)$ が可逆である。
\end{enumerate}
二つの射影がともに階数三なら前補題の noncubic 型であり、
少なくとも一方の階数が四なら一方の ruling foundation が geometric cubic 型である。
この命題は幾何点での有限 foundation の認識であり、相対 heart の存在を結論に先取りしない。
\end{proposition}
\begin{proof}
全証明は相対モジュライ論文 \cite[noncubic有限型効果性]{QuadricModuli2026} に置く。
\end{proof}

\begin{corollary}[cubic 二開部分と noncubic 閉部分]
\label{cor:lagrangian-open-closed-recognition}
前命題の $W$ において、二つの行列式を零とする閉部分スキームを
$D=V(\det M,\det N)$ と置く。
$D$ は非被約でもよい。$D$ 上では命題~\ref{prop:lagrangian-relative-double-rank-three} の正規形が
\texorpdfstring{\'etale}{etale} 局所に成立し、対応する有限 foundation は
相対中央二重平方族のものに一致する。
開補集合は $D(\det M)$ と $D(\det N)$ で覆われ、
前者は元の ruling の geometric cubic 表示、後者は交換した ruling の
geometric cubic 表示を持つ。
\end{corollary}
\begin{proof}
全証明は相対モジュライ論文 \cite[noncubic有限型効果性]{QuadricModuli2026} に置く。
\end{proof}

\subsection{全幾何 fiber の有限 foundation 認識からの一様化}

\begin{theorem}[normal foundation の全 fiber 認識による全軌道の一様化]
\label{thm:uniform-orbit-from-normal-fibers}
$R$ を本稿の基礎体 $k$ 上の可換 Noetherian 代数、$\Lambda$ を有限 projective $R$-代数とし、
$\operatorname{Perf}\Lambda$ は相対 smooth proper とする。
$F_0,F_1,F_2,F_3$ を相対 full strong exceptional foundation とし、
$\operatorname{End}(\bigoplus_{a=0}^3F_a)\simeq\Lambda$ の同定を固定する。
実際の可逆 perfect Morita 双加群による自己同値 $\mathsf s$ と同型
\[
 \mathsf s^2\simeq S^{-1}[2],\qquad
 \mathsf sF_0\simeq F_2,\qquad \mathsf sF_1\simeq F_3
\]
を固定する。ここで $S$ は canonical な相対 Serre 関手である。
各幾何点 $x$ において、$\operatorname{Perf}\Lambda_x$ とある体上の
$\HH$ 対象 $(\mathcal C_x,\mathcal O_x)$ の導来圏との同値があり、
$(F_0,F_1,F_2,F_3)_x$ を定理~\ref{thm:normal-helix} の
normal helix の最初の四項に送ると仮定する。
この仮定は各 fiber の有限 foundation に関するものであり、
$R$ 上の heart や相対正規形を先に仮定しない。

\[
 E_{2a}=\mathsf s^aF_0,\qquad E_{2a+1}=\mathsf s^aF_1,
 \qquad C_{ij}=\operatorname{Hom}_{\Lambda}(E_{-j},E_{-i})\quad(i\le j)
\]
と置き、$i>j$ では $C_{ij}=0$ とする。この同じ $R$ 上で次が成立する。
\begin{enumerate}
\item 全 $i\le j$ に対して $\mathbf R\operatorname{Hom}(E_i,E_j)$ は
  次数零の有限 projective $R$-加群に集中し、その階数は
  \[
  h_{2n}=(n+1)^2,\qquad h_{2n+1}=(n+1)(n+2)
  \]
  である。全 Hom と合成は任意の導来基底変更と整合する。
\item 全軌道環
  \[
  B=\bigoplus_{n\ge0}\operatorname{Hom}(F_0,\mathsf s^nF_0)
  \]
  は一次生成・二次・相対 Koszul であり、全 $B_n$ は階数 $(n+1)^2$ の
  有限 projective 加群である。二次関係は階数七の直和因子となり、
  双対 $B^!$ の全次数、左右および双加群 Koszul 解消は任意基底変更と両立する。
\item $C$ の各頂点単純加群には、距離四以内の五頂点 bar 解消があり、
  これは任意基底変更後にも完全である。
  任意の可換 $R$-代数 $R'$ に対し、canonical な比較
  \[
  D(\operatorname{QGr}B_{R'})\simeq D(\operatorname{Mod}\Lambda_{R'}),
  \qquad D(\operatorname{QGr}B_{R'})^c\simeq\operatorname{Perf}\Lambda_{R'}
  \]
  があり、$\pi B$ を $F_0$ に、$(1)$ を最初に選んだ実際の $\mathsf s$ に送る。
  この比較とその Serre 平方同型は、非平坦なものを含む導来基底変更と整合する。
\end{enumerate}
さらに $B$ が左右 Noetherian なら、任意の有限型可換 $R$-代数 $R'$ 上で、
$\operatorname{qgr}B_{R'}$ は Noetherian heart、module-object 基底変更と fppf descent、
一般 coherent Hom の有限性と平坦基底変更、全 $\chi_q$、AZ 豊富性を持つ。
構造対象は平坦かつ exceptional であり、全幾何 fiber は仮定した $\HH$ の heart と
構造対象を保って exact 同値である。
同じ根の相対 AS4 被覆の存在は、この定理の結論とは別に確認する。
\end{theorem}

\begin{proof}
全証明は相対モジュライ論文 \cite[noncubic有限型効果性]{QuadricModuli2026} に置く。
\end{proof}

\begin{remark}[有限型認識で必要な入力の順序]
この定理では、有限 foundation の全幾何 fiber を既知の normal または cubic 型へ
認識する有限型開条件が、全次数の一様性を供給している。
各次数の平坦 locus を無限に交叉する必要はない。
ここでの大域的な結論は有限 projective 性であり、無限個の Hom 束を
一つの開上で同時に自由化するという、不要な強い主張は含めていない。
また $B$ の Koszul 性を使って全 forward Hom を証明してはいない。
全 fiber の geometric helix から Hom を先に平坦化し、その後に相対 Koszul 原理を使う。
Noetherian 性と同じ根の AS4 被覆は、この順序の後に別途供給すべき入力である。
\end{remark}

\subsection{閉部分と二つの cubic 開部分からの Noetherian 性}

\begin{lemma}[標準 noncubic 族の相対 normal tilting]
\label{lem:nc-relative-normal-tilting}
$R$ を可換 Noetherian $k$-代数とし、
$f_1=uv$、$f_2=(\alpha u+\beta v)(\gamma u-\alpha v)$、
$\alpha^2+\beta\gamma=1$ とする。
\[
 B_0=R[u,v]\langle c,d\rangle/(c^2-f_1,d^2-f_2),\qquad
 Z=R[u,v,h],\quad h=cd+dc
\]
と置く。式~\ref{eq:nchh-plane} の行列分解の cokernel を $M_c,M_d$ とすれば、
両者は $Z$-加群として
\[
 M_c\simeq Z^{\oplus2}\oplus Z(-1)^{\oplus2},\qquad
 M_d\simeq Z^{\oplus2}\oplus Z(-1)^{\oplus2}
\]
であり、これらの表示は任意の基底変更と両立する。
特に四項
\[
 T=(\pi B_0,\pi M_c,\pi B_0(1),\pi M_d(1))
\]
は相対 full strong exceptional foundation で、その End 代数は
式~\ref{eq:nchh-plane} の有限 foundation 代数に一致する。
対応する相対 Morita 同値は四対象、合成、相対 Serre 関手を保ち、
任意の Noetherian 基底変更と整合する。
\end{lemma}
\begin{proof}
全証明は相対モジュライ論文 \cite[noncubic有限型効果性]{QuadricModuli2026} に置く。
\end{proof}

\paragraph{四周期からの Hom $\mathbb Z$-代数の比較。}
標識付き四項 foundation の同型から得た相対 Morita 同値を $F$、
両側の反 Serre 関手を $t,t_0$ と書く。
上の比較は自然同型 $\eta:Ft\simeq t_0F$ を与える。
最初の四項の同型 $\varphi_a:F(E_a)\simeq E_a^0$ ($0\le a<4$) と、
選んだ実際の平方同型による四周期を用い、
$i=a+4n$ に対して $\eta$ の反復から
$\varphi_i:F(E_i)\simeq E_i^0$ を定める。負の $n$ には逆同型を使う。
\[
 f\longmapsto\varphi_jF(f)\varphi_i^{-1}
 \qquad(f:E_i\longrightarrow E_j)
\]
は合成を保つ全 Hom $\mathbb Z$-代数の同型であり、任意基底変更と両立する。
この比較に、二つの根そのものを同一視する仮定は要らない。
異なる $\varphi_i$ の選択は頂点ごとの共役を与えるだけなので、
偶数制限の row と column の Noetherian 性、および線型解消の存在は変わらない。



\begin{lemma}[二つの cubic 開部分と noncubic 閉部分からの普遍 Noetherian 性]
\label{lem:lg-two-rulings-universal-noetherian}
$R$ を可換 Noetherian 環とし、有限 LG 認識開上の foundation と
実際の Morita 根 $s$ を取る。四項を $E_0,E_1,E_2,E_3$ と書き、
指定された同型
\[
 sE_0\simeq E_2,\qquad sE_1\simeq E_3,
 \qquad s^2\simeq S^{-1}[2]
\]
を用いて $E_{2n}=s^nE_0$, $E_{2n+1}=s^nE_1$ と置く。
全ての前向き導来 Hom は有限 projective 加群に次数零で集中し、
任意の底変更と両立することを既に確認したとする。
このとき同じ根による偶数軌道環
\[
 C_n=\operatorname{Hom}(E_0,s^nE_0),\qquad
 a*b=s^{\deg b}(a)\circ b
\]
は、任意の可換 Noetherian $R$-代数 $T$ に対して
$C_T=C\otimes_RT$ が左右 Noetherian となる。

ここで有限 LG 認識開について用いる入力は次の三つである。
二つの射影行列を $M,N$ として
\[
 a=\det M,\quad b=\det N,\quad I=(a,b),
 \qquad Z=\operatorname{Spec}(R/I)
\]
と置く。
\begin{enumerate}
\item $D(a)$ では選んだ ruling の helix が geometric cubic 型である。
\item $D(b)$ では反対の ruling の helix が geometric cubic 型である。
\item $Z$ 上では、非被約構造を含めて、foundation が \'etale 局所的に
  判別式の可逆な二形式による標準 noncubic foundation と同型である。
\end{enumerate}
これらは追加の全次数条件ではなく、有限 LG 認識で確認する条件である。
\end{lemma}

\begin{proof}
全証明は相対モジュライ論文 \cite[noncubic有限型効果性]{QuadricModuli2026} に置く。
\end{proof}

\subsection{同じ根に対する中央 AS4 被覆の延長}

\begin{lemma}[局所中央被覆の有限型近傍への延長]
\label{lem:spread-local-central-hull}
$R$ を可換 Noetherian 環、$\mathfrak p\in\operatorname{Spec}R$ とする。
$B=T_R(V)/(J)$ は連結二次 $R$-代数であり、
$V,J$ はそれぞれ階数四、七の有限 projective 加群、
$J\subset V^{\otimes2}$ は直和因子とする。
$B$ は相対 Koszul であり、全ての $B_n$ は階数 $(n+1)^2$ の
有限 projective 加群として任意の底変更と両立すると仮定する。
$E=B^!$ と置く。

局所環 $R_{\mathfrak p}$ 上で、中央正則元
$\phi_{\mathfrak p}\in(E_{\mathfrak p})_2$ があり、
$P_{\mathfrak p}=E_{\mathfrak p}/(\phi_{\mathfrak p})$ は
階数列 $(1,4,6,4,1)$ の有限自由 Frobenius 代数であるとする。
このとき $\mathfrak p\in D(f)$ となる一つの主開部分へ制限すれば、
$\phi_{\mathfrak p}$ は中央元 $\phi\in(E_f)_2$ に延長し、次が成り立つ。
\begin{enumerate}
\item $\phi$ は全次数で正則であり、この性質は任意の底変更に保たれる。
\item $P=E_f/(\phi)$ は階数列 $(1,4,6,4,1)$ の相対 Koszul
  Frobenius 代数であり、$A=P^!$ の全次数は有限 projective である。
\item $A$ は相対 quadratic Koszul AS4 型の代数であり、
  ある中央正則二次元 $g$ によって $A/(g)\cong B_f$ と書ける。
  この表示、正則性、左右および双加群 Koszul 解消は任意の底変更と両立する。
\end{enumerate}
さらに任意の可換 Noetherian $R_f$-代数 $T$ に対して $B_T$ が左右 Noetherian
ならば、$A_T$ も左右 Noetherian である。
この主張には、$D(f)$ の全ての noncubic fiber において指定した根を
別の ``良い根'' と比較する仮定は要らない。
\end{lemma}

\begin{proof}
全証明は相対モジュライ論文 \cite[noncubic有限型効果性]{QuadricModuli2026} に置く。
\end{proof}

\begin{remark}[一般の Koszul 開性との区別]
上の証明は、指定点の Frobenius 代数 $P_{\mathfrak p}$ が Koszul だから
近傍全体でも Koszul になる、という推論を使っていない。
先に全幾何 fiber の helix を H2 モデルとして認識し、
その同じ even Hom $\mathbb Z$-代数の線形解消から $B_{\bar s}$ の
Koszul 性を得る。この全 fiber の入力と、中央元の一様正則性を経由して
$P_{\bar s}$ の Koszul 性を導く点が重要である。
一般の二次代数では、三次成分の次元を固定したパラメータ空間における
$n$-Koszul locus は開であり、Koszul locus はこれらの可算共通部分である
\cite[Chapter 6, Theorem 2.1 and Corollary 2.2]{PolishchukPositselski2005}。
また同書 Chapter 6, \S6 は、Koszul locus が constructible でない族を与える。
これらを本研究の特別な中央延長に対する否定結果と解するべきではないが、
局所 Koszul 変形原理だけから単一の有限型開近傍での Koszul 性を
推論できるわけではない。本証明は、全幾何 fiber の明示的認識という
別の有限型入力を使って、この問題を回避する。
\end{remark}

\subsection{noncubic 点の有限型近傍での全相対認識}

\begin{theorem}[noncubic 境界を横断する有限型効果性]
\label{thm:noncubic-finite-type-effectivity}
任意の noncubic $\HH$ 対象の normal foundation に対応する点 $x$ と、
定理~\ref{thm:noncubic-hh-trivial-root} の良い根を取る。
有限 foundation の smooth affine chart $(S,x)$ の有限型 affine
\texorpdfstring{\'etale}{etale} 近傍 $(U,u)\to(S,x)$ 上に、
構造対象だけを指定した実際の相対族
\[
 q_U:U\longrightarrow\mathfrak H_2
\]
が存在する。この族の canonical compact 圏は普遍 foundation 代数
$\Lambda_U$ の perfect 圏であり、四項窓は普遍 projective に対応する。

同じ $U$ 上で、選んだ実際の Morita 根 $\mathsf s$ は heart の exact 自己同値となり、
全軌道環 $B$ は二次・相対 Koszul・全次数有限 projective である。
任意の Noetherian 基底変更後に $B$ は左右 Noetherian であり、
同じ根に対する中央正則二次被覆
\[
 A/(g)\simeq B,
 \qquad \rank A_n=\binom{n+3}{3}
\]
を持つ。$A$ は相対 quadratic Koszul AS4 型で、同じ普遍 Noetherian 性を持つ。
全 Hom、canonical tails の比較、相対 Serre 平方同型、被覆の正則性と
左右・双加群 Koszul 解消は任意の導来基底変更と整合する。
根・根同型・被覆はこの族を構成する補助選択であり、$q_U$ の標識には含めない。
\end{theorem}

\begin{proof}
全証明は相対モジュライ論文 \cite[noncubic有限型効果性]{QuadricModuli2026} に置く。
\end{proof}

% !TeX root = ../nc-quadric-classification.tex
% 2026-11-22 / v1.41
\section{内在的な全モジュライの有限型代数性}
\label{sec:global-h2-smooth-atlas}

本節では、全ての noncubic 点で得た有限型の効果的近傍を用いる。
入力は実際の相対 $\mathfrak H_2$ 族であり、全幾何 fiber の分類だけではない。
これらの近傍から指定された pointed 同型を保つ smooth atlas を作り、
既存の noncubic 表示族の準コンパクト性を使って有限個の chart に縮約する。
有限 foundation スタックと内在的モジュライの位相を先に同一視する必要はない。

\subsection{noncubic 点を含む実際の smooth chart}

\begin{corollary}[効果的 noncubic 近傍の smooth atlas]
\label{cor:noncubic-effective-smooth-neighborhood}
$x\in\mathfrak H_2(k)$ を noncubic 点とする。
有限 foundation の標準表示空間 $S_0=\mathbb A^{64}$ の開部分 $S$ と、
$x$ の normal foundation を表す点 $s_x\in S(k)$ を取る。
定理~\ref{thm:noncubic-finite-type-effectivity} で構成した有限型 affine étale 近傍
\[
 a_x:(U_x,u_x)\longrightarrow(S,s_x)
\]
および実際の族 $q_x:U_x\to\mathfrak H_2$ を取る。
全軌道と canonical tails の比較により、この族には次の性質を持つ列 $(E_i)$ がある。
各 $E_i$ は平坦な noetherian compact 対象で、$E_0=\OO$ であり、
各四項窓は full strong exceptional collection、
全前向き derived Hom は次数零の有限 projective 加群である。
さらに $E_{i+4}=S^{-1}E_i[2]$ であり、最初の四項窓の End 代数は
頂点・積・任意基底変更を保って $a_x^*\Lambda$ と同定される。

このとき $q_x$ はスキームで表現可能な smooth 射である。
その fppf 像 $\mathfrak V_x\subset\mathfrak H_2$ は $x$ を含む開部分スタックであり、
有限型で滑らかな Artin スタックである。
\end{corollary}

\begin{proof}
定理~\ref{thm:atlas-from-effective-foundation} の仮定を確認する。
tails の row は次数成分ごとに基底上平坦であり、商関手は tensor と整合するので、
各 $E_i$ は平坦である。Noetherian 性と compact 性は実際の相対族の構成による。
四項の完全性・strong 性と前向き Hom の消滅は、五頂点 bar 比較により任意基底変更上で成立する。
full exceptional collection の Serre rotation は mutation による helix の延長と一致するため、
$E_{i+4}=S^{-1}E_i[2]$ は De Deken--Lowen の geometric $(4,3)$-helix 条件を与える。
最後の End 代数の同定も全軌道比較の一部である。
従って同定理を適用できる。

指定された同型に関する点を明記する。
small extension $A'\twoheadrightarrow A$ と
\[
 u_A\in U_x(A),\qquad c_{A'}\in\mathfrak H_2(A'),\qquad
 \alpha_A:q_x(u_A)\xrightarrow{\sim}c_{A'}|_A
\]
を取ると、$\alpha_A$ により四項窓とその全 Hom の基底を還元先へ移す。
De Deken--Lowen の持上げを基底 $A$ に直接適用し、
例外対象の $\Ext^1,\Ext^2$ の消滅と $A'^\times\to A^\times$ の全射性により、
$E'_0=\OO_{c_{A'}}$ とその指定された pointing を保つ。
五つの Hom 空間の基底も指定還元を延長して選ぶと、実際の合成から $S(A')$ の点が得られる。
$a_x$ の全点での étale 性により、これは $u_A$ を保つ $u_{A'}\in U_x(A')$ へ持ち上がる。

両方の四項 tilting 窓の End 代数をこの基底で同定する。
得られる標準 dg Morita 同値の還元は、対象だけでなく全ての射上で $D(\alpha_A)$ と一致する。
補題~\ref{lem:artin-two-heart-comparison}、または反 Serre 再構成の全忠実性を使えば、
これは $\alpha_A$ 自体を延長する pointed exact 同型となる。
module-object 基底変更により、その同型は全有限型 test 上の相対系にも延長する。
この議論は根や被覆を同型の標識へ加えず、指定した pointed 同型をそのまま保つ。

任意の有限型 affine $T\to\mathfrak H_2$ に対して
\[
 U_x\times_{\mathfrak H_2}T\simeq U_x\times_{\mathfrak P}T,
 \qquad\mathfrak P=[X_3/\operatorname{PGL}_9]
\]
は affine 有限表示対角の基底変更なので有限型スキームである。
上の指定 Artin 持上げと \cite[Tag 02HY]{StacksSmoothAtlas} により、その $T$ への射は smooth である。
同 tag の有限型 Artin test は、$k$ が代数閉なので剰余体 $k$ の test に含まれる。
開像と smooth groupoid の商の構成は定理~\ref{thm:atlas-from-effective-foundation} の通りである。
\end{proof}

\subsection{有限個の chart による全体の表示}

\begin{theorem}[全 $\mathfrak H_2$ の有限型代数性と滑らかさ]
\label{thm:global-h2-smooth-finite-type}
局所有限型 $k$-スキームの site 上の内在的スタック $\mathfrak H_2$ は、
$k$ 上有限型で滑らかな純次元三の Artin スタックの制限により表される。
その対角は affine かつ有限表示である。
有限型の滑らかなスキーム $U$ と representable smooth 全射
\[
 q:U\longrightarrow\mathfrak H_2
\]
を取れ、全ての pointed exact 同値を用いる relation scheme
\[
 R=U\times_{\mathfrak H_2}U
   \simeq U\times_{\mathfrak P}U
\]
に対して $\mathfrak H_2\simeq[U/R]$ である。
ここで選ばれた foundation、Serre 根、平方同型、AS4 被覆を保存することは、射の条件に加えない。
\end{theorem}

\begin{proof}
まず有限個の noncubic 近傍で足りることを示す。
定理~\ref{thm:noncubic-relative-effectivity} による実際の射を、
表示スタックの atlas へ引き戻すと
\[
 p:P=Q^{\mathrm{nd}}\times Q^{\mathrm{nd}}
   \longrightarrow\mathfrak N^{\mathrm{pres}}
   \longrightarrow\mathfrak H_2
\]
を得る。$P$ は有限型の準コンパクトな $k$-スキームである。
命題~\ref{prop:noncubic-presentation-stack} により、全 noncubic $k$-対象はその像に入る。
各 noncubic 点 $x$ に前系の開部分 $\mathfrak V_x$ を取る。
引戻し $p^{-1}(\mathfrak V_x)$ は $P$ の開部分スキームであり、これらは $P$ の全閉点を覆う。
有限型 $k$-スキームは Jacobson なので、全体も覆う。
従って準コンパクト性により有限個
\[
 p^{-1}(\mathfrak V_{x_1}),\ldots,p^{-1}(\mathfrak V_{x_m})
\]
で $P$ を覆える。この段階で使ったのは $P\to\mathfrak H_2$ の存在と
各 $\mathfrak V_x$ の内在的な開性だけであり、$\mathfrak N^{\mathrm{pres}}\to\mathfrak H_2$
の full faithfulness や、任意の相対 noncubic 族の正規形は使わない。

定理~\ref{thm:cubic-smooth-open} の有限型 smooth atlas $U_{\mathrm{cub}}$ と合わせ、
\[
 U=U_{\mathrm{cub}}\amalg U_{x_1}\amalg\cdots\amalg U_{x_m}
\]
と置く。有限和なので $U$ は有限型で滑らかである。
各成分の射は representable smooth なので $q:U\to\mathfrak H_2$ も同様であり、
その開像は全ての $k$-対象を含む。

全ての相対族を覆うことも確認する。
有限型 affine $T\to\mathfrak H_2$ に対して、$U\times_{\mathfrak H_2}T\to T$ の像は開である。
$T$ の各閉点の剰余体は $k$ であり、その fiber は上の分類と有限被覆によって像に入る。
従ってこの開像は全閉点を含み、Jacobson 性により $T$ 全体である。
ゆえにこの射は smooth 全射であり、fppf 局所に section を持つ。
任意の局所有限型 $T$ については有限型 affine 開被覆へ制限すればよい。
これで $q$ が元の site 上の全対象を fppf に覆うことが示された。

相対反 Serre 認識の全忠実性から、relation は表示した $U\times_{\mathfrak P}U$ と等しい。
$\mathfrak P$ の対角が affine 有限表示なので、$R$ は有限型スキームである。
両射影 $R\rightrightarrows U$ は $q$ の基底変更として smooth であり、
\cite[Tags 04X0, 04TK]{StacksSmoothAtlas} により $[U/R]$ は通常の全スキーム上の Artin スタックとなる。
その元の site への制限は $\mathfrak H_2$ と同値である。
有限型 atlas $U$ によって有限型性が従い、$U$ の滑らかさによって $k$ 上の滑らかさが従う。
対角の性質は命題~\ref{prop:relative-h2-algebraicity-reduction} による。

最後に、系~\ref{cor:cubic-formal-h2-abelian} と系~\ref{cor:noncubic-all-artin-h2} により、
各 $k$-点の形式変形 groupoid は pointed Abel 変形と同値である。
その接空間は $\operatorname{HH}^2(\Lambda)$、無限小自己同型空間は
$\operatorname{HH}^1(\Lambda)$ であり、
命題~\ref{prop:foundation-hochschild} によって両次元の差は三である。
滑らかな Artin スタックの局所次元はこの差に等しいから、全閉点で三となる。
有限型 smooth atlas の各成分は閉点を持ち、次元は局所一定なので、全体は純次元三である。
\end{proof}

\subsection{framed 実現空間と ambient 内の非開性}

\begin{corollary}[framed 実現空間のスキーム性]
\label{cor:global-xas-smooth-scheme}
\[
 X_{\mathrm{as}}=X_3\times_{\mathfrak P}\mathfrak H_2
\]
は有限型で滑らかな純次元八十三のスキームにより表される。
$X_{\mathrm{as}}\to X_3$ は有限型の monomorphism で、separated かつ unramified、quasi-finite である。
\end{corollary}

\begin{proof}
全スキーム上では前定理の $[U/R]$ を用いる。
$R=U\times_{\mathfrak P}U$ は全ての同型を記録するため、
$[U/R]\to\mathfrak P$ は全スキーム上でも全忠実である。
実際、対象を $U$ へ fppf 局所に持ち上げると Isom はまさに $R$ であり、
この一致は降下する。従って $X_{\mathrm{as}}\to X_3$ は monomorphism である。

命題~\ref{prop:relative-h2-algebraicity-reduction} により、$X_{\mathrm{as}}$ は有限型代数空間である。
また $X_{\mathrm{as}}\to\mathfrak H_2$ は $\operatorname{PGL}_9$-torsor なので smooth で相対次元八十である。
従って $X_{\mathrm{as}}$ は $k$ 上 smooth で純次元 $3+80=83$ である。

$X_{\mathrm{as}}$ と $X_3$ はともに有限型 $k$-空間なので、両者の間の射は有限型である。
局所有限型 monomorphism は separated、unramified、locally quasi-finite かつスキームで表現可能である
\cite[Tag 05W6, Lemma 76.14.8]{StacksRepresentableMonomorphism}。
特に $X_3$ がスキームであるため $X_{\mathrm{as}}$ 自身もスキームである。
同じ結論は separated locally quasi-finite な代数空間からスキームへの射に関する
\cite[Tag 03XX, Proposition 67.50.2]{StacksRepresentableMonomorphism} からも従う。
有限型性は準コンパクト性を含むので、射は quasi-finite である。
\end{proof}

\begin{proposition}[canonical 実現像は ambient 全体で開ではない]
\label{prop:global-h2-not-ambient-open}
$\Phi:\mathfrak H_2\to\mathfrak P=[X_3/\operatorname{PGL}_9]$ は、
noncubic の二つの根対が一致する点で open immersion ではない。
従って $X_{\mathrm{as}}\to X_3$ も、その点の任意の frame の近傍で open immersion ではない。
\end{proposition}

\begin{proof}
一致点 $x$ において、系~\ref{cor:noncubic-all-artin-h2} は
\[
 \dim T_x\mathfrak H_2=4,\qquad
 \dim\operatorname{HH}^1(\Lambda)=1
\]
を与える。ここで $T_x$ は閉 fiber の同定を保つ一次変形の同型類の空間である。
一方、系~\ref{cor:artin-polarization-retraction} の分裂完全列は
\[
 0\longrightarrow\operatorname{HH}^1(\Lambda)
 \longrightarrow T_x\mathfrak P
 \longrightarrow\operatorname{HH}^2(\Lambda)\longrightarrow0
\]
を与え、$\dim T_x\mathfrak P=5$ である。
$\Phi$ は canonical 反 Serre 環を取る射なので、その微分の像は
この分裂における $\operatorname{HH}^2$ の四次元の像である。
従って一次変形の射は全射ではなく、$\Phi$ はこの点で open immersion にならない。
開 immersion なら形式変形 groupoid を同一にすることに反するからである。

framed 空間については、その任意の frame $J$ で
\[
 \dim T_JX_{\mathrm{as}}=83,\qquad
 \dim T_JX_3=5+(80-1)=84
\]
となる。後者の $80-1$ は $\operatorname{PGL}_9$ の軌道接空間の次元である。
open immersion ならこの二つの接空間は同型になるので、framed 射についても結論を得る。
\end{proof}

\begin{remark}[表現可能性から自動的には出ない比較]
上で構成したのは内在的な族のスタックそのものの有限型 smooth atlas と、
framed 部分層のスキームによる表現である。
これによって $X_{\mathrm{as}}$ を $X_3$ の開部分と同一視することはできず、
実際に前命題が一致根対でその同一視を否定する。
また有限 foundation の ruling 商 $\mathfrak Q_{\mathrm{fnd}}$ との相対同値や、
その表示からの内在位相の比較は、この atlas の構成だけでは結論しない。
全スキーム上の $[U/R]$ は元の有限型 site 上の $\mathfrak H_2$ を延長するが、
任意の Noetherian 基底上で素朴に同じ Abel 公理を課した関手との一致は別の連続性の問題である。
\end{remark}

% !TeX root = ../nc-quadric-classification.tex
% 2026-11-22 / v1.43
\section{noncubic 境界の内在的商表示とモジュライ写像}
\label{sec:noncubic-boundary-comparison}

標数零の代数閉体 $k$ 上で考える。
本節では、二形式の表示スタックと前節の内在的被約境界を、相対的な対象・射まで比較する。
幾何点の全単射だけでなく、境界の smooth chart 上の実際の相対族の比較を用いる。
その結果、表示側の不変式と慣性計算が内在的境界へ移る。
以下の相対同定とその組合せについて、包括的な新規性は未確認である。

定理~\ref{thm:noncubic-relative-effectivity} の射は、有限型 atlas $P$ と
有限型作用 groupoid $G\times P\rightrightarrows P$ 上の実際の族と同型から作られる。
従ってその降下データは、定理~\ref{thm:global-h2-smooth-finite-type} の
通常の代数スタック $[U/R]$ への射も定める。
以下で代数スタックの一般判定を使う際はこの射を用い、元の有限型 site へ制限する。
任意の Noetherian 基底に素朴に同じ Abel 公理を課す関手との一致は仮定しない。

\subsection{幾何点の慣性から相対的な忠実性へ}

\begin{lemma}[標数零における相対慣性の消滅]
\label{lem:charzero-geometric-inertia-faithful}
$k$ を標数零の体、$f:\mathcal X\to\mathcal Y$ を
有限型の代数的スタックの間の射とし、両対角は affine かつ有限表示とする。
任意の代数閉拡大体 $K/k$ と $x\in\mathcal X(K)$ に対して
\[
 \operatorname{Aut}_{\mathcal X(K)}(x)
 \longrightarrow\operatorname{Aut}_{\mathcal Y(K)}(f(x))
\]
が単射ならば、$f$ は代数空間で表現可能である。
この仮定は幾何点における抽象群の単射でよい。
\end{lemma}

\begin{proof}
有限型 smooth atlas $V\to\mathcal X$ を取り、対応する対象を $x_V$ と書く。
二つの自己同型群スキームの射の核を
\[
 K_V=\ker\bigl(
 \underline{\operatorname{Aut}}_{\mathcal X}(x_V)
 \longrightarrow
 \underline{\operatorname{Aut}}_{\mathcal Y}(f(x_V))\bigr)
\]
とする。対角の仮定と $V$ の Noetherian 性から、$K_V\to V$ は
affine かつ有限表示の群スキームである。
各幾何点 $\bar v\to V$ の上で、その fiber はただ一つの
幾何点を持つ有限型群スキームである。
Cartier の定理により標数零の有限型群スキームは smooth なので、
この fiber は $\operatorname{Spec}\kappa(\bar v)$ に等しい
\cite[Tag 0BF6, Lemma 39.8.2]{StacksBoundaryComparison}。

相対微分は fiber への基底変更と整合する。
従って有限型加群 $\Omega_{K_V/V}$ は全ての幾何点で零となり、
Nakayama の補題により零である。
ゆえに $K_V\to V$ は unramified である
\cite[Tag 02G3, Lemma 29.36.2]{StacksBoundaryComparison}。
その単位切断 $e:V\to K_V$ は unramified 射の切断だから開埋込みである
\cite[Tag 024T, Proposition 41.6.1]{StacksBoundaryComparison}。
全ての幾何 fiber が単位点だけからなるので、$e$ は全射でもある。
従って $e$ は同型であり、$K_V$ は自明である。

これは相対慣性 $\mathcal I_{\mathcal X/\mathcal Y}$ が atlas $V$ 上で自明であることを意味する。
fppf 降下により相対慣性は $\mathcal X$ 上で自明となり、
相対慣性による表現可能性判定を適用できる
\cite[Tag 04YY, Lemma 101.6.2]{StacksBoundaryComparison}。
\end{proof}

\begin{lemma}[幾何点の自己同型群のスキームとしての比較]
\label{lem:nc-geometric-inertia-scheme-iso}
自然な射
\[
 f:\mathfrak N^{\mathrm{pres}}\longrightarrow\mathfrak N
\]
は代数空間で表現可能であり、各幾何点で自己同型群スキームの同型を誘導する。
\end{lemma}

\begin{proof}
まず $\mathfrak N^{\mathrm{pres}}=[P/G]$ は smooth、従って reduced である。
定理~\ref{thm:noncubic-relative-effectivity} の射は全幾何点を
noncubic 部分へ送るので、定理~\ref{thm:noncubic-boundary-smooth-codim-two}
の reduced closed substack $\mathfrak N$ を経由する。
この因数分解は reduced atlas $P$ 上で確認して降下させればよい。

命題~\ref{prop:noncubic-presentation-stack} の慣性計算は、
この射が誘導する幾何点の自己同型群の写像を同型と同定している。
実際、$G$ の元は $B_f$ の次数付き同型を与え、その tails 関手が
相対族の構成で用いる射である。次数スカラーは、$\mathcal O$ 上で恒等となる
自然同型によって消え、他の核はない。
したがって前補題より $f$ は表現可能である。

幾何点 $p$ の自己同型群スキーム間の射を $a:G_p\to H_{f(p)}$ と書く。
両群は affine 有限型であり、Cartier の定理によって smooth である。
そのスキームとしての核は前補題の議論により自明なので、
$a$ は群スキームの monomorphism であり、その接写像は単射である。
一方、$a$ は幾何点で全射である。
各 fiber は核の移送で零次元だから両群の次元は等しい。
従って単位元における接写像は同型であり、
smooth な $k$-スキームの間の Jacobian 判定により $a$ は単位元で étale である。
群の平行移動により全ての点で étale となる。
étale な monomorphism は開埋込みであり、幾何点で全射なので同型である
\cite[Tag 025G, Theorem 41.14.1]{StacksBoundaryComparison}。
\end{proof}

\subsection{二形式表示と内在的境界の相対比較}

\begin{lemma}[境界の smooth chart 上での実際の二形式表示]
\label{lem:noncubic-boundary-chart-factorization}
任意の noncubic 幾何点の近傍で、定理
\ref{thm:noncubic-boundary-smooth-codim-two} の smooth chart
\[
 q_D:D_U\longrightarrow\mathfrak N
\]
のある \texorpdfstring{\'etale}{etale} 被覆 $D'_U\to D_U$ と射
$h:D'_U\to P$ を取れて、実際の相対族の同型
\[
 q_D|_{D'_U}\simeq q_{\mathrm{nc}}\circ h
 \tag{NC-chart-factor}\label{eq:noncubic-boundary-chart-factorization}
\]
が成立する。この同型は構造対象を保ち、任意の有限型基底変更と整合する。
根・根同型・AS4 被覆を保存する条件は加えない。
\end{lemma}
\begin{proof}
共通 chart 上で二つの行列式は実際に零である。
命題~\ref{prop:lagrangian-relative-double-rank-three} を適用し、
必要な \texorpdfstring{\'etale}{etale} 被覆上で有限 foundation を
\[
 f_1=uv,\qquad f_2=(\alpha u+\beta v)(\gamma u-\alpha v),
 \qquad\alpha^2+\beta\gamma=1
\]
の標準 foundation と同定する。両形式の判別式は単元なので、これは
$h:D'_U\to P$ を定める。

この段階で得たものは有限代数の同型であり、次に相対 Abel 圏を比較する。
補題~\ref{lem:nc-relative-normal-tilting} により、二形式族の標準四項は
相対 tilting foundation であり、その End 代数は上の有限 foundation に一致する。
一方、$q_D$ の canonical compact 圏も同じ有限代数の perfect 圏に同値である。
従って両側の間には、構造対象 $\mathcal O$ と四項を保つ相対 Morita 同値 $F$ がある。
この同値は perfect 双加群により与えられ、任意の基底変更に整合する。

相対 Serre 双対の Yoneda 一意性により、反 Serre 関手
$t=S^{-1}[2]$ と $t_f=S_f^{-1}[2]$ に対して自然同型
$Ft\simeq t_fF$ を得る。その反復と構造対象の同型を用いると
\[
 \bigoplus_{n\ge0}\operatorname{Hom}(\mathcal O,t^n\mathcal O)
 \simeq
 \bigoplus_{n\ge0}\operatorname{Hom}(\mathcal O_f,t_f^n\mathcal O_f)
\]
が積を保つ graded 環の同型となる。積との整合性は、同じ自然同型の反復を
用いて各次数を比較することから従う。各 Hom は次数零の有限 projective 加群であり、
比較は全有限型基底変更に整合する。

相対反 Serre 再構成（命題~\ref{prop:relative-h2-antiserre-reconstruction}）と
その全忠実性（命題~\ref{prop:relative-h2-stack-recognition}）によって、
この graded 同型は両相対 Abel 系の pointed exact 同値を与える。
特に canonical compact 圏だけでなく、coherent heart と module-object 基底変更の
系が一致する。これが~\eqref{eq:noncubic-boundary-chart-factorization} である。
\end{proof}

\begin{lemma}[smooth 同次元スタックの接判定]
\label{lem:smooth-stack-etale-tangent-comparison}
$k$ を標数零の代数閉体とし、$f:\mathcal X\to\mathcal Y$ を
$k$ 上 smooth、有限型、同じ純次元の Artin スタックの間の射とする。
両対角は affine 有限表示とする。全ての代数閉拡大体 $K/k$ と
$x\in\mathcal X(K)$ について、次を仮定する。
\begin{enumerate}
\item $\operatorname{Aut}_{\mathcal X}(x)(K)\to
  \operatorname{Aut}_{\mathcal Y}(f(x))(K)$ は同型である。
\item 同型類の一次変形の接空間の射
  $T^0_x\mathcal X\to T^0_{f(x)}\mathcal Y$ は全射である。
\end{enumerate}
このとき $f$ は algebraic space で表現可能な \texorpdfstring{\'etale}{etale} 射である。
\end{lemma}
\begin{proof}
補題~\ref{lem:charzero-geometric-inertia-faithful} により $f$ は表現可能である。
各幾何点の慣性群スキームの同型は、補題~\ref{lem:nc-geometric-inertia-scheme-iso}
の群スキームの議論と同じく、Cartier の定理、核の消滅、同次元性から従う。

smooth なスタックでは
\[
 \dim T^0_x\mathcal X
  =\dim_x\mathcal X+\dim\operatorname{Lie}\operatorname{Aut}(x)
\]
である。前段の慣性群の同型と同次元性から、仮定の全射は同型となる。
相対一次変形の完全列
\[
 \operatorname{Lie}\operatorname{Aut}(x)
 \longrightarrow\operatorname{Lie}\operatorname{Aut}(f(x))
 \longrightarrow T_x(\mathcal X/\mathcal Y)
 \longrightarrow T^0_x\mathcal X
 \longrightarrow T^0_{f(x)}\mathcal Y
\]
によって $T_x(\mathcal X/\mathcal Y)=0$ である。
この列は、$f(x)$ の自明な一次変形との同型を付けた $x$ の一次変形を考えれば得られる。
その同型の変更が二つの無限小自己同型空間の余核を与える。

最後に smooth scheme atlas $V\to\mathcal Y$ を取り、
$Z=\mathcal X\times_{\mathcal Y}V$ とする。表現可能性により $Z$ は代数空間で、
$Z\to\mathcal X$ は smooth なので $Z$ 自身も $k$ 上 smooth である。
各点で $Z$ と $V$ の次元は等しく、相対接空間は零である。
両者の \texorpdfstring{\'etale}{etale} scheme chart 上の Jacobian 判定により
$Z\to V$ は \texorpdfstring{\'etale}{etale} となる。
従って $f$ も \texorpdfstring{\'etale}{etale} である。
\end{proof}

\begin{theorem}[二形式表示は内在的被約境界を表す]
\label{thm:noncubic-intrinsic-presentation-equivalence}
二形式の実際の相対族は、局所有限型 $k$-スキームの site 上で
\[
 q_{\mathrm{nc}}:\mathfrak N^{\mathrm{pres}}=[P/G]
  \xrightarrow{\ \sim\ }\mathfrak N
\]
という同値を与える。ここで
$\mathfrak N=(\mathfrak H_2\setminus\mathfrak V_{\mathrm{cub}})_{\mathrm{red}}$
は内在的な被約境界である。従って任意の非被約 test も含め、対象と
pointed exact 同型の両方がこの表示商で記述される。
\end{theorem}
\begin{proof}
表示スタックは smooth、従って reduced であり、その全幾何点は noncubic である。
従って既存の実際の射 $q_{\mathrm{nc}}$ は一意に $\mathfrak N$ を経由する。
この因数分解は atlas $P$ 上で reduced source の性質を使って確認し、降下させる。
両スタックは有限型、smooth、純次元一で、対角は affine 有限表示である。

幾何点の同型類と自己同型群の対応は
命題~\ref{prop:noncubic-presentation-stack} による。
接空間の全射性を確認する。
任意の幾何点 $p\in P(K)$ とその像 $x$ に、
補題~\ref{lem:noncubic-boundary-chart-factorization} の
$q_D:D'_U\to\mathfrak N$ と $h:D'_U\to P$ を取る。
$x$ の上の点 $u$ における $h(u)$ は $p$ と同じ pointed 圏を与えるので、
命題~\ref{prop:noncubic-presentation-stack} により $G(K)$ の同じ軌道に入る。
$h$ を定数の $G(K)$ 元で移せば、$h(u)=p$ としてよい。
\eqref{eq:noncubic-boundary-chart-factorization} と $q_D$ の smooth 性から
\[
 T_uD'_U\longrightarrow T^0_p\mathfrak N^{\mathrm{pres}}
 \longrightarrow T^0_x\mathfrak N
\]
の合成は全射であり、第二の射も全射である。
比較同型の $u$ 上の値を替えると接空間上で可逆な自己同型を合成するだけなので、
この結論は比較の選択に依存しない。
前補題により $q_{\mathrm{nc}}$ は表現可能な \texorpdfstring{\'etale}{etale} 射である。

最後に、全代数閉体上で幾何点 groupoid が同値である。
従って表現可能な $q_{\mathrm{nc}}$ の各幾何 fiber は一点であり、
\texorpdfstring{\'etale}{etale} 性と合わせて全射な開埋込み、すなわち同型となる。
これは \cite[Tag 0DUD, Lemma 101.45.7]{StacksBoundaryComparison} の判定とも一致する。
代数スタックとしての同型から元の site 上での同値が従い、非被約基底も含まれる。
\end{proof}

\begin{remark}[一致根対における二つの接方向]
根対が異なる点では無限小慣性は零であり、両境界スタックの接次元は一である。
一致根対では無限小慣性は一次元であり、接次元は二である。
局所的に第一形式を $uv$ と固定し、第二形式の $uv$ 係数を一とすれば
\[
 f_2=a u^2+uv+c v^2
\]
が一致点 $(a,c)=(0,0)$ における二つの接方向を与える。
前定理はこの二方向も内在的な接空間へ同型に送られることを示す。
この chart では
\[
 \chi=\frac{1}{1-4ac},\qquad
 \chi-1=\frac{4ac}{1-4ac}.
\]
従って一致点で $d\chi=0$ であり、良モジュライ座標の一次微分だけでは
この接比較を検証できない。
\end{remark}

\subsection{境界の良モジュライ空間と全体の粗モジュライ障害}

\begin{lemma}[定値な穿孔直線の特殊化]
\label{lem:constant-punctured-line-space}
$k$ を代数閉体、$Z$ を任意の $k$-代数空間とし、
$a:\mathbb A^1_k\to Z$ を射とする。
$a|_{\mathbb G_m}$ が $z:\operatorname{Spec}k\to Z$ を経由する定値射なら、
$a(0)=z$ である。$Z$ に準分離性や有限型性は仮定しない。
\end{lemma}
\begin{proof}
スキームからの \texorpdfstring{\'etale}{etale} chart $U\to Z$ と
$a(0)$ 上の $k$-点 $u_0\in U(k)$ を選び、$u_0$ の affine 近傍に制限する。
代数空間の対角はスキームで表現可能なので
\[
 W=\mathbb A^1_k\times_Z U
\]
はスキームであり、$\mathbb A^1_k$ 上 \texorpdfstring{\'etale}{etale} である。
$(0,u_0)$ を含む既約な affine 開近傍 $W_0\subset W$ を取る。
これは滑らかな曲線であり、$W_0^\times=W_0\times_{\mathbb A^1}\mathbb G_m$
は非空な稠密既約開部分である。

仮定により、$W_0^\times\to U$ は
$U_z=U\times_Z\operatorname{Spec}k$ を経由する。
$U_z$ は $k$ 上 \texorpdfstring{\'etale}{etale} なスキームなので
$k$-点の離散和である。従って既約性により $W_0^\times$ の像は
ただ一つの $k$-点 $u_1\in U(k)$ に含まれる。
この点は affine スキーム $U$ の閉点である。
連続性と $W_0^\times$ の稠密性から、$W_0$ 全体の像も $\{u_1\}$ に含まれ、
$u_0=u_1$ となる。これを $Z$ へ送れば $a(0)=z$ を得る。
ここで使った閉点性はスキーム $U$ 上のものだけであり、
任意の代数空間 $Z$ の $k$-点が閉であるとは仮定していない。
\end{proof}

\begin{corollary}[全 $\mathfrak H_2$ の粗モジュライ代数空間の不存在]
\label{cor:global-h2-no-coarse-space}
一根を共有する根対に対応する点を $x_{\mathrm{sh}}$、
二つの根対が一致する点を $x_{\mathrm{eq}}$ とする。
任意の代数空間 $Z$ と射 $\psi:\mathfrak H_2\to Z$ に対して
\[
 \psi(x_{\mathrm{sh}})=\psi(x_{\mathrm{eq}})
\]
である。従って $\mathfrak H_2$ は、幾何点の同型類を全単射に分類する
粗モジュライ代数空間を持たない。
同じ結論は $\mathfrak N^{\mathrm{pres}}$ および $\mathfrak N$ にも成り立つ。
\end{corollary}
\begin{proof}
$\mathbb A^1_t$ 上の実際の表示族
\[
 f_1=uv,\qquad f_2=u(v-tu)
\]
を使う。両判別式は一であり、これは $P$ への射を定める。
$t=0$ の fiber は $x_{\mathrm{eq}}$、$t\ne0$ の fiber は $x_{\mathrm{sh}}$ である。
穿孔直線上で同型類が一定であるだけでなく、族自身が fppf 局所的に定値となることを確認する。
$t=r^2$ の有限 \texorpdfstring{\'etale}{etale} 被覆を取り、
\[
 A_r=\operatorname{diag}(r^2,1),\qquad a_1=a_2=r
\]
と置く。$G$ の作用式を使うと、$(A_r,r,r)$ は
$(uv,u(v-u))$ を $(uv,u(v-r^2u))$ へ送る。
よってこの族は穿孔直線の被覆上で一定の対象と同型である。

任意の $\psi$ と合成して得る $a:\mathbb A^1\to Z$ は、
穿孔直線上では fppf 局所的に値 $\psi(x_{\mathrm{sh}})$ の定値射である。
代数空間は fppf sheaf なので、これは穿孔直線上の定値射そのものである。
前補題から $\psi(x_{\mathrm{eq}})=\psi(x_{\mathrm{sh}})$ が従う。
体上の分類により二点は異なる pointed exact 同値類であるから、
幾何点の同型類に対して単射となる射は存在しない。
表示スタックについても同じ族と議論を使える。
境界については同値定理、または実際の表示族の $\mathfrak N$ への因数分解を使う。
\end{proof}

\begin{corollary}[内在的境界の良モジュライ空間]
\label{cor:intrinsic-boundary-good-moduli}
定理~\ref{thm:noncubic-intrinsic-presentation-equivalence} の同値の下で、
命題~\ref{prop:noncubic-good-moduli} の不変式は
\[
 \chi:\mathfrak N\longrightarrow\mathbb A^1_\chi
\]
という good moduli space を与える。
一根共有と根対一致はともに $\chi=1$ へ写る。
この結論は境界の good moduli space を記述するものであり、
$\mathfrak H_2$ 全体の good moduli space の存在については主張しない。
\end{corollary}
\begin{proof}
命題~\ref{prop:noncubic-good-moduli} の
$k[P]^G=k[\chi]$ と linearly reductive 群による affine 商の結果を、
表示スタックと内在的境界の同値に沿って移せばよい。
\end{proof}


\subsection{良モジュライ射の平坦性と二つの特殊 fiber}

以下では $k$ は標数零の代数閉体とする。
定理~\ref{thm:noncubic-intrinsic-presentation-equivalence} により $\mathfrak N=[P/G]$ と同定し、
良モジュライ射 $\chi:\mathfrak N\to\mathbb A^1$ を用いる。
以下の計算は表示商の smooth atlas 上で行い、この同値に沿って内在的境界へ移す。

\begin{proposition}[平坦な良モジュライ射と特殊 fiber]
\label{prop:noncubic-good-moduli-fibers}
$\chi$ は平坦かつ有限表示である。$\xi\in k$ に対し
$\mathfrak N_\xi=\chi^{-1}(\xi)$ を scheme-theoretic fiber とすると、次が成り立つ。
\begin{enumerate}
\item $\xi\ne0,1$ の fiber は smooth である。
\item $\xi=0$ の fiber は smooth な被約 Cartier divisor の二重構造である。
\item $\xi=1$ の fiber は reduced である。その幾何点は、一根共有の類と二根対一致の類の二つである。
一根共有の類は smooth な開部分をなし、その閉包は一致類を含む。
一致類の近傍で、smooth chart 上の fiber 方程式は $ac=0$ となる。
\end{enumerate}
従って $\chi$ が smooth でない点は、$\xi=0$ の全点と $\xi=1$ の一致類に限られる。
\end{proposition}

\begin{proof}
まず固定した第一形式への帰着を smooth な被覆として実現する。
$\operatorname{GL}(W)$ は $Q^{\mathrm{nd}}$ に推移的に作用し、
$uv$ の安定化群は $O(uv)\simeq\mathbb G_m\rtimes C_2$ で smooth である。
従って軌道射 $\operatorname{GL}(W)\to Q^{\mathrm{nd}}$ は smooth 全射である。
例えば単位元における微分は、二つの非対角成分から $u^2,v^2$ を、
対角成分の和から $uv$ を得るため全射であり、群作用により全点で同様である。
その基底変更に同型な射
\[
 \operatorname{GL}(W)\times Q^{\mathrm{nd}}\longrightarrow P,
 \qquad (A,g)\longmapsto
 \bigl(\operatorname{Sym}^2(A)(uv),\operatorname{Sym}^2(A)g\bigr)
\]
も smooth 全射である。
さらに $P\to[P/G]$ は smooth 全射なので、これらの合成上で fiber の性質を調べられる。

$g=au^2+buv+cv^2$、$\Delta=b^2-4ac$ と書けば、この被覆上で
\[
 \chi=\frac{b^2}{\Delta},\qquad \Delta\in\OO^\times.
\]
従って fiber の方程式は
\begin{equation}\label{eq:noncubic-fiber-quadric}
 (1-\xi)b^2+4\xi ac=0,\qquad \Delta\ne0
\end{equation}
である。$\xi\ne0,1$ なら、この二次形式は三変数の非退化二次形式であり、
その唯一の affine 特異点である原点は $\Delta\ne0$ により除かれる。
従って fiber は smooth である。

$\xi=0$ では式は $b^2=0$ となる。
被約部分 $b=0$ 上では $ac$ が可逆なので、これは smooth な Cartier divisor の二重構造である。
この divisor は $P$ 上では判別式の極化 $D(f_1,f_2)=0$ で与えられる。
$D$ は $G$ の半不変式であり、分母 $\Delta(f_1)\Delta(f_2)$ は可逆であるため、
商スタック上でも $\operatorname{div}(\chi)=2\mathfrak E_0$ と書ける。
ここで $\mathfrak E_0$ は smooth な被約 Cartier divisor である。

$\xi=1$ では $ac=0$ となり、$\Delta\ne0$ は $b$ が可逆であることと同値である。
これは二つの座標軸の通常交差と可逆な $b$ の積なので reduced である。
$(a,c)\ne(0,0)$ なら二形式は一根だけを共有し、fiber は smooth である。
$a=c=0$ は二形式が比例する場合であり、fiber はそこで交差特異点を持つ。
$b$ を一に正規化した横断座標では
\begin{equation}\label{eq:noncubic-chi-coincident-node}
 \chi-1=\frac{4ac}{1-4ac}.
\end{equation}
$a$ 軸または $c$ 軸の原点への特殊化が、一根共有の類から一致類への特殊化である。
二つの軸は $u,v$ の交換で同じ軌道へ移るので、幾何点の類は記述した二つで尽くされる。

平坦性は $P$ 上で確認できる。
$P$ は整 affine スキームであり、$\chi$ は非定数なので $k[\chi]\to k[P]$ は単射である。
従って $k[P]$ は PID $k[\chi]$ 上 torsion-free、従って flat である。
有限表示性は有限型 $k$-スキーム・スタックの間の射であることから従う。
平坦性は smooth 全射 $P\to\mathfrak N$ から降下する。
最後の smooth locus の記述は、平坦有限表示性と上で判定した幾何 fiber の smooth locus から従う。
\end{proof}

\begin{remark}[resultant による共有根の方程式]
二つの二元二次形式に対し、標準的な Sylvester resultant を用いると
\[
 D(f_1,f_2)^2-\Delta(f_1)\Delta(f_2)=4\operatorname{Res}(f_1,f_2).
\]
従って $\chi=1$ は共有根の resultant divisor である。
二つの形式自身は非退化であるが、共有根が二つになる一致類では、この divisor の局所方程式が $ac=0$ となる。
\end{remark}

\subsection{全配置の慣性と一致類の接表現}

以下で $D_8$ は位数八の二面体群を表す。
$c,d$ の独立の符号変更を $e_1,e_2$、交換を $\tau$ とする。
これらは $D_8=\langle e_1,e_2,\tau\rangle$ を生成し、その中心の非自明元は $e_1e_2$ である。

\begin{proposition}[特殊配置も含む全慣性]
\label{prop:noncubic-all-stabilizers}
内在的境界の幾何点の慣性は次の通りである。
\begin{enumerate}
\item 四根が相異なり $\chi\ne0,1$ なら $D_8\times C_2$ で、位数は十六である。
\item $\chi=0$ では位数三十二であり、根配置の慣性 $D_8$ を商とする完全列
\[
 1\longrightarrow\mu_2^2\longrightarrow H_0\longrightarrow D_8\longrightarrow1
\]
を持つ。この完全列は分裂しない。
\item 一根共有では $D_8$ である。
\item 二根対一致では
\begin{equation}\label{eq:noncubic-coincident-full-stabilizer}
 H_{\mathrm{eq}}\simeq
 \bigl((\mathbb G_m\rtimes C_2)\times D_8\bigr)/\mu_2,
\end{equation}
である。$C_2$ は $\mathbb G_m$ に逆元で作用し、割る中央 $\mu_2$ は
$(-1,e_1e_2)$ で生成される。この群の単位成分は $\mathbb G_m$ で、連結成分は八個である。
\end{enumerate}
従って Deligne--Mumford locus は、一致類を除いた開部分である。
\end{proposition}

\begin{proof}
四根が相異なる場合には、根対を $\{0,\infty\}$ と $\{1,\lambda\}$ と書ける。
順序付き四点の交比に対する $S_4$ の作用の核は Klein 四群 $V_4$ である。
二対への分割を保つ順列群は位数八の $D_8$ であり、$V_4$ を含む。
その商の非自明元は $\lambda\mapsto\lambda^{-1}$ と作用するため、
根配置の自己同型群が $V_4$ より大きくなるのは $\lambda=-1$ の場合に限られる。
$\lambda=1$ は四根相異の条件で除かれている。
$\chi=(1+\lambda)^2/(1-\lambda)^2$ なので、この例外はちょうど $\chi=0$ である。
一般慣性の既存の生成元と関係は $\lambda\ne0,1$ の全範囲で定義されるので、
$\lambda\ne-1$ では慣性は $D_8\times C_2$ となる。
特に別の pairing に関する harmonic 値や equianharmonic 値からは、追加の慣性は生じない。

$\lambda=-1$ とし、$\iota^2=-1$ を選ぶ。$f_1=uv$, $f_2=u^2-v^2$ に対して、
列 $(u,v,c,d)^{\mathsf t}$ への置換行列
\[
 R=\operatorname{diag}\!\left(\begin{pmatrix}0&-1\\1&0\end{pmatrix},\iota,\iota\right),
 \quad
 T=\operatorname{diag}\!\left(\begin{pmatrix}-1&1\\1&1\end{pmatrix},
              \begin{pmatrix}0&\iota\\-2\iota&0\end{pmatrix}\right),
 \quad S=\operatorname{diag}(1,-1,\iota,1)
\]
を取る。$R,T$ は既存の一般慣性を生成する変換であり、$S$ はその外の根配置の対称性を持ち上げる。
これらと $e_1,e_2$ は $H_0$ を生成する。
根配置の群が位数八、核が $\mu_2^2$ なので位数は三十二である。
さらに $S^2=e_1$ であり、$S$ は $e_1,e_2$ と可換である。
従って $S$ の根配置への像の任意の持上げは二乗が $e_1$ となり、位数二ではない。
根配置のこの元は位数二なので、群拡大に section は存在しない。
一根共有の場合の $D_8$ は既存の生成元の計算による。

一致類では $f_1=f_2=uv$ と選ぶ。
$O(uv)\simeq\mathbb G_m\rtimes C_2$ は $W=ku\oplus kv$ に作用し、
$D_8$ は $kc\oplus kd$ に符号付き置換行列として作用する。
二作用を block diagonal に並べて射影化する写像
\[
 O(uv)\times D_8\longrightarrow\operatorname{Aut}_{\mathrm{gr}}(B_{uv,uv})/D
\]
を考える。
右辺の任意の元は、$uv$ を定数倍する $A\in\operatorname{GL}(W)$ と、
同じ定数の平方根を係数に持つ $c,d$ の monomial 変換で表される。
その係数の一つを用いて全体を次数スカラーで割ると、$A$ は $O(uv)$ に入り、
$c,d$ の係数は $\pm1$ となる。従ってこの写像は fppf 全射である。
核は block diagonal 行列がスカラーになる条件で決まり、
$(I_2,I_2)$ と $(-I_2,-I_2)$ の二元だけである。
これで式~\eqref{eq:noncubic-coincident-full-stabilizer} を得る。
$O(uv)\times D_8$ は十六個の連結成分を持ち、核の非自明元はその単位成分に入らないので、
商の連結成分数は八である。単位成分にはこの核との交わりがなく、$\mathbb G_m$ そのものである。

標数零の有限型群スキームは smooth である。
従って一致類以外で慣性の Lie 代数は零、そこでスタックは Deligne--Mumford である。
一致類では一次元の Lie 代数があるので Deligne--Mumford ではない。
\end{proof}

\begin{proposition}[一致類の閉 gerbe と接表現]
\label{prop:noncubic-coincident-gerbe-tangent}
一致類の reduced closed substack は $BH_{\mathrm{eq}}$ と同値であり、
$\mathfrak N$ への埋込みは余次元二の正則閉埋込みである。
その法空間、すなわち一致点での一次変形空間は二次元で、
$T(t):(u,v,c,d)\mapsto(tu,t^{-1}v,c,d)$ の作用の重みは $2,-2$ である。
$u,v$ の交換はこの二つの重み空間を交換し、$\tau$ は両者に $-1$ で作用し、
$e_1,e_2$ は恒等に作用する。
\end{proposition}

\begin{proof}
$P$ 上で二形式が比例する locus は
\[
 \{(f,\rho f): f\in Q^{\mathrm{nd}},\ \rho\in\mathbb G_m\}
 \simeq Q^{\mathrm{nd}}\times\mathbb G_m
\]
であり、smooth な四次元閉部分である。
閉性は二つの係数ベクトルを並べた $2\times3$ 行列の二次小行列式で与えられる。
$P$ は smooth な六次元スキームなので、この包含は余次元二の正則閉埋込みである。
$G$ はこの locus に推移的に作用し、その安定化群が $H_{\mathrm{eq}}$ であるため、
商は $BH_{\mathrm{eq}}$ となる。正則埋込みは商へ降下する。

$(uv,uv)$ において、二形式の個別のスカラー変化と共通の
$\operatorname{GL}(W)$ による変化を除くと、接空間は
第二形式と第一形式の差の $u^2,v^2$ 係数で表される。
これを $(a,c)$ と書く。
$T(t)$ は $(a,c)\mapsto(t^2a,t^{-2}c)$ と作用し、
$u,v$ の交換は $(a,c)\mapsto(c,a)$ と作用する。
二形式の交換は差を反転するので $\tau$ は $-1$ となる。
符号変更 $e_1,e_2$ は形式自体に作用しないので接空間でも恒等である。
これは商に用いた中央 $\mu_2$ とも整合する。
\end{proof}


\begin{proposition}[一致根対における連結慣性の接表現]
\label{prop:noncubic-coincident-torus-weights}
$x\in\mathfrak N$ を一致根対の点とし、
$T=\operatorname{Aut}(x)^\circ\simeq\mathbb G_m$ とする。
指標 $\varepsilon$ は、中央二重平方表示における
\[
 u\longmapsto z u,\qquad v\longmapsto z^{-1}v,
 \qquad c\longmapsto c,\qquad d\longmapsto d
\]
の原始指標として取る。その逆指標に取り替えても以下の表示は変わらない。
このとき $T$-表現として
\[
 \begin{split}
 \operatorname{HH}^1(\Lambda_x)&\simeq k,\\
 T^0_x\mathfrak H_2\simeq\operatorname{HH}^2(\Lambda_x)
   &\simeq k_{2\varepsilon}\oplus k_{-2\varepsilon}\oplus k^{\oplus2},\\
 T^0_x\mathfrak N&\simeq k_{2\varepsilon}\oplus k_{-2\varepsilon},\\
 N_{\mathfrak N/\mathfrak H_2,x}&\simeq k^{\oplus2}.
 \end{split}
\]
特に二つの横断方向に対する連結慣性の作用は自明である。
ここでは自己同型群の非連結成分の作用を主張しない。
\end{proposition}
\begin{proof}
$f_1=uv$, $f_2=-uv$ の良い対角根を使う。
第~\ref{sec:noncubic-hh-roots}節の Grassmann 模型において
\[
 (g_A,g_B,g_C,g_K)=(T_z,T_z^{-1},T_z,T_z),
 \qquad T_z=\operatorname{diag}(z,1)
\]
が安定化トーラスである。その $H$ 上の行列は
\[
 \rho_z=\operatorname{diag}
 (1,z^{-1},z,1,\ 1,z,z^{-1},1).
\]
式~\eqref{eq:nchh-plane} の係数行列を $L$ とすれば
\[
 \rho_zL=L\operatorname{diag}(z,z^{-1},1,1),
 \qquad\rho_z^tJ^{-1}\rho_z=J^{-1}.
\]
従って交代形式は正確に保たれ、conformal 係数による指標の補正はない。
tautological 四平面の指標は $(\varepsilon,-\varepsilon,0,0)$ である。
このトーラスは有効群への写像の核を持たず、$D=\operatorname{Hom}(E_0,E_2)$
の projective 表現で原始指標が現れるため、上の $z$ は本文の内在的トーラスの
パラメータ、またはその逆に一致する。非自明な冪による isogeny ではない。

Lagrangian 接空間 $\operatorname{Sym}^2L^*$ の指標分解は
\[
 k_{2\varepsilon}\oplus k_{-2\varepsilon}
 \oplus k_\varepsilon^{\oplus2}\oplus k_{-\varepsilon}^{\oplus2}
 \oplus k^{\oplus4}.
\]
$\mathfrak g^+$ は二つの $\mathfrak{gl}_2$ から自明な一次元中心を除いたものであり、
その分解は
\[
 \mathfrak g^+
 \simeq k_\varepsilon^{\oplus2}\oplus k_{-\varepsilon}^{\oplus2}\oplus k^{\oplus3}.
\]
軌道写像の核 $\operatorname{HH}^1(\Lambda_x)$ は安定化トーラスの Lie 代数であり、
トーラスの共役作用は自明である。従って軌道接空間は
\[
 k_\varepsilon^{\oplus2}\oplus k_{-\varepsilon}^{\oplus2}\oplus k^{\oplus2}
\]
となる。Lagrangian chart の \texorpdfstring{\'etale}{etale} 性から、
この余核は内在的な全 $\operatorname{HH}^2$ に一致し、表示した分解を得る。

交点二線 $\ell=L\cap U$, $m=L\cap U^*$ は上の四基底の最後の二つで生成され、
両者に対するトーラス作用は自明である。
二つの行列式の微分はそれぞれ $q|_{m\times m}$ と $q|_{\ell\times\ell}$ であり、
上で確認した交代形式の不変性により、この同定はトーラス同変である。
従って境界の法表現は自明な二次元表現となる。
接空間の完全列とトーラスの線型簡約性から、境界の接表現は
$k_{2\varepsilon}\oplus k_{-2\varepsilon}$ となる。
\end{proof}

% !TEX root = ../nc-quadric-classification.tex
% 2026-11-21 / v1.34
\section{可換 quadric の全 Artin 変形と相対 AS4 被覆}
\label{sec:quadric-artin-hull}

以下では $(A,\mathfrak m,k)$ を標数零の体上の有限次元可換局所 Artin 環とする。
本節は一般の相対 canonical hull 定理~\ref{thm:universal-cubic-hull} を
可換 quadric の Artin 変形へ適用する。
主眼は、任意の pointed Abel 変形から、同じ構造対象を持つ exact ample な
Serre 平方根と相対 AS4 被覆を回復することである。
二周期の $\mathbb Z$-代数を用い、cubic 代数の一周期性は仮定しない。

\paragraph{入力を有限自由な tensor データとして固定する。}
$D$ を可換 quadric の cubic $\mathbb Z$-代数の $A$-平坦変形とし、
$V_i=D_{i,i+1}$、$R_i\subset V_i\otimes V_{i+1}\otimes V_{i+2}$ を
それぞれ階数二の一次生成元・三次関係の自由加群とする。
三次の標準解消は相対完全であり、
\[
 W_i=(R_i\otimes V_{i+3})\cap(V_i\otimes R_{i+1})
\]
は階数一で、任意の生成元 $w_i$ は両方の表示で非退化とする。
この非退化性は閉ファイバーでの非退化性と Nakayama から従う。
次数二の周期 $\beta:D\to D(2)$ が与えられたとし、
その偶数部分の graded 化を $B$ とする。
特に $B_1$ は階数四、$B_n$ は階数 $(n+1)^2$ である。

$B$ は二次環で相対 Koszul、二次双対 $E=B^!$ の各次数は有限自由で、
その階数は $1,4,7,8,8,\ldots$ と仮定する。
これらは後述の helix からの入力では既存の相対 Koszul 補題により成立する。
双対の積は tensor 部分加群の包含を双対化する規約で書く。
全ての $i$ に対して
\[
 S_i=R_i\otimes V_{i+3}+V_i\otimes R_{i+1}
 \subset V_i\otimes V_{i+1}\otimes V_{i+2}\otimes V_{i+3}
\]
と置く。これは階数七であり、偶数 $i$ では $\check B$ の二次関係加群である。

\begin{lemma}[平坦な cubic 変形の syzygy と四周期]
\label{lem:relative-cubic-tensors-artin}
$D$ の各成分が有限自由で、閉ファイバーが可換 quadric の cubic
$\mathbb Z$-代数であるとする。
このとき上記の三次提示、$R_i,W_i$ の自由性と底変更、非退化性、
三次の標準解消の相対完全性が成立する。
また生成元 $w_i\in W_i$ を選ぶと、$w_i$ の巡回から
$\theta_i:V_i\simeq V_{i+4}$ が定まり、$D\simeq D(4)$ を誘導する。
\end{lemma}
\begin{proof}
証明は相対モジュライ別稿 \cite[\texttt{quadric_artin_hull} 節]{QuadricModuli2026} に置く。
ここでは後続の分類側から参照するため、仮定と主張を保持する。
\end{proof}

\begin{lemma}[canonical functional の相対的な正規性]
\label{lem:relative-canonical-hull-functional}
$w_i$ が定める二つの coevaluation に対する評価写像を
$R_i\otimes V_{i+3}$ と $V_i\otimes R_{i+1}$ 上に取る。
両者は $W_i$ 上で一致し、全射
\[
 \phi_i:S_i\longrightarrow A,\qquad \phi_i(w_i)=2
\]
を定める。その核 $H_i$ は階数六の直和因子である。
$\phi_i$ の張る直線は $w_i$ の単元倍によらず、$\beta$ と整合する。
さらに $E$ の二次元に対応するこれらの元は、正規化列をなす。
従って $\beta$ に沿って graded 化すれば、$E_2$ に正規元
$\phi$ が得られる。
\end{lemma}
\begin{proof}
証明は相対モジュライ別稿 \cite[\texttt{quadric_artin_hull} 節]{QuadricModuli2026} に置く。
ここでは後続の分類側から参照するため、仮定と主張を保持する。
\end{proof}

\begin{proposition}[正規 functional からの相対 AS4 被覆]
\label{prop:relative-artin-hull-from-functional}
前補題の入力を満たすとき、$B$ は全次数平坦な二次 AS4 被覆
$\widetilde B\twoheadrightarrow B$ を持ち、その核は二次の左右正則正規元で生成される。
\end{proposition}
\begin{proof}
証明は相対モジュライ別稿 \cite[\texttt{quadric_artin_hull} 節]{QuadricModuli2026} に置く。
ここでは後続の分類側から参照するため、仮定と主張を保持する。
\end{proof}

\begin{remark}[可換ファイバーでの正規化]
標準二周期を持つ Segre 環では、$\phi_0$ は六つの交換関係を零に送り、
$q=x_{00}x_{11}-x_{01}x_{10}$ を $1$ に送る。
これは \cite[(5.3)]{VdB2011} の
\[
 w=u_0v_0u_1v_1-u_1v_0u_0v_1-u_0v_1u_1v_0+u_1v_1u_0v_0
\]
を逆 coevaluation の定義へ代入して確かめられる。
従って $E_k/(\phi_0)$ は四変数の外積代数であり、
その双対 hull は四変数多項式環となる。
\end{remark}

\begin{theorem}[可換 quadric の全 Artin 変形の相対実現]
\label{thm:quadric-all-artin-h2}
$Q=\mathbb P^1_k\times\mathbb P^1_k$ とし、
$(\mathcal C_A,\mathcal O_A)$ を $(\operatorname{coh}Q,\mathcal O_Q)$ の
任意の flat pointed Abel 変形とする。
この変形には $s_0=(-)\otimes\mathcal O_Q(1,1)$ を還元とする
exact ample な相対 Serre 平方根 $s_A$ が存在し、その全軌道環 $B_A$ は
\[
 B_A\simeq\widetilde B_A/(g),\qquad
 \operatorname{rank}_A(B_A)_n=(n+1)^2
\]
と表示できる。ここで $\widetilde B_A$ は型 $(1,4,6,4,1)$ の相対 quadratic
AS 正則環であり、$g$ は二次の左右正則正規元である。
従って $(\mathcal C_A,\mathcal O_A)$ は、任意の有限型 $A$-代数への基底変更を
備えた定義~\ref{def:relative-h2} の相対 $\mathfrak H_2$ 族へ昇格する。
根と被覆は存在条件であり、新しい標識として加えない。
\end{theorem}
\begin{proof}
証明は相対モジュライ別稿 \cite[\texttt{quadric_artin_hull} 節]{QuadricModuli2026} に置く。
ここでは後続の分類側から参照するため、仮定と主張を保持する。
\end{proof}

\begin{corollary}[可換点での形式変形問題]
\label{cor:quadric-formal-deformation-equivalence}
$Q$ の相対 $\mathfrak H_2$ 変形から flat pointed Abel 変形への忘却は、
局所 Artin $k$-代数上の変形 groupoid の同値を与える。
この変形問題は障害を持たず、その接空間は
\[
 \operatorname{HH}^2(Q)=H^0(Q,\mathcal O_Q(2,2)),\qquad \dim=9
\]
である。無限小自己同型の空間は $\operatorname{HH}^1(Q)=H^0(Q,T_Q)$、次元六である。
\end{corollary}
\begin{proof}
証明は相対モジュライ別稿 \cite[\texttt{quadric_artin_hull} 節]{QuadricModuli2026} に置く。
ここでは後続の分類側から参照するため、仮定と主張を保持する。
\end{proof}

この定理は可換点における全 Artin 基底の認識を与える。
一般有限型基底への代数化、一般の非可換点での同様の認識、
内在的 $\mathfrak H_2$ 全体の代数性は別に証明する必要がある。

\subsection{cubic 点での条件付き Artin 認識}

\begin{proposition}[適合 cubic helix と Hochschild 自明作用からの全 Artin 認識]
\label{prop:cubic-artin-hhtrivial-conditional}
標数零の代数閉体 $k$ 上の $(\mathcal C,\mathcal O)\in\mathcal H_2$ について、
次のデータと条件を仮定する。
\begin{enumerate}
\item $E_0=\mathcal O$ を含む ample geometric four-helix $(E_i)_{i\in\mathbb Z}$ があり、
  $D_{ij}=\operatorname{Hom}(E_{-j},E_{-i})$ は標準的な三次三次元 AS 正則
  $\mathbb Z$-代数である。
\item exact な Serre 平方根 $s_0^2\simeq S^{-1}[2]$ があり、
  $(\mathcal O,s_0)$ は Artin--Zhang ample で、
  全ての $i$ について $s_0E_i\simeq E_{i+2}$ である。
  ここでは対象ごとの同型で十分であり、別途 coherent な周期の標識は要求しない。
\item $(s_0)_*$ は $\operatorname{HH}^2(\mathcal C)$ 上恒等に作用する。
\end{enumerate}
このとき、任意の局所 Artin $k$-代数 $A$ 上の flat pointed Abel 変形
$(\mathcal C_A,\mathcal O_A)$ は、$s_0$ の持上げを用いる全軌道環の
相対 quadratic AS4 被覆を持つ。
従って全有限型基底変更を含む相対 $\mathfrak H_2$ 族へ昇格する。
この点での相対 $\mathfrak H_2$ 変形と flat pointed Abel 変形の groupoid は、
局所 Artin 基底上で同値となる。

さらに $(s_0)_*$ が $\operatorname{HH}^1(\mathcal C)$ 上も恒等なら、
この還元を持つ Serre 平方根の持上げは Morita 自己同値の同型類として一意である。
\end{proposition}
\begin{proof}
証明は相対モジュライ別稿 \cite[\texttt{quadric_artin_hull} 節]{QuadricModuli2026} に置く。
ここでは後続の分類側から参照するため、仮定と主張を保持する。
\end{proof}

\begin{remark}[Hochschild 作用条件は対象ごとの周期性から従わない]
\label{rem:cubic-helix-root-action}
この命題は、全ての cubic 点または全ての平方根について
$(s_0)_*=1$ on $\operatorname{HH}^2$ と主張していない。
$s_0^2=S^{-1}[2]$ から得られるのは作用が対合であることまでである。
実際、可換 $Q$ 上で $\varphi(z,w)=(-z,w)$ とすると
$s_0=(-\otimes\mathcal O_Q(1,1))\circ\varphi^*$ は Serre 平方根であり、
標準 helix の各線束を対象ごとに二項先へ送る。
しかし $\partial_z\wedge\partial_w$ を負に送るので、
$\operatorname{HH}^2$ 上恒等ではない。
従って cubic helix との対象ごとの適合性から作用条件を除くことはできない。
一般 cubic 点に前命題を適用するには、適合する根で Hochschild 二次に
恒等に作用するものが存在する範囲を別に確定する必要がある。
\end{remark}

% !TeX root = ../nc-quadric-classification.tex
% 2026-11-21 / v1.12
\section{可換 quadric の全一次方向を覆う二次 AS 被覆}
\label{sec:quadric-first-order-as}

$k$ は標数零、$A=k[\varepsilon]/(\varepsilon^2)$、
$P=k[x_0,x_1,x_2,x_3]$、$q=x_0x_3-x_1x_2$ とする。
$\epsilon_{ijkl}$ は $\epsilon_{0123}=1$ で定まる交代記号である。

\begin{proposition}[一次の Jacobian 被覆]
\label{prop:quadric-first-order-jacobian-cover}
$f\in P_2/kq$ を取り、$q_i=\partial_iq$、$f_i=\partial_if$ と書く。
$P$ 上の交代二重導分 $\gamma_f$ を
\[
 \gamma_f(x_i,x_j)=\sum_{k,l=0}^3\epsilon_{ijkl}q_kf_l
\]
で定める。$f$ の $kq$ による変更はこの式を変えない。
$\mathcal A_f=P\otimes_kA$ に積
\[
 a\star b=ab+\frac{\varepsilon}{2}\gamma_f(a,b)
\]
を入れると、これは次数別有限自由の二次 AS4 被覆となり、
ベクトル空間の同一視で指定した $q\in(\mathcal A_f)_2$ は中心正則である。
商 $\mathcal B_f=\mathcal A_f/(q)$ も全次数有限自由であり、
階数は $(n+1)^2$ である。
これらの平坦性・相対 Koszul 完全性・相対 Gorenstein 双対は底変更と両立する。
任意の Noetherian $A$-代数への底変更後にも左右 noetherian である。
\end{proposition}
\begin{proof}
証明は相対モジュライ別稿 \cite[\texttt{quadric_first_order} 節]{QuadricModuli2026} に置く。
ここでは後続の分類側から参照するため、仮定と主張を保持する。
\end{proof}

ここで $q$ は次数別加群 $P\otimes A$ 内の元である。
非可換の積記号を使って $x_0\star x_3-x_1\star x_2$ と書いた元とは、
一般に $\varepsilon$ の二次補正だけ異なる。

\begin{lemma}[全 bivector 方向の回復]
\label{lem:quadric-jacobian-bivectors}
$Q=\operatorname{Proj}(P/(q))$ の $x_0\ne0$ の chart で
$z=x_2/x_0$、$w=x_1/x_0$ と置く。
$\mathcal B_f$ の次数零局所化の一次交換子は
\[
 \{z,w\}=2\frac{f}{x_0^2}
        =2f(1,w,z,zw)
\]
である。従って $f\mapsto\gamma_f|_Q$ は
$P_2/kq\simeq H^0(Q,\bigwedge^2T_Q)$ という線形同型を与える。
\end{lemma}
\begin{proof}
証明は相対モジュライ別稿 \cite[\texttt{quadric_first_order} 節]{QuadricModuli2026} に置く。
ここでは後続の分類側から参照するため、仮定と主張を保持する。
\end{proof}

\begin{proposition}[Abel 変形の Hochschild 類と係数の規約]
\label{prop:quadric-jacobian-abelian-class}
$\operatorname{qgr}\mathcal B_f$ を閉 fiber $\operatorname{coh}Q$ の
pointed flat Abel 一次変形とみなしたとき、その Hochschild 類は、
各標準 affine chart 上で二コサイクル $\gamma_f/2$ により表される。
HKR の規約を
\[
 v\wedge w\longmapsto
 \left[(a,b)\longmapsto
   \frac{v(a)w(b)-w(a)v(b)}2\right]
\]
と固定すると、対応する bivector は
\[
 \pi_f=2f(1,w,z,zw)\,\partial_z\wedge\partial_w
 \quad (x_0\ne0)
\]
である。従って任意の指定された一次類
$u\leftrightarrow h(z,w)\partial_z\wedge\partial_w$
に対して、$f$ の制限を $h/2$ と取ればこの類を実現できる。
\end{proposition}
\begin{proof}
証明は相対モジュライ別稿 \cite[\texttt{quadric_first_order} 節]{QuadricModuli2026} に置く。
ここでは後続の分類側から参照するため、仮定と主張を保持する。
\end{proof}

\begin{lemma}[近恒等の補正と右 Zhang twist の規約]
\label{lem:quadric-first-order-root-twist}
$\delta\in\mathfrak{sl}_2\oplus\mathfrak{sl}_2$ を Segre 表示の $P_1$ に作用させ、
$\alpha=1+\varepsilon\delta$ と置く。
これは $\mathcal A_f$ と $\mathcal B_f$ の次数付き自己同型で、$\alpha(q)=q$ である。
$\mathcal C=\operatorname{qgr}\mathcal B_f$、$F=(1)$ とし、
その全軌道環が $\mathcal B_f$ に一致すると仮定する。
右加群に作用を $\alpha^{-1}$ で twist する自己同値を $G_\alpha$ とすれば、
標識 $G_\alpha(\pi\mathcal B_f)\simeq\pi\mathcal B_f$ を
$b\mapsto\alpha(b)$ で取ることにより、
$G_\alpha$ の全軌道環上の作用は $\alpha$ となる。
従って $s=G_\alpha F$ の右規約の全軌道環は
\[
 a*_{s}b=\alpha^{\deg b}(a)*_Fb
\]
という右 Zhang twist $\mathcal B_f^\alpha$ であり、
$\mathcal A_f^\alpha$ はその相対 AS4 被覆を与える。
\end{lemma}
\begin{proof}
証明は相対モジュライ別稿 \cite[\texttt{quadric_first_order} 節]{QuadricModuli2026} に置く。
ここでは後続の分類側から参照するため、仮定と主張を保持する。
\end{proof}


\subsection{可換 quadric の全一次変形に対する適合表示}
\label{subsec:quadric-first-order-bridge}

$k$ は標数零、$A=k[\varepsilon]/(\varepsilon^2)$ とする。
$Q=\mathbf P^1_k\times\mathbf P^1_k$、$L=\mathcal O_Q(1,1)$ と置き、
命題~\ref{prop:quadric-first-order-jacobian-cover} の
$\mathcal B_f=\mathcal A_f/(q)$ を用いる。
以下の最初の補題では、変形された圏が $\mathfrak H_2$ に属するとは仮定しない。

\begin{lemma}[商圏の平坦性と全次数の飽和]
\label{lem:quadric-first-order-flat-saturated}
$\mathcal G_A=\operatorname{QGr}\mathcal B_f$、
$\mathcal C_A=\operatorname{qgr}\mathcal B_f$、
$\mathcal O_A=\pi\mathcal B_f$、$F=(1)$ と置く。
このとき $(\mathcal C_A,\mathcal O_A)$ は
$(\operatorname{coh}Q,\mathcal O_Q)$ の flat pointed Abel 変形であり、
$\mathcal O_A$ は $A$ 上平坦である。
自然な射は全次数で同型
\[
 \mathcal B_f\xrightarrow{\ \sim\ }
 \bigoplus_{n\ge0}\operatorname{Hom}_{\mathcal C_A}
       (\mathcal O_A,F^n\mathcal O_A)
\]
となる。また $\mathcal O_A$ は相対例外的である。
\end{lemma}
\begin{proof}
証明は相対モジュライ別稿 \cite[\texttt{quadric_first_order} 節]{QuadricModuli2026} に置く。
ここでは後続の分類側から参照するため、仮定と主張を保持する。
\end{proof}

\begin{lemma}[平方根補正の全近恒等方向は次数付き自己同型から来る]
\label{lem:quadric-first-order-correction-geometric}
$\mathcal G_A$ の finite foundation 表示を取り、
その近恒等な derived Picard 群の核を $K_A$ と書く。
命題~\ref{prop:artin-serre-square-correction} の補正元は、
ある $\delta\in\mathfrak{sl}_2\oplus\mathfrak{sl}_2$ と
$\alpha=1+\varepsilon\delta$ に対する $G_\alpha$ で表される。
ここで $G_\alpha$ は右加群の作用を $\alpha^{-1}$ で twist する関手である。
\end{lemma}
\begin{proof}
証明は相対モジュライ別稿 \cite[\texttt{quadric_first_order} 節]{QuadricModuli2026} に置く。
ここでは後続の分類側から参照するため、仮定と主張を保持する。
\end{proof}

\begin{theorem}[可換 quadric の全一次方向は相対 $\mathfrak H_2$ に入る]
\label{thm:quadric-all-first-order-h2}
$u\in\operatorname{HH}^2(Q)$ を任意に取る。
この一次類を持つ $(\operatorname{coh}Q,\mathcal O_Q)$ の flat pointed Abel 変形には、
$s_0=(-)\otimes\mathcal O_Q(1,1)$ の持上げ $s_A$ と相対 AS4 被覆が存在し、
定義~\ref{def:relative-h2} の全有限型基底変更を含む相対族を構成できる。
従って、この点での $\mathfrak H_2$ から pointed Abel 変形への一次忘却写像は全射である。
\end{theorem}
\begin{proof}
証明は相対モジュライ別稿 \cite[\texttt{quadric_first_order} 節]{QuadricModuli2026} に置く。
ここでは後続の分類側から参照するため、仮定と主張を保持する。
\end{proof}


\begin{corollary}[可換点の内在的な接空間]
\label{cor:quadric-intrinsic-tangent}
可換対象 $x=(\operatorname{coh}Q,\mathcal O_Q)$ において、
一次変形の忘却により自然な同定
\[
 T_x\mathfrak H_2\simeq\operatorname{HH}^2(Q)
      \simeq H^0(Q,\mathcal O_Q(2,2))
\]
を得る。従って内在的接空間は九次元である。
反 Serre 関係式を表す点を $\bar x\in[X_3/\operatorname{PGL}_9]$ と書けば、
\[
 \dim H^0T_{\bar x}[X_3/\operatorname{PGL}_9]=15,
 \qquad \dim\ker(T_{\bar x}\to\operatorname{HH}^2(Q))=6.
\]
後者の六次元は canonical 反 Serre 偏極からのずれであり、
$\mathfrak H_2$ の追加変形方向ではない。
\end{corollary}
\begin{proof}
証明は相対モジュライ別稿 \cite[\texttt{quadric_first_order} 節]{QuadricModuli2026} に置く。
ここでは後続の分類側から参照するため、仮定と主張を保持する。
\end{proof}

全 Artin 変形の認識は系~\ref{cor:quadric-formal-deformation-equivalence} で得られる。
これらの形式的結果だけから内在的 $\mathfrak H_2$ 全体の代数性は従わない。

% !TEX root = ../nc-quadric-classification.tex
% 2026-11-21 / v1.38
\section{構造対象の忘却と楕円型の無限ファイバー}
\label{sec:unpointed-elliptic-fiber}

全ての Abel 圏を一つの代数スタックに集めることは仮定しない。
ここでは、$\HH$ の同型類から、その基礎 Abel 圏の exact 同値類への
集合論的な忘却写像を考える。固定 $\cC$ でのこの写像のファイバーは、
許容される $\OO$ の同型類を
$\operatorname{Aut}^{\mathrm{ex}}(\cC)$ で割った集合である。
固定 $\cC$ の上のスタックの \emph{2-fiber} はこの商ではなく、
許容 $\OO$ の同型群oidであることに注意する。

\begin{lemma}[許容構造対象の数値類と一意性]
\label{lem:forgetting-numerical-rigidity}
$(\cC,\OO)\in\HH$ を固定し、
定理~\ref{thm:normal-helix} の foundation によって
$K_0(\cC)\simeq K_0(Q)$、$Q=\PP^1\times\PP^1$ と同定する。
$\cC$ 内の Gabriel--Krull rank-one torsion-free exceptional 対象 $L$ は
ある $a,b\in\mathbb Z$ に対して
\[
 \operatorname{ch}[L]
 =1+aH_1+bH_2+abH_1H_2
 =\operatorname{ch}[\OO_Q(a,b)]
\]
を満たす。この数値類を指定したとき、該当する $L$ の同型類は
高々一つである。特に、固定 $\cC$ 上の許容構造対象の
同型類の集合は高々可算である。
\end{lemma}
\begin{proof}
命題~\ref{prop:first-orbit-object} の Euler 形式の計算と
$r=1$ という rank 条件を使う。
$\operatorname{ch}[L]=r+aH_1+bH_2+cH_1H_2$ と書くと
\[
 1=\chi(L,L)=r^2+2rc-2ab=1+2c-2ab.
\]
従って $c=ab$。上の $K_0$ は捩れを持たず、
この Chern 指標で整数類が一意に定まる。

同じ類を持つ二つの対象 $L,L'$ を取る。
$t=S^{-1}[2]$ と置く。もし $\Ext^2(L,L')\ne0$ なら
Serre 双対性により $\Hom(L',t^{-1}L)\ne0$ であり、
rank-one torsion-free 性から $L'\hookrightarrow t^{-1}L$ が得られる。
命題~\ref{prop:first-orbit-object} の証明で用いた AZ 豊富性と
自由被覆・負切断消滅を任意の coherent 対象に適用すると、
$n\gg0$ で
$\dim\Hom(\OO,t^n M)=\chi(\OO,t^nM)$ が成り立つ。
両辺に $\Hom(\OO,t^n-)$ を施すと、Riemann--Roch により
\[
 (2n+a+1)(2n+b+1)
 \le(2n+a-1)(2n+b-1)
\]
となり、矛盾する。従って $\Ext^2(L,L')=0$。
$\chi(L,L')=1$ から $\Hom(L,L')\ne0$ であり、同様に単射
$L\hookrightarrow L'$ を得る。余核は $K_0$ 類が零なので
Hilbert 多項式も零である。
$t$ の AZ 豊富性により非零 coherent tails 対象には
十分高い次数に非零切断が存在するから、余核は零である。
\end{proof}

滑らかな楕円曲線 $E$ と次数二の線束 $L_0,L_1,L_2$ に対し、
線束の類を加法的に $\ell_i=[L_i]\in\operatorname{Pic}^2(E)$ と書く。
差を
\[
 \eta=\ell_1-\ell_0,\qquad \delta=\ell_2-\ell_0
 \quad\in\operatorname{Pic}^0(E)
\]
とする。$(E,L_0,L_1,L_2)$ が VdB 許容であることは、
この滑らかな場合には $\eta\ne0,\eta\ne\delta,\delta\ne0$
と同値である。実際、相異なる次数二の線束による二つの
二重被覆 $E\to\PP^1$ は異なる被覆対合を持つので、
積写像は閉埋込みとなる。

\begin{lemma}[楕円四つ組の二つの反射]
\label{lem:elliptic-two-reflections}
許容楕円四つ組の helix の添字を一つ進める操作 $T$ は
\[
 T:(\eta,\delta)\longmapsto(\delta-\eta,\delta)
\]
を誘導する。また、
\[
 \eta\ne2\delta
 \tag{R-open}\label{eq:elliptic-ruling-open}
\]
の下で、命題~\ref{prop:ruling-involution-mutation} による
もう一方の ruling の foundation も cubic 型であり、
対応する四つ組は
\[
 R:(L_0,L_1,L_2)\longmapsto
 (L_0L_1L_2^{-1},L_2,L_1).
\]
従って
\[
 R:(\eta,\delta)\longmapsto(2\delta-\eta,\delta),
 \qquad RT:(\eta,\delta)\longmapsto(\eta+\delta,\delta).
\]
二つの操作は、それぞれ underlying Abel 圏の exact 同値類を保つ。
$R$ は $\OO$ を保ち、$T$ は一般に $\OO$ を変更する。
\end{lemma}
\begin{proof}
helix の関係
$L_{i+3}=L_i^{-1}L_{i+1}L_{i+2}$ により、
$\ell_{i+2}-\ell_i=\delta$ は $i$ に依存しない。
従って $(L_1,L_2,L_3)$ の差は
$(\delta-\eta,\delta)$ である。
添字の変更は同じ $\mathbb Z$-代数の reindexing なので、
基礎 tails Abel 圏を変えない。

後半は ruling 交換の二写像を楕円曲線の切断で計算する。
$V_i=H^0(E,L_i)$、$W_2=H^0(E,L_0L_1)$、
$W_3=H^0(E,L_0L_1L_2)$ とすると、
元の foundation の二写像は
\[
 V_1\otimes V_0\longrightarrow W_2,\qquad
 V_2\otimes W_2\longrightarrow W_3
\]
という切断の積である。
次数二の線束の評価列を $L_0L_1$ で捩ると、
\[
 \ker(V_2\otimes W_2\longrightarrow W_3)
 \simeq H^0(E,L_0L_1L_2^{-1})
\]
を得る。ここでは $\deg(L_0L_1L_2^{-1})=2$ による
$H^1$ 消滅を用いた。
一方、$L_1$ の評価列を $L_0L_1$ で捩ると、
$V_0\to V_1^\vee\otimes W_2$ の余核は
$H^0(E,L_0L_1^2)$ に同定される。
この同定では二次元空間の
$V_1^\vee\simeq V_1\otimes(\det V_1)^{-1}$ に由来する
一次元のスカラー因子を許す。
従って交換後の二写像は、
$(L_0L_1L_2^{-1},L_2,L_1)$ から得る
隣接二つの切断乗法と一致する。

新しい隣接次数二線束が相異なる条件は
$\eta\ne2\delta$ および $\eta\ne\delta$ であり、
残る非周期性条件は $\delta\ne0$ である。
ゆえに式~\eqref{eq:elliptic-ruling-open} の下で
新しい四つ組も VdB 許容である。
命題~\ref{prop:ruling-involution-mutation} と
定理~\ref{thm:normal-helix} による qgr の再構成により、
これは元の $(\cC,\OO)$ のもう一方の normalized foundation である。
差の計算は直接に
$\eta'=2\delta-\eta,\ \delta'=\delta$ を与える。
\end{proof}

\begin{theorem}[一般楕円型における忘却の可算無限ファイバー]
\label{thm:generic-elliptic-unpointed-fiber}
標数零の代数閉体上で、滑らかな楕円許容四つ組
$U=(E,L_0,L_1,L_2)$ に対し、上記の $\eta,\delta$ が
\[
 \delta\text{ は非捩れ},\qquad
 2\eta\notin\mathbb Z\delta
 \tag{G}\label{eq:elliptic-unpointed-generic}
\]
を満たすとする。
$\cC=\qgr D(U)$ に対して、
\[
 \operatorname{Aut}^{\mathrm{ex}}(\cC)\backslash
 \bigl\{[L]:(\cC,L)\in\HH\bigr\}
\]
は可算無限集合である。
特に滑らかな楕円型 cubic locus の幾何学的一般点では、
構造対象を忘れる写像は同型類の意味で可算無限ファイバーを持つ。
これはスタック射が étale な $\mathbb Z$-被覆であるという主張ではない。
\end{theorem}
\begin{proof}
補題~\ref{lem:elliptic-two-reflections} の $R,T$ を交互に適用する。
$U_n$ を $(RT)^n U$ と書けば、その差のパラメータは
\[
 (\eta_n,\delta_n)=(\eta+n\delta,\delta),
 \qquad n\in\mathbb Z.
\]
仮定~\eqref{eq:elliptic-unpointed-generic} により
全ての途中段階でも隣接線束の一致と
$\eta_n=2\delta$ が起こらない。
従って全 $U_n$ は VdB 許容四つ組であり、
定理~\ref{thm:all-admissible-inverse} により
$(\qgr D(U_n),\pi e_0D(U_n))\in\HH$ である。
また補題~\ref{lem:elliptic-two-reflections} により
これらの underlying Abel 圏は全て $\cC$ と exact 同値である。

これらの pointed 同値類が異なることを示す。
系~\ref{cor:two-intrinsic-foundations} により、pointed exact 同値は
二つの内在的 foundation の非順序対を保つ。
$U_n$ のその二つは $U_n$ と $RU_n$ の geometric cubic foundation である。
従って二つの pointed 対象が同値なら、
Van den Bergh の四つ組の復元
\cite[arXiv v3, Theorem 5.5.10]{VdB2011} により、
$U_n$ または $RU_n$ が $U_m$ または $RU_m$ と、
曲線の自己同型を介して同型になる。
楕円曲線の自己同型は平行移動と原点を保つ有限群の積であり、
前者は線束の差 $\eta,\delta$ に作用しない。
後者の $\sigma$ が $\delta$ を固定し、$\delta$ が非捩れなら、
$\sigma=1$ である。
従って
\[
 \eta+n\delta=\eta+m\delta
 \quad\text{または}\quad
 \eta+n\delta=2\delta-\eta-m\delta
\]
が必要となる。
第一式は $n=m$ を、第二式は
$2\eta=(2-n-m)\delta$ という仮定への矛盾を与える。
よって異なる $n$ は異なる pointed 同値類を与える。

一方、補題~\ref{lem:forgetting-numerical-rigidity} により
固定 $\cC$ 上の許容構造対象は高々可算個である。
その自己同値による商も高々可算なので、結論を得る。
幾何学的一般点では $\eta,\delta$ が独立な Picard パラメータであり、
全ての等式 $m\delta=0$ と $2\eta=n\delta$ を避ける。
\end{proof}

\begin{remark}[既存研究と残る問題]
上の定理は幾何点でのファイバーの濃度を決めるが、
自然な $\mathbb Z$-torsor 構造や相対ファイバーの代数性は与えない。
Kitamura--Okawa の2022年 Oberwolfach 報告\cite{OkawaKitamura2022OWR}は
cubic $\mathbb Z$-代数の無標識同値に $D_\infty$ が現れることと
rank-one torsion-free line bundle の $\mathbb Z^2$ 分類を述べる。
同報告で詳細論文は当時準備中とされているので、
本節の証明はその未照合の全分類定理に依存させず、
本稿の二つの foundation と楕円評価列から構成した。
\end{remark}


\section{同値類の分類後に残る問題}
\label{sec:open-problems}

定理~\ref{thm:complete-h2-classification} は、体上の全対象を分類し、
構成するデータと必要十分な同値判定を与える。
以下では、この結論に含めない構造を区別する。

\begin{problem}[幾何データ上での同値判定の簡約]
全許容四つ組について成立する有限判定を、曲線・直線束の変換による
より小さい正規形や幾何的な対応で記述せよ。
特殊な自己同型を持つ点と全ての退化型を含め、既存の tensor の GIT 商と比較せよ。
この簡約を待たずに、主定理の有限判定は全同値類を区別する。
\end{problem}

\begin{problem}[平方根と表示の忘却]
固定 $(\cC,\OO)$ 上の適合する平方根を分類し、その自己同値群の作用を求めよ。
さらに固定平方根に対する ambient 表示を分類し、
これらの二段階の忘却が既存 tensor モジュライとどう関係するかを記述せよ。
可換点の四つの平方根や固定切断環の非一意な ambient を、
全非可換点での有限性・一般次数の結論に外挿してはならない。
\end{problem}

\begin{remark}[全体の Deligne--Mumford 性]
可換対象 $(\coh Q,\mathcal O_Q)$ における pointed 自己同値群は
\[
 \operatorname{Aut}(Q)=(\operatorname{PGL}_2\times\operatorname{PGL}_2)\rtimes S_2
\]
である。実際、\cite[Theorem 3.1]{BondalOrlov2001} による幾何自己同型、
線束、cohomological shift の記述で、exact 性と $\mathcal O_Q$ の保存により
最後の二つの因子が消える。幾何自己同型は族でも作用するので、
この自己同値を通常の慣性として保持する全体のモジュライは、
可換点で六次元の慣性を持つ。従って目標とする代数スタックは全体では
Deligne--Mumford スタックにはならない。
\end{remark}

\begin{problem}[境界の相対表示と framed 局所構造]
定理~\ref{thm:global-h2-smooth-finite-type} により $\HHstack$ の有限型代数性は成立し、
系~\ref{cor:global-xas-smooth-scheme} により framed 実現空間は smooth スキームとなる。
有限 foundation の ruling 商 $\mathfrak Q_{\mathrm{fnd}}$ との相対比較を具体化し、
境界で確立した相対商表示（定理~\ref{thm:noncubic-intrinsic-presentation-equivalence}）を踏まえて、
全体の相対 groupoid を比較せよ。
全体の既約性と被約境界の滑らかさ・余次元二は、
命題~\ref{prop:global-h2-irreducible} と定理~\ref{thm:noncubic-boundary-smooth-codim-two} で成立した。
既存の smooth proper dg 特殊化族そのものの相対 $\HHstack$ 認識は、各族について別に確認する。
この内在的な既約性の証明は、実際の相対族の連続性を使い、
命題~\ref{prop:h2-component-density} の誘導位相との同一視を必要としない。

また $X_{\mathrm{as}}\to X_3$ の局所閉性と具体的な認識方程式を調べよ。
一致根対では命題~\ref{prop:global-h2-not-ambient-open} が開性を否定する。
余分な一次方向は系~\ref{cor:artin-polarization-retraction} の
$\operatorname{HH}^1$ に対応する偏極変形である。
全スキーム上の商による延長と、任意 Noetherian 基底上で同じ Abel 公理を課す関手との一致も、
極限保存と canonical 導来圏の比較を伴う別の課題である。
\end{problem}

\bibliographystyle{alpha}
\bibliography{references}
\end{document}
